%%==========================================================================
%% This is the very first version of a Streetwise LaTeX template. By
%% chosing the document class, it is easily convertible to a two-column IEEE
%% paper format, provided any conflicting packages and (new)commmands are
%% removed.
%%
%% Erwin de Gelder, January 6, 2018
%%==========================================================================

%%==========================================================================
%% Title:
%% Author(s):
%% To be published in:
%% Date of submission:
%% Date of submission R1:
%% Date of submission R2:
%% Date of final submission:
%%==========================================================================


%% Generic
\documentclass[10pt,final,a4paper,oneside,onecolumn]{article}

%% IEEE
%% Document class options: replace "draftcls" by "final" for final document.
%% All other options may just as well be omitted because they are the default values.
%\documentclass[10pt,final,journal,letterpaper,twoside,twocolumn]{IEEEtran}


%%==========================================================================
%% Document automation
%%==========================================================================

\def\reptitle{Storage of real-world scenarios}
\def\repauthor{Erwin de Gelder, Jan-Pieter Paardekooper}


%%==========================================================================
%% Packages
%%==========================================================================

\usepackage[a4paper,left=3.5cm,right=3.5cm,top=3cm,bottom=3cm]{geometry} %% change page layout; remove for IEEE paper format
\usepackage[T1]{fontenc}                        %% output font encoding for international characters (e.g., accented)
\usepackage[cmex10]{amsmath}                    %% math typesetting; consider using the [cmex10] option
\usepackage{amssymb}                            %% special (symbol) fonts for math typesetting
\usepackage{amsthm}                             %% theorem styles
\usepackage{dsfont}                             %% double stroke roman fonts: the real numbers R: $\mathds{R}$
\usepackage{mathrsfs}                           %% formal script fonts: the Laplace transform L: $\mathscr{L}$
\usepackage[pdftex]{graphicx}                   %% graphics control; use dvips for TeXify; use pdftex for PDFTeXify
\usepackage{array}                              %% array functionality (array, tabular)
\usepackage{upgreek}                            %% upright Greek letters; add the prefix 'up', e.g. \upphi
\usepackage{stfloats}                           %% improved handling of floats
\usepackage{multirow}                           %% cells spanning multiple rows in tables
%\usepackage{subfigure}                         %% subfigures and corresponding captions (for use with IEEEconf.cls)
%\usepackage{subfig}                             %% subfigures (IEEEtran.cls: set caption=false)
\usepackage{fancyhdr}                           %% page headers and footers
\usepackage[official,left]{eurosym}             %% the euro symbol; command: \euro
\usepackage{appendix}                           %% appendix layout
\usepackage{xspace}                             %% add space after macro depending on context
\usepackage{verbatim}                           %% provides the comment environment
\usepackage[dutch,USenglish]{babel}             %% language support
\usepackage{wrapfig}                            %% wrapping text around figures
\usepackage{longtable}                          %% tables spanning multiple pages
\usepackage{pgfplots}                           %% support for TikZ figures (Matlab)
\usepackage[breaklinks=true,hidelinks,          %% implement hyperlinks (dvips yields minor problems with breaklinks;
            bookmarksnumbered=true]{hyperref}   %% IEEEtran: set bookmarks=false)
%\usepackage[hyphenbreaks]{breakurl}            %% allow line breaks in URLs (don't use with PDFTeX)
\usepackage{lmodern} 
\usepackage{etoolbox}							%% Needed for apptocmd later
\usepackage[capitalize]{cleveref}
\crefname{figure}{Figure}{Figures}
\crefname{equation}{}{}
\usepackage{units}
\usepackage{subcaption}

\usepackage[utf8]{inputenc}   				 	%% utf8 support (required for biblatex)
\usepackage[style=ieee,doi=false,isbn=false,url=false,date=year,minbibnames=15,maxbibnames=15,backend=biber]{biblatex}
%\renewcommand*{\bibfont}{\footnotesize}		%% Use this for papers
\setlength{\biblabelsep}{\labelsep}
\bibliography{../bib}

\usepackage{csquotes}							%% Quoted texts are typeset according to rules of main language

%%==========================================================================
%% Fancy headers and footers
%%==========================================================================

\newtoggle{standalone}
\toggletrue{standalone}
\pagestyle{fancy}                                       %% set page style
\fancyhf{}                                              %% clear all header & footer fields
\iftoggle{standalone}{%
	\fancyhead[L]{\includegraphics[width=15mm]{Streetwise}} %% define headers (LE: left field/even pages, etc.)
	\fancyhead[R]{\small\emph{\reptitle}}                   %% similar
	\fancyfoot[C]{\thepage}                                 %% define footer
	\setlength{\headheight}{18pt}                           %% increase head height to accommodate an oversized picture in the header
	\renewcommand*{\headrulewidth}{0.25pt}                  %% header line width
	%% Redefine the default "plain" page style (automatically activated by \maketitle, \section, ...)
	\fancypagestyle{plain}{%
		\fancyhf{}
		\fancyhead[C]{\includegraphics[width=30mm]{Streetwise}}
		\renewcommand{\headrulewidth}{0pt}
		\renewcommand{\footrulewidth}{0pt}
		\setlength{\headheight}{36pt}
	}
}{%
	\renewcommand*{\headrulewidth}{0pt}                  	%% No line in this case
	%% Redefine the default "plain" page style (automatically activated by \maketitle, \section, ...)
	\fancypagestyle{plain}{%
		\fancyhf{}
	}
}
\renewcommand*{\footrulewidth}{0pt}                     %% footer line width


%%==========================================================================
%% TikZ figures
%%==========================================================================

\newlength\figurewidth
\setlength\figurewidth{0.35\textwidth}              %% set figure width
\newlength\figureheight
\setlength\figureheight{0.3\textwidth}              %% set figure height
\pgfplotsset{every axis/.append style={
    scaled y ticks=false,
    scaled x ticks=false,
    y tick label style={/pgf/number format/fixed},
    x tick label style={/pgf/number format/fixed},
    legend style={font=\small}},
    compat=1.9}                                     %% PGFPlots package options
\usetikzlibrary{shapes.geometric, arrows, arrows.meta}
\usepgfplotslibrary{groupplots}
\usetikzlibrary{external}                           %% Create pdf figures from TikZ. Use PDFTeXify ...
\tikzexternalize[prefix=./tikz/]                    %% ... with --tex-option=--shell-escape switch.
%\tikzset{external/force remake}                    %% force pdf figure update


%%==========================================================================
%% User-defined commands
%%==========================================================================

\newcommand*{\mat}[1]{\mathbf{#1}}                              %% matrix/vector notation
\newcommand*{\matsym}[1]{\boldsymbol{#1}}                       %% matrix/vector notation for Greek letters
\newcommand*{\T}{^{\scriptscriptstyle\mathsf{T}}}               %% transpose operator
\newcommand*{\Hr}{^{\scriptscriptstyle\mathsf{H}}}              %% conjugate transpose operator
\newcommand*{\ud}{\mathrm{\,d}}                                 %% differential operator (upright d)
\newcommand*{\defeq}{\mathrel{\mathop:}=}                       %% definition sign :=
\newcommand*{\eqdef}{=\mathrel{\mathop:}}                       %% definition sign =:
\newcommand*{\ip}[2]{\left\langle#1\,{,}\,#2\right\rangle}      %% inner product
\newcommand*{\real}[1]{\mathrm{Re}(#1)}                         %% real part
\newcommand*{\imag}[1]{\mathrm{Im}(#1)}                         %% imaginary part
\newcommand*{\lsup}[1]{{}^{#1}\!}                               %% left superscript
\newcommand*{\hi}[1]{$^\text{#1}$}                              %% superscript in normal text
\newcommand*{\lo}[1]{$_\text{#1}$}                              %% subscript in normal text
\newcommand*{\w}[1]{\mathrm{#1}}                                %% multiple character super-/subscript in math mode
\newcommand*{\Ln}{{\ensuremath{\mathcal{L}}}\xspace}            %% L norm
\newcommand*{\Lone}{{\ensuremath{\mathcal{L}_1}}\xspace}        %% L_1 norm
\newcommand*{\Ltwo}{{\ensuremath{\mathcal{L}_2}}\xspace}        %% L_2 norm
\newcommand*{\Linf}{{\ensuremath{\mathcal{L}_\infty}}\xspace}   %% L_inf norm
\newcommand*{\Lp}{{\ensuremath{\mathcal{L}_p}}\xspace}          %% L_p norm
\newcommand*{\Htwo}{{\ensuremath{\mathcal{H}_2}}\xspace}        %% H_2 norm
\newcommand*{\Hinf}{{\ensuremath{\mathcal{H}_\infty}}\xspace}   %% H_inf norm
\newcommand*{\Hp}{{\ensuremath{\mathcal{H}_p}}\xspace}          %% H_p norm
\newcommand*{\capskip}{\vspace{-12pt}}                          %% caption skip for figures with subfloats
\newcommand*{\etal}{et al.}                                     %% may be required for Natbib bibliography styles
\renewcommand*{\qedsymbol}{$\blacksquare$}                      %% redefine the end-of-proof symbol
\renewcommand*{\labelitemi}{$\bullet$}                          %% first level item list bullet
\renewcommand*{\labelitemii}{$-$}                               %% second level item list bullet
%\renewcommand*{\theenumi}{\textit{\roman{enumi}}}               %% first level enumerator
\renewcommand*{\labelenumi}{\theenumi.}
\renewcommand*{\theenumii}{\textit{\alph{enumii}}}              %% second level enumerator
\renewcommand*{\labelenumii}{\theenumii.}
\DeclareMathOperator{\tr}{tr}                                   %% trace of a matrix
\DeclareMathOperator{\sgn}{sgn}                                 %% signum function
\DeclareMathOperator{\atan}{atan}                               %% arc tangent

\newcommand{\deuclid}{d_{\textup{Euclidian}}}
\newcommand{\ddtw}{d_{\textup{DTW}}}
\newcommand{\imagi}{\textbf{j}}
\newcommand{\profile}[1]{\textbf{v}_{#1}}
\newcommand{\barprofile}[1]{\bar{\textbf{v}}_{#1}}
\newcommand{\parvstart}{v_{\textup{start}}}
\newcommand{\parvend}{v_{\textup{end}}}
\newcommand{\pardv}{v_{\Delta}}
\newcommand{\pardt}{t_{\textup{duration}}}


%%==========================================================================
%% User-defined environments
%%==========================================================================

\theoremstyle{plain}\newtheorem{definition}{Definition}[section]    %% definition
                    \newtheorem{theorem}{Theorem}[section]          %% theorem
                    \newtheorem{lemma}[theorem]{Lemma}              %% lemma
                    \newtheorem{corollary}[theorem]{Corollary}      %% corollary
                    \newtheorem{assumption}{Assumption}[section]    %% assumption
                    \newtheorem{condition}{Condition}[section]      %% condition
\theoremstyle{definition}\newtheorem{example}{Example}[section]     %% examples
\theoremstyle{remark}\newtheorem{remarkenv}{Remark}[section]        %% remarks
\newenvironment{remark}{\begin{remarkenv}}%
                       {\hfill$\blacklozenge$\end{remarkenv}}       %% end remark with a lozenge
\newenvironment{reviewer}{\itshape}{\upshape}                       %% environment for reviewer's comments


%%==========================================================================
%% Miscellaneous
%%==========================================================================

\graphicspath{{./../}{./figures/}{./more_figures/}} %% (graphicx) directory path for figures
%\setlength{\parindent}{0pt}                        %% no paragraph indentation
%\setlength{\parskip}{2ex}                          %% create empty line between paragraphs
%\interdisplaylinepenalty=2500                      %% (amsmath) allow for page breaks within multiline equations
%\numberwithin{equation}{section}                   %% (amsmath) include section number in equation numbering
%\numberwithin{figure}{section}                     %% (amsmath) include section number in figure numbering
%\numberwithin{table}{section}                      %% (amsmath) include section number in table numbering
\addtolength{\arraycolsep}{-0.5mm}                  %% squeeze matrix columns a little
\fnbelowfloat                                       %% (stfloats) put footnote below a float at the page bottom
\urlstyle{same}                                     %% (hyperref) use current font for URLs
\hypersetup{pdftitle={\reptitle},
            pdfauthor={\repauthor}}                 %% (hyperref) pdf properties title and author
%\raggedbottom                                      %% don't add inter-paragraph spacing to achieve \textheight
%\setlength\subfigcapskip{-3pt}                     %% (subfigure) distance between subfloat and subcaption
%\setlength\subfigbottomskip{-3pt}                  %% (subfigure) distance between subcaption and caption
%\renewcommand*{\subcapsize}{\small}                %% (subfigure) subcaption font size
\renewcommand*{\thesubfigure}{(\alph{subfigure})}   %% (subfig) implement 1(a) instead of 1a ...as subfigure reference
\captionsetup[subfloat]{labelformat=simple}        %% ... as subfigure reference
%\captionsetup[subfloat]{
%    farskip=0pt,
%    nearskip=8pt,
%    captionskip=1.5pt,
%    labelfont={small,bf},
%    textfont=small}                                %% (subfig) subfloat caption format
%\captionsetup[figure]{
%    labelfont={small,bf},
%    textfont=small}                                %% (subfig/caption) figure caption format
%\captionsetup[table]{
%    aboveskip=2pt,
%    labelfont={small,bf},
%    textfont={small,sc}}                           %% (subfig/caption) table caption format


% Table stuff
\usepackage{booktabs}
\usepackage{tabularx}
\setlength{\heavyrulewidth}{0.1em}
\newcommand{\otoprule}{\midrule[\heavyrulewidth]}

%%==========================================================================
%% Begin document
%%==========================================================================

\begin{document}

\selectlanguage{USenglish}

\title{\textbf{\reptitle}}
\author{\repauthor}
\date{\today}
\maketitle

\tableofcontents

\newpage

\section{Introduction}
\label{sec:introduction}

The goal of Streetwise is the storage of real-world scenarios mined from naturalistic driving data. This document provides an example (an extension of the example described in \cite{degelder2018ontology}) of several scenarios found in a approximately 40 seconds of naturalistic driving data.


\section{Scenario description}

A schematic overview of the example is shown in \cref{fig:example schematic}. Three vehicles are present in the example: a pickup truck, a sedan, and a station wagon, see \cref{fig:example schematic}(a) from left to right, respectively. The ego vehicle (the sedan), is displayed in red. The example can be described in words as follows. The ego vehicle accelerates towards its predecessor (the pickup truck), such that the distance between the ego vehicle and its predecessor becomes smaller. In the meantime, the car behind the ego vehicle (the station wagon) accelerates and takes over both other vehicles (the ego vehicle and the pickup truck). When the station wagon has passed both vehicles, the ego vehicle accelerates and performs a lane change after which it overtakes the pickup truck.

\setlength\figureheight{106pt}
\setlength\figurewidth{260pt}
\begin{figure}[b]
	\begin{minipage}[b]{\linewidth}
		\centering
		% This file was created by matplotlib2tikz v0.6.14.
\begin{tikzpicture}

\begin{axis}[
xmin=-20, xmax=20,
ymin=-5, ymax=5,
width=\figurewidth,
height=\figureheight,
tick align=outside,
tick pos=left,
x grid style={lightgray!92.026143790849673!black},
y grid style={lightgray!92.026143790849673!black},
axis background/.style={fill=white!90.0!black},
ticks=none,
hide axis
]
\path [draw=white, fill=white] (axis cs:-20,3.5)
--(axis cs:20,3.5)
--(axis cs:20,-3.5)
--(axis cs:-20,-3.5)
--cycle;

\addplot [semithick, black, forget plot]
table {%
-20 3.5
20 3.5
};
\addplot [semithick, black, forget plot]
table {%
-20 -3.5
20 -3.5
};
\addplot [semithick, black, forget plot]
table {%
-20 0
-15.5555555555556 0
};
\addplot [semithick, black, forget plot]
table {%
-6.66666666666666 0
-2.22222222222222 0
};
\addplot [semithick, black, forget plot]
table {%
6.66666666666667 0
11.1111111111111 0
};
\addplot [red, forget plot]
table {%
7.25 1.75
7.25 1.318
7.12614678899083 1.318
7.0848623853211 1.354
7.0848623853211 1.534
7.12614678899083 1.57
7.25 1.57
7.25 1.282
7.0848623853211 1.066
6.79587155963303 1.21
6.79587155963303 1.75
6.79587155963303 1.21
6.67201834862385 1.138
4.81422018348624 1.138
4.73165137614679 1.282
4.73165137614679 1.75
};
\addplot [red, forget plot]
table {%
7.25 1.75
7.25 2.182
7.12614678899083 2.182
7.0848623853211 2.146
7.0848623853211 1.966
7.12614678899083 1.93
7.25 1.93
7.25 2.218
7.0848623853211 2.434
6.79587155963303 2.29
6.79587155963303 1.75
6.79587155963303 2.29
6.67201834862385 2.362
4.81422018348624 2.362
4.73165137614679 2.218
4.73165137614679 1.75
};
\addplot [red, forget plot]
table {%
4.81422018348624 1.138
4.56651376146789 0.958
4.19495412844037 1.246
4.02981651376147 1.534
4.02981651376147 1.75
};
\addplot [red, forget plot]
table {%
4.81422018348624 2.362
4.56651376146789 2.542
4.19495412844037 2.254
4.02981651376147 1.966
4.02981651376147 1.75
};
\addplot [red, forget plot]
table {%
4.56651376146789 0.958
6.7545871559633 0.886
7.0848623853211 1.066
6.7545871559633 0.886
4.56651376146789 0.958
4.69036697247706 0.85
4.56651376146789 0.958
3.4105504587156 0.922
3.16284403669725 0.994
2.9151376146789 1.246
2.99770642201835 1.282
3.2454128440367 1.138
3.16284403669725 0.994
2.9151376146789 1.246
2.83256880733945 1.39
2.75 1.642
2.75 1.75
};
\addplot [red, forget plot]
table {%
4.56651376146789 2.542
6.7545871559633 2.614
7.0848623853211 2.434
6.7545871559633 2.614
4.56651376146789 2.542
4.69036697247706 2.65
4.56651376146789 2.542
3.4105504587156 2.578
3.16284403669725 2.506
2.9151376146789 2.254
2.99770642201835 2.218
3.2454128440367 2.362
3.16284403669725 2.506
2.9151376146789 2.254
2.83256880733945 2.11
2.75 1.858
2.75 1.75
};
\addplot [blue, forget plot]
table {%
-3.22093023255814 1.75
-3.22093023255814 1.73767123287671
-4.50290697674419 1.73767123287671
-4.50290697674419 1.75
};
\addplot [blue, forget plot]
table {%
-3.22093023255814 1.75
-3.22093023255814 1.76232876712329
-4.50290697674419 1.76232876712329
-4.50290697674419 1.75
};
\addplot [blue, forget plot]
table {%
-3.22093023255814 1.49109589041096
-4.50290697674419 1.49109589041096
-4.50290697674419 1.51575342465753
-3.22093023255814 1.51575342465753
-3.22093023255814 1.49109589041096
};
\addplot [blue, forget plot]
table {%
-3.22093023255814 2.00890410958904
-4.50290697674419 2.00890410958904
-4.50290697674419 1.98424657534247
-3.22093023255814 1.98424657534247
-3.22093023255814 2.00890410958904
};
\addplot [blue, forget plot]
table {%
-3.22093023255814 1.21986301369863
-4.50290697674419 1.21986301369863
-4.50290697674419 1.24452054794521
-3.22093023255814 1.24452054794521
-3.22093023255814 1.21986301369863
};
\addplot [blue, forget plot]
table {%
-3.22093023255814 2.28013698630137
-4.50290697674419 2.28013698630137
-4.50290697674419 2.25547945205479
-3.22093023255814 2.25547945205479
-3.22093023255814 2.28013698630137
};
\addplot [blue, forget plot]
table {%
-2.75 1.75
-2.75 1.07191780821918
-2.82848837209302 0.997945205479452
-2.9593023255814 0.997945205479452
-2.9593023255814 1.75
-2.9593023255814 0.923972602739726
-3.03779069767442 0.85
-5.99418604651163 0.874657534246575
-5.83720930232558 0.726712328767123
-5.99418604651163 0.874657534246575
-6.75290697674419 0.899315068493151
-7.09302325581395 1.02260273972603
-7.25 1.63904109589041
-7.25 1.75
};
\addplot [blue, forget plot]
table {%
-2.75 1.75
-2.75 2.42808219178082
-2.82848837209302 2.50205479452055
-2.9593023255814 2.50205479452055
-2.9593023255814 1.75
-2.9593023255814 2.57602739726027
-3.03779069767442 2.65
-5.99418604651163 2.62534246575343
-5.83720930232558 2.77328767123288
-5.99418604651163 2.62534246575343
-6.75290697674419 2.60068493150685
-7.09302325581395 2.47739726027397
-7.25 1.86095890410959
-7.25 1.75
};
\addplot [blue, forget plot]
table {%
-4.65988372093023 1.75
-4.65988372093023 0.997945205479452
-4.8953488372093 1.09657534246575
-4.8953488372093 1.75
};
\addplot [blue, forget plot]
table {%
-4.65988372093023 1.75
-4.65988372093023 2.50205479452055
-4.8953488372093 2.40342465753425
-4.8953488372093 1.75
};
\addplot [blue, forget plot]
table {%
-5.83720930232558 1.75
-5.83720930232558 1.44178082191781
-5.75872093023256 1.07191780821918
-6.28197674418605 0.973287671232877
-6.36046511627907 1.46643835616438
-6.36046511627907 1.75
};
\addplot [blue, forget plot]
table {%
-5.83720930232558 1.75
-5.83720930232558 2.05821917808219
-5.75872093023256 2.42808219178082
-6.28197674418605 2.52671232876712
-6.36046511627907 2.03356164383562
-6.36046511627907 1.75
};
\addplot [blue, forget plot]
table {%
-3.11627906976744 1.75
-3.11627906976744 1.02260273972603
-4.52906976744186 1.02260273972603
};
\addplot [blue, forget plot]
table {%
-3.11627906976744 1.75
-3.11627906976744 2.47739726027397
-4.52906976744186 2.47739726027397
};
\addplot [blue, forget plot]
table {%
17.25 1.75
17.25 1.40714285714286
17.1590909090909 1.10714285714286
14.5227272727273 1.10714285714286
14.3863636363636 1.40714285714286
14.3409090909091 1.53571428571429
14.3409090909091 1.75
};
\addplot [blue, forget plot]
table {%
17.25 1.75
17.25 2.09285714285714
17.1590909090909 2.39285714285714
14.5227272727273 2.39285714285714
14.3863636363636 2.09285714285714
14.3409090909091 1.96428571428571
14.3409090909091 1.75
};
\addplot [blue, forget plot]
table {%
14.5227272727273 1.10714285714286
14.2954545454545 0.935714285714286
14.0227272727273 0.935714285714286
13.7954545454545 1.15
13.6590909090909 1.49285714285714
13.6590909090909 1.75
};
\addplot [blue, forget plot]
table {%
14.5227272727273 2.39285714285714
14.2954545454545 2.56428571428571
14.0227272727273 2.56428571428571
13.7954545454545 2.35
13.6590909090909 2.00714285714286
13.6590909090909 1.75
};
\addplot [blue, forget plot]
table {%
14.2954545454545 0.935714285714286
15.7954545454545 0.978571428571429
15.7954545454545 1.10714285714286
15.7954545454545 0.978571428571429
17.1136363636364 0.978571428571429
17.1590909090909 1.10714285714286
17.1136363636364 0.978571428571429
16.9772727272727 0.85
13.1590909090909 0.85
12.8409090909091 1.15
13.0227272727273 1.15
13.2045454545455 0.978571428571429
13.2045454545455 0.85
};
\addplot [blue, forget plot]
table {%
14.2954545454545 2.56428571428571
15.7954545454545 2.52142857142857
15.7954545454545 2.39285714285714
15.7954545454545 2.52142857142857
17.1136363636364 2.52142857142857
17.1590909090909 2.39285714285714
17.1136363636364 2.52142857142857
16.9772727272727 2.65
13.1590909090909 2.65
12.8409090909091 2.35
13.0227272727273 2.35
13.2045454545455 2.52142857142857
13.2045454545455 2.65
};
\addplot [blue, forget plot]
table {%
12.8409090909091 1.15
12.75 1.45
12.75 1.75
};
\addplot [blue, forget plot]
table {%
12.8409090909091 2.35
12.75 2.05
12.75 1.75
};
\addplot [semithick, blue, dashed, forget plot]
table {%
12.75 1.75
12.6717171717172 1.74996447297214
12.5934343434343 1.74972010306519
12.5151515151515 1.74906980619819
12.4368686868687 1.74782908517264
12.3585858585859 1.74582576468553
12.280303030303 1.74289972634234
12.2020202020202 1.73890264367003
12.1237373737374 1.73369771713003
12.0454545454545 1.72715940913127
11.9671717171717 1.71917317904316
11.8888888888889 1.70963521820861
11.8106060606061 1.69845218495698
11.7323232323232 1.68554093961714
11.6540404040404 1.67082827953046
11.5757575757576 1.65425067406376
11.4974747474747 1.63575399962237
11.4191919191919 1.6152932746631
11.3409090909091 1.59283239470727
11.2626262626263 1.56834386735364
11.1843434343434 1.54180854729151
11.1060606060606 1.51321537131363
11.0277777777778 1.48256109332927
10.9494949494949 1.44985001937716
10.8712121212121 1.41509374263854
10.7929292929293 1.37831087845013
10.7146464646465 1.33952679931715
10.6363636363636 1.2987733699263
10.5580808080808 1.25608868215878
10.479797979798 1.21151679010327
10.4015151515152 1.16510744506897
10.3232323232323 1.11691583059853
10.2449494949495 1.06700229748114
10.1666666666667 1.01543209876543
10.0883838383838 0.962275124772575
10.010101010101 0.907605638109211
9.93181818181818 0.85150200868048
9.85353535353535 0.794046448703017
9.77525252525253 0.735324747717952
9.6969696969697 0.67542600760391
9.61868686868687 0.614442377590011
9.54040404040404 0.552468789268875
9.46212121212121 0.489602691609613
9.38383838383838 0.425943785970835
9.30555555555556 0.361593761113651
9.22727272727273 0.296656028214665
9.1489898989899 0.231235455878979
9.07070707070707 0.165438105153197
8.99242424242424 0.0993709645384162
8.91414141414141 0.033141685003238
8.83585858585858 -0.0331416850032387
8.75757575757576 -0.0993709645384178
8.67929292929293 -0.165438105153197
8.6010101010101 -0.23123545587898
8.52272727272727 -0.296656028214666
8.44444444444444 -0.361593761113652
8.36616161616162 -0.425943785970836
8.28787878787879 -0.489602691609613
8.20959595959596 -0.552468789268874
8.13131313131313 -0.614442377590012
8.0530303030303 -0.67542600760391
7.97474747474747 -0.735324747717952
7.89646464646465 -0.794046448703018
7.81818181818182 -0.851502008680481
7.73989898989899 -0.907605638109211
7.66161616161616 -0.962275124772574
7.58333333333333 -1.01543209876543
7.50505050505051 -1.06700229748113
7.42676767676768 -1.11691583059853
7.34848484848485 -1.16510744506897
7.27020202020202 -1.21151679010327
7.19191919191919 -1.25608868215878
7.11363636363636 -1.2987733699263
7.03535353535353 -1.33952679931715
6.95707070707071 -1.37831087845013
6.87878787878788 -1.41509374263854
6.80050505050505 -1.44985001937716
6.72222222222222 -1.48256109332927
6.64393939393939 -1.51321537131363
6.56565656565657 -1.54180854729151
6.48737373737374 -1.56834386735365
6.40909090909091 -1.59283239470727
6.33080808080808 -1.6152932746631
6.25252525252525 -1.63575399962237
6.17424242424242 -1.65425067406375
6.0959595959596 -1.67082827953045
6.01767676767677 -1.68554093961715
5.93939393939394 -1.69845218495698
5.86111111111111 -1.7096352182086
5.78282828282828 -1.71917317904316
5.70454545454545 -1.72715940913127
5.62626262626263 -1.73369771713003
5.5479797979798 -1.73890264367003
5.46969696969697 -1.74289972634234
5.39141414141414 -1.74582576468553
5.31313131313131 -1.74782908517264
5.23484848484848 -1.7490698061982
5.15656565656566 -1.74972010306519
5.07828282828283 -1.74996447297215
5 -1.75
-10 -1.75
};
\addplot [semithick, blue, forget plot]
table {%
-10 -1.75
-9.25 -1
};
\addplot [semithick, blue, forget plot]
table {%
-10 -1.75
-9.25 -2.5
};
\addplot [semithick, red, dashed, forget plot]
table {%
2.75 1.75
0.75 1.75
};
\addplot [semithick, red, forget plot]
table {%
1.5 2.5
0.75 1.75
1.5 1
};
\addplot [semithick, blue, dashed, forget plot]
table {%
-7.25 1.75
-9.25 1.75
};
\addplot [semithick, blue, forget plot]
table {%
-8.5 2.5
-9.25 1.75
-8.5 1
};
\end{axis}

\end{tikzpicture}%
		\subcaption{Last vehicle overtakes other vehicles.}
	\end{minipage}
	\setlength\figureheight{100pt}
	\begin{minipage}[b]{\linewidth}
		\centering
		% This file was created by matplotlib2tikz v0.6.14.
\begin{tikzpicture}

\begin{axis}[
xmin=-20, xmax=20,
ymin=-5, ymax=5,
width=\figurewidth,
height=\figureheight,
tick align=outside,
tick pos=left,
x grid style={lightgray!92.026143790849673!black},
y grid style={lightgray!92.026143790849673!black},
axis background/.style={fill=white!90.0!black},
ticks=none,
hide axis
]
\path [draw=white, fill=white] (axis cs:-20,3.5)
--(axis cs:20,3.5)
--(axis cs:20,-3.5)
--(axis cs:-20,-3.5)
--cycle;

\addplot [semithick, black, forget plot]
table {%
-20 3.5
20 3.5
};
\addplot [semithick, black, forget plot]
table {%
-20 -3.5
20 -3.5
};
\addplot [semithick, black, forget plot]
table {%
-20 0
-15.5555555555556 0
};
\addplot [semithick, black, forget plot]
table {%
-6.66666666666666 0
-2.22222222222222 0
};
\addplot [semithick, black, forget plot]
table {%
6.66666666666667 0
11.1111111111111 0
};
\addplot [red, forget plot]
table {%
7.25 1.77571428571429
7.25 1.36428571428571
7.13235294117647 1.08142857142857
6.48529411764706 0.927142857142857
4.33823529411765 0.927142857142857
4.48529411764706 0.798571428571428
4.33823529411765 0.927142857142857
3.30882352941176 0.927142857142857
2.89705882352941 1.08142857142857
2.75 1.44142857142857
2.75 1.75
2.75 1.46714285714286
3.72058823529412 1.13285714285714
};
\addplot [red, forget plot]
table {%
7.25 1.72428571428571
7.25 2.13571428571429
7.13235294117647 2.41857142857143
6.48529411764706 2.57285714285714
4.33823529411765 2.57285714285714
4.48529411764706 2.70142857142857
4.33823529411765 2.57285714285714
3.30882352941176 2.57285714285714
2.89705882352941 2.41857142857143
2.75 2.05857142857143
2.75 1.75
2.75 2.03285714285714
3.72058823529412 2.36714285714286
};
\addplot [red, forget plot]
table {%
6.16176470588235 1.03
5.77941176470588 1.13285714285714
4.89705882352941 1.13285714285714
4.48529411764706 1.03
};
\addplot [red, forget plot]
table {%
6.16176470588235 2.47
5.77941176470588 2.36714285714286
4.89705882352941 2.36714285714286
4.48529411764706 2.47
};
\addplot [red, forget plot]
table {%
6.16176470588235 1.75
6.16176470588235 1.57
6.04411764705882 1.21
6.39705882352941 1.21
6.63235294117647 1.28714285714286
6.75 1.39
6.83823529411765 1.59571428571429
6.83823529411765 1.75
};
\addplot [red, forget plot]
table {%
6.16176470588235 1.75
6.16176470588235 1.93
6.04411764705882 2.29
6.39705882352941 2.29
6.63235294117647 2.21285714285714
6.75 2.11
6.83823529411765 1.90428571428571
6.83823529411765 1.75
};
\addplot [red, forget plot]
table {%
4.69117647058824 1.75
4.69117647058824 1.41571428571429
4.72058823529412 1.21
4.01470588235294 1.03
3.86764705882353 1.62142857142857
3.86764705882353 1.75
};
\addplot [red, forget plot]
table {%
4.69117647058824 1.75
4.69117647058824 2.08428571428571
4.72058823529412 2.29
4.01470588235294 2.47
3.86764705882353 1.87857142857143
3.86764705882353 1.75
};
\addplot [blue, forget plot]
table {%
-3.22093023255814 1.75
-3.22093023255814 1.73767123287671
-4.50290697674419 1.73767123287671
-4.50290697674419 1.75
};
\addplot [blue, forget plot]
table {%
-3.22093023255814 1.75
-3.22093023255814 1.76232876712329
-4.50290697674419 1.76232876712329
-4.50290697674419 1.75
};
\addplot [blue, forget plot]
table {%
-3.22093023255814 1.49109589041096
-4.50290697674419 1.49109589041096
-4.50290697674419 1.51575342465753
-3.22093023255814 1.51575342465753
-3.22093023255814 1.49109589041096
};
\addplot [blue, forget plot]
table {%
-3.22093023255814 2.00890410958904
-4.50290697674419 2.00890410958904
-4.50290697674419 1.98424657534247
-3.22093023255814 1.98424657534247
-3.22093023255814 2.00890410958904
};
\addplot [blue, forget plot]
table {%
-3.22093023255814 1.21986301369863
-4.50290697674419 1.21986301369863
-4.50290697674419 1.24452054794521
-3.22093023255814 1.24452054794521
-3.22093023255814 1.21986301369863
};
\addplot [blue, forget plot]
table {%
-3.22093023255814 2.28013698630137
-4.50290697674419 2.28013698630137
-4.50290697674419 2.25547945205479
-3.22093023255814 2.25547945205479
-3.22093023255814 2.28013698630137
};
\addplot [blue, forget plot]
table {%
-2.75 1.75
-2.75 1.07191780821918
-2.82848837209302 0.997945205479452
-2.9593023255814 0.997945205479452
-2.9593023255814 1.75
-2.9593023255814 0.923972602739726
-3.03779069767442 0.85
-5.99418604651163 0.874657534246575
-5.83720930232558 0.726712328767123
-5.99418604651163 0.874657534246575
-6.75290697674419 0.899315068493151
-7.09302325581395 1.02260273972603
-7.25 1.63904109589041
-7.25 1.75
};
\addplot [blue, forget plot]
table {%
-2.75 1.75
-2.75 2.42808219178082
-2.82848837209302 2.50205479452055
-2.9593023255814 2.50205479452055
-2.9593023255814 1.75
-2.9593023255814 2.57602739726027
-3.03779069767442 2.65
-5.99418604651163 2.62534246575343
-5.83720930232558 2.77328767123288
-5.99418604651163 2.62534246575343
-6.75290697674419 2.60068493150685
-7.09302325581395 2.47739726027397
-7.25 1.86095890410959
-7.25 1.75
};
\addplot [blue, forget plot]
table {%
-4.65988372093023 1.75
-4.65988372093023 0.997945205479452
-4.8953488372093 1.09657534246575
-4.8953488372093 1.75
};
\addplot [blue, forget plot]
table {%
-4.65988372093023 1.75
-4.65988372093023 2.50205479452055
-4.8953488372093 2.40342465753425
-4.8953488372093 1.75
};
\addplot [blue, forget plot]
table {%
-5.83720930232558 1.75
-5.83720930232558 1.44178082191781
-5.75872093023256 1.07191780821918
-6.28197674418605 0.973287671232877
-6.36046511627907 1.46643835616438
-6.36046511627907 1.75
};
\addplot [blue, forget plot]
table {%
-5.83720930232558 1.75
-5.83720930232558 2.05821917808219
-5.75872093023256 2.42808219178082
-6.28197674418605 2.52671232876712
-6.36046511627907 2.03356164383562
-6.36046511627907 1.75
};
\addplot [blue, forget plot]
table {%
-3.11627906976744 1.75
-3.11627906976744 1.02260273972603
-4.52906976744186 1.02260273972603
};
\addplot [blue, forget plot]
table {%
-3.11627906976744 1.75
-3.11627906976744 2.47739726027397
-4.52906976744186 2.47739726027397
};
\addplot [semithick, red, dashed, forget plot]
table {%
2.75 1.75
2.67171717171717 1.74996447297214
2.59343434343434 1.74972010306519
2.51515151515152 1.74906980619819
2.43686868686869 1.74782908517264
2.35858585858586 1.74582576468553
2.28030303030303 1.74289972634234
2.2020202020202 1.73890264367003
2.12373737373737 1.73369771713003
2.04545454545455 1.72715940913127
1.96717171717172 1.71917317904316
1.88888888888889 1.70963521820861
1.81060606060606 1.69845218495698
1.73232323232323 1.68554093961714
1.6540404040404 1.67082827953046
1.57575757575758 1.65425067406376
1.49747474747475 1.63575399962237
1.41919191919192 1.6152932746631
1.34090909090909 1.59283239470727
1.26262626262626 1.56834386735364
1.18434343434343 1.54180854729151
1.10606060606061 1.51321537131363
1.02777777777778 1.48256109332927
0.949494949494949 1.44985001937716
0.871212121212121 1.41509374263854
0.792929292929293 1.37831087845013
0.714646464646465 1.33952679931715
0.636363636363636 1.2987733699263
0.558080808080808 1.25608868215878
0.47979797979798 1.21151679010327
0.401515151515151 1.16510744506897
0.323232323232323 1.11691583059853
0.244949494949495 1.06700229748114
0.166666666666667 1.01543209876543
0.0883838383838382 0.962275124772575
0.0101010101010099 0.907605638109211
-0.0681818181818183 0.85150200868048
-0.146464646464647 0.794046448703017
-0.224747474747475 0.735324747717952
-0.303030303030303 0.67542600760391
-0.381313131313131 0.614442377590011
-0.45959595959596 0.552468789268875
-0.537878787878788 0.489602691609613
-0.616161616161616 0.425943785970835
-0.694444444444445 0.361593761113651
-0.772727272727273 0.296656028214665
-0.851010101010101 0.231235455878979
-0.92929292929293 0.165438105153197
-1.00757575757576 0.0993709645384162
-1.08585858585859 0.033141685003238
-1.16414141414141 -0.0331416850032387
-1.24242424242424 -0.0993709645384178
-1.32070707070707 -0.165438105153197
-1.3989898989899 -0.23123545587898
-1.47727272727273 -0.296656028214666
-1.55555555555556 -0.361593761113652
-1.63383838383838 -0.425943785970836
-1.71212121212121 -0.489602691609613
-1.79040404040404 -0.552468789268874
-1.86868686868687 -0.614442377590012
-1.9469696969697 -0.67542600760391
-2.02525252525253 -0.735324747717952
-2.10353535353535 -0.794046448703018
-2.18181818181818 -0.851502008680481
-2.26010101010101 -0.907605638109211
-2.33838383838384 -0.962275124772574
-2.41666666666667 -1.01543209876543
-2.49494949494949 -1.06700229748113
-2.57323232323232 -1.11691583059853
-2.65151515151515 -1.16510744506897
-2.72979797979798 -1.21151679010327
-2.80808080808081 -1.25608868215878
-2.88636363636364 -1.2987733699263
-2.96464646464647 -1.33952679931715
-3.04292929292929 -1.37831087845013
-3.12121212121212 -1.41509374263854
-3.19949494949495 -1.44985001937716
-3.27777777777778 -1.48256109332927
-3.35606060606061 -1.51321537131363
-3.43434343434343 -1.54180854729151
-3.51262626262626 -1.56834386735365
-3.59090909090909 -1.59283239470727
-3.66919191919192 -1.6152932746631
-3.74747474747475 -1.63575399962237
-3.82575757575758 -1.65425067406375
-3.9040404040404 -1.67082827953045
-3.98232323232323 -1.68554093961715
-4.06060606060606 -1.69845218495698
-4.13888888888889 -1.7096352182086
-4.21717171717172 -1.71917317904316
-4.29545454545455 -1.72715940913127
-4.37373737373737 -1.73369771713003
-4.4520202020202 -1.73890264367003
-4.53030303030303 -1.74289972634234
-4.60858585858586 -1.74582576468553
-4.68686868686869 -1.74782908517264
-4.76515151515152 -1.7490698061982
-4.84343434343434 -1.74972010306519
-4.92171717171717 -1.74996447297215
-5 -1.75
-11 -1.75
};
\addplot [semithick, red, forget plot]
table {%
-10.25 -1
-11 -1.75
-10.25 -2.5
};
\addplot [semithick, blue, dashed, forget plot]
table {%
-7.25 1.75
-9.25 1.75
};
\addplot [semithick, blue, forget plot]
table {%
-8.5 2.5
-9.25 1.75
-8.5 1
};
\end{axis}

\end{tikzpicture}%
		\subcaption{Another overtaking.}
	\end{minipage}
	\caption{Schematic overview of the traffic scenarios. The sedan (red, the middle car in (a)) is the defined as the ego vehicle. First, the station wagon (i.e., the last car in (a)) overtakes both vehicles. Next, the ego vehicle overtakes the pickup truck (first vehicle in (a)).}
	\label{fig:example schematic}
\end{figure}

Several different scenarios can be identified:
\begin{itemize} 
    \item \textbf{Ego gap closing}: Starts at the beginning of the example and ends when the lane change starts.
    \item \textbf{Ego overtaking 1}: Starts at the lane change, ends when ego vehicle completes overtaking of the pickup truck.
    \item \textbf{Ego overtaking 2}: Starts at the lane change of the station wagon, ends when ego vehicle completes overtaking of the pickup truck.
    \item \textbf{Pickup free driving}: Starts at beginning of example, ends when the station wagon makes a cut-in.
    \item \textbf{Pickup cut-in 1}: Starts either when the station wagon starts changing lane, or when the station wagon comes into sensor range of the pickup. Ends when the station wagon in the same lane in front of the pickup.
    \item \textbf{Pickup cut-in 2}: Starts either when the sedan starts changing lane, or when the sedan comes into sensor range of the pickup. Ends when the sedan in the same lane in front of the pickup.
    \item \textbf{Station car following}: Starts at the beginning of the example and ends when the lane change starts.
    \item \textbf{Station overtaking}: Starts at the lane change, ends when the station completes overtaking both cars and is back in the right lane.
\end{itemize}
The scenarios include all 'relevant' events. In case of the 'ego overtaking' scenarios, it might be relevant that another car is already overtaking, because it means that the ego cannot change lane as long as it is occupied. In the end, during test case generation, what is 'relevant' depends on the system under test. In this particular example, suppose that we are testing an automatic lane change function. If the function needs to decide whether it is safe to overtake, the second scenario is the most relevant. If not, the first might be more relevant.

\section{Ego vehicle}

The activities of the ego vehicle are shown in \cref{fig:example ego states}, in which the longitudinal velocity and lateral position are shown with respect to time $t$. Here, the longitudinal and lateral positions are expressed in the Road Coordinate System \cite{zofka2015datadrivetrafficscenarios}. Regarding the longitudinal velocity, four events are identified. At the first event, the ego vehicle's longitudinal acceleration transits from being positive to negative, i.e., the ego vehicle starts braking. As there are four events identified regarding the longitudinal velocity, five different activities are to be modeled, denoted by A, B, C, D, and E in \cref{fig:example ego states lon}. During the time intervals A, B, C, D, and E, the ego vehicle accelerates, brakes, cruises, accelerates, and cruises, respectively. Regarding the lateral position, see \cref{fig:example ego states lat}, also four events are identified, such that at each event time the ego vehicle starts or completes a lane change. During the second and fourth activity (i.e., time intervals G and I in \cref{fig:example ego states lat}) the ego vehicle performs the lane change, while in the remaining time (i.e., time intervals F, H and J in \cref{fig:example ego states lat}) the ego vehicles drives straight ahead. An complete overview of the events is presented in \cref{fig:example events}.

\setlength\figureheight{170pt}
\setlength\figurewidth{0.49\linewidth}
\begin{figure}
	\centering
	\begin{minipage}[b]{0.49\linewidth}
		\centering
		% This file was created by matplotlib2tikz v0.6.14.
\begin{tikzpicture}

\definecolor{color0}{rgb}{0.12156862745098,0.466666666666667,0.705882352941177}
\definecolor{color1}{rgb}{1,0.498039215686275,0.0549019607843137}

\begin{axis}[
xlabel={$t$ [s]},
ylabel={Longitudinal velocity [km/h]},
xmin=0, xmax=40,
ymin=35, ymax=75,
width=\figurewidth,
height=\figureheight,
tick align=outside,
tick pos=left,
xmajorgrids,
x grid style={lightgray!92.026143790849673!black},
ymajorgrids,
y grid style={lightgray!92.026143790849673!black},
clip marker paths
]
\addplot [semithick, color0, mark=*, mark size=1, mark options={solid}, only marks, forget plot]
table {%
0 37.4401428403431
0.025 37.43996431043
0.05 37.5405880740954
0.075 37.6778728877673
0.1 37.7752498731806
0.125 37.947297130267
0.15 38.0982448203184
0.175 38.1850927617531
0.2 38.2780279942145
0.225 38.3680283893234
0.25 38.5280185061743
0.275 38.6715175936735
0.3 38.735646326606
0.325 38.8407305620651
0.35 38.9451844399066
0.375 39.0746025315269
0.4 39.2238734143283
0.425 39.3318738159309
0.45 39.4515929033386
0.475 39.5713397583653
0.5 39.7198979366809
0.525 39.857654656494
0.55 39.9252401503896
0.575 40.045546607177
0.6 40.1770879125208
0.625 40.2492796093048
0.65 40.334136515428
0.675 40.4223605060126
0.7 40.5341846525105
0.725 40.6913299918757
0.75 40.8088933015721
0.775 40.9053299048099
0.8 41.00402973893
0.825 41.0760292998492
0.85 41.1940432205534
0.875 41.3077606771305
0.9 41.3874812448841
0.925 41.5179665642182
0.95 41.6360342159925
0.975 41.734028838642
1 41.8337834953435
1.025 41.9055837346934
1.05 42.0404560371994
1.075 42.1720032536764
1.1 42.2515078594697
1.125 42.3279265718808
1.15 42.3800800635402
1.175 42.5143154611608
1.2 42.6876253021909
1.225 42.8136257771312
1.25 42.9094030407382
1.275 42.9621593843763
1.3 43.0189202582297
1.325 43.1112414661543
1.35 43.1947722340595
1.375 43.3016956459513
1.4 43.4073985153599
1.425 43.4613987161613
1.45 43.521240558735
1.475 43.5922678559857
1.5 43.7031937252783
1.525 43.8445395749978
1.55 43.9390395170369
1.575 43.9964574853556
1.6 44.0933483123004
1.625 44.21955873256
1.65 44.3801994442414
1.675 44.5047833899699
1.7 44.566540786837
1.725 44.6737601059195
1.75 44.770701919897
1.775 44.8446857578458
1.8 44.9188918714645
1.825 44.9728915608496
1.85 45.0667189688517
1.875 45.1857465641861
1.9 45.3347688823756
1.925 45.445780538885
1.95 45.477767418483
1.975 45.5899440289695
2 45.7020716822103
2.025 45.7378536399392
2.05 45.805842881925
2.075 45.8890953661698
2.1 45.983209401329
2.125 46.1398926475155
2.15 46.272538144231
2.175 46.3568218528836
2.2 46.4229133579704
2.225 46.4589134751927
2.25 46.5240883402058
2.275 46.6031512862001
2.3 46.6992407606259
2.325 46.8422144501083
2.35 46.9534191280526
2.375 47.019383559678
2.4 47.0767304568906
2.425 47.1307306950836
2.45 47.2240863239652
2.475 47.3287553076495
2.5 47.4628496312203
2.525 47.5816148892767
2.55 47.6235723903986
2.575 47.7059115184451
2.6 47.7978169178488
2.625 47.8520449210575
2.65 47.9432891415744
2.675 48.0403459980348
2.7 48.1390102552415
2.725 48.217004921237
2.75 48.2429142853792
2.775 48.319633010916
2.8 48.4398553847838
2.825 48.5478547794312
2.85 48.6792105868379
2.875 48.7694426431567
2.9 48.7799534558937
2.925 48.8851228457686
2.95 49.0344217297474
2.975 49.127253118404
3 49.1761875914911
3.025 49.1759531068561
3.05 49.2879029298287
3.075 49.4289577573252
3.1 49.4970086474854
3.125 49.5492388424953
3.15 49.5803611927873
3.175 49.693311827111
3.2 49.8327080832055
3.225 49.9227084091675
3.25 49.968190603683
3.275 49.9983920299798
3.3 50.1063924213416
3.325 50.222701578948
3.35 50.2801120327383
3.375 50.331602189975
3.4 50.4227356701119
3.425 50.5487361903824
3.45 50.7474882860562
3.475 50.7971960584297
3.5 50.7559012183913
3.525 50.8611725748315
3.55 50.9568341250207
3.575 51.035612269255
3.6 51.1095128809523
3.625 51.1861793929761
3.65 51.2779398889559
3.675 51.3325291083389
3.7 51.3589626015638
3.725 51.4641136298429
3.75 51.5923669874695
3.775 51.6439586083864
3.8 51.6950575553225
3.825 51.7670571241807
3.85 51.867743376261
3.875 51.9652841357822
3.9 52.0632433991985
3.925 52.1669860591707
3.95 52.2250627151082
3.975 52.2760272860995
4 52.3444875486562
4.025 52.4702374315002
4.05 52.5599499079724
4.075 52.5599498841554
4.1 52.6135418839279
4.125 52.6861733479854
4.15 52.7223301749962
4.175 52.759398916919
4.2 52.8313985231674
4.225 52.9212591403127
4.25 52.9935115758849
4.275 53.0981185746318
4.3 53.1881189594989
4.325 53.2967376877398
4.35 53.386738008067
4.375 53.3882036589361
4.4 53.3931299272678
4.425 53.3980561899644
4.45 53.5787034911732
4.475 53.7930221112129
4.5 53.877449579409
4.525 53.96412173262
4.55 54.0361219910955
4.575 54.1006440357069
4.6 54.1366441800421
4.625 54.2336541261238
4.65 54.327590349064
4.675 54.3461431601438
4.7 54.4150195186551
4.725 54.5230199688979
4.75 54.6275274151875
4.775 54.6992667384909
4.8 54.7544987786583
4.825 54.8046345517998
4.85 54.8780840126093
4.875 54.951674205129
4.9 54.9874120185631
4.925 55.0989517542321
4.95 55.2327075666356
4.975 55.35770416905
5 55.4490195855873
5.025 55.4667551387826
5.05 55.5630784796288
5.075 55.6890777123185
5.1 55.7624466656812
5.125 55.8014714954698
5.15 55.8207055905182
5.175 55.8730126836845
5.2 55.9630122169316
5.225 56.0281261441424
5.25 56.0463933003361
5.275 56.1342324248767
5.3 56.2296002534621
5.325 56.2733136563573
5.35 56.3246117695189
5.375 56.4148805827708
5.4 56.5190536512721
5.425 56.5910539176536
5.45 56.7070878889653
5.475 56.7920996150776
5.5 56.7852279877083
5.525 56.8183304726681
5.55 56.9083308223607
5.575 57.0374876434223
5.6 57.1454881080465
5.625 57.168218132892
5.65 57.1802680137966
5.675 57.2882684166073
5.7 57.3952114205372
5.725 57.4312115425834
5.75 57.4858461984894
5.775 57.5395719276308
5.8 57.6210190772231
5.825 57.6943916269583
5.85 57.7633181069556
5.875 57.8519822427408
5.9 57.9777058442677
5.925 58.1074768405993
5.95 58.1663019839616
5.975 58.2347701310876
6 58.3007316705235
6.025 58.3544534518366
6.05 58.4315280282154
6.075 58.5035276427362
6.1 58.5501881207208
6.125 58.5696327132425
6.15 58.6271092278341
6.175 58.7044366856197
6.2 58.7944361714351
6.225 58.8690968769898
6.25 58.8873775914665
6.275 58.938695228133
6.3 58.9926954493027
6.325 59.1116650036734
6.35 59.2324200384071
6.375 59.2869807263217
6.4 59.3799537490738
6.425 59.4879541375027
6.45 59.6073981876099
6.475 59.6950828194742
6.5 59.7579122326932
6.525 59.8142050058052
6.55 59.8502051580456
6.575 59.9240459709409
6.6 59.9960462373215
6.625 60.0829877388286
6.65 60.150375457012
6.675 60.1336168635415
6.7 60.1659033073796
6.725 60.25590364126
6.75 60.3382956052396
6.775 60.3920077710487
6.8 60.4600340615284
6.825 60.5279451801729
6.85 60.6018540023132
6.875 60.6909828221255
6.9 60.7626931566213
6.925 60.8687441987167
6.95 60.962399361411
6.975 61.0235976604566
7 61.0736101986657
7.025 61.1093188414708
7.05 61.1661220863327
7.075 61.2201217742212
7.1 61.3050352855201
7.125 61.3895883021609
7.15 61.4642018070261
7.175 61.5472603510805
7.2 61.6012600231587
7.225 61.6575853814513
7.25 61.6938793673994
7.275 61.7366731995669
7.3 61.7781778316458
7.325 61.8174453545316
7.35 61.8831590743799
7.375 61.9734543480539
7.4 62.0622331317266
7.425 62.1162333133687
7.45 62.2158662794525
7.475 62.3314664809303
7.5 62.43183793168
7.525 62.535691027315
7.55 62.5910555794157
7.575 62.6233045487439
7.6 62.6627691559728
7.625 62.7350680914819
7.65 62.8232890129898
7.675 62.8952893095645
7.7 62.9279399838726
7.725 62.9309234812944
7.75 63.0212238893649
7.775 63.1140557513263
7.8 63.1680559408736
7.825 63.2299950222877
7.85 63.2659951443339
7.875 63.3309878540051
7.9 63.4026856522013
7.925 63.4619791945439
7.95 63.5093139994851
7.975 63.6412748885102
8 63.791192107757
8.025 63.8628876923541
8.05 63.9159634757538
8.075 63.9407124258254
8.1 64.0146883710464
8.125 64.0995574127227
8.15 64.1532515168544
8.175 64.1987625360974
8.2 64.2167624390157
8.225 64.2540498544336
8.25 64.3002062039178
8.275 64.3957881249178
8.3 64.4870878974871
8.325 64.5230876889385
8.35 64.5676430838157
8.375 64.6039510111216
8.4 64.6768756980325
8.425 64.748875986701
8.45 64.8436346004948
8.475 64.9497771245957
8.5 65.0143901258389
8.525 65.091593793126
8.55 65.1524020977151
8.575 65.2387455329103
8.6 65.3149934098203
8.625 65.3513049464533
8.65 65.3859606204344
8.675 65.4039606972677
8.7 65.4828294912576
8.725 65.574318451085
8.75 65.6657391091642
8.775 65.731017503914
8.8 65.7130174587021
8.825 65.727564188074
8.85 65.7815643934314
8.875 65.849081060225
8.9 65.9027669287695
8.925 65.8626460339765
8.95 65.787569775521
8.975 65.6968603144749
9 65.6194437557189
9.025 65.6011309354335
9.05 65.4941357753864
9.075 65.3552878098024
9.1 65.2370886784668
9.125 65.1879335943838
9.15 65.24162256412
9.175 65.2264866544567
9.2 65.1724869665696
9.225 65.1064351316943
9.25 65.057034232462
9.275 65.0742264865203
9.3 65.0422511216257
9.325 64.9522516279064
9.35 64.8699075703824
9.375 64.8162168887036
9.4 64.7148903310581
9.425 64.6172661335099
9.45 64.5043151939642
9.475 64.3716315670094
9.5 64.2639382414043
9.525 64.0937045770665
9.55 63.9559696859885
9.575 63.8162052494388
9.6 63.6917696168598
9.625 63.6020730850762
9.65 63.500460693574
9.675 63.4284604350976
9.7 63.3817278563165
9.725 63.360822223028
9.75 63.2493357244163
9.775 63.1391080241197
9.8 63.1211079551924
9.825 63.0716505847499
9.85 62.9996503341829
9.875 62.9094369636823
9.9 62.8371370880973
9.925 62.7491119155905
9.95 62.6812749625699
9.975 62.698102932251
10 62.7122392812123
10.025 62.5836394668384
10.05 62.4735825406772
10.075 62.4717265447862
10.1 62.4017827364773
10.125 62.3341910795011
10.15 62.2978939631617
10.175 62.2522001290863
10.2 62.2162003586809
10.225 62.1545352627713
10.25 62.1099400481191
10.275 62.1450489158731
10.3 62.1477503634083
10.325 62.0937506926162
10.35 62.0639408379295
10.375 62.0642367642481
10.4 62.0639407991323
10.425 62.0534686226306
10.45 61.9994689518365
10.475 61.9347946619586
10.5 61.8810898740016
10.525 61.8464478014171
10.55 61.8941015318815
10.575 61.8775948671253
10.6 61.7940381358435
10.625 61.7583327271913
10.65 61.7654942943462
10.675 61.8014944242996
10.7 61.7779586415034
10.725 61.7419584892652
10.75 61.7402354941969
10.775 61.7345339458243
10.8 61.716533868992
10.825 61.6946985384773
10.85 61.6766984867888
10.875 61.6682352594377
10.9 61.667941189995
10.925 61.6500880616359
10.95 61.6403928046989
10.975 61.676245992656
11 61.6604016472154
11.025 61.5701079510887
11.05 61.5402610076713
11.075 61.5477174607125
11.1 61.5133977104494
11.125 61.5065673092131
11.15 61.5422738950947
11.175 61.4924336661739
11.2 61.425219624241
11.225 61.4210619561799
11.25 61.4165693077322
11.275 61.4159414780552
11.3 61.3974343003866
11.325 61.3614344760022
11.35 61.3692673955308
11.375 61.4055600461513
11.4 61.3986235384701
11.425 61.3772027920661
11.45 61.358520883654
11.475 61.3439414835249
11.5 61.3442340112233
11.525 61.3174026622738
11.55 61.2835729760352
11.575 61.2921119356945
11.6 61.280600528785
11.625 61.2268926480571
11.65 61.1741642713912
11.675 61.1381641479084
11.7 61.2142178646218
11.725 61.3010113825532
11.75 61.2870269891186
11.775 61.2719415744526
11.8 61.2719415744526
11.825 61.2597980610531
11.85 61.24179800325
11.875 61.2199758951098
11.9 61.2016839769579
11.925 61.1999416496397
11.95 61.1929230943468
11.975 61.174923029232
12 61.1642333482673
12.025 61.1639416678347
12.05 61.1958682834782
12.075 61.2264161056628
12.1 61.2124895478268
12.125 61.1869811981589
12.15 61.1664362616025
12.175 61.1256883072383
12.2 61.0919417283986
12.225 61.0919417514551
12.25 61.112225098732
12.275 61.1479335496527
12.3 61.1639416757071
12.325 61.1629271604948
12.35 61.1269270060915
12.375 61.0974640615395
12.4 61.1151726739227
12.425 61.1279417253855
12.45 61.1279416865905
12.475 61.1124135560862
12.5 61.0767049246905
12.525 61.0954832904215
12.55 61.1531762732287
12.575 61.2036342481922
12.6 61.2359416083904
12.625 61.2362335867999
12.65 61.248646593167
12.675 61.2666466740246
12.7 61.3075489527894
12.725 61.352601489209
12.75 61.4072865975735
12.775 61.4288165546996
12.8 61.3748163503592
12.825 61.3439414913931
12.85 61.3439415144496
12.875 61.3113910721387
12.9 61.2860404612302
12.925 61.3548835232369
12.95 61.3946324003974
12.975 61.3406322039295
13 61.3082338798701
13.025 61.3079415719981
13.05 61.2756568394762
13.075 61.2396570297151
13.1 61.2359415769071
13.125 61.2362335710572
13.15 61.2359415847773
13.175 61.2851869569079
13.2 61.3418966457189
13.225 61.3646515750917
13.25 61.3802341577584
13.275 61.37994148838
13.3 61.4040561243881
13.325 61.4400558633102
13.35 61.4366662699946
13.375 61.4169679781143
13.4 61.3959502587339
13.425 61.3861911467457
13.45 61.4041910628751
13.475 61.4374296740249
13.5 61.4737226467941
13.525 61.4578536027885
13.55 61.4193799907219
13.575 61.3954680575994
13.6 61.3987864878963
13.625 61.435079288065
13.65 61.3868295326684
13.675 61.3127894162413
13.7 61.2879018015893
13.725 61.2777986838921
13.75 61.2960909084479
13.775 61.3186451364314
13.8 61.3366452094206
13.825 61.3439414913931
13.85 61.3439414756526
13.875 61.3442339797379
13.9 61.3548282699174
13.925 61.3908284164525
13.95 61.4159414156452
13.975 61.4159414471328
14 61.4449312557874
14.025 61.4812436375244
14.05 61.5065468982412
14.075 61.5239413217044
14.1 61.5239412986543
14.125 61.5508679433721
14.15 61.5785334337761
14.175 61.5339001146483
14.2 61.4879413792551
14.225 61.4879413792573
14.25 61.4882345743094
14.275 61.4879413247175
14.3 61.5223354249973
14.325 61.5551515928396
14.35 61.5387575582376
14.375 61.5164857717254
14.4 61.4981924958116
14.425 61.5126348468798
14.45 61.5486346403416
14.475 61.5599412883275
14.5 61.5557733713392
14.525 61.5194800376221
14.55 61.5021418879076
14.575 61.538141673495
14.6 61.5599412883254
14.625 61.5602348600659
14.65 61.5599413192499
14.675 61.5655606516869
14.7 61.5835605526322
14.725 61.5426706039213
14.75 61.4756792441307
14.775 61.5286561896096
14.8 61.5797851767
14.825 61.5257854749832
14.85 61.487941356203
14.875 61.4882346209767
14.9 61.5164986431935
14.925 61.5528513119446
14.95 61.5782938200032
14.975 61.5907356457848
15 61.5730292462997
15.025 61.5506406905713
15.05 61.5326405939712
15.075 61.5107094786318
15.1 61.4927094365714
15.125 61.5061024558208
15.15 61.5279659489898
15.175 61.5460981516056
15.2 61.5504131875803
15.225 61.5324131612605
15.25 61.5077162095512
15.275 61.4948733745501
15.3 61.5113918701007
15.325 61.533141341067
15.35 61.5511414303578
15.375 61.5602348443253
15.4 61.5599412725826
15.425 61.55994125684
15.45 61.5553256483591
15.475 61.5373255826857
15.5 61.5242346505642
15.525 61.5239413143907
15.55 61.5369112156693
15.575 61.5549111160561
15.6 61.5775661343914
15.625 61.5925813603528
15.65 61.5746629813412
15.675 61.5599413035073
15.7 61.559941280455
15.725 61.5733097346137
15.75 61.5858959874688
15.775 61.5542340454422
15.8 61.5152794477645
15.825 61.4972796097878
15.85 61.475129090807
15.875 61.4574222044409
15.9 61.5034233059405
15.925 61.5507661252816
15.95 61.5172841250504
15.975 61.4879413635167
16 61.4882345506942
16.025 61.5007967247787
16.05 61.5187967904606
16.075 61.5403899119083
16.1 61.5528621813556
16.125 61.5187068936914
16.15 61.4809192133878
16.175 61.4629191634552
16.2 61.4746765957021
16.225 61.5106767034422
16.25 61.5242347287126
16.275 61.5239413531877
16.3 61.5239413222674
16.325 61.5104616165218
16.35 61.4924615508462
16.375 61.4882345664369
16.4 61.4819670530806
16.425 61.463966987966
16.45 61.462319046736
16.475 61.4803190809303
16.5 61.5023145541818
16.525 61.519736682174
16.55 61.5056566477682
16.575 61.4776111268221
16.6 61.4416109960276
16.625 61.4346683269805
16.65 61.4703752554092
16.675 61.4482901530103
16.7 61.3942904586029
16.725 61.3977901104736
16.75 61.4118313578684
16.775 61.3936900276076
16.8 61.3616094333337
16.825 61.3256095932002
16.85 61.3079415405128
16.875 61.2974442884243
16.9 61.2611521010718
16.925 61.2359416083926
16.95 61.2359415926476
16.975 61.2359415926476
17 61.2362335946701
17.025 61.2180434867383
17.05 61.2000434289374
17.075 61.199941681125
17.1 61.2084698017393
17.125 61.2267616653494
17.15 61.247268866242
17.175 61.2652689004406
17.2 61.2009506404698
17.225 61.135875545454
17.25 61.1711584990936
17.275 61.2076765342568
17.3 61.2256765847568
17.325 61.2087871705726
17.35 61.1617849504501
17.375 61.0812309506828
17.4 61.0376215735046
17.425 61.0736216964331
17.45 61.1300727077019
17.475 61.1840729277829
17.5 61.1655152689196
17.525 61.1360667728696
17.55 61.1707849520893
17.575 61.2077633854788
17.6 61.225763466895
17.625 61.321264505426
17.65 61.4469717762991
17.675 61.4550304413163
17.7 61.4176989369699
17.725 61.3966096154723
17.75 61.4081187931825
17.775 61.4798257112885
17.8 61.5466926253169
17.825 61.5826924423832
17.85 61.6448091389039
17.875 61.7049340065587
17.9 61.7277719961021
17.925 61.7805100187814
17.95 61.8445411039443
17.975 62.0199730447762
18 62.1494487779542
18.025 62.0771527154252
18.05 62.0878005323529
18.075 62.1957999127317
18.1 62.3438285405205
18.125 62.470125789022
18.15 62.511444515013
18.175 62.5314829034927
18.2 62.5519785676202
18.225 62.6101803822453
18.25 62.7184791263603
18.275 62.805337385641
18.3 62.8413375243103
18.325 62.9471787671305
18.35 63.070580302448
18.375 63.1956421163545
18.4 63.3244884023529
18.425 63.3964886723734
18.45 63.5136614661615
18.475 63.601625761677
18.5 63.5922065197024
18.525 63.59499996671
18.55 63.6490001395696
18.575 63.7269397313644
18.6 63.7989400637903
18.625 63.8880880645761
18.65 63.9650670670867
18.675 64.0492230624955
18.7 64.1296008961608
18.725 64.1836010538359
18.75 64.2543426897982
18.775 64.3080361445177
18.8 64.3888134643589
18.825 64.466024059443
18.85 64.5531497477809
18.875 64.6597152146328
18.9 64.7854063728359
18.925 64.9124627322581
18.95 64.9692265516305
18.975 65.0007022182763
19 65.0248096410059
19.025 65.0424995039048
19.05 65.1410267514407
19.075 65.2670260636873
19.1 65.3380832159678
19.125 65.3811771141881
19.15 65.4187198047675
19.175 65.4943005424434
19.2 65.6022999228225
19.225 65.6880949432076
19.25 65.7244082121265
19.275 65.7685457440657
19.3 65.8098773881943
19.325 65.8492685167357
19.35 65.9008048312782
19.375 65.955119209354
19.4 66.0446299653376
19.425 66.1346303083387
19.45 66.2497275609189
19.475 66.3509270677661
19.5 66.3954524153032
19.525 66.4743416235428
19.55 66.5502966458634
19.575 66.6038910688594
19.6 66.652503478808
19.625 66.7068214123619
19.65 66.7939593121259
19.675 66.8839596866164
19.7 66.9525452719872
19.725 66.9911882933904
19.75 67.0708987255093
19.775 67.1544312158321
19.8 67.2084314432227
19.825 67.2296031344849
19.85 67.193602995829
19.875 67.2022921503848
19.9 67.2379716092909
19.925 67.3373405179759
19.95 67.4190878915224
19.975 67.3836834505298
20 67.4283437680351
20.025 67.5720217207008
20.05 67.6274774373864
20.075 67.615011917393
20.1 67.6314699489665
20.125 67.6626275248827
20.15 67.6983047630224
20.175 67.7444518298366
20.2 67.780451591811
20.225 67.8238159719519
20.25 67.8601393810745
20.275 67.8599352973706
20.3 67.8854068924524
20.325 67.9394065947299
20.35 67.9806590951164
20.375 67.9989832549067
20.4 68.0382818498646
20.425 68.0810691170918
20.45 68.1187220163269
20.475 68.1715048063586
20.5 68.2258299762534
20.525 68.2559349021225
20.55 68.2617215867276
20.575 68.2977213959279
20.6 68.3279348421181
20.625 68.3282606549364
20.65 68.3405375449751
20.675 68.3585376027804
20.7 68.3998638557736
20.725 68.4421931490727
20.75 68.4785903338611
20.775 68.507934665115
20.8 68.5079346881652
20.825 68.5441128223883
20.85 68.5981130036833
20.875 68.6488430606418
20.9 68.6818470106122
20.925 68.6672657658382
20.95 68.6730916298897
20.975 68.7270918184902
21 68.7602623132743
21.025 68.7599344438522
21.05 68.7928872300922
21.075 68.8057402283493
21.1 68.7367872640937
21.125 68.7280908805091
21.15 68.7805066537134
21.175 68.7766779084581
21.2 68.7768603927834
21.225 68.830860605003
21.25 68.8789589477332
21.275 68.8966304142208
21.3 68.9039342693062
21.325 68.9050003176257
21.35 68.9410004253667
21.375 68.9874715784171
21.4 69.0231424662829
21.425 69.0386802349343
21.45 69.0206803424179
21.475 69.0119341753642
21.5 69.0122632456854
21.525 69.0474838682066
21.55 69.0834840299257
21.575 69.0839341310988
21.6 69.0929451365117
21.625 69.1112746098392
21.65 69.1365706426522
21.675 69.1489714674007
21.7 69.1323348994119
21.725 69.1470217454601
21.75 69.2013515379294
21.775 69.2563555921569
21.8 69.2923557460016
21.825 69.2641452741372
21.85 69.2430364025593
21.875 69.2971555738249
21.9 69.3190946164658
21.925 69.2830944929851
21.95 69.2767120786706
21.975 69.2947121443511
22 69.2832818832407
22.025 69.2713668696464
22.05 69.3063494087591
22.075 69.2977795421661
22.1 69.2077792064795
22.125 69.1968780334272
22.15 69.2535479232287
22.175 69.2849339916433
22.2 69.2999339207153
22.225 69.2999338971093
22.25 69.338976757745
22.275 69.3926460010056
22.3 69.4079338194628
22.325 69.4146680959552
22.35 69.4506682110099
22.375 69.4958731805108
22.4 69.5315417059282
22.425 69.5519336994566
22.45 69.5519336685278
22.475 69.5519336915871
22.5 69.5522653593693
22.525 69.5519336837119
22.55 69.5519336606616
22.575 69.5519336685319
22.6 69.5633457210136
22.625 69.5816774387308
22.65 69.5549070333326
22.675 69.5189068631825
22.7 69.5159336940387
22.725 69.5296405187758
22.75 69.5492011408978
22.775 69.5711626611059
22.8 69.598777842922
22.825 69.6347776369425
22.85 69.6704783882654
22.875 69.6888106237861
22.9 69.6959335406472
22.925 69.6961525833471
22.95 69.7141524528092
22.975 69.7319334679195
23 69.7322660301956
23.025 69.7140278580572
23.05 69.6780276805934
23.075 69.700084293051
23.1 69.7540845361954
23.125 69.7326725255611
23.15 69.6963400124018
23.175 69.6959335249066
23.2 69.7050492831944
23.225 69.7230493331343
23.25 69.7157355340536
23.275 69.6914466389989
23.3 69.6719769576811
23.325 69.6599335824593
23.35 69.6599335976442
23.375 69.6744598960626
23.4 69.6921276108587
23.425 69.6959335485216
23.45 69.7055617092896
23.475 69.7415618322146
23.5 69.7516816124306
23.525 69.715349138627
23.55 69.7212985696896
23.575 69.757298355279
23.6 69.767933480647
23.625 69.7682661536764
23.65 69.7679334654621
23.675 69.7825864686546
23.7 69.8185862621143
23.725 69.8265036704214
23.75 69.8088366072356
23.775 69.8039334079093
23.8 69.803933400039
23.825 69.8039334157797
23.85 69.8039334393955
23.875 69.8042662923371
23.9 69.8039334545746
23.925 69.8004981291328
23.95 69.7824982051351
23.975 69.7755907754966
24 69.7939235215112
24.025 69.7800199674788
24.05 69.7440198130734
24.075 69.7813111647678
24.1 69.8301948561137
24.125 69.7991497636012
24.15 69.7821652262554
24.175 69.8181653418659
24.2 69.8399334054569
24.225 69.8399333897172
24.25 69.8402664152613
24.275 69.8457177118729
24.3 69.8637177544942
24.325 69.8759333327202
24.35 69.8759333636498
24.375 69.8762665932817
24.4 69.8731976105577
24.425 69.8551975291374
24.45 69.8468256084892
24.475 69.8648256662995
24.5 69.876266538742
24.525 69.8759333484649
24.55 69.8450080465404
24.575 69.8144159197049
24.6 69.845341214202
24.625 69.8624914138262
24.65 69.8261584530086
24.675 69.8143594459079
24.7 69.8323593478416
24.725 69.8100419827714
24.75 69.7827325450416
24.775 69.8302908551099
24.8 69.8851093578905
24.825 69.9031092755362
24.85 69.9284585181284
24.875 69.9446794717191
24.9 69.9278209679232
24.925 69.918242065476
24.95 69.9362419438684
24.975 69.969136881762
25 70.005470580309
25.025 69.9904713435298
25.05 69.9544712021965
25.075 69.9479332890654
25.1 69.9423477575195
25.125 69.9246812401479
25.15 69.921728825865
25.175 69.9397289004627
25.2 69.9336017945888
25.225 69.9151094246621
25.25 69.8937743036403
25.275 69.875933351735
25.3 69.8759333438851
25.325 69.8759333360394
25.35 69.8788594980223
25.375 69.8971927609645
25.4 69.9119333204043
25.425 69.9119333125422
25.45 69.924433872079
25.475 69.9424339466817
25.5 69.9308996679087
25.525 69.915747626268
25.55 69.9331147889601
25.575 69.9479332890654
25.6 69.947933304766
25.625 69.9361501275352
25.65 69.9178166840742
25.675 69.944690113908
25.7 69.9780748043117
25.725 69.9633180775362
25.75 69.9482668227387
25.775 69.9479333283147
25.8 69.9591901567645
25.825 69.9771900351447
25.85 69.9506292470504
25.875 69.9209228999365
25.9 69.9358932945051
25.925 69.9343945276595
25.95 69.9176711231375
25.975 69.9312098688364
26 69.9317662213461
26.025 69.8774330054333
26.05 69.8704405390494
26.075 69.9244401977434
26.1 69.9055402932794
26.125 69.8531175147799
26.15 69.8771770121611
26.175 69.9168508977354
26.2 69.9348508232146
26.225 69.9479332890663
26.25 69.948266861992
26.275 69.9479332969112
26.3 69.9479332890614
26.325 69.9479332812156
26.35 69.9522210388521
26.375 69.9705546628137
26.4 69.9554226969239
26.425 69.9014225045507
26.45 69.9092192016437
26.475 69.9394915074898
26.5 69.9245389597603
26.525 69.9211141218351
26.55 69.9391141885748
26.575 69.9479333126068
26.6 69.9479333126118
26.625 69.9304604187013
26.65 69.9075899424237
26.675 69.8893962522269
26.7 69.8673734936664
26.725 69.8493733955446
26.75 69.8660471993066
26.775 69.9017140031596
26.8 69.9297777255269
26.825 69.9513608895644
26.85 69.9695166336177
26.875 69.984266964072
26.9 69.9839332734304
26.925 69.9839332577348
26.95 69.9839332734304
26.975 69.9839332734354
27 69.9753544775935
27.025 69.9570208928803
27.05 69.9730789566624
27.075 70.0090787291141
27.1 69.9859755138785
27.125 69.9503092907907
27.15 69.947933304757
27.175 69.9403791569283
27.2 69.9223792550069
27.225 69.9247873010107
27.25 69.9431207837037
27.275 69.9479333283024
27.3 69.9479333126068
27.325 69.9479333126068
27.35 69.947933304766
27.375 69.9482668384507
27.4 69.9479332812156
27.425 69.9479332969112
27.45 69.9311949415172
27.475 69.9190548507653
27.5 69.9361267869307
27.525 69.9615960131181
27.55 69.9783067270745
27.575 69.9646440265632
27.6 69.9535157127278
27.625 69.9718493288437
27.65 69.9938060053882
27.675 70.0118060878316
27.7 69.9911744499126
27.725 69.9551742928887
27.75 69.9482668306008
27.775 69.9527118430332
27.8 69.9707119019189
27.825 69.9747392477342
27.85 69.9567391888535
27.875 69.9625336086487
27.9 69.9801999768928
27.925 69.9661335168234
27.95 69.9481335834984
27.975 69.9479332812206
28 69.9393074824037
28.025 69.9209740546425
28.05 69.9286162438088
28.075 69.9408761141107
28.1 69.9241932220901
28.125 69.9215532376237
28.15 69.9392197235992
28.175 69.9618107090318
28.2 69.9798106580523
28.225 69.966110882567
28.25 69.9484444986223
28.275 69.9479332969112
28.3 69.9395682395192
28.325 69.9215683062065
28.35 69.9244302697647
28.375 69.9427637132043
28.4 69.9479333126027
28.425 69.9479333440103
28.45 69.9479332969112
28.475 69.9811248583009
28.5 70.0714587570709
28.525 70.0731505038292
28.55 69.9986115022961
28.575 69.963394636918
28.6 69.9295325863617
28.625 69.8938659787131
28.65 69.8887113144746
28.675 69.9067113733644
28.7 69.8957518276547
28.725 69.8846134210011
28.75 69.9191283414988
28.775 69.9479333047611
28.8 69.9479333047611
28.825 69.9359331242325
28.85 69.9179330418064
28.875 69.9285832891409
28.9 69.9435755017543
28.925 69.9272588783396
28.95 69.9185896710304
28.975 69.9365897770243
29 69.9482668227387
29.025 69.947933304757
29.05 69.9645496203242
29.075 69.9825494987043
29.1 69.9839332734354
29.125 69.984266964072
29.15 69.9826016893816
29.175 69.9646016226297
29.2 69.9479333126068
29.225 69.9479333047783
29.25 69.9582065235776
29.275 69.9758728839637
29.3 69.9551679593737
29.325 69.9197681284476
29.35 69.9305333600659
29.375 69.9532079290624
29.4 69.9708742973066
29.425 69.9567111785417
29.45 69.9027114963064
29.475 69.9043783485062
29.5 69.9407116546426
29.525 69.9123808259377
29.55 69.8844038185702
29.575 69.9199562659982
29.6 69.9572879756773
29.625 69.9756215917932
29.65 69.9839332498849
29.675 69.9778845869768
29.7 69.9598846614976
29.725 69.9718882560195
29.75 70.0082219624163
29.775 69.9913227577398
29.8 69.9553226242687
29.825 69.930009773875
29.85 69.9204664395805
29.875 69.9567235473413
29.9 69.9839332577348
29.925 69.9839332655846
29.95 69.9839332498849
29.975 69.9839332420401
30 69.984266964072
30.025 69.9805515614816
30.05 69.9625514476428
30.075 69.9479332812115
30.1 69.9479333126241
30.125 69.9482668070431
30.15 69.9479332969112
30.175 69.9479333204689
30.2 69.9576129234119
30.225 69.9756129901565
30.25 69.968533848013
30.275 69.9483719104118
30.3 69.9281048772918
30.325 69.9180302263138
30.35 69.9360303009001
30.375 69.9589218574874
30.4 69.9765882100563
30.425 69.9390925667311
30.45 69.8879047057443
30.475 69.9327453810614
30.5 69.9741492342045
30.525 69.9378157396995
30.55 69.9259704281048
30.575 69.9619702084227
30.6 69.9617965037654
30.625 69.9261301001679
30.65 69.9454303031443
30.675 69.9756456152035
30.7 69.9601485500806
30.725 69.9617740186062
30.75 69.9801076190305
30.775 69.9839332420433
30.8 69.9839332734345
30.825 69.9839332969758
30.85 69.9740249576383
30.875 69.9563585972391
30.9 69.9192407829577
30.925 69.8862852568135
30.95 69.9329776569929
30.975 69.9789209052718
31 69.9612545762801
31.025 69.9572464585998
31.05 69.9752465253313
31.075 69.9981617757138
31.1 70.0161618189171
31.125 69.984520328181
31.15 69.9564652074677
31.175 69.9922119532966
31.2 70.0104523717891
31.225 69.9924523364456
31.25 70.0177986124071
31.275 70.0353748438419
31.3 69.9838431523166
31.325 69.9600237319763
31.35 69.9780237987128
31.375 69.9842669797635
31.4 69.9811026085953
31.425 69.963102534014
31.45 69.9550575898492
31.475 69.9730576408941
31.5 69.9726219626263
31.525 69.9542883700468
31.55 69.963649217504
31.575 69.9816491272876
31.6 69.9839332577397
31.625 69.9842669326808
31.65 69.9839332498849
31.675 69.9839332655846
31.7 69.9839332577348
31.725 69.9968446088382
31.75 70.0109582209961
31.775 69.9797133087176
31.8 69.9579303603849
31.825 69.9759302230695
31.85 69.9839332577397
31.875 69.9822638705708
31.9 69.9639302701416
31.925 69.9542176340282
31.95 69.9722175281203
31.975 69.9732179721972
32 69.9555516196446
32.025 69.9630010486676
32.05 69.983592864713
32.075 70.0225247873472
32.1 70.0390700294178
32.125 69.985403880872
32.15 69.9479332812065
32.175 69.9479333125978
32.2 69.9621968022524
32.225 69.9801968611422
32.25 69.9842670190259
32.275 69.9839332812712
32.3 69.9839332420137
32.325 69.9699464728086
32.35 69.9553765572449
32.375 69.9696970720546
32.4 69.9754080495526
32.425 69.9574079671215
32.45 69.9732118035012
32.475 70.0092119448221
32.5 69.9861765617455
32.525 69.9566406544395
32.55 69.9907312137997
32.575 70.0044929745585
32.6 69.9684927939719
32.625 69.9601639796878
32.65 69.9778303400803
32.675 69.9866049725446
32.7 69.9896104643634
32.725 69.9869387273333
32.75 69.9842669671629
32.775 69.9839332541773
32.8 69.9951418871873
32.825 70.0131417866836
32.85 70.0199332179085
32.875 70.0014027050705
32.9 69.9470691048122
32.925 69.9400040322057
32.95 69.9760038783937
32.975 69.9839332851938
33 69.9842669828953
33.025 69.9839332463111
33.05 69.9839332851838
33.075 69.9839332851938
33.1 69.9835995328871
33.125 69.9676553913654
33.15 69.9856552907599
33.175 70.0150273115082
33.2 69.997027386043
33.225 69.9839332184905
33.25 69.9842669640769
33.275 69.983933289121
33.3 69.983933289121
33.325 69.9839332577397
33.35 69.9880138302036
33.375 70.0063476111221
33.4 70.0102441427911
33.425 69.9922440839013
33.45 70.0175727406262
33.475 70.0413691264324
33.5 70.0080633420333
33.525 69.9750015038712
33.55 69.9570014606811
33.575 69.9621444781727
33.6 69.9801445292127
33.625 69.9842669719268
33.65 69.9839332969709
33.675 69.9839332734295
33.7 69.9788383498529
33.725 69.9428382320733
33.75 69.9192658737314
33.775 69.9369323753803
33.8 69.9479332890646
33.825 69.9479332969162
33.85 69.9479332733666
33.875 69.9544301481595
33.9 69.9720965242436
33.925 69.9787525687911
33.95 69.9739055783579
33.975 69.9790863000629
34 69.9842669750291
34.025 69.9839332305787
34.05 69.9945870927654
34.075 70.0125869769876
34.1 70.0056271116968
34.125 69.9891557232544
34.15 70.0031279438658
34.175 70.0199332414972
34.2 70.0199332183669
34.225 70.0018023367214
34.25 69.9661361169823
34.275 69.9750310715748
34.3 70.0110311847706
34.325 69.9659417599736
34.35 69.9201963993067
34.375 69.9565216785971
34.4 69.9934798746739
34.425 70.0114799236348
34.45 70.0199331947682
34.475 70.0137075843753
34.5 69.9960414057308
34.525 69.9959397927354
34.55 70.0139398648366
34.575 70.0034830677845
34.6 69.9854830109574
34.625 69.9842669750291
34.65 69.9979971516681
34.675 70.0339973116028
34.7 70.0318220076464
34.725 69.9958218629857
34.75 69.999885738723
34.775 70.0155266431306
34.8 69.9999079183094
34.825 69.990939118716
34.85 70.0089391672185
34.875 70.0095624699883
34.9 69.9912286770354
34.925 69.9996078430981
34.95 70.0206782335657
34.975 70.041003780981
35 70.0465592568484
35.025 70.0282253222871
35.05 70.0199332100323
35.075 70.016444008825
35.1 69.9804439034955
35.125 69.967149751133
35.15 70.0208158939553
35.175 70.0453930032922
35.2 70.0273931038076
35.225 70.0199332105007
35.25 70.0189132109657
35.275 70.000579472618
35.3 69.9895989685733
35.325 70.0075988602115
35.35 70.0298235272943
35.375 70.0481574618374
35.4 70.0271857573275
35.425 69.9911855973928
35.45 69.9839332541873
35.475 69.9839332773275
35.5 69.9842669671629
35.525 69.9839332541773
35.55 69.9873043608863
35.575 70.0053042446501
35.6 70.0199332179085
35.625 70.0202670960911
35.65 70.032748053944
35.675 70.0507481339113
35.7 70.0204545703294
35.725 69.9914176973929
35.75 70.0272303594905
35.775 70.0328129517542
35.8 69.9788127664573
35.825 69.9601020235735
35.85 69.9781021035409
35.875 69.9842669592884
35.9 69.9817955437097
35.925 69.9637955183557
35.95 69.9620576043138
35.975 69.998057733242
36 70.0089635681092
36.025 69.9906297825641
36.05 69.9839332773176
36.075 69.9778695009885
36.1 69.9598694052887
36.125 69.962394190154
36.15 69.9787173147765
36.175 69.9645899471247
36.2 69.9533318533571
36.225 69.9713319411907
36.25 69.9945333708332
36.275 70.0121995494678
36.3 69.9915091859052
36.325 69.9566344521849
36.35 69.9850584999107
36.375 70.0105297906697
36.4 69.9741960977639
36.425 69.9479332899895
36.45 69.9479333057319
36.475 69.9608607811646
36.5 69.9791943862268
36.525 69.966075873222
36.55 69.9522748021986
36.575 69.9701321595652
36.6 69.9738386984321
36.625 69.9561723707715
36.65 69.9988312536302
36.675 70.0403324538906
36.7 70.0074343818882
36.725 69.9839332615951
36.75 69.9842669750291
36.775 69.9839332541773
36.8 69.978299325538
36.825 69.9422995497156
36.85 69.9260922987554
36.875 69.96242573295
36.9 69.9717614126443
36.925 69.9537613715497
36.95 69.9479332742588
36.975 69.9459823864676
37 69.928315846861
37.025 69.9237813826007
37.05 69.9597815504098
37.075 69.9839332773275
37.1 69.983933246321
37.125 69.9397141286648
37.15 69.8916124009897
37.175 69.9181651285057
37.2 69.9479332821332
37.225 69.9479332821332
37.25 69.9313638440211
37.275 69.9141066208944
37.3 70.0381342253126
37.325 70.0581638727703
37.35 69.9690393013356
37.375 69.9842670020165
37.4 69.9839332711621
37.425 69.9690635767333
37.45 69.9523943933858
37.475 69.9672640480088
37.5 69.975909825885
37.525 69.9395763036541
37.55 69.9119333031608
37.575 69.9119333188131
37.6 69.8976938374357
37.625 69.881409691436
37.65 69.8953157952187
37.675 69.9057820024322
37.7 69.8877821191958
37.725 69.8901381697364
37.75 69.9139851651875
37.775 69.9534468526167
37.8 69.9839332789792
37.825 69.9839332711721
37.85 69.9711312302492
37.875 69.9534649442115
37.9 69.9126102225341
37.925 69.8834239501216
37.95 69.9187470666506
37.975 69.9479332758059
38 69.9482668253917
38.025 69.9479332750756
38.05 69.9479332594414
38.075 69.9314824637571
38.1 69.913482391435
38.125 69.9122666807556
38.15 69.9119333507919
38.175 69.911933358609
38.2 69.9119333266302
38.225 69.9119333429766
38.25 69.9122667205614
38.275 69.9056788321399
38.3 69.887678768347
38.325 69.9040717673664
38.35 69.9374904484051
38.375 69.909685249841
38.4 69.8873546835507
38.425 69.9233548040332
38.45 69.9379764114077
38.475 69.9199763234514
38.5 69.8980347578476
38.525 69.8829567226862
38.55 69.9331887667789
38.575 69.9790431556652
38.6 69.9610430911601
38.625 69.9482668331988
38.65 69.9479332594414
38.675 69.9479332516343
38.7 69.9479332992572
38.725 69.9479333070743
38.75 69.9523319339382
38.775 69.9699983074047
38.8 69.9740194724685
38.825 69.9560195643384
38.85 69.9479332829109
38.875 69.9359219592829
38.9 69.8995886254174
38.925 69.8878262276024
38.95 69.9058261591921
38.975 69.8954924755157
39 69.8833254503911
39.025 69.917432839868
39.05 69.940937597827
39.075 69.9229376904171
39.1 69.8878186534888
39.125 69.8521519340424
39.15 69.8555204480176
39.175 69.871523880562
39.2 69.8559368394984
39.225 69.8461390983124
39.25 69.8644721949417
39.275 69.8759333788689
39.3 69.8759333788788
39.325 69.8759333468901
39.35 69.8792222523483
39.375 69.8975555060646
39.4 69.9217458314974
39.425 69.9397458881953
39.45 69.9305863798312
39.475 69.9125863785915
39.5 69.9122667205515
39.525 69.9211577195997
39.55 69.9391577841146
39.575 69.905221007671
39.6 69.851220806347
39.625 69.8580328738812
39.65 69.8756997413397
39.675 69.8759333312559
39.7 69.8592334285262
39.725 69.8232332838901
39.75 69.8299657194561
39.775 69.8656326748628
39.8 69.8759333312559
39.825 69.8826284783967
39.85 69.9186286543012
39.875 69.9405846772356
39.9 69.9222512195197
39.925 69.9396054607363
39.95 69.9692699407873
39.975 69.9415977910388
40 69.9219796417478
40.025 69.9396461794062
40.05 69.947933290728
40.075 69.9479333148914
40.1 69.9479333227085
40.125 69.9545036369768
40.15 69.9721699862716
40.175 69.962907127799
40.2 69.9269073364225
40.225 69.9119333188049
40.25 69.9134393022791
40.275 69.9311058328326
40.3 69.936933598215
40.325 69.9009337826668
40.35 69.885746903737
40.375 69.904080196547
40.4 69.8976814996408
40.425 69.880412164422
40.45 69.8946639992186
40.475 69.8995294469809
40.5 69.8638626642855
40.525 69.8684511415608
40.55 69.9044512783815
40.575 69.9119333351595
40.6 69.920123304339
40.625 69.938456714422
40.65 69.9479332992473
40.675 69.9479332992572
40.7 69.9479332672685
40.725 69.9479332516343
40.75 69.9482668019304
40.775 69.9479333148815
40.8 69.9479333305256
40.825 69.960043146057
40.85 69.9780431871026
40.875 69.9842669622007
40.9 69.9811626487603
40.925 69.9631625608139
40.95 69.9409240275503
40.975 69.9229239786877
41 69.9122666885628
41.025 69.9119333429666
41.05 69.9119333507937
41.075 69.9058307423169
41.1 69.8878306621858
41.125 69.8903008944803
41.15 69.9053051133064
41.175 69.873270698978
41.2 69.851126749017
41.225 69.8871268851238
41.25 69.8916330074075
41.275 69.8552998548007
41.3 69.8399334147892
41.325 69.8435573688761
41.35 69.8975575467406
41.375 69.9387588232025
41.4 69.9024255206136
41.425 69.8759333867041
41.45 69.8759333710617
41.475 69.9044957251257
41.5 69.9408290589813
41.525 69.9121552937899
41.55 69.8803462754598
41.575 69.8981243806041
41.6 69.9218229046541
41.625 69.9401563154574
41.65 69.9133222247887
41.675 69.871144871356
41.7 69.8517562548453
41.725 69.8520396678644
41.75 69.8703727246879
41.775 69.8759333234288
41.8 69.8785832966608
41.825 69.8965832282505
41.85 69.89792919436
41.875 69.8622623953019
41.9 69.8641561860681
41.925 69.9001563548858
41.95 69.9423410940866
41.975 69.9713072196431
42 69.9052327100012
42.025 69.8469519714576
42.05 69.882952091932
42.075 69.9197104674989
42.1 69.937710532004
42.125 69.9623864742748
42.15 69.9772952774605
42.175 69.9631756051261
42.2 69.9283130617035
42.225 69.8923129405188
42.25 69.8919249839568
42.275 69.9074936356241
42.3 69.8918351877868
42.325 69.8762665204853
};
\addplot [line width=2.4000000000000004pt, color1, dashed, forget plot]
table {%
0 37.6080419972848
0.025 37.7152844154791
0.05 37.822339777017
0.075 37.9292080818984
0.1 38.0358893301232
0.125 38.1423835216916
0.15 38.2486906566035
0.175 38.3548107348589
0.2 38.4607437564578
0.225 38.5664897214002
0.25 38.6720486296862
0.275 38.7774204813156
0.3 38.8826052762886
0.325 38.987603014605
0.35 39.092413696265
0.375 39.1970373212685
0.4 39.3014738896155
0.425 39.405723401306
0.45 39.50978585634
0.475 39.6136612547175
0.5 39.7173495964385
0.525 39.8208508815031
0.55 39.9241651099112
0.575 40.0272922816627
0.6 40.1302323967578
0.625 40.2329854551964
0.65 40.3355514569785
0.675 40.4379304021041
0.7 40.5401222905732
0.725 40.6421271223858
0.75 40.743944897542
0.775 40.8455756160416
0.8 40.9470192778848
0.825 41.0482758830714
0.85 41.1493454316016
0.875 41.2502279234753
0.9 41.3509233586925
0.925 41.4514317372532
0.95 41.5517530591574
0.975 41.6518873244051
1 41.7518345329964
1.025 41.8515946849311
1.05 41.9511677802094
1.075 42.0505538188312
1.1 42.1497528007964
1.125 42.2487647261052
1.15 42.3475895947575
1.175 42.4462274067534
1.2 42.5446781620927
1.225 42.6429418607755
1.25 42.7410185028019
1.275 42.8389080881717
1.3 42.9366106168851
1.325 43.0341260889419
1.35 43.1314545043423
1.375 43.2285958630862
1.4 43.3255501651736
1.425 43.4223174106045
1.45 43.518897599379
1.475 43.6152907314969
1.5 43.7114968069583
1.525 43.8075158257633
1.55 43.9033477879118
1.575 43.9989926934037
1.6 44.0944505422392
1.625 44.1897213344182
1.65 44.2848050699407
1.675 44.3797017488067
1.7 44.4744113710163
1.725 44.5689339365693
1.75 44.6632694454659
1.775 44.7574178977059
1.8 44.8513792932895
1.825 44.9451536322166
1.85 45.0387409144872
1.875 45.1321411401013
1.9 45.2253543090589
1.925 45.31838042136
1.95 45.4112194770046
1.975 45.5038714759928
2 45.5963364183244
2.025 45.6886143039996
2.05 45.7807051330182
2.075 45.8726089053804
2.1 45.9643256210861
2.125 46.0558552801353
2.15 46.147197882528
2.175 46.2383534282642
2.2 46.329321917344
2.225 46.4201033497672
2.25 46.510697725534
2.275 46.6011050446442
2.3 46.691325307098
2.325 46.7813585128953
2.35 46.8712046620361
2.375 46.9608637545204
2.4 47.0503357903482
2.425 47.1396207695195
2.45 47.2287186920344
2.475 47.3176295578927
2.5 47.4063533670946
2.525 47.4948901196399
2.55 47.5832398155288
2.575 47.6714024547612
2.6 47.7593780373371
2.625 47.8471665632565
2.65 47.9347680325194
2.675 48.0221824451259
2.7 48.1094098010758
2.725 48.1964501003692
2.75 48.2833033430062
2.775 48.3699695289867
2.8 48.4564486583106
2.825 48.5427407309781
2.85 48.6288457469891
2.875 48.7147637063436
2.9 48.8004946090417
2.925 48.8860384550832
2.95 48.9713952444682
2.975 49.0565649771968
3 49.1415476532689
3.025 49.2263432726844
3.05 49.3109518354435
3.075 49.3953733415461
3.1 49.4796077909922
3.125 49.5636551837818
3.15 49.6475155199149
3.175 49.7311887993916
3.2 49.8146750222117
3.225 49.8979741883754
3.25 49.9810862978825
3.275 50.0640113507332
3.3 50.1467493469274
3.325 50.2293002864651
3.35 50.3116641693463
3.375 50.393840995571
3.4 50.4758307651392
3.425 50.557633478051
3.45 50.6392491343062
3.475 50.720677733905
3.5 50.8019192768472
3.525 50.882973763133
3.55 50.9638411927623
3.575 51.0445215657351
3.6 51.1250148820514
3.625 51.2053211417112
3.65 51.2854403447146
3.675 51.3653724910614
3.7 51.4451175807517
3.725 51.5246756137856
3.75 51.604046590163
3.775 51.6832305098839
3.8 51.7622273729483
3.825 51.8410371793561
3.85 51.9196599291076
3.875 51.9980956222025
3.9 52.0763442586409
3.925 52.1544058384228
3.95 52.2322803615483
3.975 52.3099678280173
4 52.3874682378297
4.025 52.4647815909857
4.05 52.5419078874852
4.075 52.6188471273282
4.1 52.6955993105147
4.125 52.7721644370448
4.15 52.8485425069183
4.175 52.9247335201353
4.2 53.0007374766959
4.225 53.0765543766
4.25 53.1521842198475
4.275 53.2276270064386
4.3 53.3028827363732
4.325 53.3779514096513
4.35 53.4528330262729
4.375 53.5275275862381
4.4 53.6020350895467
4.425 53.6763555361988
4.45 53.7504889261945
4.475 53.8244352595337
4.5 53.8982111571136
4.525 53.9719163443169
4.55 54.0455674420409
4.575 54.1191644502856
4.6 54.1927073690509
4.625 54.2661961983369
4.65 54.3396309381435
4.675 54.4130115884708
4.7 54.4863381493188
4.725 54.5596106206874
4.75 54.6328290025767
4.775 54.7059932949866
4.8 54.7791034979172
4.825 54.8521596113684
4.85 54.9251616353404
4.875 54.9981095698329
4.9 55.0710034148462
4.925 55.1438431703801
4.95 55.2166288364346
4.975 55.2893604130098
5 55.3620379001057
5.025 55.4346612977223
5.05 55.5072306058594
5.075 55.5797458245173
5.1 55.6522069536958
5.125 55.724613993395
5.15 55.7969669436148
5.175 55.8692658043553
5.2 55.9415105756165
5.225 56.0137012573983
5.25 56.0858378497007
5.275 56.1579203525239
5.3 56.2299487658677
5.325 56.3019230897321
5.35 56.3738433241172
5.375 56.445709469023
5.4 56.5175215244494
5.425 56.5892794903965
5.45 56.6609833668643
5.475 56.7326331538527
5.5 56.8042288513618
5.525 56.8757704593915
5.55 56.9472579779419
5.575 57.0186914070129
5.6 57.0900707466046
5.625 57.161395996717
5.65 57.23266715735
5.675 57.3038842285037
5.7 57.3750472101781
5.725 57.4461561023731
5.75 57.5172109050887
5.775 57.5882116183251
5.8 57.6591582420821
5.825 57.7300507763597
5.85 57.800889221158
5.875 57.871673576477
5.9 57.9424038423166
5.925 58.0130800186769
5.95 58.0837021055578
5.975 58.1542701029594
6 58.2247840108817
6.025 58.2952438293246
6.05 58.3656495582882
6.075 58.4360011977725
6.1 58.5062987477774
6.125 58.5765422083029
6.15 58.6467315793492
6.175 58.716866860916
6.2 58.7869480530036
6.225 58.8569751556118
6.25 58.9269481687407
6.275 58.9968670923902
6.3 59.0667319265604
6.325 59.1365426712512
6.35 59.2062993264627
6.375 59.2760018921949
6.4 59.3456503684477
6.425 59.4152447552212
6.45 59.4847850525153
6.475 59.5542712603301
6.5 59.6237033786656
6.525 59.6930814075217
6.55 59.7624053468985
6.575 59.8316751967959
6.6 59.900890957214
6.625 59.9700526281528
6.65 60.0391602096122
6.675 60.1082137015923
6.7 60.177213104093
6.725 60.2461584171144
6.75 60.3150496406565
6.775 60.3838867747192
6.8 60.4526698193026
6.825 60.5213987744066
6.85 60.5900736400313
6.875 60.6586944161767
6.9 60.7272611028427
6.925 60.7957737000294
6.95 60.8642322077367
6.975 60.9326366259647
7 61.0009869547134
7.025 61.0692831939827
7.05 61.1375253437727
7.075 61.2057134040833
7.1 61.2738473749146
7.125 61.3419272562666
7.15 61.4099530481392
7.175 61.4779247505325
7.2 61.5458423634464
7.225 61.613705886881
7.25 61.6815153208363
7.275 61.7492706653122
7.3 61.8169719203087
7.325 61.884619085826
7.35 61.9522121618639
7.375 62.0197511484224
7.4 62.0872360455017
7.425 62.1546668531015
7.45 62.2220435712221
7.475 62.2893661998633
7.5 62.3566347390251
7.525 62.4238491887076
7.55 62.4910095489108
7.575 62.5581158196346
7.6 62.6251680008791
7.625 62.6921660926443
7.65 62.7591100949301
7.675 62.8260000077366
7.7 62.8928358310637
7.725 62.9596175649115
7.75 63.02634520928
7.775 63.0930187641691
7.8 63.1596382295788
7.825 63.2262036055093
7.85 63.2927148919603
7.875 63.3591720889321
7.9 63.4255751964245
7.925 63.4919242144376
7.95 63.5582191429713
7.975 63.6244599820257
8 63.6906467316008
8.025 63.7567793916965
8.05 63.8228579623129
8.075 63.8888824434499
8.1 63.9548528351076
8.125 64.0207691372859
8.15 64.0866313499849
8.175 64.1524394732046
8.2 64.2181935069449
8.225 64.2838934512059
8.25 64.3495393059876
8.275 64.4151310712899
8.3 64.4806687471128
8.325 64.5461523334565
8.35 64.6115818303207
8.375 64.6769572377057
8.4 64.7422785556113
8.425 64.8075457840376
8.45 64.8727589229845
8.475 64.9379179724521
8.5 65.0030229324403
8.525 65.0680738029492
8.55 65.1330705839788
8.575 65.198013275529
8.6 65.2629018775999
8.625 65.3277363901914
8.65 65.3925168133037
8.675 65.4572431469365
8.7 65.52191539109
8.725 65.5865335457642
8.75 65.6510976109591
8.775 65.7156075866746
8.8 65.7800634729108
8.825 65.8444652696676
8.85 65.9088129769451
8.875 65.9731065947432
8.9 66.037346123062
8.925 66.1015315619015
8.95 66.1656629112616
8.975 66.2297401711424
9 65.9163525027803
9.025 65.816737139492
9.05 65.7180682141708
9.075 65.6203457268167
9.1 65.5235696774297
9.125 65.4277400660097
9.15 65.3328568925569
9.175 65.2389201570712
9.2 65.1459298595526
9.225 65.053886000001
9.25 64.9627885784166
9.275 64.8726375947992
9.3 64.783433049149
9.325 64.6951749414658
9.35 64.6078632717497
9.375 64.5214980400008
9.4 64.4360792462189
9.425 64.3516068904041
9.45 64.2680809725564
9.475 64.1855014926758
9.5 64.1038684507623
9.525 64.0231818468159
9.55 63.9434416808365
9.575 63.8646479528243
9.6 63.7868006627792
9.625 63.7098998107011
9.65 63.6339453965902
9.675 63.5589374204463
9.7 63.4848758822695
9.725 63.4117607820599
9.75 63.3395921198173
9.775 63.2683698955418
9.8 63.1980941092334
9.825 63.1287647608922
9.85 63.0603818505179
9.875 62.9929453781108
9.9 62.9264553436708
9.925 62.8609117471979
9.95 62.7963145886921
9.975 62.7326638681533
10 62.6700920870255
10.025 62.6093942539702
10.05 62.5507028704314
10.075 62.4940179364089
10.1 62.4393394519028
10.125 62.3866674169131
10.15 62.3360018314397
10.175 62.2873426954828
10.2 62.2406900090422
10.225 62.196043772118
10.25 62.1534039847103
10.275 62.1127706468189
10.3 62.0741437584438
10.325 62.0375233195852
10.35 62.002909330243
10.375 61.9703017904171
10.4 61.9397007001076
10.425 61.9111060593145
10.45 61.8845178680378
10.475 61.8599361262775
10.5 61.8373608340335
10.525 61.816791991306
10.55 61.7982295980948
10.575 61.7816736544
10.6 61.7671241602216
10.625 61.7545811155596
10.65 61.744044520414
10.675 61.7355143747847
10.7 61.7289906786719
10.725 61.7244734320754
10.75 61.7219626349953
10.775 61.7214582874316
10.8 61.7229603893843
10.825 61.7264689408534
10.85 61.7319839418388
10.875 61.7395053923406
10.9 61.7490332923588
10.925 61.7605676418935
10.95 61.7741084409445
10.975 61.7896556895118
11 61.392702927607
11.025 61.392702927607
11.05 61.392702927607
11.075 61.392702927607
11.1 61.392702927607
11.125 61.392702927607
11.15 61.392702927607
11.175 61.392702927607
11.2 61.392702927607
11.225 61.392702927607
11.25 61.392702927607
11.275 61.392702927607
11.3 61.392702927607
11.325 61.392702927607
11.35 61.392702927607
11.375 61.392702927607
11.4 61.392702927607
11.425 61.392702927607
11.45 61.392702927607
11.475 61.392702927607
11.5 61.392702927607
11.525 61.392702927607
11.55 61.392702927607
11.575 61.392702927607
11.6 61.392702927607
11.625 61.392702927607
11.65 61.392702927607
11.675 61.392702927607
11.7 61.392702927607
11.725 61.392702927607
11.75 61.392702927607
11.775 61.392702927607
11.8 61.392702927607
11.825 61.392702927607
11.85 61.392702927607
11.875 61.392702927607
11.9 61.392702927607
11.925 61.392702927607
11.95 61.392702927607
11.975 61.392702927607
12 61.392702927607
12.025 61.392702927607
12.05 61.392702927607
12.075 61.392702927607
12.1 61.392702927607
12.125 61.392702927607
12.15 61.392702927607
12.175 61.392702927607
12.2 61.392702927607
12.225 61.392702927607
12.25 61.392702927607
12.275 61.392702927607
12.3 61.392702927607
12.325 61.392702927607
12.35 61.392702927607
12.375 61.392702927607
12.4 61.392702927607
12.425 61.392702927607
12.45 61.392702927607
12.475 61.392702927607
12.5 61.392702927607
12.525 61.392702927607
12.55 61.392702927607
12.575 61.392702927607
12.6 61.392702927607
12.625 61.392702927607
12.65 61.392702927607
12.675 61.392702927607
12.7 61.392702927607
12.725 61.392702927607
12.75 61.392702927607
12.775 61.392702927607
12.8 61.392702927607
12.825 61.392702927607
12.85 61.392702927607
12.875 61.392702927607
12.9 61.392702927607
12.925 61.392702927607
12.95 61.392702927607
12.975 61.392702927607
13 61.392702927607
13.025 61.392702927607
13.05 61.392702927607
13.075 61.392702927607
13.1 61.392702927607
13.125 61.392702927607
13.15 61.392702927607
13.175 61.392702927607
13.2 61.392702927607
13.225 61.392702927607
13.25 61.392702927607
13.275 61.392702927607
13.3 61.392702927607
13.325 61.392702927607
13.35 61.392702927607
13.375 61.392702927607
13.4 61.392702927607
13.425 61.392702927607
13.45 61.392702927607
13.475 61.392702927607
13.5 61.392702927607
13.525 61.392702927607
13.55 61.392702927607
13.575 61.392702927607
13.6 61.392702927607
13.625 61.392702927607
13.65 61.392702927607
13.675 61.392702927607
13.7 61.392702927607
13.725 61.392702927607
13.75 61.392702927607
13.775 61.392702927607
13.8 61.392702927607
13.825 61.392702927607
13.85 61.392702927607
13.875 61.392702927607
13.9 61.392702927607
13.925 61.392702927607
13.95 61.392702927607
13.975 61.392702927607
14 61.392702927607
14.025 61.392702927607
14.05 61.392702927607
14.075 61.392702927607
14.1 61.392702927607
14.125 61.392702927607
14.15 61.392702927607
14.175 61.392702927607
14.2 61.392702927607
14.225 61.392702927607
14.25 61.392702927607
14.275 61.392702927607
14.3 61.392702927607
14.325 61.392702927607
14.35 61.392702927607
14.375 61.392702927607
14.4 61.392702927607
14.425 61.392702927607
14.45 61.392702927607
14.475 61.392702927607
14.5 61.392702927607
14.525 61.392702927607
14.55 61.392702927607
14.575 61.392702927607
14.6 61.392702927607
14.625 61.392702927607
14.65 61.392702927607
14.675 61.392702927607
14.7 61.392702927607
14.725 61.392702927607
14.75 61.392702927607
14.775 61.392702927607
14.8 61.392702927607
14.825 61.392702927607
14.85 61.392702927607
14.875 61.392702927607
14.9 61.392702927607
14.925 61.392702927607
14.95 61.392702927607
14.975 61.392702927607
15 61.392702927607
15.025 61.392702927607
15.05 61.392702927607
15.075 61.392702927607
15.1 61.392702927607
15.125 61.392702927607
15.15 61.392702927607
15.175 61.392702927607
15.2 61.392702927607
15.225 61.392702927607
15.25 61.392702927607
15.275 61.392702927607
15.3 61.392702927607
15.325 61.392702927607
15.35 61.392702927607
15.375 61.392702927607
15.4 61.392702927607
15.425 61.392702927607
15.45 61.392702927607
15.475 61.392702927607
15.5 61.392702927607
15.525 61.392702927607
15.55 61.392702927607
15.575 61.392702927607
15.6 61.392702927607
15.625 61.392702927607
15.65 61.392702927607
15.675 61.392702927607
15.7 61.392702927607
15.725 61.392702927607
15.75 61.392702927607
15.775 61.392702927607
15.8 61.392702927607
15.825 61.392702927607
15.85 61.392702927607
15.875 61.392702927607
15.9 61.392702927607
15.925 61.392702927607
15.95 61.392702927607
15.975 61.392702927607
16 61.392702927607
16.025 61.392702927607
16.05 61.392702927607
16.075 61.392702927607
16.1 61.392702927607
16.125 61.392702927607
16.15 61.392702927607
16.175 61.392702927607
16.2 61.392702927607
16.225 61.392702927607
16.25 61.392702927607
16.275 61.392702927607
16.3 61.392702927607
16.325 61.392702927607
16.35 61.392702927607
16.375 61.392702927607
16.4 61.392702927607
16.425 61.392702927607
16.45 61.392702927607
16.475 61.392702927607
16.5 61.392702927607
16.525 61.392702927607
16.55 61.392702927607
16.575 61.392702927607
16.6 61.392702927607
16.625 61.392702927607
16.65 61.392702927607
16.675 61.392702927607
16.7 61.392702927607
16.725 61.392702927607
16.75 61.392702927607
16.775 61.392702927607
16.8 61.392702927607
16.825 61.392702927607
16.85 61.392702927607
16.875 61.392702927607
16.9 61.392702927607
16.925 61.392702927607
16.95 61.392702927607
16.975 61.392702927607
17 61.392702927607
17.025 61.392702927607
17.05 61.392702927607
17.075 61.392702927607
17.1 61.392702927607
17.125 61.392702927607
17.15 61.392702927607
17.175 61.392702927607
17.2 61.392702927607
17.225 61.392702927607
17.25 61.392702927607
17.275 61.392702927607
17.3 61.392702927607
17.325 61.392702927607
17.35 61.392702927607
17.375 61.392702927607
17.4 61.392702927607
17.425 61.392702927607
17.45 61.392702927607
17.475 61.392702927607
17.5 61.392702927607
17.525 61.392702927607
17.55 61.392702927607
17.575 61.392702927607
17.6 61.392702927607
17.625 61.392702927607
17.65 61.392702927607
17.675 61.392702927607
17.7 61.392702927607
17.725 61.392702927607
17.75 61.392702927607
17.775 61.392702927607
17.8 61.392702927607
17.825 61.392702927607
17.85 61.392702927607
17.875 61.392702927607
17.9 61.392702927607
17.925 61.392702927607
17.95 61.392702927607
17.975 61.392702927607
18 61.7326653372238
18.025 61.8311539114564
18.05 61.9289491671942
18.075 62.026051104437
18.1 62.1224597231851
18.125 62.2181750234382
18.15 62.3131970051965
18.175 62.4075256684599
18.2 62.5011610132285
18.225 62.5941030395022
18.25 62.686351747281
18.275 62.7779071365649
18.3 62.868769207354
18.325 62.9589379596482
18.35 63.0484133934476
18.375 63.1371955087521
18.4 63.2252843055617
18.425 63.3126797838764
18.45 63.3993819436963
18.475 63.4853907850213
18.5 63.5707063078515
18.525 63.6553285121868
18.55 63.7392573980272
18.575 63.8224929653727
18.6 63.9050352142234
18.625 63.9868841445793
18.65 64.0680397564402
18.675 64.1485020498063
18.7 64.2282710246775
18.725 64.3073466810539
18.75 64.3857290189354
18.775 64.463418038322
18.8 64.5404137392137
18.825 64.6167161216106
18.85 64.6923251855127
18.875 64.7672409309198
18.9 64.8414633578321
18.925 64.9149924662495
18.95 64.9878282561721
18.975 65.0599707275998
19 65.1314198805326
19.025 65.2021757149706
19.05 65.2722382309137
19.075 65.3416074283619
19.1 65.4102833073152
19.125 65.4782658677737
19.15 65.5455551097374
19.175 65.6121510332061
19.2 65.67805363818
19.225 65.743262924659
19.25 65.8077788926432
19.275 65.8716015421325
19.3 65.9347308731269
19.325 65.9971668856265
19.35 66.0589095796312
19.375 66.119958955141
19.4 66.180315012156
19.425 66.239977750676
19.45 66.2989471707013
19.475 66.3572232722316
19.5 66.4148060552671
19.525 66.4716955198078
19.55 66.5278916658535
19.575 66.5833944934044
19.6 66.6382040024605
19.625 66.6923201930216
19.65 66.7457430650879
19.675 66.7984726186594
19.7 66.8505088537359
19.725 66.9018517703176
19.75 66.9525013684045
19.775 67.0024576479964
19.8 67.0517206090935
19.825 67.1002902516958
19.85 67.1481665758031
19.875 67.1953495814156
19.9 67.2418392685333
19.925 67.287635637156
19.95 67.3327386872839
19.975 67.3771484189169
20 67.4208648320551
20.025 67.4638879266984
20.05 67.5062177028469
20.075 67.5478541605004
20.1 67.5887972996591
20.125 67.629047120323
20.15 67.6686036224919
20.175 67.707466806166
20.2 67.7456366713453
20.225 67.7831132180296
20.25 67.8198964462192
20.275 67.8559863559138
20.3 67.8913829471136
20.325 67.9260862198185
20.35 67.9600961740285
20.375 67.9934128097437
20.4 68.026036126964
20.425 68.0579661256894
20.45 68.08920280592
20.475 68.1197461676557
20.5 68.1496431270512
20.525 68.1791751810337
20.55 68.2083892457577
20.575 68.2372853212234
20.6 68.2658634074307
20.625 68.2941235043795
20.65 68.32206561207
20.675 68.3496897305021
20.7 68.3769958596758
20.725 68.4039839995912
20.75 68.4306541502481
20.775 68.4570063116466
20.8 68.4830404837868
20.825 68.5087566666685
20.85 68.5341548602919
20.875 68.5592350646568
20.9 68.5839972797634
20.925 68.6084415056116
20.95 68.6325677422014
20.975 68.6563759895328
21 68.6798662476058
21.025 68.7030385164204
21.05 68.7258927959766
21.075 68.7484290862745
21.1 68.7706473873139
21.125 68.792547699095
21.15 68.8141300216176
21.175 68.8353943548819
21.2 68.8563406988878
21.225 68.8769690536352
21.25 68.8972794191243
21.275 68.917271795355
21.3 68.9369461823274
21.325 68.9563025800413
21.35 68.9753409884968
21.375 68.9940614076939
21.4 69.0124638376327
21.425 69.030548278313
21.45 69.048314729735
21.475 69.0657631918986
21.5 69.0828936648037
21.525 69.0997061484505
21.55 69.1162006428389
21.575 69.1323771479689
21.6 69.1482356638405
21.625 69.1637761904538
21.65 69.1789987278086
21.675 69.193903275905
21.7 69.2084898347431
21.725 69.2227584043227
21.75 69.236708984644
21.775 69.2503415757069
21.8 69.2636561775114
21.825 69.2766527900574
21.85 69.2893314133451
21.875 69.3016920473744
21.9 69.3137346921453
21.925 69.3254593476579
21.95 69.336866013912
21.975 69.3479546909077
22 69.3587253786451
22.025 69.369178077124
22.05 69.3793127863446
22.075 69.3891295063068
22.1 69.3986282370105
22.125 69.4078089784559
22.15 69.4166717306429
22.175 69.4252164935716
22.2 69.4334432672418
22.225 69.4413520516536
22.25 69.448942846807
22.275 69.4562156527021
22.3 69.4631704693387
22.325 69.469807296717
22.35 69.4761261348368
22.375 69.4821269836983
22.4 69.4878098433014
22.425 69.4931747136461
22.45 69.4982215947324
22.475 69.5029504865603
22.5 69.5073613891298
22.525 69.5114543024409
22.55 69.5152292264937
22.575 69.518686161288
22.6 69.521825106824
22.625 69.5246460631015
22.65 69.5271490301207
22.675 69.5293340078815
22.7 69.5312009963838
22.725 69.5327499956278
22.75 69.5339810056134
22.775 69.5348940263407
22.8 69.5354890578095
22.825 69.5357661000199
22.85 69.5357251529719
22.875 69.5353662166656
22.9 69.5346892911008
22.925 69.5336943762777
22.95 69.5323814721961
22.975 69.5307505788562
23 69.935051266673
23.025 69.935051266673
23.05 69.935051266673
23.075 69.935051266673
23.1 69.935051266673
23.125 69.935051266673
23.15 69.935051266673
23.175 69.935051266673
23.2 69.935051266673
23.225 69.935051266673
23.25 69.935051266673
23.275 69.935051266673
23.3 69.935051266673
23.325 69.935051266673
23.35 69.935051266673
23.375 69.935051266673
23.4 69.935051266673
23.425 69.935051266673
23.45 69.935051266673
23.475 69.935051266673
23.5 69.935051266673
23.525 69.935051266673
23.55 69.935051266673
23.575 69.935051266673
23.6 69.935051266673
23.625 69.935051266673
23.65 69.935051266673
23.675 69.935051266673
23.7 69.935051266673
23.725 69.935051266673
23.75 69.935051266673
23.775 69.935051266673
23.8 69.935051266673
23.825 69.935051266673
23.85 69.935051266673
23.875 69.935051266673
23.9 69.935051266673
23.925 69.935051266673
23.95 69.935051266673
23.975 69.935051266673
24 69.935051266673
24.025 69.935051266673
24.05 69.935051266673
24.075 69.935051266673
24.1 69.935051266673
24.125 69.935051266673
24.15 69.935051266673
24.175 69.935051266673
24.2 69.935051266673
24.225 69.935051266673
24.25 69.935051266673
24.275 69.935051266673
24.3 69.935051266673
24.325 69.935051266673
24.35 69.935051266673
24.375 69.935051266673
24.4 69.935051266673
24.425 69.935051266673
24.45 69.935051266673
24.475 69.935051266673
24.5 69.935051266673
24.525 69.935051266673
24.55 69.935051266673
24.575 69.935051266673
24.6 69.935051266673
24.625 69.935051266673
24.65 69.935051266673
24.675 69.935051266673
24.7 69.935051266673
24.725 69.935051266673
24.75 69.935051266673
24.775 69.935051266673
24.8 69.935051266673
24.825 69.935051266673
24.85 69.935051266673
24.875 69.935051266673
24.9 69.935051266673
24.925 69.935051266673
24.95 69.935051266673
24.975 69.935051266673
25 69.935051266673
25.025 69.935051266673
25.05 69.935051266673
25.075 69.935051266673
25.1 69.935051266673
25.125 69.935051266673
25.15 69.935051266673
25.175 69.935051266673
25.2 69.935051266673
25.225 69.935051266673
25.25 69.935051266673
25.275 69.935051266673
25.3 69.935051266673
25.325 69.935051266673
25.35 69.935051266673
25.375 69.935051266673
25.4 69.935051266673
25.425 69.935051266673
25.45 69.935051266673
25.475 69.935051266673
25.5 69.935051266673
25.525 69.935051266673
25.55 69.935051266673
25.575 69.935051266673
25.6 69.935051266673
25.625 69.935051266673
25.65 69.935051266673
25.675 69.935051266673
25.7 69.935051266673
25.725 69.935051266673
25.75 69.935051266673
25.775 69.935051266673
25.8 69.935051266673
25.825 69.935051266673
25.85 69.935051266673
25.875 69.935051266673
25.9 69.935051266673
25.925 69.935051266673
25.95 69.935051266673
25.975 69.935051266673
26 69.935051266673
26.025 69.935051266673
26.05 69.935051266673
26.075 69.935051266673
26.1 69.935051266673
26.125 69.935051266673
26.15 69.935051266673
26.175 69.935051266673
26.2 69.935051266673
26.225 69.935051266673
26.25 69.935051266673
26.275 69.935051266673
26.3 69.935051266673
26.325 69.935051266673
26.35 69.935051266673
26.375 69.935051266673
26.4 69.935051266673
26.425 69.935051266673
26.45 69.935051266673
26.475 69.935051266673
26.5 69.935051266673
26.525 69.935051266673
26.55 69.935051266673
26.575 69.935051266673
26.6 69.935051266673
26.625 69.935051266673
26.65 69.935051266673
26.675 69.935051266673
26.7 69.935051266673
26.725 69.935051266673
26.75 69.935051266673
26.775 69.935051266673
26.8 69.935051266673
26.825 69.935051266673
26.85 69.935051266673
26.875 69.935051266673
26.9 69.935051266673
26.925 69.935051266673
26.95 69.935051266673
26.975 69.935051266673
27 69.935051266673
27.025 69.935051266673
27.05 69.935051266673
27.075 69.935051266673
27.1 69.935051266673
27.125 69.935051266673
27.15 69.935051266673
27.175 69.935051266673
27.2 69.935051266673
27.225 69.935051266673
27.25 69.935051266673
27.275 69.935051266673
27.3 69.935051266673
27.325 69.935051266673
27.35 69.935051266673
27.375 69.935051266673
27.4 69.935051266673
27.425 69.935051266673
27.45 69.935051266673
27.475 69.935051266673
27.5 69.935051266673
27.525 69.935051266673
27.55 69.935051266673
27.575 69.935051266673
27.6 69.935051266673
27.625 69.935051266673
27.65 69.935051266673
27.675 69.935051266673
27.7 69.935051266673
27.725 69.935051266673
27.75 69.935051266673
27.775 69.935051266673
27.8 69.935051266673
27.825 69.935051266673
27.85 69.935051266673
27.875 69.935051266673
27.9 69.935051266673
27.925 69.935051266673
27.95 69.935051266673
27.975 69.935051266673
28 69.935051266673
28.025 69.935051266673
28.05 69.935051266673
28.075 69.935051266673
28.1 69.935051266673
28.125 69.935051266673
28.15 69.935051266673
28.175 69.935051266673
28.2 69.935051266673
28.225 69.935051266673
28.25 69.935051266673
28.275 69.935051266673
28.3 69.935051266673
28.325 69.935051266673
28.35 69.935051266673
28.375 69.935051266673
28.4 69.935051266673
28.425 69.935051266673
28.45 69.935051266673
28.475 69.935051266673
28.5 69.935051266673
28.525 69.935051266673
28.55 69.935051266673
28.575 69.935051266673
28.6 69.935051266673
28.625 69.935051266673
28.65 69.935051266673
28.675 69.935051266673
28.7 69.935051266673
28.725 69.935051266673
28.75 69.935051266673
28.775 69.935051266673
28.8 69.935051266673
28.825 69.935051266673
28.85 69.935051266673
28.875 69.935051266673
28.9 69.935051266673
28.925 69.935051266673
28.95 69.935051266673
28.975 69.935051266673
29 69.935051266673
29.025 69.935051266673
29.05 69.935051266673
29.075 69.935051266673
29.1 69.935051266673
29.125 69.935051266673
29.15 69.935051266673
29.175 69.935051266673
29.2 69.935051266673
29.225 69.935051266673
29.25 69.935051266673
29.275 69.935051266673
29.3 69.935051266673
29.325 69.935051266673
29.35 69.935051266673
29.375 69.935051266673
29.4 69.935051266673
29.425 69.935051266673
29.45 69.935051266673
29.475 69.935051266673
29.5 69.935051266673
29.525 69.935051266673
29.55 69.935051266673
29.575 69.935051266673
29.6 69.935051266673
29.625 69.935051266673
29.65 69.935051266673
29.675 69.935051266673
29.7 69.935051266673
29.725 69.935051266673
29.75 69.935051266673
29.775 69.935051266673
29.8 69.935051266673
29.825 69.935051266673
29.85 69.935051266673
29.875 69.935051266673
29.9 69.935051266673
29.925 69.935051266673
29.95 69.935051266673
29.975 69.935051266673
30 69.935051266673
30.025 69.935051266673
30.05 69.935051266673
30.075 69.935051266673
30.1 69.935051266673
30.125 69.935051266673
30.15 69.935051266673
30.175 69.935051266673
30.2 69.935051266673
30.225 69.935051266673
30.25 69.935051266673
30.275 69.935051266673
30.3 69.935051266673
30.325 69.935051266673
30.35 69.935051266673
30.375 69.935051266673
30.4 69.935051266673
30.425 69.935051266673
30.45 69.935051266673
30.475 69.935051266673
30.5 69.935051266673
30.525 69.935051266673
30.55 69.935051266673
30.575 69.935051266673
30.6 69.935051266673
30.625 69.935051266673
30.65 69.935051266673
30.675 69.935051266673
30.7 69.935051266673
30.725 69.935051266673
30.75 69.935051266673
30.775 69.935051266673
30.8 69.935051266673
30.825 69.935051266673
30.85 69.935051266673
30.875 69.935051266673
30.9 69.935051266673
30.925 69.935051266673
30.95 69.935051266673
30.975 69.935051266673
31 69.935051266673
31.025 69.935051266673
31.05 69.935051266673
31.075 69.935051266673
31.1 69.935051266673
31.125 69.935051266673
31.15 69.935051266673
31.175 69.935051266673
31.2 69.935051266673
31.225 69.935051266673
31.25 69.935051266673
31.275 69.935051266673
31.3 69.935051266673
31.325 69.935051266673
31.35 69.935051266673
31.375 69.935051266673
31.4 69.935051266673
31.425 69.935051266673
31.45 69.935051266673
31.475 69.935051266673
31.5 69.935051266673
31.525 69.935051266673
31.55 69.935051266673
31.575 69.935051266673
31.6 69.935051266673
31.625 69.935051266673
31.65 69.935051266673
31.675 69.935051266673
31.7 69.935051266673
31.725 69.935051266673
31.75 69.935051266673
31.775 69.935051266673
31.8 69.935051266673
31.825 69.935051266673
31.85 69.935051266673
31.875 69.935051266673
31.9 69.935051266673
31.925 69.935051266673
31.95 69.935051266673
31.975 69.935051266673
32 69.935051266673
32.025 69.935051266673
32.05 69.935051266673
32.075 69.935051266673
32.1 69.935051266673
32.125 69.935051266673
32.15 69.935051266673
32.175 69.935051266673
32.2 69.935051266673
32.225 69.935051266673
32.25 69.935051266673
32.275 69.935051266673
32.3 69.935051266673
32.325 69.935051266673
32.35 69.935051266673
32.375 69.935051266673
32.4 69.935051266673
32.425 69.935051266673
32.45 69.935051266673
32.475 69.935051266673
32.5 69.935051266673
32.525 69.935051266673
32.55 69.935051266673
32.575 69.935051266673
32.6 69.935051266673
32.625 69.935051266673
32.65 69.935051266673
32.675 69.935051266673
32.7 69.935051266673
32.725 69.935051266673
32.75 69.935051266673
32.775 69.935051266673
32.8 69.935051266673
32.825 69.935051266673
32.85 69.935051266673
32.875 69.935051266673
32.9 69.935051266673
32.925 69.935051266673
32.95 69.935051266673
32.975 69.935051266673
33 69.935051266673
33.025 69.935051266673
33.05 69.935051266673
33.075 69.935051266673
33.1 69.935051266673
33.125 69.935051266673
33.15 69.935051266673
33.175 69.935051266673
33.2 69.935051266673
33.225 69.935051266673
33.25 69.935051266673
33.275 69.935051266673
33.3 69.935051266673
33.325 69.935051266673
33.35 69.935051266673
33.375 69.935051266673
33.4 69.935051266673
33.425 69.935051266673
33.45 69.935051266673
33.475 69.935051266673
33.5 69.935051266673
33.525 69.935051266673
33.55 69.935051266673
33.575 69.935051266673
33.6 69.935051266673
33.625 69.935051266673
33.65 69.935051266673
33.675 69.935051266673
33.7 69.935051266673
33.725 69.935051266673
33.75 69.935051266673
33.775 69.935051266673
33.8 69.935051266673
33.825 69.935051266673
33.85 69.935051266673
33.875 69.935051266673
33.9 69.935051266673
33.925 69.935051266673
33.95 69.935051266673
33.975 69.935051266673
34 69.935051266673
34.025 69.935051266673
34.05 69.935051266673
34.075 69.935051266673
34.1 69.935051266673
34.125 69.935051266673
34.15 69.935051266673
34.175 69.935051266673
34.2 69.935051266673
34.225 69.935051266673
34.25 69.935051266673
34.275 69.935051266673
34.3 69.935051266673
34.325 69.935051266673
34.35 69.935051266673
34.375 69.935051266673
34.4 69.935051266673
34.425 69.935051266673
34.45 69.935051266673
34.475 69.935051266673
34.5 69.935051266673
34.525 69.935051266673
34.55 69.935051266673
34.575 69.935051266673
34.6 69.935051266673
34.625 69.935051266673
34.65 69.935051266673
34.675 69.935051266673
34.7 69.935051266673
34.725 69.935051266673
34.75 69.935051266673
34.775 69.935051266673
34.8 69.935051266673
34.825 69.935051266673
34.85 69.935051266673
34.875 69.935051266673
34.9 69.935051266673
34.925 69.935051266673
34.95 69.935051266673
34.975 69.935051266673
35 69.935051266673
35.025 69.935051266673
35.05 69.935051266673
35.075 69.935051266673
35.1 69.935051266673
35.125 69.935051266673
35.15 69.935051266673
35.175 69.935051266673
35.2 69.935051266673
35.225 69.935051266673
35.25 69.935051266673
35.275 69.935051266673
35.3 69.935051266673
35.325 69.935051266673
35.35 69.935051266673
35.375 69.935051266673
35.4 69.935051266673
35.425 69.935051266673
35.45 69.935051266673
35.475 69.935051266673
35.5 69.935051266673
35.525 69.935051266673
35.55 69.935051266673
35.575 69.935051266673
35.6 69.935051266673
35.625 69.935051266673
35.65 69.935051266673
35.675 69.935051266673
35.7 69.935051266673
35.725 69.935051266673
35.75 69.935051266673
35.775 69.935051266673
35.8 69.935051266673
35.825 69.935051266673
35.85 69.935051266673
35.875 69.935051266673
35.9 69.935051266673
35.925 69.935051266673
35.95 69.935051266673
35.975 69.935051266673
36 69.935051266673
36.025 69.935051266673
36.05 69.935051266673
36.075 69.935051266673
36.1 69.935051266673
36.125 69.935051266673
36.15 69.935051266673
36.175 69.935051266673
36.2 69.935051266673
36.225 69.935051266673
36.25 69.935051266673
36.275 69.935051266673
36.3 69.935051266673
36.325 69.935051266673
36.35 69.935051266673
36.375 69.935051266673
36.4 69.935051266673
36.425 69.935051266673
36.45 69.935051266673
36.475 69.935051266673
36.5 69.935051266673
36.525 69.935051266673
36.55 69.935051266673
36.575 69.935051266673
36.6 69.935051266673
36.625 69.935051266673
36.65 69.935051266673
36.675 69.935051266673
36.7 69.935051266673
36.725 69.935051266673
36.75 69.935051266673
36.775 69.935051266673
36.8 69.935051266673
36.825 69.935051266673
36.85 69.935051266673
36.875 69.935051266673
36.9 69.935051266673
36.925 69.935051266673
36.95 69.935051266673
36.975 69.935051266673
37 69.935051266673
37.025 69.935051266673
37.05 69.935051266673
37.075 69.935051266673
37.1 69.935051266673
37.125 69.935051266673
37.15 69.935051266673
37.175 69.935051266673
37.2 69.935051266673
37.225 69.935051266673
37.25 69.935051266673
37.275 69.935051266673
37.3 69.935051266673
37.325 69.935051266673
37.35 69.935051266673
37.375 69.935051266673
37.4 69.935051266673
37.425 69.935051266673
37.45 69.935051266673
37.475 69.935051266673
37.5 69.935051266673
37.525 69.935051266673
37.55 69.935051266673
37.575 69.935051266673
37.6 69.935051266673
37.625 69.935051266673
37.65 69.935051266673
37.675 69.935051266673
37.7 69.935051266673
37.725 69.935051266673
37.75 69.935051266673
37.775 69.935051266673
37.8 69.935051266673
37.825 69.935051266673
37.85 69.935051266673
37.875 69.935051266673
37.9 69.935051266673
37.925 69.935051266673
37.95 69.935051266673
37.975 69.935051266673
38 69.935051266673
38.025 69.935051266673
38.05 69.935051266673
38.075 69.935051266673
38.1 69.935051266673
38.125 69.935051266673
38.15 69.935051266673
38.175 69.935051266673
38.2 69.935051266673
38.225 69.935051266673
38.25 69.935051266673
38.275 69.935051266673
38.3 69.935051266673
38.325 69.935051266673
38.35 69.935051266673
38.375 69.935051266673
38.4 69.935051266673
38.425 69.935051266673
38.45 69.935051266673
38.475 69.935051266673
38.5 69.935051266673
38.525 69.935051266673
38.55 69.935051266673
38.575 69.935051266673
38.6 69.935051266673
38.625 69.935051266673
38.65 69.935051266673
38.675 69.935051266673
38.7 69.935051266673
38.725 69.935051266673
38.75 69.935051266673
38.775 69.935051266673
38.8 69.935051266673
38.825 69.935051266673
38.85 69.935051266673
38.875 69.935051266673
38.9 69.935051266673
38.925 69.935051266673
38.95 69.935051266673
38.975 69.935051266673
39 69.935051266673
39.025 69.935051266673
39.05 69.935051266673
39.075 69.935051266673
39.1 69.935051266673
39.125 69.935051266673
39.15 69.935051266673
39.175 69.935051266673
39.2 69.935051266673
39.225 69.935051266673
39.25 69.935051266673
39.275 69.935051266673
39.3 69.935051266673
39.325 69.935051266673
39.35 69.935051266673
39.375 69.935051266673
39.4 69.935051266673
39.425 69.935051266673
39.45 69.935051266673
39.475 69.935051266673
39.5 69.935051266673
39.525 69.935051266673
39.55 69.935051266673
39.575 69.935051266673
39.6 69.935051266673
39.625 69.935051266673
39.65 69.935051266673
39.675 69.935051266673
39.7 69.935051266673
39.725 69.935051266673
39.75 69.935051266673
39.775 69.935051266673
39.8 69.935051266673
39.825 69.935051266673
39.85 69.935051266673
39.875 69.935051266673
39.9 69.935051266673
39.925 69.935051266673
39.95 69.935051266673
39.975 69.935051266673
40 69.935051266673
40.025 69.935051266673
40.05 69.935051266673
40.075 69.935051266673
40.1 69.935051266673
40.125 69.935051266673
40.15 69.935051266673
40.175 69.935051266673
40.2 69.935051266673
40.225 69.935051266673
40.25 69.935051266673
40.275 69.935051266673
40.3 69.935051266673
40.325 69.935051266673
40.35 69.935051266673
40.375 69.935051266673
40.4 69.935051266673
40.425 69.935051266673
40.45 69.935051266673
40.475 69.935051266673
40.5 69.935051266673
40.525 69.935051266673
40.55 69.935051266673
40.575 69.935051266673
40.6 69.935051266673
40.625 69.935051266673
40.65 69.935051266673
40.675 69.935051266673
40.7 69.935051266673
40.725 69.935051266673
40.75 69.935051266673
40.775 69.935051266673
40.8 69.935051266673
40.825 69.935051266673
40.85 69.935051266673
40.875 69.935051266673
40.9 69.935051266673
40.925 69.935051266673
40.95 69.935051266673
40.975 69.935051266673
41 69.935051266673
41.025 69.935051266673
41.05 69.935051266673
41.075 69.935051266673
41.1 69.935051266673
41.125 69.935051266673
41.15 69.935051266673
41.175 69.935051266673
41.2 69.935051266673
41.225 69.935051266673
41.25 69.935051266673
41.275 69.935051266673
41.3 69.935051266673
41.325 69.935051266673
41.35 69.935051266673
41.375 69.935051266673
41.4 69.935051266673
41.425 69.935051266673
41.45 69.935051266673
41.475 69.935051266673
41.5 69.935051266673
41.525 69.935051266673
41.55 69.935051266673
41.575 69.935051266673
41.6 69.935051266673
41.625 69.935051266673
41.65 69.935051266673
41.675 69.935051266673
41.7 69.935051266673
41.725 69.935051266673
41.75 69.935051266673
41.775 69.935051266673
41.8 69.935051266673
41.825 69.935051266673
41.85 69.935051266673
41.875 69.935051266673
41.9 69.935051266673
41.925 69.935051266673
41.95 69.935051266673
41.975 69.935051266673
42 69.935051266673
42.025 69.935051266673
42.05 69.935051266673
42.075 69.935051266673
42.1 69.935051266673
42.125 69.935051266673
42.15 69.935051266673
42.175 69.935051266673
42.2 69.935051266673
42.225 69.935051266673
42.25 69.935051266673
42.275 69.935051266673
42.3 69.935051266673
42.325 69.935051266673
};
\addplot [ultra thick, black, forget plot]
table {%
9 35
9 75
};
\addplot [ultra thick, black, forget plot]
table {%
11 35
11 75
};
\addplot [ultra thick, black, forget plot]
table {%
18 35
18 75
};
\addplot [ultra thick, black, forget plot]
table {%
23 35
23 75
};
\node at (axis cs:4.5,40)[
  anchor=base,
  text=black,
  rotate=0.0
]{ A};
\node at (axis cs:10,40)[
  anchor=base,
  text=black,
  rotate=0.0
]{ B};
\node at (axis cs:14.5,40)[
  anchor=base,
  text=black,
  rotate=0.0
]{ C};
\node at (axis cs:20.5,40)[
  anchor=base,
  text=black,
  rotate=0.0
]{ D};
\node at (axis cs:31.5,40)[
  anchor=base,
  text=black,
  rotate=0.0
]{ E};
\end{axis}

\end{tikzpicture}%
		\subcaption{Longitudinal velocity ego vehicle. During the time intervals A and D, the ego vehicle accelerates, during the time interval B, the ego vehicle brakes and during the time intervals C and E, the ego vehicles cruises.}
		\label{fig:example ego states lon}
	\end{minipage}%
	\hfill
	\begin{minipage}[b]{0.49\linewidth}
		\centering
		% This file was created by matplotlib2tikz v0.6.14.
\begin{tikzpicture}

\definecolor{color0}{rgb}{0.12156862745098,0.466666666666667,0.705882352941177}
\definecolor{color1}{rgb}{1,0.498039215686275,0.0549019607843137}

\begin{axis}[
xlabel={Time [s]},
ylabel={Lateral distance [m]},
xmin=0, xmax=40,
ymin=-0.2, ymax=3.8,
width=\figurewidth,
height=\figureheight,
tick align=outside,
tick pos=left,
xmajorgrids,
x grid style={white!69.019607843137251!black},
ymajorgrids,
y grid style={white!69.019607843137251!black},
clip marker paths
]
\addplot [semithick, color0, mark=*, mark size=1, mark options={solid}, only marks, forget plot]
table {%
0 0.0937500001330113
0.025 0.0828903467114039
0.05 0.0720307975106381
0.075 0.0594505826611748
0.1 0.0470313523314523
0.125 0.0352070005756295
0.15 0.0244521743136864
0.175 0.0141362518893688
0.2 0.00462722711661243
0.225 0.00453474629515783
0.25 0.0133658820522649
0.275 0.0213851986889686
0.3 0.0291163107499829
0.325 0.0364014834129398
0.35 0.0429244132663706
0.375 0.0491345539530386
0.4 0.0548429105440976
0.425 0.0601911966537087
0.45 0.0652041258725517
0.475 0.0696687390350636
0.5 0.0739920138997101
0.525 0.0781016767740133
0.55 0.0816518148104859
0.575 0.0850028583115035
0.6 0.088117750405292
0.625 0.0910301416922158
0.65 0.0937522882449332
0.675 0.0961850404614553
0.7 0.0985493281102856
0.725 0.100813504571676
0.75 0.102534635330981
0.775 0.10407642126486
0.8 0.105491911585923
0.825 0.106791091813252
0.85 0.107974379674953
0.875 0.109051950579307
0.9 0.110088055486755
0.925 0.111059705738672
0.95 0.111661493011074
0.975 0.112139340593796
1 0.112576020228859
1.025 0.113057246471211
1.05 0.113579168666344
1.075 0.11404941537109
1.1 0.114504751148708
1.125 0.114937356924901
1.15 0.11532093219849
1.175 0.115685619215265
1.2 0.116021209177081
1.225 0.116281559693602
1.25 0.116468251785197
1.275 0.116635078240359
1.3 0.116717619744347
1.325 0.116619707210667
1.35 0.116533168531571
1.375 0.11645101545443
1.4 0.116375313806244
1.425 0.116114653609613
1.45 0.115674965077235
1.475 0.115279710096107
1.5 0.114695281157758
1.525 0.113709191016803
1.55 0.112829119821645
1.575 0.111990893108854
1.6 0.111213904632034
1.625 0.109772873949879
1.65 0.107716350809616
1.675 0.105882793814197
1.7 0.103975405517058
1.725 0.101862565534966
1.75 0.100020295504264
1.775 0.0982752398255858
1.8 0.0966631034448312
1.825 0.0947964198609719
1.85 0.0926739868378996
1.875 0.0907821927001157
1.9 0.0889127628526968
1.925 0.0870522073238416
1.95 0.0853984032767185
1.975 0.0838269942876814
2 0.0823894248763972
2.025 0.0808140384456236
2.05 0.079099268625621
2.075 0.0775656048233277
2.1 0.076086772923478
2.125 0.0746907480690876
2.15 0.0734400192582771
2.175 0.0722478364252017
2.2 0.0711418981985308
2.225 0.0698556671427727
2.25 0.0684073354334131
2.275 0.0670636932730139
2.3 0.0657246863498245
2.325 0.0643634418861496
2.35 0.0631438887690473
2.375 0.0619455750479519
2.4 0.0604136997882718
2.425 0.0589270664471708
2.45 0.0574937685129643
2.475 0.0562327671067012
2.5 0.0547799308865657
2.525 0.0527186809506036
2.55 0.0508721580700069
2.575 0.0491245498441236
2.6 0.0475146199581872
2.625 0.045745858812044
2.65 0.0438043114353348
2.675 0.0420354549408974
2.7 0.0401852899595268
2.725 0.038063209910462
2.75 0.0361506247736728
2.775 0.0343337928341486
2.8 0.0326535212618332
2.825 0.0307125776379243
2.85 0.0285287293014327
2.875 0.026491970559177
2.9 0.0244117545764509
2.925 0.0221727595746638
2.95 0.0201668845164022
2.975 0.0182391643286841
3 0.0162383935200944
3.025 0.0143364172765091
3.05 0.0125359013132222
3.075 0.0108996168460826
3.1 0.00917425190141665
3.125 0.00716508175155698
3.15 0.00534044137446139
3.175 0.00359567884903606
3.2 0.00192079282468207
3.225 0.000319731292477796
3.25 0.00120673187629871
3.275 0.00259468278899574
3.3 0.00400747389285228
3.325 0.00552669007709789
3.35 0.00690145186290519
3.375 0.00828607035493733
3.4 0.0103080235717949
3.425 0.0122555140476516
3.45 0.0141179014503684
3.475 0.000930918439993354
3.5 0.0175522470781492
3.525 0.0194731746699773
3.55 0.0212042644105164
3.575 0.0228468448453827
3.6 0.0243492893767263
3.625 0.0264419763743761
3.65 0.0280585013359203
3.675 0.0292408108612145
3.7 0.0307596789300958
3.725 0.0322062012946317
3.75 0.0335040589346945
3.775 0.0347331501306652
3.8 0.0358492920727369
3.825 0.0369178637475248
3.85 0.0379364235635357
3.875 0.0388630287450428
3.9 0.0397238192799848
3.925 0.0404465008230731
3.95 0.0411019096593025
3.975 0.0416683372308041
4 0.0417400270532535
4.025 0.0418106174927982
4.05 0.04187848515461
4.075 0.0419412357176481
4.1 0.0419998455300424
4.125 0.0419716689790136
4.15 0.0419138240906501
4.175 0.0418624661587798
4.2 0.0418125024525442
4.225 0.0417663039192073
4.25 0.0417223343723764
4.275 0.041680434890548
4.3 0.0416426985763516
4.325 0.0416061372973319
4.35 0.0415723637368928
4.375 0.0415923524081005
4.4 0.041784710314458
4.425 0.00164794447865664
4.45 0.0516592945093573
4.475 0.0423404451808666
4.5 0.0425339771702229
4.525 0.0427801873454714
4.55 0.0430197272993023
4.575 0.0432526660631627
4.6 0.0434579587432346
4.625 0.0436562176801357
4.65 0.0438390465335492
4.675 0.0439655257806944
4.7 0.0439068862458425
4.725 0.0438523456735977
4.75 0.043801742134144
4.775 0.0435951035304992
4.8 0.0433255862228156
4.825 0.0430833061191727
4.85 0.0428469853755793
4.875 0.0426303743477673
4.9 0.042272724946372
4.925 0.0417277971738127
4.95 0.0412067956072244
4.975 0.0406878973328025
5 0.040093655647362
5.025 0.0395577483770125
5.05 0.0390827200850529
5.075 0.0386295127953756
5.1 0.038186044023931
5.125 0.0378003001903472
5.15 0.037421008544834
5.175 0.0370634993887123
5.2 0.0368129493063491
5.225 0.0367151518971414
5.25 0.0366248888599905
5.275 0.0365382279796054
5.3 0.0363584424842424
5.325 0.036174769407985
5.35 0.0360187432551956
5.375 0.0357342915206352
5.4 0.0350189631163555
5.425 0.0343124121647861
5.45 0.0336150704181417
5.475 0.0327625881217262
5.5 0.0319327652384821
5.525 0.031217584919035
5.55 0.030531870579018
5.575 0.029873853531433
5.6 0.0291683197489171
5.625 0.0284503043149913
5.65 0.0277127579234485
5.675 0.0269585857146695
5.7 0.0261070702570375
5.725 0.0252772038499333
5.75 0.0244948365871197
5.775 0.023753557505925
5.8 0.0230460843499689
5.825 0.0223979162322381
5.85 0.0217752138803479
5.875 0.021533871892373
5.9 0.021291835677909
5.925 0.0210480526931526
5.95 0.0208310101839832
5.975 0.0206660651452854
6 0.020932283427054
6.025 0.021191517262773
6.05 0.0214348885778857
6.075 0.0219700335073666
6.1 0.0227441101400666
6.125 0.0234534179537148
6.15 0.024144304683699
6.175 0.024774533606781
6.2 0.0253863749094266
6.225 0.0259490148158042
6.25 0.0266057727940669
6.275 0.027353499013669
6.3 0.0280190187655241
6.325 0.028662917118124
6.35 0.0292534321193369
6.375 0.0298566331537929
6.4 0.0305012927367248
6.425 0.0311418514455273
6.45 0.0317780100884759
6.475 0.0324692739917272
6.5 0.0331279763543994
6.525 0.033688985984918
6.55 0.0342870333305839
6.575 0.0349210658862872
6.6 0.0354887616495046
6.625 0.036037667951111
6.65 0.037318795235879
6.675 0.0385837310907457
6.7 0.0397569237063454
6.725 0.0408474636107899
6.75 0.0418588276710913
6.775 0.0429812804023961
6.8 0.0441502248987133
6.825 0.0451885506834772
6.85 0.0462667532756874
6.875 0.0475315083447857
6.9 0.0487590588740332
6.925 0.0499366969521752
6.95 0.0511582219185159
6.975 0.0523841925021744
7 0.0534986567010675
7.025 0.0545165725381367
7.05 0.0554310091905185
7.075 0.05597409335529
7.1 0.0563633011666637
7.125 0.0567072831274513
7.15 0.0569876783514218
7.175 0.0569675172651814
7.2 0.056948141198807
7.225 0.0569297961077755
7.25 0.0566046628889646
7.275 0.0561556385875553
7.3 0.055737655122808
7.325 0.0553123279776291
7.35 0.0548291758802594
7.375 0.0543644103701772
7.4 0.0539282516224672
7.425 0.053518469143687
7.45 0.0531238401582796
7.475 0.0530605457293891
7.5 0.0530394288594526
7.525 0.053019357649522
7.55 0.0529150676665243
7.575 0.0527821205134967
7.6 0.0526678610896973
7.625 0.0524874880338097
7.65 0.0521647159187152
7.675 0.0518578892393511
7.7 0.0515642886907682
7.725 0.0512951811798999
7.75 0.0510305438768691
7.775 0.0507118143112995
7.8 0.0503929811792123
7.825 0.0500733145311238
7.85 0.0495585458152529
7.875 0.0487391536104326
7.9 0.047959084994135
7.925 0.0472059337774543
7.95 0.0462305519051314
7.975 0.0452130625384504
8 0.0442155357655194
8.025 0.0432044234726499
8.05 0.0421789199111946
8.075 0.0413367922618765
8.1 0.0405225574915033
8.125 0.0397593290962135
8.15 0.0389864532469252
8.175 0.0381962602598893
8.2 0.0374556527648801
8.225 0.0367479346479021
8.25 0.0358985166317877
8.275 0.0350344800540938
8.3 0.0342769839193291
8.325 0.0335384619323765
8.35 0.0328422537331374
8.375 0.0321830142175985
8.4 0.0315543515498977
8.425 0.0308482055002559
8.45 0.0301215527815056
8.475 0.0293429299096416
8.5 0.0285780811056558
8.525 0.0278613130010608
8.55 0.0268690573284208
8.575 0.0257951433186002
8.6 0.0248935990315108
8.625 0.0239459006066366
8.65 0.0228940850288927
8.675 0.0219064704440019
8.7 0.0209699343650884
8.725 0.0203675727359765
8.75 0.0198562266564151
8.775 0.0193899700825987
8.8 0.0189373957888283
8.825 0.0185211848341238
8.85 0.0182394348138724
8.875 0.0180860362942818
8.9 0.0179450462233644
8.925 0.0178080603438566
8.95 0.0179033891437041
8.975 0.0179995966627722
9 0.0180885681231342
9.025 0.0184185301379545
9.05 0.0190000670889748
9.075 0.0195183403231304
9.1 0.0200040430279172
9.125 0.0204248051166734
9.15 0.0208270977417118
9.175 0.0212118303342786
9.2 0.0215622841787689
9.225 0.0219023576751054
9.25 0.0223233588355407
9.275 0.0227362104508156
9.3 0.0231205924108513
9.325 0.0238082561075532
9.35 0.0247900713506881
9.375 0.0257183384453364
9.4 0.0266241159348821
9.425 0.0278893233302005
9.45 0.0291711379576596
9.475 0.0303959234370244
9.5 0.0317815441695558
9.525 0.0335439661376161
9.55 0.0351108821304218
9.575 0.036582819955911
9.6 0.0378160641167113
9.625 0.0391207433788545
9.65 0.0405821943450011
9.675 0.0419659180787217
9.7 0.0432856383563183
9.725 0.0447615527219445
9.75 0.0462751249480967
9.775 0.0476441720664905
9.8 0.0489753439759153
9.825 0.0502033378038996
9.85 0.0513763121926828
9.875 0.0524951675847219
9.9 0.0532630399350445
9.925 0.053920768691217
9.95 0.0545220964003146
9.975 0.0551087505161445
10 0.0556497570124183
10.025 0.0225843739299107
10.05 0.049305133199644
10.075 0.0569215149804316
10.1 0.0571696086872812
10.125 0.0571398186182733
10.15 0.0571109339581273
10.175 0.0570829489792512
10.2 0.0570576744193212
10.225 0.0570332072008517
10.25 0.0570105877212156
10.275 0.0570435108377658
10.3 0.0572403563770903
10.325 0.0574308603500816
10.35 0.0576135432224254
10.375 0.0577841927589238
10.4 0.0579496216908954
10.425 0.057991210965166
10.45 0.058026498458009
10.475 0.0580583137429669
10.5 0.0579972776213882
10.525 0.0577919812371441
10.55 0.0594468382743136
10.575 0.0573647758340038
10.6 0.0568378731545828
10.625 0.0563575245425989
10.65 0.0559242413337808
10.675 0.0553613071305664
10.7 0.0547453619704262
10.725 0.0542284642226387
10.75 0.053717473670496
10.775 0.0532076274499891
10.8 0.0526519787960712
10.825 0.0519896464705851
10.85 0.0513550955405798
10.875 0.0507468565494754
10.9 0.0501053962876794
10.925 0.0494535794672638
10.95 0.0488633557408286
10.975 0.0482343395346191
11 0.0474742413903244
11.025 0.0467167423495823
11.05 0.045961991336905
11.075 0.0452884000988201
11.1 0.0446272863021616
11.125 0.0439760936436389
11.15 0.0433495420689695
11.175 0.0427467439862248
11.2 0.0186546207986583
11.225 0.0421486505671494
11.25 0.0411884302250756
11.275 0.0407425874392129
11.3 0.0404157315108041
11.325 0.0401120812830512
11.35 0.0398303693104289
11.375 0.0397191937274682
11.4 0.0396798954675589
11.425 0.0396449669874875
11.45 0.0396524984988282
11.475 0.0398402642034397
11.5 0.0400252480763952
11.525 0.0402061436817667
11.55 0.0403774934606725
11.575 0.0405288539746703
11.6 0.0405592342031647
11.625 0.0405867173438491
11.65 0.0406113720256855
11.675 0.0407141363084883
11.7 0.0408478787361333
11.725 0.0409673029414032
11.75 0.0411240022199016
11.775 0.0414555006088702
11.8 0.041773023874273
11.825 0.0420716844440943
11.85 0.0423480845870751
11.875 0.0426137535714801
11.9 0.0427575111214031
11.925 0.0428859396901374
11.95 0.0430021999745392
11.975 0.0430708054431429
12 0.0430648517736617
12.025 0.043059271001645
12.05 0.0430539024785132
12.075 0.0430485617247672
12.1 0.0430433698042012
12.125 0.0430385130654874
12.15 0.0430342195914248
12.175 0.0430580342293255
12.2 0.0434998076854582
12.225 0.0439280870091652
12.25 0.0443260752601994
12.275 0.0447621182766088
12.3 0.0452286315919089
12.325 0.0456585740548785
12.35 0.0460873165065422
12.375 0.0468082521249068
12.4 0.0474884574507906
12.425 0.0480692356331978
12.45 0.0487030851365653
12.475 0.0494357376866815
12.5 0.0501734522620378
12.525 0.0509178335593052
12.55 0.0518041107829693
12.575 0.0527349466159507
12.6 0.0536015817320545
12.625 0.0543745074144306
12.65 0.055059105276401
12.675 0.0556384395708129
12.7 0.0561739133889976
12.725 0.0566290123719952
12.75 0.0570788498908901
12.775 0.0575069328796744
12.8 0.0579185273753245
12.825 0.0583081143924991
12.85 0.0589677034929508
12.875 0.0597386059486222
12.9 0.0604519316805775
12.925 0.0611829031081757
12.95 0.0620432750730388
12.975 0.0628687542441736
13 0.0636407259012504
13.025 0.0643709137382594
13.05 0.0650767506075755
13.075 0.0654418862401948
13.1 0.0657589135550939
13.125 0.0660563920168168
13.15 0.0663012881457246
13.175 0.0665287192000699
13.2 0.0667436120564862
13.225 0.0669373701808377
13.25 0.0670187288399794
13.275 0.0670976784848815
13.3 0.0671728205799505
13.325 0.0672403579218194
13.35 0.0673041471910935
13.375 0.0674602225329597
13.4 0.0676306043218253
13.425 0.0677836479075624
13.45 0.0677767629690307
13.475 0.0674691872315906
13.5 0.0671809819956593
13.525 0.0669058622683222
13.55 0.0666460893053094
13.575 0.0663911969919777
13.6 0.066148028492708
13.625 0.065665374836244
13.65 0.0649647122644543
13.675 0.0643821447457815
13.7 0.063818424008728
13.725 0.0633044708302146
13.75 0.0628022633074374
13.775 0.0623244178933219
13.8 0.0618786852304042
13.825 0.0614548332145521
13.85 0.0609721043269834
13.875 0.0604830874226082
13.9 0.060046837290523
13.925 0.0593009513132076
13.95 0.0578407876879546
13.975 0.05640240648227
14 0.0549871407067252
14.025 0.0529337578939867
14.05 0.0507241905310118
14.075 0.0486924571786
14.1 0.0466627431544503
14.125 0.0446376171267591
14.15 0.0420594290732335
14.175 0.0393494552190748
14.2 0.036997816441832
14.225 0.0346334153623516
14.25 0.0322540265713424
14.275 0.0300525501259409
14.3 0.0279227646148946
14.325 0.02633288595748
14.35 0.0247913250028745
14.375 0.0233572094618425
14.4 0.0218878628412901
14.425 0.0203719225292919
14.45 0.0189770933126975
14.475 0.0176377354396754
14.5 0.0161265169179041
14.525 0.0146115336150731
14.55 0.0131918523870971
14.575 0.0118372369455962
14.6 0.0105825432032133
14.625 0.00953752378201326
14.65 0.00861394467033343
14.675 0.00770781694864297
14.7 0.00690638149283064
14.725 0.00633150032643828
14.75 0.00582248429097655
14.775 0.00535578434683633
14.8 0.00516242502016178
14.825 0.00497341955936286
14.85 0.00479461859340614
14.875 0.00481994467359351
14.9 0.00499304487841015
14.925 0.00515332843695988
14.95 0.00531003467781676
14.975 0.00545481655105457
15 0.00565825151189213
15.025 0.0060058434947606
15.05 0.00633464526942826
15.075 0.00664587968349768
15.1 0.00718605158106193
15.125 0.00780872963419373
15.15 0.00836351137270913
15.175 0.00896207023926997
15.2 0.00970811483589581
15.225 0.0104344962777902
15.25 0.0111434056427396
15.275 0.013006982735649
15.3 0.0148616416770716
15.325 0.0164573676593931
15.35 0.0179965315752547
15.375 0.0194816235576886
15.4 0.0211628639409793
15.425 0.0228999076657065
15.45 0.024503271503943
15.475 0.0261492302020307
15.5 0.0279205082522297
15.525 0.0295826320369752
15.55 0.0311395924308808
15.575 0.033015947579894
15.6 0.0350161521327148
15.625 0.0367659482534773
15.65 0.0385251870801428
15.675 0.0403507541256717
15.7 0.0421320074331953
15.725 0.0438526998026403
15.75 0.0455151994286943
15.775 0.0470630959818971
15.8 0.0481097569008054
15.825 0.0490824657096764
15.85 0.0499754294676562
15.875 0.0503340540938549
15.9 0.0504763179045789
15.925 0.0506043250394593
15.95 0.0506268533340045
15.975 0.0501875513532176
16 0.0497616751560622
16.025 0.0493541042375007
16.05 0.0489889389262982
16.075 0.0486407399153824
16.1 0.047912264748648
16.125 0.0471598430194179
16.15 0.0464731654373724
16.175 0.0455959808584437
16.2 0.0444209548265299
16.225 0.0433072962796525
16.25 0.0422294250745643
16.275 0.041068218450726
16.3 0.039926521255374
16.325 0.0388478873867961
16.35 0.0379059822695427
16.375 0.0369827433928762
16.4 0.0356286188487179
16.425 0.0343258951636832
16.45 0.033126594469852
16.475 0.0316961781448336
16.5 0.0300959832406023
16.525 0.0286031572440493
16.55 0.0271349531868337
16.575 0.0253378293550722
16.6 0.0236043302055387
16.625 0.0220288761150735
16.65 0.0203102417040775
16.675 0.0184551841650977
16.7 0.0167116421639243
16.725 0.0150084776103158
16.75 0.0134851901464039
16.775 0.0118604822710618
16.8 0.00990289040125165
16.825 0.00797321822543921
16.85 0.0060704936950458
16.875 0.00379998437945094
16.9 0.00162408895184712
16.925 0.000245865153594583
16.95 0.00240660724797677
16.975 0.00484476589169738
17 0.00700347978626173
17.025 0.00909097880452181
17.05 0.0110490655101411
17.075 0.0134289769863933
17.1 0.0169859010803222
17.125 0.0203298323714284
17.15 0.0234371628231518
17.175 0.0275366954781984
17.2 0.0322207771845197
17.225 0.0360946233634679
17.25 0.040214177071323
17.275 0.0454134974249623
17.3 0.0504940301999144
17.325 0.0554170296601922
17.35 0.0638945718094495
17.375 0.0733348308017881
17.4 0.0816925501081944
17.425 0.0893846255933985
17.45 0.0963870996339837
17.475 0.104509966617721
17.5 0.113099211968271
17.525 0.120723317736903
17.55 0.128625448197616
17.575 0.137864406533366
17.6 0.146693268925681
17.625 0.154988684624959
17.65 0.16624007817892
17.675 0.17891559817192
17.7 0.18998163010242
17.725 0.201130304091572
17.75 0.215205779736734
17.775 0.228544401617479
17.8 0.240717148713732
17.825 0.254235333349373
17.85 0.268536806821565
17.875 0.282619816047361
17.9 0.296246177092236
17.925 0.308967321872066
17.95 0.322831380807163
17.975 0.336959959361066
18 0.349366311306964
18.025 0.363556222374015
18.05 0.38159758738865
18.075 0.398643013982791
18.1 0.414887186816871
18.125 0.43138938990489
18.15 0.447959006167044
18.175 0.464448236773233
18.2 0.481187275945912
18.225 0.499982277460327
18.25 0.517767613526006
18.275 0.533980797074274
18.3 0.554374110347912
18.325 0.577620478494071
18.35 0.599394526666116
18.375 0.620705177661819
18.4 0.640854364312325
18.425 0.662229533663766
18.45 0.684794260392932
18.475 0.703303281326228
18.5 0.721685241353057
18.525 0.74229032131788
18.55 0.762184291334003
18.575 0.780775298601419
18.6 0.800707293464191
18.625 0.821543956046575
18.65 0.840722136652283
18.675 0.860035719166599
18.7 0.885315043432724
18.725 0.910378009987315
18.75 0.934803116271472
18.775 0.959894458605611
18.8 0.985514654623191
18.825 1.00964031491779
18.85 1.03315995314325
18.875 1.05401039049577
18.9 1.07643387801769
18.925 1.10230589982071
18.95 1.12586339217794
18.975 1.1495252625463
19 1.17800672438313
19.025 1.20450845108865
19.05 1.22882745579058
19.075 1.25379271018141
19.1 1.27902538517729
19.125 1.30184632927338
19.15 1.32482214297243
19.175 1.34902926496669
19.2 1.37220204797256
19.225 1.39399921058012
19.25 1.42512155897684
19.275 1.46081521962101
19.3 1.49340049171301
19.325 1.52517537765532
19.35 1.55414597141253
19.375 1.58426119207533
19.4 1.61619334445924
19.425 1.64648587358356
19.45 1.67581559583443
19.475 1.70666592771266
19.5 1.73694043521999
19.525 1.76560107700804
19.55 1.79389576885965
19.575 1.82151993731505
19.6 1.84466668616142
19.625 1.86849133225063
19.65 1.89384718226811
19.675 1.91772054760153
19.7 1.94040957520715
19.725 1.96132408954908
19.75 1.98201360974035
19.775 2.00882225542413
19.8 2.03494448148737
19.825 2.05919691305911
19.85 2.08394041833894
19.875 2.10915762546669
19.9 2.13308434371337
19.925 2.15651695414636
19.95 2.18128365422603
19.975 2.20559612711451
20 2.22835239344105
20.025 2.25514264469042
20.05 2.28595690345595
20.075 2.31362891343604
20.1 2.33936209987282
20.125 2.36139228918606
20.15 2.38456210882462
20.175 2.40882549247525
20.2 2.4304224555356
20.225 2.45131939249065
20.25 2.47488305184092
20.275 2.49762930602905
20.3 2.51809428234759
20.325 2.54160187725095
20.35 2.56851677140092
20.375 2.59364663847068
20.4 2.61803620638292
20.425 2.64353596369807
20.45 2.6686410441614
20.475 2.69181605965899
20.5 2.7160234694134
20.525 2.74156349687695
20.55 2.76616529269913
20.575 2.78993284910355
20.6 2.81094738337209
20.625 2.83187468575636
20.65 2.85271585895476
20.675 2.87288061758859
20.7 2.89275597180972
20.725 2.9152473475692
20.75 2.93751568501055
20.775 2.95871327898736
20.8 2.98002706505988
20.825 3.00149205290701
20.85 3.0215216279289
20.875 3.04084439709682
20.9 3.06090638651476
20.925 3.08074764931588
20.95 3.09887650015776
20.975 3.11629333253814
21 3.13260531129625
21.025 3.15125704631818
21.05 3.17125397687198
21.075 3.19071981627751
21.1 3.21034024552285
21.125 3.23039122409911
21.15 3.24829414060972
21.175 3.26539062500419
21.2 3.28192785877827
21.225 3.29804735680378
21.25 3.31325235953081
21.275 3.32720722587874
21.3 3.34030831946944
21.325 3.35523976459716
21.35 3.37059949780006
21.375 3.38486138139377
21.4 3.39884199533578
21.425 3.4122003497547
21.45 3.42496440782774
21.475 3.43716654030764
21.5 3.44935813923867
21.525 3.46154606254447
21.55 3.47255119517105
21.575 3.483651383132
21.6 3.49510922557977
21.625 3.5065112447136
21.65 3.51785746826286
21.675 3.52793942753779
21.7 3.53747237326957
21.725 3.54601855255526
21.75 3.55435540690651
21.775 3.56248421252031
21.8 3.56987520541131
21.825 3.57706956123351
21.85 3.58383704096053
21.875 3.59033227131328
21.9 3.59611894232224
21.925 3.60160042993954
21.95 3.60673477772107
21.975 3.61148867083501
22 3.61608715698155
22.025 3.61994122212596
22.05 3.62368658925056
22.075 3.62713659240528
22.1 3.63040615715488
22.125 3.63343061046277
22.15 3.63633913763346
22.175 3.63914459298881
22.2 3.64162587932398
22.225 3.64388395849529
22.25 3.64593274949318
22.275 3.6478783576382
22.3 3.64978319960428
22.325 3.65145866574443
22.35 3.65310263893143
22.375 3.65463508452507
22.4 3.65585572440297
22.425 3.6566689789565
22.45 3.65743317200153
22.475 3.65817214511083
22.5 3.65884513222553
22.525 3.65945065050725
22.55 3.65971589308469
22.575 3.65997010301879
22.6 3.6602075589015
22.625 3.65956800076006
22.65 3.65842215395071
22.675 3.65730509655788
22.7 3.65622287932382
22.725 3.65520647072733
22.75 3.65370430315619
22.775 3.65206553382523
22.8 3.65063888244137
22.825 3.64903776406497
22.85 3.64703197368698
22.875 3.64519131175082
22.9 3.64346741205732
22.925 3.64091219060076
22.95 3.63813811807561
22.975 3.63559439923262
23 3.63312600310578
23.025 3.63086557141359
23.05 3.62832136863562
23.075 3.6254903599518
23.1 3.62287093264515
23.125 3.62032445082831
23.15 3.61705321593707
23.175 3.61383150655683
23.2 3.61080720005285
23.225 3.60782162452561
23.25 3.60487373961777
23.275 3.59970978127363
23.3 3.59473280814744
23.325 3.59053210297617
23.35 3.58608223601134
23.375 3.58139059005117
23.4 3.57720955981001
23.425 3.57316547956068
23.45 3.56937946135319
23.475 3.56553454928535
23.5 3.5615237046182
23.525 3.55764964217185
23.55 3.55393574619875
23.575 3.55024901998613
23.6 3.54657373404605
23.625 3.54324296958245
23.65 3.53992099023105
23.675 3.53655704476945
23.7 3.53323667512785
23.725 3.52997986866198
23.75 3.52686586464606
23.775 3.52382829109783
23.8 3.52093911950931
23.825 3.51831340136208
23.85 3.51594477531742
23.875 3.51337588416051
23.9 3.51071823965078
23.925 3.50831825265762
23.95 3.5059816600345
23.975 3.50384971038052
24 3.50147619688393
24.025 3.49877642652999
24.05 3.49624561178001
24.075 3.49380015257154
24.1 3.49140224820227
24.125 3.48903832127165
24.15 3.48685276695079
24.175 3.48471640279898
24.2 3.48265522550026
24.225 3.48068945749553
24.25 3.47877883635321
24.275 3.47693677504693
24.3 3.475183824117
24.325 3.47360694280698
24.35 3.47232607900201
24.375 3.47113539386022
24.4 3.47010352634718
24.425 3.46910284528941
24.45 3.46819184404049
24.475 3.46786961278951
24.5 3.46849625104405
24.525 3.46907018005431
24.55 3.46961518183578
24.575 3.4703851541371
24.6 3.47064062350245
24.625 3.47121442386473
24.65 3.47175750407042
24.675 3.47225101306139
24.7 3.47291801923278
24.725 3.47371104958367
24.75 3.47446418817823
24.775 3.47519956812986
24.8 3.47588250191083
24.825 3.47735430006816
24.85 3.47958468220453
24.875 3.48139643026348
24.9 3.48325735979681
24.925 3.48576919585787
24.95 3.48819543855676
24.975 3.49046299162382
25 3.49311059290947
25.025 3.49602338563937
25.05 3.49871641027344
25.075 3.50135127818903
25.1 3.50384610835494
25.125 3.50625826682239
25.15 3.50848662451288
25.175 3.51060500134037
25.2 3.51263130116714
25.225 3.51456705098991
25.25 3.51647604761824
25.275 3.5182570417433
25.3 3.52022634179652
25.325 3.52260000727595
25.35 3.5247714895654
25.375 3.52682433860096
25.4 3.52860960616397
25.425 3.53037332293543
25.45 3.53211336711615
25.475 3.53376528156973
25.5 3.53537844840413
25.525 3.53710634896218
25.55 3.5388147688877
25.575 3.54041307649245
25.6 3.54184059137809
25.625 3.54304428077412
25.65 3.54416544858249
25.675 3.54524653788518
25.7 3.54612786212184
25.725 3.54697723668819
25.75 3.54775507340155
25.775 3.54850259784782
25.8 3.54920194180864
25.825 3.5495123861067
25.85 3.54958986244664
25.875 3.54966447294611
25.9 3.54936797764758
25.925 3.54832228966534
25.95 3.54739830459951
25.975 3.54652339047888
26 3.54576580823803
26.025 3.5449255504019
26.05 3.54389741835788
26.075 3.54293889437659
26.1 3.54203399112837
26.125 3.54110179562472
26.15 3.54016555657064
26.175 3.53932507174916
26.2 3.53835602165534
26.225 3.5370449782747
26.25 3.53580246172816
26.275 3.53462525299166
26.3 3.53319721081728
26.325 3.53168413149619
26.35 3.53032453749652
26.375 3.52899030906867
26.4 3.52773451387777
26.425 3.52592642833538
26.45 3.52362455024413
26.475 3.52153914082137
26.5 3.51943137978876
26.525 3.51723795589095
26.55 3.51510471490144
26.575 3.51302924673528
26.6 3.51084150760613
26.625 3.50862393471973
26.65 3.50655324415787
26.675 3.50454671154637
26.7 3.50272574522079
26.725 3.50087824365245
26.75 3.4990014423737
26.775 3.49723474732509
26.8 3.49551171293634
26.825 3.49401737621531
26.85 3.49255100032006
26.875 3.49118902101999
26.9 3.48986456177379
26.925 3.48858951865825
26.95 3.48734146442899
26.975 3.48613210016993
27 3.48507660417368
27.025 3.4838393327183
27.05 3.48241672585918
27.075 3.48107236706797
27.1 3.47976180120253
27.125 3.4783639918866
27.15 3.47698960538083
27.175 3.47573801487414
27.2 3.47447312786023
27.225 3.47318988023389
27.25 3.47204393075587
27.275 3.47095295390784
27.3 3.46993327702716
27.325 3.46907744070463
27.35 3.46905874974472
27.375 3.46904076013782
27.4 3.46902386970421
27.425 3.46944176940057
27.45 3.47011550803733
27.475 3.47076007921161
27.5 3.47158904298136
27.525 3.47280455273786
27.55 3.47386829502034
27.575 3.47487384668433
27.6 3.47574941934699
27.625 3.47666659257971
27.65 3.47767628730032
27.675 3.47863401448282
27.7 3.47954897946335
27.725 3.48048408349574
27.75 3.48141728073758
27.775 3.48223841605713
27.8 3.48307160146108
27.825 3.48393815594595
27.85 3.4847689051995
27.875 3.48556537427003
27.9 3.48627177371304
27.925 3.48695459721023
27.95 3.48890305362923
27.975 3.49082686472783
28 3.4926205257857
28.025 3.49484060544286
28.05 3.49749095325152
28.075 3.49984477990067
28.1 3.50202850428343
28.125 3.50389898504672
28.15 3.50616382246591
28.175 3.50879865946545
28.2 3.51119354149216
28.225 3.51351710954457
28.25 3.51587074132355
28.275 3.51815921938188
28.3 3.52024277369514
28.325 3.52218165886842
28.35 3.52395387158399
28.375 3.5256460008925
28.4 3.5273028590165
28.425 3.52838104472281
28.45 3.52943786766786
28.475 3.53041703320523
28.5 3.53139840007313
28.525 3.5323835200528
28.55 3.53328871386237
28.575 3.53415569899106
28.6 3.53494287722621
28.625 3.5356553813859
28.65 3.53629645710441
28.675 3.53688983072464
28.7 3.53746371857548
28.725 3.53821943070524
28.75 3.53897711129404
28.775 3.53965634181858
28.8 3.54030460340061
28.825 3.54091243341605
28.85 3.54155511462708
28.875 3.54221522852116
28.9 3.54281621375726
28.925 3.54335237955789
28.95 3.54355208490384
28.975 3.5437460070105
29 3.54393008028148
29.025 3.54352574353661
29.05 3.54278019467555
29.075 3.54206031757581
29.1 3.54136347052101
29.125 3.54071006414418
29.15 3.53978076324413
29.175 3.5387679461035
29.2 3.53789560519088
29.225 3.53694605371515
29.25 3.53582396278707
29.275 3.53476262084749
29.3 3.53375055359337
29.325 3.5330896639109
29.35 3.53251895115579
29.375 3.53200163442398
29.4 3.53149935202163
29.425 3.53103681189138
29.45 3.53053323122158
29.475 3.52998928964645
29.5 3.52950733436717
29.525 3.52904184289721
29.55 3.52773904186508
29.575 3.52644544214957
29.6 3.52522072320101
29.625 3.52357461403269
29.65 3.52153907612803
29.675 3.5196875240058
29.7 3.51793932297582
29.725 3.516366234688
29.75 3.51479013741149
29.775 3.51321253605616
29.8 3.51171535446417
29.825 3.5102389258294
29.85 3.5085495241046
29.875 3.50689936239307
29.9 3.50538001025895
29.925 3.5039331558997
29.95 3.50256743045042
29.975 3.50127584299958
30 3.5000150314059
30.025 3.49892623197581
30.05 3.49790066646353
30.075 3.49713751769503
30.1 3.4963797784626
30.125 3.49563024867335
30.15 3.49473942269398
30.175 3.49381903764517
30.2 3.49298608801487
30.225 3.49201952562981
30.25 3.49093809715048
30.275 3.49000340076219
30.3 3.48909867766884
30.325 3.48827198978768
30.35 3.48728162567061
30.375 3.48597169236353
30.4 3.48474515411297
30.425 3.48357611606908
30.45 3.4821475646733
30.475 3.48067344877016
30.5 3.47933456698784
30.525 3.47799676846436
30.55 3.47666175322514
30.575 3.47534813249428
30.6 3.47406801528608
30.625 3.4726146135487
30.65 3.47105269515924
30.675 3.469574881804
30.7 3.46813509796942
30.725 3.46676239541105
30.75 3.46460183729979
30.775 3.46222840655845
30.8 3.46015967865129
30.825 3.45791183519143
30.85 3.45525323906634
30.875 3.45274117817281
30.9 3.45034880884389
30.925 3.44813451051851
30.95 3.44592432851682
30.975 3.44302451396429
31 3.44019846652735
31.025 3.43756363446017
31.05 3.43392907578462
31.075 3.42936196444322
31.1 3.42528350988953
31.125 3.42133452992544
31.15 3.41649852993306
31.175 3.41173247921026
31.2 3.40722152019811
31.225 3.40280349462376
31.25 3.39846889378193
31.275 3.393173178719
31.3 3.38805158564211
31.325 3.38341594261604
31.35 3.37701777644293
31.375 3.36975805134091
31.4 3.36311887176154
31.425 3.35633695295517
31.45 3.34845562091313
31.475 3.34086606317875
31.5 3.33372172938866
31.525 3.32567349611873
31.55 3.31713205615053
31.575 3.30931933138679
31.6 3.30122433501184
31.625 3.28977484928827
31.65 3.27840504322446
31.675 3.26736221887234
31.7 3.25657966683319
31.725 3.24616305896993
31.75 3.23389099563917
31.775 3.22100652254282
31.8 3.20901710103864
31.825 3.19663378041891
31.85 3.18393534129983
31.875 3.17242361154262
31.9 3.160561083381
31.925 3.14414139603487
31.95 3.12827035239722
31.975 3.11342256174471
32 3.0994782337967
32.025 3.08614822053823
32.05 3.07110210312872
32.075 3.05582953318763
32.1 3.04194499036854
32.125 3.02729472461917
32.15 3.01095842831005
32.175 2.99547503381106
32.2 2.98069552414978
32.225 2.96042748743337
32.25 2.93887849489365
32.275 2.91945239237163
32.3 2.89859086661435
32.325 2.87488734798859
32.35 2.85331748217955
32.375 2.83303202492041
32.4 2.8156030268306
32.425 2.79755720943338
32.45 2.7788255147909
32.475 2.76129995655423
32.5 2.74428571489395
32.525 2.72543499228396
32.55 2.70684277523795
32.575 2.68980201226006
32.6 2.67030015617786
32.625 2.64752078665386
32.65 2.62611423618532
32.675 2.60541210054091
32.7 2.60918833424665
32.725 2.56003338519708
32.75 2.53953486639948
32.775 2.51825358370813
32.8 2.49576538581225
32.825 2.47408167702896
32.85 2.45288574983436
32.875 2.43244999679766
32.9 2.40946837310165
32.925 2.38164177258743
32.95 2.35753690403403
32.975 2.33441501816898
33 2.31044176405235
33.025 2.28697913213889
33.05 2.26467710803421
33.075 2.24212126028671
33.1 2.21939238959634
33.125 2.1987378816953
33.15 2.17621323208626
33.175 2.15357970655093
33.2 2.13149049968689
33.225 2.11085488541089
33.25 2.08672410467926
33.275 2.05910833937282
33.3 2.03446841653969
33.325 2.01063958718687
33.35 1.97664419353156
33.375 1.94305252980103
33.4 1.91094366478981
33.425 1.87974332094757
33.45 1.84932171008278
33.475 1.81943472927451
33.5 1.79181920471497
33.525 1.76856029430243
33.55 1.74227714094924
33.575 1.71292302945731
33.6 1.68849128603188
33.625 1.66537156377182
33.65 1.64146960901813
33.675 1.6174180951244
33.7 1.5935362686859
33.725 1.56949587280498
33.75 1.5443779369111
33.775 1.5202576823738
33.8 1.49770574830296
33.825 1.4736256787401
33.85 1.44865597996983
33.875 1.42472630617282
33.9 1.40151046443805
33.925 1.37953348667978
33.95 1.30435020088538
33.975 1.32813528307277
34 1.3050788580828
34.025 1.28154783903337
34.05 1.256944084154
34.075 1.23439657967048
34.1 1.21326730799578
34.125 1.19097127836411
34.15 1.16857266703384
34.175 1.14896819320139
34.2 1.12799564993756
34.225 1.10338708930238
34.25 1.0798979254598
34.275 1.05751224178821
34.3 1.03448030162168
34.325 1.01123576577745
34.35 0.99052429399778
34.375 0.970107762668252
34.4 0.950681911271828
34.425 0.929780607048498
34.45 0.90757244633141
34.475 0.887440868019283
34.5 0.867333588262886
34.525 0.847012386892439
34.55 0.828563816364588
34.575 0.8110497585659
34.6 0.789883127878166
34.625 0.768676720083586
34.65 0.749223143781047
34.675 0.729276373026001
34.7 0.708559863349167
34.725 0.689793369604557
34.75 0.671988106353826
34.775 0.655356560980265
34.8 0.638765109398681
34.825 0.621077004727509
34.85 0.60432341572922
34.875 0.589035921315935
34.9 0.571155561847501
34.925 0.551508294175912
34.95 0.533357381054118
34.975 0.515189236516231
35 0.495856741415309
35.025 0.4774021684172
35.05 0.459697646160834
35.075 0.445565954674604
35.1 0.431542793647069
35.125 0.416261067372524
35.15 0.401559708000822
35.175 0.387937812197724
35.2 0.372944267921576
35.225 0.356980176402963
35.25 0.342477646594658
35.275 0.328243159131682
35.3 0.313863266332165
35.325 0.299947039296751
35.35 0.287040652395106
35.375 0.274239253511002
35.4 0.261524043606119
35.425 0.249913034247832
35.45 0.23857048691983
35.475 0.226910870945729
35.5 0.215310078378534
35.525 0.203827297605311
35.55 0.192025902204609
35.575 0.180504892295186
35.6 0.170587049783392
35.625 0.161213236187488
35.65 0.152461151413508
35.675 0.143998485931415
35.7 0.135652808474393
35.725 0.128400980859076
35.75 0.121285965311757
35.775 0.114674750060279
35.8 0.108327522527182
35.825 0.102332737174582
35.85 0.096908612009724
35.875 0.0917578658289242
35.9 0.0871698810824815
35.925 0.0827341584724207
35.95 0.0790478584480782
35.975 0.0755342231077944
36 0.072287924536184
36.025 0.0695275688784588
36.05 0.0670551213032032
36.075 0.0646988023297341
36.1 0.0624067662943787
36.125 0.0602260957938459
36.15 0.0586279563068764
36.175 0.0572037616470194
36.2 0.055956188984885
36.225 0.054940470509046
36.25 0.0544659123367125
36.275 0.0540254285371651
36.3 0.0536106146211718
36.325 0.0536632995743997
36.35 0.0538400287111955
36.375 0.0539993668183689
36.4 0.054154347859754
36.425 0.054297577983898
36.45 0.0545987799409128
36.475 0.0550532282187999
36.5 0.0554858103614232
36.525 0.0559098119648317
36.55 0.0568355519594502
36.575 0.0577536918217105
36.6 0.0586336551025853
36.625 0.059749903198182
36.65 0.0610511889225896
36.675 0.062245645010947
36.7 0.0634141583346353
36.725 0.0645460946881252
36.75 0.0655711447516819
36.775 0.0665430106064401
36.8 0.0674310442299998
36.825 0.0683479000912137
36.85 0.0694661923889504
36.875 0.0705511394001264
36.9 0.0715846328239481
36.925 0.0728103839485714
36.95 0.0741282073676508
36.975 0.0753556572452934
37 0.076591239092667
37.025 0.0780033018606415
37.05 0.079378330177419
37.075 0.0806779052132153
37.1 0.082217604297365
37.125 0.0839247723030692
37.15 0.0855771665958319
37.175 0.0871933995846432
37.2 0.0886924649618426
37.225 0.0901854606757357
37.25 0.0916734331708779
37.275 0.092873027373224
37.3 0.0947894788077575
37.325 0.096158064709031
37.35 0.096193548954935
37.375 0.0971880834702321
37.4 0.0979895384690288
37.425 0.0986571355731403
37.45 0.0992730818351006
37.475 0.0997854767967812
37.5 0.0991370415547667
37.525 0.0984874551709728
37.55 0.0978340656561057
37.575 0.0969926725329499
37.6 0.0960190750824526
37.625 0.095089433608319
37.65 0.0941792126476118
37.675 0.0933628930928086
37.7 0.0923299998323924
37.725 0.0908799690167937
37.75 0.0895968678126332
37.775 0.0883120393145822
37.8 0.0867710532876744
37.825 0.0853238190984908
37.85 0.0839854317278287
37.875 0.0828118334402272
37.9 0.0817051502384574
37.925 0.0807408251696595
37.95 0.0797732812046076
37.975 0.0787803183215203
38 0.0778307561493616
38.025 0.0769385538238412
38.05 0.0761042664845148
38.075 0.0752985702970274
38.1 0.0727521948422332
38.125 0.0700803793887392
38.15 0.0676316894056146
38.175 0.0648017944105832
38.2 0.0613683585179966
38.225 0.0580755007505
38.25 0.0548645544579951
38.275 0.0507561958995987
38.3 0.0467134971751228
38.325 0.0428971249983246
38.35 0.0396052846761258
38.375 0.0363679146983763
38.4 0.0319390847478411
38.425 0.0276577249798974
38.45 0.0236935858805228
38.475 0.0196404875869742
38.5 0.0155154592219613
38.525 0.0116620781651776
38.55 0.00789038267005615
38.575 0.00427027508948957
38.6 0.000734468133665922
38.625 0.00257546078165048
38.65 0.00571741458153111
38.675 0.00869864891170328
38.7 0.0117750497198474
38.725 0.014882821104266
38.75 0.0177080984344937
38.775 0.0204920258160864
38.8 0.0231523092463983
38.825 0.0257415324296114
38.85 0.02827274878742
38.875 0.0306227817731826
38.9 0.0328608595710501
38.925 0.0348964874960468
38.95 0.03675251259068
38.975 0.0385162835014938
39 0.0396014037561721
39.025 0.0406062890210453
39.05 0.0415231396193089
39.075 0.0421920027061725
39.1 0.0424707628024162
39.125 0.0427229227070186
39.15 0.0429619720779232
39.175 0.0428964010574425
39.2 0.0427849951254011
39.225 0.0426841805235613
39.25 0.0424325079156703
39.275 0.0418941217783383
39.3 0.0413911269691768
39.325 0.0409121125673162
39.35 0.0404457003479271
39.375 0.0399246998334258
39.4 0.039147760447454
39.425 0.0383996467843651
39.45 0.037675573372633
39.475 0.0368473336350497
39.5 0.0360403783011563
39.525 0.0353067329517679
39.55 0.0342575272535809
39.575 0.0329080973134943
39.6 0.0317093226357303
39.625 0.0305503465021417
39.65 0.0294566067701922
39.675 0.0282604530392112
39.7 0.0267782257533704
39.725 0.0253466089045804
39.75 0.0239734502882566
39.775 0.0226232823989178
39.8 0.0212823421760754
39.825 0.0200765974840904
39.85 0.0189396187124743
39.875 0.018075924651015
39.9 0.0172288693117017
39.925 0.0164041670077174
39.95 0.0156194109108004
39.975 0.0148845564389049
40 0.014327025389697
40.025 0.0138056953552514
40.05 0.0133152549456281
40.075 0.0130768535061336
40.1 0.0128802283054872
40.125 0.0127031076265278
40.15 0.0124550872512488
40.175 0.0120733440500074
40.2 0.0117107605957766
40.225 0.0113618248699746
40.25 0.0109567742677314
40.275 0.0105440946113783
40.3 0.010165459401873
40.325 0.00979843201522314
40.35 0.00945780522813607
40.375 0.00919581477650779
40.4 0.00899942634591127
40.425 0.00881840966010503
40.45 0.00866531721333953
40.475 0.00907664704651874
40.5 0.00948561246915795
40.525 0.00989007009494812
40.55 0.0105398722362369
40.575 0.0112335105398985
40.6 0.0118338742306321
40.625 0.0124032078476351
40.65 0.012935368798304
40.675 0.0138355015552595
40.7 0.0148910871491464
40.725 0.0158025103120952
40.75 0.016862523966703
40.775 0.0185874570023137
40.8 0.0202404749274741
40.825 0.0217969076720808
40.85 0.0230964151424298
40.875 0.0242709778256596
40.9 0.025337959547241
40.925 0.0263780894300366
40.95 0.0273266052714329
40.975 0.028196803062957
41 0.0289442168654491
41.025 0.029663598717992
41.05 0.0303668867062307
41.075 0.0306782415881494
41.1 0.0309657289564923
41.125 0.0312404373842086
41.15 0.0313879076926793
41.175 0.0314983581805248
41.2 0.0315954615324617
41.225 0.0316899391324363
41.25 0.03177860088574
41.275 0.0318611801422544
41.3 0.0319392294269222
41.325 0.032012592251298
41.35 0.0320913654348882
41.375 0.0322619893879679
41.4 0.0324293614969951
41.425 0.0325876758193751
41.45 0.032738986017293
41.475 0.0328837026771028
41.5 0.0330124773656732
41.525 0.0331370995950733
41.55 0.033251776394515
41.575 0.0333648307950029
41.6 0.0334727372981813
41.625 0.033577924053141
41.65 0.0336808773200302
41.675 0.033424153669018
41.7 0.0331649268960064
41.725 0.0329285371775062
41.75 0.0325669812404885
41.775 0.0321444920077049
41.8 0.0317590695936748
41.825 0.0313666830491258
41.85 0.0309118984607449
41.875 0.0304708982810854
41.9 0.0300515541760494
41.925 0.0296759932256767
41.95 0.02932171572635
41.975 0.0288904093714641
42 0.0284524904131624
42.025 0.0280843836640323
42.05 0.0276751326518466
42.075 0.0270959745393589
42.1 0.0265260917923815
42.125 0.0259683889999963
42.15 0.0252626248699349
42.175 0.0245261389217737
42.2 0.0238447085055249
42.225 0.0232137737273139
42.25 0.0226249827965509
42.275 0.0214296636592117
42.3 0.0201595660206856
42.325 0.019010067001732
};
\addplot [line width=2.4000000000000004pt, color1, dashed, forget plot]
table {%
0 0.0937500001330113
0.025 0.0936190645892096
0.05 0.093488129045408
0.075 0.0933571935016064
0.1 0.0932262579578047
0.125 0.0930953224140031
0.15 0.0929643868702014
0.175 0.0928334513263998
0.2 0.0927025157825982
0.225 0.0925715802387965
0.25 0.0924406446949949
0.275 0.0923097091511932
0.3 0.0921787736073916
0.325 0.09204783806359
0.35 0.0919169025197883
0.375 0.0917859669759867
0.4 0.091655031432185
0.425 0.0915240958883834
0.45 0.0913931603445818
0.475 0.0912622248007801
0.5 0.0911312892569785
0.525 0.0910003537131768
0.55 0.0908694181693752
0.575 0.0907384826255735
0.6 0.0906075470817719
0.625 0.0904766115379703
0.65 0.0903456759941686
0.675 0.090214740450367
0.7 0.0900838049065654
0.725 0.0899528693627637
0.75 0.0898219338189621
0.775 0.0896909982751604
0.8 0.0895600627313588
0.825 0.0894291271875571
0.85 0.0892981916437555
0.875 0.0891672560999539
0.9 0.0890363205561522
0.925 0.0889053850123506
0.95 0.0887744494685489
0.975 0.0886435139247473
1 0.0885125783809457
1.025 0.088381642837144
1.05 0.0882507072933424
1.075 0.0881197717495407
1.1 0.0879888362057391
1.125 0.0878579006619375
1.15 0.0877269651181358
1.175 0.0875960295743342
1.2 0.0874650940305325
1.225 0.0873341584867309
1.25 0.0872032229429293
1.275 0.0870722873991276
1.3 0.086941351855326
1.325 0.0868104163115243
1.35 0.0866794807677227
1.375 0.0865485452239211
1.4 0.0864176096801194
1.425 0.0862866741363178
1.45 0.0861557385925161
1.475 0.0860248030487145
1.5 0.0858938675049128
1.525 0.0857629319611112
1.55 0.0856319964173096
1.575 0.0855010608735079
1.6 0.0853701253297063
1.625 0.0852391897859047
1.65 0.085108254242103
1.675 0.0849773186983014
1.7 0.0848463831544997
1.725 0.0847154476106981
1.75 0.0845845120668964
1.775 0.0844535765230948
1.8 0.0843226409792932
1.825 0.0841917054354915
1.85 0.0840607698916899
1.875 0.0839298343478882
1.9 0.0837988988040866
1.925 0.083667963260285
1.95 0.0835370277164833
1.975 0.0834060921726817
2 0.08327515662888
2.025 0.0831442210850784
2.05 0.0830132855412768
2.075 0.0828823499974751
2.1 0.0827514144536735
2.125 0.0826204789098718
2.15 0.0824895433660702
2.175 0.0823586078222686
2.2 0.0822276722784669
2.225 0.0820967367346653
2.25 0.0819658011908636
2.275 0.081834865647062
2.3 0.0817039301032604
2.325 0.0815729945594587
2.35 0.0814420590156571
2.375 0.0813111234718554
2.4 0.0811801879280538
2.425 0.0810492523842521
2.45 0.0809183168404505
2.475 0.0807873812966489
2.5 0.0806564457528472
2.525 0.0805255102090456
2.55 0.080394574665244
2.575 0.0802636391214423
2.6 0.0801327035776407
2.625 0.080001768033839
2.65 0.0798708324900374
2.675 0.0797398969462357
2.7 0.0796089614024341
2.725 0.0794780258586325
2.75 0.0793470903148308
2.775 0.0792161547710292
2.8 0.0790852192272275
2.825 0.0789542836834259
2.85 0.0788233481396243
2.875 0.0786924125958226
2.9 0.078561477052021
2.925 0.0784305415082193
2.95 0.0782996059644177
2.975 0.0781686704206161
3 0.0780377348768144
3.025 0.0779067993330128
3.05 0.0777758637892111
3.075 0.0776449282454095
3.1 0.0775139927016079
3.125 0.0773830571578062
3.15 0.0772521216140046
3.175 0.0771211860702029
3.2 0.0769902505264013
3.225 0.0768593149825997
3.25 0.076728379438798
3.275 0.0765974438949964
3.3 0.0764665083511947
3.325 0.0763355728073931
3.35 0.0762046372635915
3.375 0.0760737017197898
3.4 0.0759427661759882
3.425 0.0758118306321865
3.45 0.0756808950883849
3.475 0.0755499595445833
3.5 0.0754190240007816
3.525 0.07528808845698
3.55 0.0751571529131783
3.575 0.0750262173693767
3.6 0.074895281825575
3.625 0.0747643462817734
3.65 0.0746334107379718
3.675 0.0745024751941701
3.7 0.0743715396503685
3.725 0.0742406041065669
3.75 0.0741096685627652
3.775 0.0739787330189636
3.8 0.0738477974751619
3.825 0.0737168619313603
3.85 0.0735859263875586
3.875 0.073454990843757
3.9 0.0733240552999554
3.925 0.0731931197561537
3.95 0.0730621842123521
3.975 0.0729312486685504
4 0.0728003131247488
4.025 0.0726693775809472
4.05 0.0725384420371455
4.075 0.0724075064933439
4.1 0.0722765709495422
4.125 0.0721456354057406
4.15 0.072014699861939
4.175 0.0718837643181373
4.2 0.0717528287743357
4.225 0.071621893230534
4.25 0.0714909576867324
4.275 0.0713600221429308
4.3 0.0712290865991291
4.325 0.0710981510553275
4.35 0.0709672155115258
4.375 0.0708362799677242
4.4 0.0707053444239226
4.425 0.0705744088801209
4.45 0.0704434733363193
4.475 0.0703125377925176
4.5 0.070181602248716
4.525 0.0700506667049143
4.55 0.0699197311611127
4.575 0.0697887956173111
4.6 0.0696578600735094
4.625 0.0695269245297078
4.65 0.0693959889859062
4.675 0.0692650534421045
4.7 0.0691341178983029
4.725 0.0690031823545012
4.75 0.0688722468106996
4.775 0.0687413112668979
4.8 0.0686103757230963
4.825 0.0684794401792947
4.85 0.068348504635493
4.875 0.0682175690916914
4.9 0.0680866335478897
4.925 0.0679556980040881
4.95 0.0678247624602865
4.975 0.0676938269164848
5 0.0675628913726832
5.025 0.0674319558288815
5.05 0.0673010202850799
5.075 0.0671700847412783
5.1 0.0670391491974766
5.125 0.066908213653675
5.15 0.0667772781098733
5.175 0.0666463425660717
5.2 0.0665154070222701
5.225 0.0663844714784684
5.25 0.0662535359346668
5.275 0.0661226003908651
5.3 0.0659916648470635
5.325 0.0658607293032619
5.35 0.0657297937594602
5.375 0.0655988582156586
5.4 0.0654679226718569
5.425 0.0653369871280553
5.45 0.0652060515842537
5.475 0.065075116040452
5.5 0.0649441804966504
5.525 0.0648132449528487
5.55 0.0646823094090471
5.575 0.0645513738652455
5.6 0.0644204383214438
5.625 0.0642895027776422
5.65 0.0641585672338405
5.675 0.0640276316900389
5.7 0.0638966961462372
5.725 0.0637657606024356
5.75 0.063634825058634
5.775 0.0635038895148323
5.8 0.0633729539710307
5.825 0.063242018427229
5.85 0.0631110828834274
5.875 0.0629801473396258
5.9 0.0628492117958241
5.925 0.0627182762520225
5.95 0.0625873407082208
5.975 0.0624564051644192
6 0.0623254696206176
6.025 0.0621945340768159
6.05 0.0620635985330143
6.075 0.0619326629892126
6.1 0.061801727445411
6.125 0.0616707919016094
6.15 0.0615398563578077
6.175 0.0614089208140061
6.2 0.0612779852702044
6.225 0.0611470497264028
6.25 0.0610161141826012
6.275 0.0608851786387995
6.3 0.0607542430949979
6.325 0.0606233075511962
6.35 0.0604923720073946
6.375 0.060361436463593
6.4 0.0602305009197913
6.425 0.0600995653759897
6.45 0.059968629832188
6.475 0.0598376942883864
6.5 0.0597067587445848
6.525 0.0595758232007831
6.55 0.0594448876569815
6.575 0.0593139521131798
6.6 0.0591830165693782
6.625 0.0590520810255766
6.65 0.0589211454817749
6.675 0.0587902099379733
6.7 0.0586592743941716
6.725 0.05852833885037
6.75 0.0583974033065683
6.775 0.0582664677627667
6.8 0.0581355322189651
6.825 0.0580045966751634
6.85 0.0578736611313618
6.875 0.0577427255875601
6.9 0.0576117900437585
6.925 0.0574808544999569
6.95 0.0573499189561552
6.975 0.0572189834123536
7 0.0570880478685519
7.025 0.0569571123247503
7.05 0.0568261767809487
7.075 0.056695241237147
7.1 0.0565643056933454
7.125 0.0564333701495437
7.15 0.0563024346057421
7.175 0.0561714990619405
7.2 0.0560405635181388
7.225 0.0559096279743372
7.25 0.0557786924305355
7.275 0.0556477568867339
7.3 0.0555168213429323
7.325 0.0553858857991306
7.35 0.055254950255329
7.375 0.0551240147115273
7.4 0.0549930791677257
7.425 0.0548621436239241
7.45 0.0547312080801224
7.475 0.0546002725363208
7.5 0.0544693369925191
7.525 0.0543384014487175
7.55 0.0542074659049158
7.575 0.0540765303611142
7.6 0.0539455948173126
7.625 0.0538146592735109
7.65 0.0536837237297093
7.675 0.0535527881859076
7.7 0.053421852642106
7.725 0.0532909170983044
7.75 0.0531599815545027
7.775 0.0530290460107011
7.8 0.0528981104668994
7.825 0.0527671749230978
7.85 0.0526362393792962
7.875 0.0525053038354945
7.9 0.0523743682916929
7.925 0.0522434327478912
7.95 0.0521124972040896
7.975 0.051981561660288
8 0.0518506261164863
8.025 0.0517196905726847
8.05 0.051588755028883
8.075 0.0514578194850814
8.1 0.0513268839412798
8.125 0.0511959483974781
8.15 0.0510650128536765
8.175 0.0509340773098748
8.2 0.0508031417660732
8.225 0.0506722062222716
8.25 0.0505412706784699
8.275 0.0504103351346683
8.3 0.0502793995908666
8.325 0.050148464047065
8.35 0.0500175285032634
8.375 0.0498865929594617
8.4 0.0497556574156601
8.425 0.0496247218718584
8.45 0.0494937863280568
8.475 0.0493628507842552
8.5 0.0492319152404535
8.525 0.0491009796966519
8.55 0.0489700441528502
8.575 0.0488391086090486
8.6 0.048708173065247
8.625 0.0485772375214453
8.65 0.0484463019776437
8.675 0.048315366433842
8.7 0.0481844308900404
8.725 0.0480534953462388
8.75 0.0479225598024371
8.775 0.0477916242586355
8.8 0.0476606887148338
8.825 0.0475297531710322
8.85 0.0473988176272305
8.875 0.0472678820834289
8.9 0.0471369465396273
8.925 0.0470060109958256
8.95 0.046875075452024
8.975 0.0467441399082223
9 0.0466132043644207
9.025 0.0464822688206191
9.05 0.0463513332768174
9.075 0.0462203977330158
9.1 0.0460894621892141
9.125 0.0459585266454125
9.15 0.0458275911016109
9.175 0.0456966555578092
9.2 0.0455657200140076
9.225 0.0454347844702059
9.25 0.0453038489264043
9.275 0.0451729133826027
9.3 0.045041977838801
9.325 0.0449110422949994
9.35 0.0447801067511977
9.375 0.0446491712073961
9.4 0.0445182356635945
9.425 0.0443873001197928
9.45 0.0442563645759912
9.475 0.0441254290321895
9.5 0.0439944934883879
9.525 0.0438635579445863
9.55 0.0437326224007846
9.575 0.043601686856983
9.6 0.0434707513131813
9.625 0.0433398157693797
9.65 0.043208880225578
9.675 0.0430779446817764
9.7 0.0429470091379748
9.725 0.0428160735941731
9.75 0.0426851380503715
9.775 0.0425542025065698
9.8 0.0424232669627682
9.825 0.0422923314189666
9.85 0.0421613958751649
9.875 0.0420304603313633
9.9 0.0418995247875616
9.925 0.04176858924376
9.95 0.0416376536999584
9.975 0.0415067181561567
10 0.0413757826123551
10.025 0.0412448470685534
10.05 0.0411139115247518
10.075 0.0409829759809502
10.1 0.0408520404371485
10.125 0.0407211048933469
10.15 0.0405901693495452
10.175 0.0404592338057436
10.2 0.040328298261942
10.225 0.0401973627181403
10.25 0.0400664271743387
10.275 0.039935491630537
10.3 0.0398045560867354
10.325 0.0396736205429338
10.35 0.0395426849991321
10.375 0.0394117494553305
10.4 0.0392808139115288
10.425 0.0391498783677272
10.45 0.0390189428239256
10.475 0.0388880072801239
10.5 0.0387570717363223
10.525 0.0386261361925206
10.55 0.038495200648719
10.575 0.0383642651049173
10.6 0.0382333295611157
10.625 0.0381023940173141
10.65 0.0379714584735124
10.675 0.0378405229297108
10.7 0.0377095873859091
10.725 0.0375786518421075
10.75 0.0374477162983059
10.775 0.0373167807545042
10.8 0.0371858452107026
10.825 0.0370549096669009
10.85 0.0369239741230993
10.875 0.0367930385792977
10.9 0.036662103035496
10.925 0.0365311674916944
10.95 0.0364002319478927
10.975 0.0362692964040911
11 0.0361383608602895
11.025 0.0360074253164878
11.05 0.0358764897726862
11.075 0.0357455542288845
11.1 0.0356146186850829
11.125 0.0354836831412813
11.15 0.0353527475974796
11.175 0.035221812053678
11.2 0.0350908765098763
11.225 0.0349599409660747
11.25 0.0348290054222731
11.275 0.0346980698784714
11.3 0.0345671343346698
11.325 0.0344361987908681
11.35 0.0343052632470665
11.375 0.0341743277032649
11.4 0.0340433921594632
11.425 0.0339124566156616
11.45 0.0337815210718599
11.475 0.0336505855280583
11.5 0.0335196499842567
11.525 0.033388714440455
11.55 0.0332577788966534
11.575 0.0331268433528517
11.6 0.0329959078090501
11.625 0.0328649722652485
11.65 0.0327340367214468
11.675 0.0326031011776452
11.7 0.0324721656338435
11.725 0.0323412300900419
11.75 0.0322102945462402
11.775 0.0320793590024386
11.8 0.031948423458637
11.825 0.0318174879148353
11.85 0.0316865523710337
11.875 0.031555616827232
11.9 0.0314246812834304
11.925 0.0312937457396288
11.95 0.0311628101958271
11.975 0.0310318746520255
12 0.0309009391082238
12.025 0.0307700035644222
12.05 0.0306390680206206
12.075 0.0305081324768189
12.1 0.0303771969330173
12.125 0.0302462613892156
12.15 0.030115325845414
12.175 0.0299843903016124
12.2 0.0298534547578107
12.225 0.0297225192140091
12.25 0.0295915836702074
12.275 0.0294606481264058
12.3 0.0293297125826042
12.325 0.0291987770388025
12.35 0.0290678414950009
12.375 0.0289369059511992
12.4 0.0288059704073976
12.425 0.0286750348635959
12.45 0.0285440993197943
12.475 0.0284131637759927
12.5 0.028282228232191
12.525 0.0281512926883894
12.55 0.0280203571445878
12.575 0.0278894216007861
12.6 0.0277584860569845
12.625 0.0276275505131828
12.65 0.0274966149693812
12.675 0.0273656794255795
12.7 0.0272347438817779
12.725 0.0271038083379763
12.75 0.0269728727941746
12.775 0.026841937250373
12.8 0.0267110017065713
12.825 0.0265800661627697
12.85 0.0264491306189681
12.875 0.0263181950751664
12.9 0.0261872595313648
12.925 0.0260563239875631
12.95 0.0259253884437615
12.975 0.0257944528999599
13 0.0256635173561582
13.025 0.0255325818123566
13.05 0.0254016462685549
13.075 0.0252707107247533
13.1 0.0251397751809517
13.125 0.02500883963715
13.15 0.0248779040933484
13.175 0.0247469685495467
13.2 0.0246160330057451
13.225 0.0244850974619435
13.25 0.0243541619181418
13.275 0.0242232263743402
13.3 0.0240922908305385
13.325 0.0239613552867369
13.35 0.0238304197429352
13.375 0.0236994841991336
13.4 0.023568548655332
13.425 0.0234376131115303
13.45 0.0233066775677287
13.475 0.0231757420239271
13.5 0.0230448064801254
13.525 0.0229138709363238
13.55 0.0227829353925221
13.575 0.0226519998487205
13.6 0.0225210643049188
13.625 0.0223901287611172
13.65 0.0222591932173156
13.675 0.0221282576735139
13.7 0.0219973221297123
13.725 0.0218663865859106
13.75 0.021735451042109
13.775 0.0216045154983074
13.8 0.0214735799545057
13.825 0.0213426444107041
13.85 0.0212117088669024
13.875 0.0210807733231008
13.9 0.0209498377792992
13.925 0.0208189022354975
13.95 0.0206879666916959
13.975 0.0205570311478942
14 0.0204260956040926
14.025 0.020295160060291
14.05 0.0201642245164893
14.075 0.0200332889726877
14.1 0.019902353428886
14.125 0.0197714178850844
14.15 0.0196404823412828
14.175 0.0195095467974811
14.2 0.0193786112536795
14.225 0.0192476757098778
14.25 0.0191167401660762
14.275 0.0189858046222746
14.3 0.0188548690784729
14.325 0.0187239335346713
14.35 0.0185929979908696
14.375 0.018462062447068
14.4 0.0183311269032664
14.425 0.0182001913594647
14.45 0.0180692558156631
14.475 0.0179383202718614
14.5 0.0178073847280598
14.525 0.0176764491842581
14.55 0.0175455136404565
14.575 0.0174145780966549
14.6 0.0172836425528532
14.625 0.0171527070090516
14.65 0.01702177146525
14.675 0.0168908359214483
14.7 0.0167599003776467
14.725 0.016628964833845
14.75 0.0164980292900434
14.775 0.0163670937462417
14.8 0.0162361582024401
14.825 0.0161052226586385
14.85 0.0159742871148368
14.875 0.0158433515710352
14.9 0.0157124160272335
14.925 0.0155814804834319
14.95 0.0154505449396303
14.975 0.0153196093958286
15 0.015188673852027
15.025 0.0150577383082253
15.05 0.0149268027644237
15.075 0.0147958672206221
15.1 0.0146649316768204
15.125 0.0145339961330188
15.15 0.0144030605892171
15.175 0.0142721250454155
15.2 0.0141411895016139
15.225 0.0140102539578122
15.25 0.0138793184140106
15.275 0.0137483828702089
15.3 0.0136174473264073
15.325 0.0134865117826057
15.35 0.013355576238804
15.375 0.0132246406950024
15.4 0.0130937051512007
15.425 0.0129627696073991
15.45 0.0128318340635974
15.475 0.0127008985197958
15.5 0.0125699629759942
15.525 0.0124390274321925
15.55 0.0123080918883909
15.575 0.0121771563445893
15.6 0.0120462208007876
15.625 0.011915285256986
15.65 0.0117843497131843
15.675 0.0116534141693827
15.7 0.011522478625581
15.725 0.0113915430817794
15.75 0.0112606075379778
15.775 0.0111296719941761
15.8 0.0109987364503745
15.825 0.0108678009065728
15.85 0.0107368653627712
15.875 0.0106059298189696
15.9 0.0104749942751679
15.925 0.0103440587313663
15.95 0.0102131231875646
15.975 0.010082187643763
16 0.00995125209996137
16.025 0.00982031655615971
16.05 0.00968938101235808
16.075 0.00955844546855644
16.1 0.0094275099247548
16.125 0.00929657438095316
16.15 0.00916563883715152
16.175 0.00903470329334988
16.2 0.00890376774954825
16.225 0.00877283220574659
16.25 0.00864189666194497
16.275 0.00851096111814331
16.3 0.00838002557434167
16.325 0.00824909003054004
16.35 0.00811815448673839
16.375 0.00798721894293676
16.4 0.00785628339913511
16.425 0.00772534785533348
16.45 0.00759441231153184
16.475 0.0074634767677302
16.5 0.00733254122392855
16.525 0.0072016056801269
16.55 0.00707067013632527
16.575 0.00693973459252364
16.6 0.00680879904872199
16.625 0.00667786350492035
16.65 0.00654692796111871
16.675 0.00641599241731707
16.7 0.00628505687351544
16.725 0.00615412132971378
16.75 0.00602318578591216
16.775 0.0058922502421105
16.8 0.00576131469830887
16.825 0.00563037915450723
16.85 0.00549944361070558
16.875 0.00536850806690395
16.9 0.00523757252310229
16.925 0.00510663697930067
16.95 0.00497570143549902
16.975 0.00484476589169738
17 0.00700347978626206
17.025 0.00578320420254621
17.05 0.00512201173751309
17.075 0.00501416250648042
17.1 0.00545391662476623
17.125 0.00643553420768833
17.15 0.00795327537056491
17.175 0.0100014002287135
17.2 0.0125741688974521
17.225 0.0156658414920993
17.25 0.019270678127972
17.275 0.0233829389203892
17.3 0.0279968839846676
17.325 0.0331067734361258
17.35 0.0387068673900825
17.375 0.0447914259618539
17.4 0.0513547092667599
17.425 0.0583909774201164
17.45 0.0658944905372423
17.475 0.0738595087334566
17.5 0.0822802921240751
17.525 0.091151100824418
17.55 0.100466194949801
17.575 0.110219834615542
17.6 0.120406279936963
17.625 0.131019791029376
17.65 0.142054628008104
17.675 0.153505050988461
17.7 0.165365320085767
17.725 0.177629695415341
17.75 0.190292437092498
17.775 0.20334780523256
17.8 0.216790059950839
17.825 0.230613461362658
17.85 0.244812269583334
17.875 0.259380744728182
17.9 0.274313146912525
17.925 0.289603736251676
17.95 0.305246772860954
17.975 0.321236516855681
18 0.337567228351169
18.025 0.354233167462741
18.05 0.371228594305711
18.075 0.388547768995398
18.1 0.406184951647123
18.125 0.424134402376199
18.15 0.44239038129795
18.175 0.460947148527686
18.2 0.47979896418073
18.225 0.498940088372402
18.25 0.518364781218014
18.275 0.53806730283289
18.3 0.558041913332342
18.325 0.578282872831691
18.35 0.598784441446258
18.375 0.619540879291355
18.4 0.640546446482305
18.425 0.661795403134421
18.45 0.683282009363023
18.475 0.705000525283432
18.5 0.726945211010961
18.525 0.749110326660933
18.55 0.771490132348661
18.575 0.794078888189464
18.6 0.816870854298666
18.625 0.839860290791575
18.65 0.863041457783518
18.675 0.886408615389805
18.7 0.909956023725758
18.725 0.933677942906698
18.75 0.957568633047937
18.775 0.981622354264798
18.8 1.00583336667259
18.825 1.03019593038664
18.85 1.05470430552227
18.875 1.07935275219479
18.9 1.10413553051952
18.925 1.12904690061177
18.95 1.15408112258686
18.975 1.17923245656013
19 1.20449516264687
19.025 1.22986350096241
19.05 1.25533173162207
19.075 1.28089411474116
19.1 1.30654491043501
19.125 1.33227837881892
19.15 1.35808878000823
19.175 1.38397037411824
19.2 1.40991742126427
19.225 1.43592418156165
19.25 1.46198491512568
19.275 1.4880938820717
19.3 1.51424534251501
19.325 1.54043355657094
19.35 1.56665278435479
19.375 1.5928972859819
19.4 1.61916132156758
19.425 1.64543915122713
19.45 1.6717250350759
19.475 1.69801323322918
19.5 1.72429800580231
19.525 1.75057361291059
19.55 1.77683431466935
19.575 1.8030743711939
19.6 1.82928804259956
19.625 1.85546958900165
19.65 1.88161327051549
19.675 1.90771334725639
19.7 1.93376407933968
19.725 1.95975972688067
19.75 1.98569454999467
19.775 2.01156280879701
19.8 2.03735876340301
19.825 2.06307667392798
19.85 2.08871080048724
19.875 2.1142554031961
19.9 2.1397047421699
19.925 2.16505307752394
19.95 2.19029466937354
19.975 2.21542377783402
20 2.24043473011387
20.025 2.26532339656456
20.05 2.29008719068056
20.075 2.31472359304947
20.1 2.33923008425894
20.125 2.36360414489659
20.15 2.38784325555004
20.175 2.41194489680692
20.2 2.43590654925487
20.225 2.4597256934815
20.25 2.48339981007444
20.275 2.50692637962133
20.3 2.53030288270978
20.325 2.55352679992742
20.35 2.57659561186189
20.375 2.5995067991008
20.4 2.62225784223179
20.425 2.64484622184248
20.45 2.66726941852051
20.475 2.68952491285348
20.5 2.71161018542904
20.525 2.73352271683481
20.55 2.75525998765842
20.575 2.77681947848749
20.6 2.79819866990966
20.625 2.81939504251253
20.65 2.84040607688376
20.675 2.86122925361096
20.7 2.88186205328176
20.725 2.90230195648378
20.75 2.92254644380465
20.775 2.94259299583201
20.8 2.96243909315347
20.825 2.98208221635667
20.85 3.00151984602923
20.875 3.02074946275877
20.9 3.03976854713293
20.925 3.05857457973933
20.95 3.0771650411656
20.975 3.09553741199937
21 3.11368917282826
21.025 3.1316178042399
21.05 3.14932078682191
21.075 3.16679560116193
21.1 3.18403972784758
21.125 3.20105064746649
21.15 3.21782584060628
21.175 3.23436278785458
21.2 3.25065896979903
21.225 3.26671186702724
21.25 3.28251896012684
21.275 3.29807772968546
21.3 3.31338565629072
21.325 3.32844022053026
21.35 3.3432389029917
21.375 3.35777918426267
21.4 3.3720585449308
21.425 3.3860744655837
21.45 3.39982442680901
21.475 3.41330590919436
21.5 3.42651639332737
21.525 3.43945335979567
21.55 3.45211428918689
21.575 3.46449666208865
21.6 3.47659795908858
21.625 3.48841566077431
21.65 3.49994724773346
21.675 3.51119020055366
21.7 3.52214199982255
21.725 3.53280012612773
21.75 3.54316206005685
21.775 3.55322528219753
21.8 3.56298727313739
21.825 3.57244551346407
21.85 3.58159748376519
21.875 3.59044066462837
21.9 3.59897253664125
21.925 3.60719058039145
21.95 3.6150922764666
21.975 3.62267510545433
22 3.62993654794225
22.025 3.63687408451801
22.05 3.64348519576922
22.075 3.64976736228351
22.1 3.65571806464852
22.125 3.66133478345186
22.15 3.66661499928116
22.175 3.67155619272406
22.2 3.67615584436817
22.225 3.68041143480113
22.25 3.68432044461056
22.275 3.68788035438409
22.3 3.69108864470935
22.325 3.69394279617395
22.35 3.69644028936554
22.375 3.69857860487174
22.4 3.70035522328017
22.425 3.70176762517845
22.45 3.70281329115423
22.475 3.70348970179512
22.5 3.70379433768876
22.525 3.70372467942276
22.55 3.70327820758476
22.575 3.70245240276238
22.6 3.70124474554325
22.625 3.699652716515
22.65 3.69767379626525
22.675 3.69530546538163
22.7 3.69254520445177
22.725 3.68939049406329
22.75 3.68583881480383
22.775 3.681887647261
22.8 3.67753447202244
22.825 3.67277676967577
22.85 3.66761202080862
22.875 3.66203770600862
22.9 3.65605130586339
22.925 3.64965030096056
22.95 3.64283217188776
22.975 3.63559439923261
23 3.63312600310578
23.025 3.63242593469094
23.05 3.6317258662761
23.075 3.63102579786126
23.1 3.63032572944642
23.125 3.62962566103158
23.15 3.62892559261674
23.175 3.62822552420191
23.2 3.62752545578707
23.225 3.62682538737223
23.25 3.62612531895739
23.275 3.62542525054255
23.3 3.62472518212771
23.325 3.62402511371287
23.35 3.62332504529803
23.375 3.62262497688319
23.4 3.62192490846836
23.425 3.62122484005352
23.45 3.62052477163868
23.475 3.61982470322384
23.5 3.619124634809
23.525 3.61842456639416
23.55 3.61772449797932
23.575 3.61702442956448
23.6 3.61632436114964
23.625 3.6156242927348
23.65 3.61492422431997
23.675 3.61422415590513
23.7 3.61352408749029
23.725 3.61282401907545
23.75 3.61212395066061
23.775 3.61142388224577
23.8 3.61072381383093
23.825 3.61002374541609
23.85 3.60932367700125
23.875 3.60862360858642
23.9 3.60792354017158
23.925 3.60722347175674
23.95 3.6065234033419
23.975 3.60582333492706
24 3.60512326651222
24.025 3.60442319809738
24.05 3.60372312968254
24.075 3.6030230612677
24.1 3.60232299285286
24.125 3.60162292443803
24.15 3.60092285602319
24.175 3.60022278760835
24.2 3.59952271919351
24.225 3.59882265077867
24.25 3.59812258236383
24.275 3.59742251394899
24.3 3.59672244553415
24.325 3.59602237711931
24.35 3.59532230870447
24.375 3.59462224028964
24.4 3.5939221718748
24.425 3.59322210345996
24.45 3.59252203504512
24.475 3.59182196663028
24.5 3.59112189821544
24.525 3.5904218298006
24.55 3.58972176138576
24.575 3.58902169297092
24.6 3.58832162455609
24.625 3.58762155614125
24.65 3.58692148772641
24.675 3.58622141931157
24.7 3.58552135089673
24.725 3.58482128248189
24.75 3.58412121406705
24.775 3.58342114565221
24.8 3.58272107723737
24.825 3.58202100882253
24.85 3.5813209404077
24.875 3.58062087199286
24.9 3.57992080357802
24.925 3.57922073516318
24.95 3.57852066674834
24.975 3.5778205983335
25 3.57712052991866
25.025 3.57642046150382
25.05 3.57572039308898
25.075 3.57502032467414
25.1 3.57432025625931
25.125 3.57362018784447
25.15 3.57292011942963
25.175 3.57222005101479
25.2 3.57151998259995
25.225 3.57081991418511
25.25 3.57011984577027
25.275 3.56941977735543
25.3 3.56871970894059
25.325 3.56801964052576
25.35 3.56731957211092
25.375 3.56661950369608
25.4 3.56591943528124
25.425 3.5652193668664
25.45 3.56451929845156
25.475 3.56381923003672
25.5 3.56311916162188
25.525 3.56241909320704
25.55 3.5617190247922
25.575 3.56101895637737
25.6 3.56031888796253
25.625 3.55961881954769
25.65 3.55891875113285
25.675 3.55821868271801
25.7 3.55751861430317
25.725 3.55681854588833
25.75 3.55611847747349
25.775 3.55541840905865
25.8 3.55471834064382
25.825 3.55401827222898
25.85 3.55331820381414
25.875 3.5526181353993
25.9 3.55191806698446
25.925 3.55121799856962
25.95 3.55051793015478
25.975 3.54981786173994
26 3.5491177933251
26.025 3.54841772491026
26.05 3.54771765649543
26.075 3.54701758808059
26.1 3.54631751966575
26.125 3.54561745125091
26.15 3.54491738283607
26.175 3.54421731442123
26.2 3.54351724600639
26.225 3.54281717759155
26.25 3.54211710917671
26.275 3.54141704076187
26.3 3.54071697234704
26.325 3.5400169039322
26.35 3.53931683551736
26.375 3.53861676710252
26.4 3.53791669868768
26.425 3.53721663027284
26.45 3.536516561858
26.475 3.53581649344316
26.5 3.53511642502832
26.525 3.53441635661349
26.55 3.53371628819865
26.575 3.53301621978381
26.6 3.53231615136897
26.625 3.53161608295413
26.65 3.53091601453929
26.675 3.53021594612445
26.7 3.52951587770961
26.725 3.52881580929477
26.75 3.52811574087993
26.775 3.5274156724651
26.8 3.52671560405026
26.825 3.52601553563542
26.85 3.52531546722058
26.875 3.52461539880574
26.9 3.5239153303909
26.925 3.52321526197606
26.95 3.52251519356122
26.975 3.52181512514638
27 3.52111505673155
27.025 3.52041498831671
27.05 3.51971491990187
27.075 3.51901485148703
27.1 3.51831478307219
27.125 3.51761471465735
27.15 3.51691464624251
27.175 3.51621457782767
27.2 3.51551450941283
27.225 3.51481444099799
27.25 3.51411437258316
27.275 3.51341430416832
27.3 3.51271423575348
27.325 3.51201416733864
27.35 3.5113140989238
27.375 3.51061403050896
27.4 3.50991396209412
27.425 3.50921389367928
27.45 3.50851382526444
27.475 3.5078137568496
27.5 3.50711368843477
27.525 3.50641362001993
27.55 3.50571355160509
27.575 3.50501348319025
27.6 3.50431341477541
27.625 3.50361334636057
27.65 3.50291327794573
27.675 3.50221320953089
27.7 3.50151314111605
27.725 3.50081307270122
27.75 3.50011300428638
27.775 3.49941293587154
27.8 3.4987128674567
27.825 3.49801279904186
27.85 3.49731273062702
27.875 3.49661266221218
27.9 3.49591259379734
27.925 3.4952125253825
27.95 3.49451245696766
27.975 3.49381238855283
28 3.49311232013799
28.025 3.49241225172315
28.05 3.49171218330831
28.075 3.49101211489347
28.1 3.49031204647863
28.125 3.48961197806379
28.15 3.48891190964895
28.175 3.48821184123411
28.2 3.48751177281928
28.225 3.48681170440444
28.25 3.4861116359896
28.275 3.48541156757476
28.3 3.48471149915992
28.325 3.48401143074508
28.35 3.48331136233024
28.375 3.4826112939154
28.4 3.48191122550056
28.425 3.48121115708572
28.45 3.48051108867089
28.475 3.47981102025605
28.5 3.47911095184121
28.525 3.47841088342637
28.55 3.47771081501153
28.575 3.47701074659669
28.6 3.47631067818185
28.625 3.47561060976701
28.65 3.47491054135217
28.675 3.47421047293733
28.7 3.4735104045225
28.725 3.47281033610766
28.75 3.47211026769282
28.775 3.47141019927798
28.8 3.47071013086314
28.825 3.4700100624483
28.85 3.46930999403346
28.875 3.46860992561862
28.9 3.46790985720378
28.925 3.46720978878895
28.95 3.46650972037411
28.975 3.46580965195927
29 3.46510958354443
29.025 3.46440951512959
29.05 3.46370944671475
29.075 3.46300937829991
29.1 3.46230930988507
29.125 3.46160924147023
29.15 3.46090917305539
29.175 3.46020910464056
29.2 3.45950903622572
29.225 3.45880896781088
29.25 3.45810889939604
29.275 3.4574088309812
29.3 3.45670876256636
29.325 3.45600869415152
29.35 3.45530862573668
29.375 3.45460855732184
29.4 3.453908488907
29.425 3.45320842049217
29.45 3.45250835207733
29.475 3.45180828366249
29.5 3.45110821524765
29.525 3.45040814683281
29.55 3.44970807841797
29.575 3.44900801000313
29.6 3.44830794158829
29.625 3.44760787317345
29.65 3.44690780475862
29.675 3.44620773634378
29.7 3.44550766792894
29.725 3.4448075995141
29.75 3.44410753109926
29.775 3.44340746268442
29.8 3.44270739426958
29.825 3.44200732585474
29.85 3.4413072574399
29.875 3.44060718902506
29.9 3.43990712061023
29.925 3.43920705219539
29.95 3.43850698378055
29.975 3.43780691536571
30 3.43710684695087
30.025 3.43640677853603
30.05 3.43570671012119
30.075 3.43500664170635
30.1 3.43430657329151
30.125 3.43360650487668
30.15 3.43290643646184
30.175 3.432206368047
30.2 3.43150629963216
30.225 3.43080623121732
30.25 3.43010616280248
30.275 3.42940609438764
30.3 3.4287060259728
30.325 3.42800595755796
30.35 3.42730588914312
30.375 3.42660582072829
30.4 3.42590575231345
30.425 3.42520568389861
30.45 3.42450561548377
30.475 3.42380554706893
30.5 3.42310547865409
30.525 3.42240541023925
30.55 3.42170534182441
30.575 3.42100527340957
30.6 3.42030520499473
30.625 3.4196051365799
30.65 3.41890506816506
30.675 3.41820499975022
30.7 3.41750493133538
30.725 3.41680486292054
30.75 3.4161047945057
30.775 3.41540472609086
30.8 3.41470465767602
30.825 3.41400458926118
30.85 3.41330452084635
30.875 3.41260445243151
30.9 3.41190438401667
30.925 3.41120431560183
30.95 3.41050424718699
30.975 3.40980417877215
31 3.40910411035731
31.025 3.40840404194247
31.05 3.40770397352763
31.075 3.40700390511279
31.1 3.40630383669796
31.125 3.40560376828312
31.15 3.40490369986828
31.175 3.40420363145344
31.2 3.4035035630386
31.225 3.40280349462376
31.25 3.39846889378193
31.275 3.39513073016018
31.3 3.39130807823856
31.325 3.38700645866187
31.35 3.38223139207489
31.375 3.37698839912243
31.4 3.37128300044927
31.425 3.36512071670021
31.45 3.35850706852006
31.475 3.35144757655359
31.5 3.34394776144561
31.525 3.33601314384091
31.55 3.32764924438429
31.575 3.31886158372054
31.6 3.30965568249445
31.625 3.30003706135082
31.65 3.29001124093445
31.675 3.27958374189013
31.7 3.26876008486265
31.725 3.25754579049681
31.75 3.2459463794374
31.775 3.23396737232922
31.8 3.22161428981706
31.825 3.20889265254572
31.85 3.19580798116
31.875 3.18236579630468
31.9 3.16857161862456
31.925 3.15443096876444
31.95 3.13994936736911
31.975 3.12513233508336
32 3.109985392552
32.025 3.09451406041981
32.05 3.07872385933159
32.075 3.06262030993214
32.1 3.04620893286624
32.125 3.0294952487787
32.15 3.01248477831431
32.175 2.99518304211786
32.2 2.97759556083415
32.225 2.95972785510797
32.25 2.94158544558413
32.275 2.9231738529074
32.3 2.90449859772258
32.325 2.88556520067449
32.35 2.8663791824079
32.375 2.84694606356761
32.4 2.82727136479841
32.425 2.8073606067451
32.45 2.78721931005249
32.475 2.76685299536535
32.5 2.74626718332849
32.525 2.72546739458669
32.55 2.70445914978476
32.575 2.68324796956749
32.6 2.66183937457967
32.625 2.6402388854661
32.65 2.61845202287157
32.675 2.59648430744088
32.7 2.57434125981882
32.725 2.5520284006502
32.75 2.52955125057979
32.775 2.5069153302524
32.8 2.48412616031281
32.825 2.46118926140584
32.85 2.43811015417627
32.875 2.41489435926889
32.9 2.39154739732851
32.925 2.3680747889999
32.95 2.34448205492788
32.975 2.32077471575723
33 2.29695829213275
33.025 2.27303830469923
33.05 2.24902027410147
33.075 2.22490972098427
33.1 2.20071216599241
33.125 2.1764331297707
33.15 2.15207813296392
33.175 2.12765269621687
33.2 2.10316234017435
33.225 2.07861258548115
33.25 2.05400895278207
33.275 2.0293569627219
33.3 2.00466213594543
33.325 1.97992999309747
33.35 1.9551660548228
33.375 1.93037584176622
33.4 1.90556487457252
33.425 1.8807386738865
33.45 1.85590276035296
33.475 1.83106265461668
33.5 1.80622387732247
33.525 1.78139194911512
33.55 1.75657239063941
33.575 1.73177072254016
33.6 1.70699246546215
33.625 1.68224314005017
33.65 1.65752826694903
33.675 1.63285336680351
33.7 1.60822396025841
33.725 1.58364556795853
33.75 1.5591236754774
33.775 1.53466296174964
33.8 1.51026729907093
33.825 1.48594052466572
33.85 1.46168647575843
33.875 1.43750898957349
33.9 1.41341190333535
33.925 1.38939905426842
33.95 1.36547427959715
33.975 1.34164141654598
34 1.31790430233932
34.025 1.29426677420162
34.05 1.2707326693573
34.075 1.24730582503081
34.1 1.22399007844657
34.125 1.20078926682902
34.15 1.17770722740259
34.175 1.15474779739171
34.2 1.13191481402081
34.225 1.10921211451434
34.25 1.08664353609673
34.275 1.06421291599239
34.3 1.04192409142577
34.325 1.01978089962131
34.35 0.997787177803433
34.375 0.975946763196573
34.4 0.954263493025165
34.425 0.932741204513635
34.45 0.911383734886428
34.475 0.890194921367971
34.5 0.869178601182698
34.525 0.848338611555041
34.55 0.827678789709427
34.575 0.807202972870303
34.6 0.786914998262093
34.625 0.766818703109233
34.65 0.746917924636155
34.675 0.727216500067285
34.7 0.70771826662707
34.725 0.688427061539935
34.75 0.669346722030314
34.775 0.650481085322641
34.8 0.631833988641342
34.825 0.613409269210862
34.85 0.595210764255629
34.875 0.577242311000075
34.9 0.559507746668633
34.925 0.542010908485732
34.95 0.524755633675816
34.975 0.507745759463311
35 0.49098512307265
35.025 0.474477561728267
35.05 0.458226912654592
35.075 0.442237013076064
35.1 0.426511700217115
35.125 0.411054811302174
35.15 0.395870183555678
35.175 0.380961654202054
35.2 0.366333060465744
35.225 0.351988239571176
35.25 0.337931028742783
35.275 0.324165265204999
35.3 0.310694786182253
35.325 0.297523428898986
35.35 0.284655030579626
35.375 0.272093428448609
35.4 0.259842459730364
35.425 0.247905961649325
35.45 0.236287771429928
35.475 0.224991726296605
35.5 0.214021663473789
35.525 0.203381420185911
35.55 0.193074833657403
35.575 0.183105741112704
35.6 0.173477979776244
35.625 0.164195386872455
35.65 0.155261799625771
35.675 0.146681055260623
35.7 0.138456991001448
35.725 0.130593444072677
35.75 0.123094251698743
35.775 0.11596325110408
35.8 0.109204279513118
35.825 0.102821174150295
35.85 0.0968177722400401
35.875 0.0911979110067889
35.9 0.0859654276749726
35.925 0.0811241594690237
35.95 0.0766779436133787
35.975 0.0726306173324662
36 0.0689860178507242
36.025 0.0657479823925828
36.05 0.0629203481824741
36.075 0.0605069524448343
36.1 0.0585116324040937
36.125 0.0569382252846871
36.15 0.0557905683110462
36.175 0.0550724987076046
36.2 0.0547878536987958
36.225 0.0549404705090528
36.25 0.0544659123367125
36.275 0.0543200035081735
36.3 0.0541740946796345
36.325 0.0540281858510955
36.35 0.0538822770225564
36.375 0.0537363681940174
36.4 0.0535904593654784
36.425 0.0534445505369394
36.45 0.0532986417084004
36.475 0.0531527328798614
36.5 0.0530068240513224
36.525 0.0528609152227834
36.55 0.0527150063942443
36.575 0.0525690975657053
36.6 0.0524231887371663
36.625 0.0522772799086273
36.65 0.0521313710800883
36.675 0.0519854622515492
36.7 0.0518395534230102
36.725 0.0516936445944712
36.75 0.0515477357659322
36.775 0.0514018269373932
36.8 0.0512559181088542
36.825 0.0511100092803152
36.85 0.0509641004517762
36.875 0.0508181916232371
36.9 0.0506722827946981
36.925 0.0505263739661591
36.95 0.0503804651376201
36.975 0.0502345563090811
37 0.0500886474805421
37.025 0.0499427386520031
37.05 0.049796829823464
37.075 0.049650920994925
37.1 0.049505012166386
37.125 0.049359103337847
37.15 0.049213194509308
37.175 0.049067285680769
37.2 0.0489213768522299
37.225 0.0487754680236909
37.25 0.0486295591951519
37.275 0.0484836503666129
37.3 0.0483377415380739
37.325 0.0481918327095349
37.35 0.0480459238809959
37.375 0.0479000150524569
37.4 0.0477541062239179
37.425 0.0476081973953788
37.45 0.0474622885668398
37.475 0.0473163797383008
37.5 0.0471704709097618
37.525 0.0470245620812228
37.55 0.0468786532526837
37.575 0.0467327444241447
37.6 0.0465868355956057
37.625 0.0464409267670667
37.65 0.0462950179385277
37.675 0.0461491091099887
37.7 0.0460032002814497
37.725 0.0458572914529106
37.75 0.0457113826243716
37.775 0.0455654737958326
37.8 0.0454195649672936
37.825 0.0452736561387546
37.85 0.0451277473102156
37.875 0.0449818384816766
37.9 0.0448359296531376
37.925 0.0446900208245985
37.95 0.0445441119960595
37.975 0.0443982031675205
38 0.0442522943389815
38.025 0.0441063855104425
38.05 0.0439604766819034
38.075 0.0438145678533644
38.1 0.0436686590248254
38.125 0.0435227501962864
38.15 0.0433768413677474
38.175 0.0432309325392084
38.2 0.0430850237106694
38.225 0.0429391148821304
38.25 0.0427932060535914
38.275 0.0426472972250523
38.3 0.0425013883965133
38.325 0.0423554795679743
38.35 0.0422095707394353
38.375 0.0420636619108963
38.4 0.0419177530823572
38.425 0.0417718442538182
38.45 0.0416259354252792
38.475 0.0414800265967402
38.5 0.0413341177682012
38.525 0.0411882089396622
38.55 0.0410423001111232
38.575 0.0408963912825841
38.6 0.0407504824540451
38.625 0.0406045736255061
38.65 0.0404586647969671
38.675 0.0403127559684281
38.7 0.0401668471398891
38.725 0.0400209383113501
38.75 0.0398750294828111
38.775 0.039729120654272
38.8 0.039583211825733
38.825 0.039437302997194
38.85 0.039291394168655
38.875 0.039145485340116
38.9 0.0389995765115769
38.925 0.0388536676830379
38.95 0.0387077588544989
38.975 0.0385618500259599
39 0.0384159411974209
39.025 0.0382700323688819
39.05 0.0381241235403429
39.075 0.0379782147118039
39.1 0.0378323058832649
39.125 0.0376863970547258
39.15 0.0375404882261868
39.175 0.0373945793976478
39.2 0.0372486705691088
39.225 0.0371027617405698
39.25 0.0369568529120308
39.275 0.0368109440834917
39.3 0.0366650352549527
39.325 0.0365191264264137
39.35 0.0363732175978747
39.375 0.0362273087693357
39.4 0.0360813999407967
39.425 0.0359354911122576
39.45 0.0357895822837186
39.475 0.0356436734551796
39.5 0.0354977646266406
39.525 0.0353518557981016
39.55 0.0352059469695626
39.575 0.0350600381410236
39.6 0.0349141293124846
39.625 0.0347682204839456
39.65 0.0346223116554065
39.675 0.0344764028268675
39.7 0.0343304939983285
39.725 0.0341845851697895
39.75 0.0340386763412505
39.775 0.0338927675127114
39.8 0.0337468586841724
39.825 0.0336009498556334
39.85 0.0334550410270944
39.875 0.0333091321985554
39.9 0.0331632233700164
39.925 0.0330173145414774
39.95 0.0328714057129384
39.975 0.0327254968843994
40 0.0325795880558603
40.025 0.0324336792273213
40.05 0.0322877703987823
40.075 0.0321418615702433
40.1 0.0319959527417043
40.125 0.0318500439131653
40.15 0.0317041350846262
40.175 0.0315582262560872
40.2 0.0314123174275482
40.225 0.0312664085990092
40.25 0.0311204997704702
40.275 0.0309745909419312
40.3 0.0308286821133921
40.325 0.0306827732848531
40.35 0.0305368644563141
40.375 0.0303909556277751
40.4 0.0302450467992361
40.425 0.0300991379706971
40.45 0.0299532291421581
40.475 0.0298073203136191
40.5 0.0296614114850801
40.525 0.029515502656541
40.55 0.029369593828002
40.575 0.029223684999463
40.6 0.029077776170924
40.625 0.028931867342385
40.65 0.0287859585138459
40.675 0.0286400496853069
40.7 0.0284941408567679
40.725 0.0283482320282289
40.75 0.0282023231996899
40.775 0.0280564143711509
40.8 0.0279105055426119
40.825 0.0277645967140728
40.85 0.0276186878855338
40.875 0.0274727790569948
40.9 0.0273268702284558
40.925 0.0271809613999168
40.95 0.0270350525713778
40.975 0.0268891437428388
41 0.0267432349142998
41.025 0.0265973260857607
41.05 0.0264514172572217
41.075 0.0263055084286827
41.1 0.0261595996001437
41.125 0.0260136907716047
41.15 0.0258677819430656
41.175 0.0257218731145266
41.2 0.0255759642859876
41.225 0.0254300554574486
41.25 0.0252841466289096
41.275 0.0251382378003706
41.3 0.0249923289718316
41.325 0.0248464201432926
41.35 0.0247005113147536
41.375 0.0245546024862145
41.4 0.0244086936576755
41.425 0.0242627848291365
41.45 0.0241168760005975
41.475 0.0239709671720585
41.5 0.0238250583435195
41.525 0.0236791495149804
41.55 0.0235332406864414
41.575 0.0233873318579024
41.6 0.0232414230293634
41.625 0.0230955142008244
41.65 0.0229496053722854
41.675 0.0228036965437464
41.7 0.0226577877152073
41.725 0.0225118788866683
41.75 0.0223659700581293
41.775 0.0222200612295903
41.8 0.0220741524010513
41.825 0.0219282435725123
41.85 0.0217823347439733
41.875 0.0216364259154343
41.9 0.0214905170868952
41.925 0.0213446082583562
41.95 0.0211986994298172
41.975 0.0210527906012782
42 0.0209068817727392
42.025 0.0207609729442001
42.05 0.0206150641156611
42.075 0.0204691552871221
42.1 0.0203232464585831
42.125 0.0201773376300441
42.15 0.0200314288015051
42.175 0.0198855199729661
42.2 0.0197396111444271
42.225 0.019593702315888
42.25 0.019447793487349
42.275 0.01930188465881
42.3 0.019155975830271
42.325 0.019010067001732
};
\addplot [ultra thick, black, forget plot]
table {%
17 -0.2
17 3.8
};
\addplot [ultra thick, black, forget plot]
table {%
23 -0.2
23 3.8
};
\addplot [ultra thick, black, forget plot]
table {%
31.25 -0.2
31.25 3.8
};
\addplot [ultra thick, black, forget plot]
table {%
36.25 -0.2
36.25 3.8
};
\node at (axis cs:8.5,0.3)[
  anchor=base,
  text=black,
  rotate=0.0
]{ F};
\node at (axis cs:20,0.3)[
  anchor=base,
  text=black,
  rotate=0.0
]{ G};
\node at (axis cs:27.125,0.3)[
  anchor=base,
  text=black,
  rotate=0.0
]{ H};
\node at (axis cs:33.75,0.3)[
  anchor=base,
  text=black,
  rotate=0.0
]{ I};
\node at (axis cs:38,0.3)[
  anchor=base,
  text=black,
  rotate=0.0
]{ J};
\end{axis}

\end{tikzpicture}%
		\subcaption{Lateral position ego vehicle. During the time intervals F, H and J, the ego vehicles drives straight ahead while during the time intervals G and I, the ego vehicle performs a lane change.}
		\label{fig:example ego states lat}
	\end{minipage}%
	\caption{Activities of the ego vehicle. The blue line shows the measured data and the orange dashed line shows the output of the model fitted to the data. The black vertical lines indicate the events.}
	\label{fig:example ego states}
\end{figure}

\begin{figure}
	\centering
	\setlength\figureheight{170pt}
	\setlength\figurewidth{\linewidth}
	% This file was created by matplotlib2tikz v0.6.14.
\begin{tikzpicture}

\begin{axis}[
xlabel={$t$ [s]},
xmin=-7, xmax=42.325,
ymin=0, ymax=6,
width=\figurewidth,
height=\figureheight,
tick align=outside,
tick pos=left,
x grid style={white!69.019607843137251!black},
clip marker paths,
ytick=\empty,
xtick={0,5,10,15,20,25,30,35,40},
y axis line style={draw opacity=0}
]
\path [draw=white!80.0!black, fill=white!80.0!black] (axis cs:0,0)
--(axis cs:42.325,0)
--(axis cs:42.325,1)
--(axis cs:0,1)
--cycle;

\path [draw=white!70.0!black, fill=white!70.0!black] (axis cs:0,1)
--(axis cs:42.325,1)
--(axis cs:42.325,2)
--(axis cs:0,2)
--cycle;

\path [draw=white!80.0!black, fill=white!80.0!black] (axis cs:0,2)
--(axis cs:17.7200000286102,2)
--(axis cs:17.7200000286102,3)
--(axis cs:0,3)
--cycle;

\path [draw=white!70.0!black, fill=white!70.0!black] (axis cs:0,3)
--(axis cs:17.7200000286102,3)
--(axis cs:17.7200000286102,4)
--(axis cs:0,4)
--cycle;

\path [draw=white!80.0!black, fill=white!80.0!black] (axis cs:0.800000190734863,4)
--(axis cs:34.1600000858307,4)
--(axis cs:34.1600000858307,5)
--(axis cs:0.800000190734863,5)
--cycle;

\path [draw=white!70.0!black, fill=white!70.0!black] (axis cs:0.800000190734863,5)
--(axis cs:34.1600000858307,5)
--(axis cs:34.1600000858307,6)
--(axis cs:0.800000190734863,6)
--cycle;

\addplot [semithick, black, forget plot]
table {%
0 0
0 1
};
\addplot [semithick, black, forget plot]
table {%
9 0
9 1
};
\addplot [semithick, black, forget plot]
table {%
11 0
11 1
};
\addplot [semithick, black, forget plot]
table {%
18 0
18 1
};
\addplot [semithick, black, forget plot]
table {%
23 0
23 1
};
\addplot [semithick, black, forget plot]
table {%
0 1
0 2
};
\addplot [semithick, black, forget plot]
table {%
17 1
17 2
};
\addplot [semithick, black, forget plot]
table {%
23 1
23 2
};
\addplot [semithick, black, forget plot]
table {%
31.25 1
31.25 2
};
\addplot [semithick, black, forget plot]
table {%
36.25 1
36.25 2
};
\addplot [semithick, black, forget plot]
table {%
42.325 0
42.325 2
};
\addplot [semithick, black, forget plot]
table {%
0 2
0 3
};
\addplot [semithick, black, forget plot]
table {%
0 3
0 4
};
\addplot [semithick, black, forget plot]
table {%
17.7200000286102 2
17.7200000286102 3
};
\addplot [semithick, black, forget plot]
table {%
17.7200000286102 3
17.7200000286102 4
};
\addplot [semithick, black, forget plot]
table {%
13.2000000476837 2
13.2000000476837 3
};
\addplot [semithick, black, forget plot]
table {%
9.20000004768372 3
9.20000004768372 4
};
\addplot [semithick, black, forget plot]
table {%
11.8000001907349 3
11.8000001907349 4
};
\addplot [semithick, black, forget plot]
table {%
0.800000190734863 4
0.800000190734863 5
};
\addplot [semithick, black, forget plot]
table {%
0.800000190734863 5
0.800000190734863 6
};
\addplot [semithick, black, forget plot]
table {%
34.1600000858307 4
34.1600000858307 5
};
\addplot [semithick, black, forget plot]
table {%
34.1600000858307 5
34.1600000858307 6
};
\addplot [semithick, black, forget plot]
table {%
0 1
42.325 1
};
\addplot [semithick, black, forget plot]
table {%
0 2
42.325 2
};
\addplot [semithick, black, forget plot]
table {%
0 3
42.325 3
};
\addplot [semithick, black, forget plot]
table {%
0 4
42.325 4
};
\addplot [semithick, black, forget plot]
table {%
0 5
42.325 5
};
\node at (axis cs:4.4875,0.5)[
  scale=0.75,
  text=black,
  rotate=0.0
]{ Accelerating};
\node at (axis cs:9.9875,0.5)[
  scale=0.75,
  text=black,
  rotate=0.0,
  align=center
]{ Bra-\\
king};
\node at (axis cs:14.4875,0.5)[
  scale=0.75,
  text=black,
  rotate=0.0
]{ Cruising};
\node at (axis cs:20.4875,0.5)[
  scale=0.75,
  text=black,
  rotate=0.0
]{ Accelerating};
\node at (axis cs:32.6625,0.5)[
  scale=0.75,
  text=black,
  rotate=0.0
]{ Cruising};
\node at (axis cs:8.4875,1.5)[
  scale=0.75,
  text=black,
  rotate=0.0
]{ Straight};
\node at (axis cs:19.9875,1.5)[
  scale=0.75,
  text=black,
  rotate=0.0,
  align=center
]{ Lane\\
Change};
\node at (axis cs:27.1125,1.5)[
  scale=0.75,
  text=black,
  rotate=0.0
]{ Straight};
\node at (axis cs:33.7375,1.5)[
  scale=0.75,
  text=black,
  rotate=0.0,
  align=center
]{ Lane\\
Change};
\node at (axis cs:39.2875,1.5)[
  scale=0.75,
  text=black,
  rotate=0.0
]{ Straight};
\node at (axis cs:6.60000002384186,2.5)[
  scale=0.75,
  text=black,
  rotate=0.0
]{ Accelerating};
\node at (axis cs:15.460000038147,2.5)[
  scale=0.75,
  text=black,
  rotate=0.0
]{ Cruising};
\node at (axis cs:4.60000002384186,3.5)[
  scale=0.75,
  text=black,
  rotate=0.0
]{ Straight};
\node at (axis cs:10.5000001192093,3.5)[
  scale=0.75,
  text=black,
  rotate=0.0,
  align=center
]{ Lane\\
Change};
\node at (axis cs:14.7600001096725,3.5)[
  scale=0.75,
  text=black,
  rotate=0.0
]{ Straight};
\node at (axis cs:17.4800001382828,4.5)[
  scale=0.75,
  text=black,
  rotate=0.0
]{ Cruising};
\node at (axis cs:17.4800001382828,5.5)[
  scale=0.75,
  text=black,
  rotate=0.0
]{ Straight};
\node at (axis cs:-3.5,0.5)[
  scale=0.75,
  text=black,
  rotate=0.0,
  align=center
]{ Ego vehicle,\\
Longitudinal state};
\node at (axis cs:-3.5,1.5)[
  scale=0.75,
  text=black,
  rotate=0.0,
  align=center
]{ Ego vehicle,\\
Lateral state};
\node at (axis cs:-3.5,2.5)[
  scale=0.75,
  text=black,
  rotate=0.0,
  align=center
]{ Station wagon,\\
Longitudinal state};
\node at (axis cs:-3.5,3.5)[
  scale=0.75,
  text=black,
  rotate=0.0,
  align=center
]{ Station wagon,\\
Lateral state};
\node at (axis cs:-3.5,4.5)[
  scale=0.75,
  text=black,
  rotate=0.0,
  align=center
]{ Pickup truck,\\
Longitudinal state};
\node at (axis cs:-3.5,5.5)[
  scale=0.75,
  text=black,
  rotate=0.0,
  align=center
]{ Pickup truck,\\
Lateral state};
\end{axis}

\end{tikzpicture}%
	\caption{Complete overview of the events and activities of the example. Every vertical black line represents an event. The end of the 'ego gap closing' scenario and the start of the 'ego overtaking 1' scenario is at the moment the ego vehicle starts the first lane change, i.e., at $t=17\ \textup{s}$ as indicated by the red line.}
	\label{fig:example events}
\end{figure}

For parametrization of the acceleration, braking, and lane change activities, splines are employed. First, the position is scaled such that it starts at $0$ and ends at $1$. For this, the starting position $x_{\textup{start}}$ and end position $x_{\textup{end}}$ are used. Next, the time is scaled, such that the scaled time is on the interval $[0, 1]$. For this, the start time $t_{\textup{start}}$ and duration $t_{\textup{duration}}$ are used. The scaled position is modeled using the splines. In total, two splines are used. The first spline for the interval $[0, 0.5]$ and the second spline for the interval $[0.5, 1]$. Each spline takes the following form:

\begin{equation}
	x = at^3 + bt^2 + ct + d,
\end{equation}
where $a$, $b$, $c$, and $d$ are the spline coefficients and $x$ and $t$ refer to the scaled position and time, respectively. 

For the cruising and straight-driving activities, a constant velocity is assumed. Hence, only the normalization parameters are enough to model these activities. The parameters for the acceleration activities A and C are presented in \cref{tab:parameters ego A} and \cref{tab:parameters ego D}, respectively. The parameters for the braking activity are presented in \cref{tab:parameters ego B}. The parameters for the cruising activities are presented in \cref{tab:parameters ego cruising}.

\begin{table}
	\centering
	\caption{Parameters of the first acceleration activity of the ego vehicle, denoted by A in \cref{fig:example ego states lon}.}
	\label{tab:parameters ego A}
	\begin{minipage}[b]{.33\linewidth}
		\centering
		\subcaption{Normalization parameters.}
		\begin{tabular}{cr}
    \toprule
    Parameter & Value \\ \otoprule
    $x_{\textup{start}}$    & $\unit[0.00]{m}$ \\
    $x_{\textup{end}}$      & $\unit[1.33 \cdot 10^{2}]{m}$ \\
    $t_{\textup{start}}$    & $\unit[0.00]{s}$ \\
    $t_{\textup{duration}}$ & $\unit[8.97]{s}$ \\
    \bottomrule
\end{tabular}
%
	\end{minipage}%
	\begin{minipage}[b]{.5\linewidth}
		\centering
		\subcaption{Spline coefficients.}
		\begin{tabular}{crr}
    \toprule
    Coefficients & Spline 1 & Spline 2 \\ \otoprule
    $a$ & $-7.55 \cdot 10^{-2}$ & $-2.18 \cdot 10^{-2}$ \\
    $b$ & $3.62 \cdot 10^{-1}$ & $2.82 \cdot 10^{-1}$ \\
    $c$ & $7.07 \cdot 10^{-1}$ & $7.47 \cdot 10^{-1}$ \\
    $d$ & $3.74 \cdot 10^{-16}$ & $-6.71 \cdot 10^{-3}$ \\
    \bottomrule
\end{tabular}
%
	\end{minipage}%
\end{table}

\begin{table}
	\centering
	\caption{Parameters of the braking activity of the ego vehicle, denoted by B in \cref{fig:example ego states lon}.}
	\label{tab:parameters ego B}
	\begin{minipage}[b]{.33\linewidth}
		\centering
		\subcaption{Normalization parameters.}
		\begin{tabular}{cr}
    \toprule
    Parameter & Value \\ \otoprule
    $x_{\textup{start}}$    & $\unit[1.33 \cdot 10^{2}]{m}$ \\
    $x_{\textup{end}}$      & $\unit[1.68 \cdot 10^{2}]{m}$ \\
    $t_{\textup{start}}$    & $\unit[9.00]{s}$ \\
    $t_{\textup{duration}}$ & $\unit[1.98]{s}$ \\
    \bottomrule
\end{tabular}
%
	\end{minipage}%
	\begin{minipage}[b]{.5\linewidth}
		\centering
		\subcaption{Spline coefficients.}
		\begin{tabular}{crr}
    \toprule
    Coefficients & Spline 1 & Spline 2 \\ \otoprule
    $a$ & $1.56 \cdot 10^{-2}$ & $3.31 \cdot 10^{-2}$ \\
    $b$ & $-6.27 \cdot 10^{-2}$ & $-8.89 \cdot 10^{-2}$ \\
    $c$ & $1.04$ & $1.06$ \\
    $d$ & $3.96 \cdot 10^{-16}$ & $-2.18 \cdot 10^{-3}$ \\
    \bottomrule
\end{tabular}
%
	\end{minipage}%
\end{table}

\begin{table}
	\centering
	\caption{Parameters of the second acceleration activity of the ego vehicle, denoted by D in \cref{fig:example ego states lon}.}
	\label{tab:parameters ego D}
	\begin{minipage}[b]{.33\linewidth}
		\centering
		\subcaption{Normalization parameters.}
		\begin{tabular}{cr}
    \toprule
    Parameter & Value \\ \otoprule
    $x_{\textup{start}}$    & $\unit[2.88 \cdot 10^{2}]{m}$ \\
    $x_{\textup{end}}$      & $\unit[3.81 \cdot 10^{2}]{m}$ \\
    $t_{\textup{start}}$    & $\unit[1.80 \cdot 10^{1}]{s}$ \\
    $t_{\textup{duration}}$ & $\unit[4.98]{s}$ \\
    \bottomrule
\end{tabular}
%
	\end{minipage}%
	\begin{minipage}[b]{.5\linewidth}
		\centering
		\subcaption{Spline coefficients.}
		\begin{tabular}{crr}
    \toprule
    Coefficients & Spline 1 & Spline 2 \\ \otoprule
    $a$ & $-6.80 \cdot 10^{-2}$ & $-3.12 \cdot 10^{-2}$ \\
    $b$ & $1.46 \cdot 10^{-1}$ & $9.09 \cdot 10^{-2}$ \\
    $c$ & $9.17 \cdot 10^{-1}$ & $9.45 \cdot 10^{-1}$ \\
    $d$ & $3.85 \cdot 10^{-16}$ & $-4.60 \cdot 10^{-3}$ \\
    \bottomrule
\end{tabular}
%
	\end{minipage}%
\end{table}

\begin{table}
	\centering
	\caption{Normalization parameters of the two cruising activities of the ego vehicle.}
	\label{tab:parameters ego cruising}
	\begin{minipage}[b]{.33\linewidth}
		\centering
		\subcaption{Event C of \cref{fig:example ego states lon}.}
		\begin{tabular}{cr}
    \toprule
    Parameter & Value \\ \otoprule
    $x_{\textup{start}}$    & $\unit[1.68 \cdot 10^{2}]{m}$ \\
    $x_{\textup{end}}$      & $\unit[2.87 \cdot 10^{2}]{m}$ \\
    $t_{\textup{start}}$    & $\unit[1.10 \cdot 10^{1}]{s}$ \\
    $t_{\textup{duration}}$ & $\unit[6.98]{s}$ \\
    \bottomrule
\end{tabular}
%
	\end{minipage}%
	\begin{minipage}[b]{.33\linewidth}
		\centering
		\subcaption{Event E of \cref{fig:example ego states lon}.}
		\begin{tabular}{cr}
    \toprule
    Parameter & Value \\ \otoprule
    $x_{\textup{start}}$    & $\unit[3.81 \cdot 10^{2}]{m}$ \\
    $x_{\textup{end}}$      & $\unit[7.56 \cdot 10^{2}]{m}$ \\
    $t_{\textup{start}}$    & $\unit[2.30 \cdot 10^{1}]{s}$ \\
    $t_{\textup{duration}}$ & $\unit[1.93 \cdot 10^{1}]{s}$ \\
    \bottomrule
\end{tabular}
%
	\end{minipage}%
\end{table}

The parameters of the two lane change activities are presented in \cref{tab:parameters ego G} and \cref{tab:parameters ego I}, respectively. The parameters for the straight-driving activities are presented in \cref{tab:parameters ego straight}.

\begin{table}
	\centering
	\caption{Parameters of the first lane change activity of the ego vehicle, denoted by G in \cref{fig:example ego states lat}.}
	\label{tab:parameters ego G}
	\begin{minipage}[b]{.33\linewidth}
		\centering
		\subcaption{Normalization parameters.}
		\begin{tabular}{cr}
    \toprule
    Parameter & Value \\ \otoprule
    $x_{\textup{start}}$    & $\unit[7.00 \cdot 10^{-3}]{m}$ \\
    $x_{\textup{end}}$      & $\unit[3.64]{m}$ \\
    $t_{\textup{start}}$    & $\unit[1.70 \cdot 10^{1}]{s}$ \\
    $t_{\textup{duration}}$ & $\unit[5.98]{s}$ \\
    \bottomrule
\end{tabular}
%
	\end{minipage}%
	\begin{minipage}[b]{.5\linewidth}
		\centering
		\subcaption{Spline coefficients.}
		\begin{tabular}{crr}
    \toprule
    Coefficients & Spline 1 & Spline 2 \\ \otoprule
    $a$ & $-3.60$ & $-1.58$ \\
    $b$ & $4.45$ & $1.42$ \\
    $c$ & $-9.89 \cdot 10^{-2}$ & $1.42$ \\
    $d$ & $9.24 \cdot 10^{-17}$ & $-2.52 \cdot 10^{-1}$ \\
    \bottomrule
\end{tabular}
%
	\end{minipage}%
\end{table}

\begin{table}
	\centering
	\caption{Parameters of the second lane change activity of the ego vehicle, denoted by I in \cref{fig:example ego states lat}.}
	\label{tab:parameters ego I}
	\begin{minipage}[b]{.33\linewidth}
		\centering
		\subcaption{Normalization parameters.}
		\begin{tabular}{cr}
    \toprule
    Parameter & Value \\ \otoprule
    $x_{\textup{start}}$    & $\unit[3.40]{m}$ \\
    $x_{\textup{end}}$      & $\unit[5.49 \cdot 10^{-2}]{m}$ \\
    $t_{\textup{start}}$    & $\unit[3.12 \cdot 10^{1}]{s}$ \\
    $t_{\textup{duration}}$ & $\unit[4.98]{s}$ \\
    \bottomrule
\end{tabular}
%
	\end{minipage}%
	\begin{minipage}[b]{.5\linewidth}
		\centering
		\subcaption{Spline coefficients.}
		\begin{tabular}{crr}
    \toprule
    Coefficients & Spline 1 & Spline 2 \\ \otoprule
    $a$ & $-2.17$ & $-1.51$ \\
    $b$ & $2.90$ & $1.91$ \\
    $c$ & $1.84 \cdot 10^{-1}$ & $6.80 \cdot 10^{-1}$ \\
    $d$ & $2.16 \cdot 10^{-16}$ & $-8.27 \cdot 10^{-2}$ \\
    \bottomrule
\end{tabular}
%
	\end{minipage}%
\end{table}

\begin{table}
	\centering
	\caption{Normalization parameters of straight-driving activities of the ego vehicle.}
	\label{tab:parameters ego straight}
	\begin{minipage}[b]{.33\linewidth}
		\centering
		\subcaption{Event F of \cref{fig:example ego states lat}.}
		\begin{tabular}{cr}
    \toprule
    Parameter & Value \\ \otoprule
    $x_{\textup{start}}$    & $\unit[9.38 \cdot 10^{-2}]{m}$ \\
    $x_{\textup{end}}$      & $\unit[4.84 \cdot 10^{-3}]{m}$ \\
    $t_{\textup{start}}$    & $\unit[0.00]{s}$ \\
    $t_{\textup{duration}}$ & $\unit[1.70 \cdot 10^{1}]{s}$ \\
    \bottomrule
\end{tabular}
%
	\end{minipage}%
	\begin{minipage}[b]{.33\linewidth}
		\centering
		\subcaption{Event H of \cref{fig:example ego states lat}.}
		\begin{tabular}{cr}
    \toprule
    Parameter & Value \\ \otoprule
    $x_{\textup{start}}$    & $\unit[3.63]{m}$ \\
    $x_{\textup{end}}$      & $\unit[3.40]{m}$ \\
    $t_{\textup{start}}$    & $\unit[2.30 \cdot 10^{1}]{s}$ \\
    $t_{\textup{duration}}$ & $\unit[8.23]{s}$ \\
    \bottomrule
\end{tabular}
%
	\end{minipage}%
	\begin{minipage}[b]{.33\linewidth}
		\centering
		\subcaption{Event J of \cref{fig:example ego states lat}.}
		\begin{tabular}{cr}
    \toprule
    Parameter & Value \\ \otoprule
    $x_{\textup{start}}$    & $\unit[5.45 \cdot 10^{-2}]{m}$ \\
    $x_{\textup{end}}$      & $\unit[1.90 \cdot 10^{-2}]{m}$ \\
    $t_{\textup{start}}$    & $\unit[3.62 \cdot 10^{1}]{s}$ \\
    $t_{\textup{duration}}$ & $\unit[6.08]{s}$ \\
    \bottomrule
\end{tabular}
%
	\end{minipage}%
\end{table}

\cref{fig:example ego states pos acc} show the longitudinal position and acceleration of the ego vehicle. \Cref{fig:example ego states lon acc} shows that by employing splines the acceleration is automatically smoothed. 

\setlength\figureheight{170pt}
\setlength\figurewidth{0.475\linewidth}
\begin{figure}
	\centering
	\begin{minipage}[b]{0.49\linewidth}
		\centering
		% This file was created by matplotlib2tikz v0.6.14.
\begin{tikzpicture}

\definecolor{color1}{rgb}{1,0.498039215686275,0.0549019607843137}
\definecolor{color0}{rgb}{0.12156862745098,0.466666666666667,0.705882352941177}

\begin{axis}[
xlabel={$t$ [s]},
ylabel={Position [m]},
xmin=0, xmax=40,
ymin=-37.8221793267088, ymax=794.265765860885,
width=\figurewidth,
height=\figureheight,
tick align=outside,
tick pos=left,
xmajorgrids,
x grid style={white!69.019607843137251!black},
ymajorgrids,
y grid style={white!69.019607843137251!black},
clip marker paths
]
\addplot [semithick, color0, mark=*, mark size=1, mark options={solid}, only marks, forget plot]
table {%
0 0
0.025 0.260000991946827
0.05 0.519999504311528
0.075 0.781398048531486
0.1 1.04330329441941
0.125 1.30605429677011
0.15 1.57034908789534
0.175 1.83519658594119
0.2 2.10069759847524
0.225 2.36683586363862
0.25 2.63358688166029
0.275 2.90194723177993
0.3 3.17069129268353
0.325 3.43994231964945
0.35 3.71014588382332
0.375 3.9808476590926
0.4 4.2528486967612
0.425 4.52562367873605
0.45 4.79912472198246
0.475 5.0735624690602
0.5 5.34872666307087
0.525 5.62522771818077
0.55 5.90230519996662
0.575 6.17974494249173
0.6 6.45849334728852
0.625 6.73776005238786
0.65 7.01751111963998
0.675 7.29795639287991
0.7 7.5789327933346
0.725 7.86093117972033
0.75 8.14409015433287
0.775 8.42772136446439
0.8 8.71221973634412
0.825 8.99722177750509
0.85 9.28272014328647
0.875 9.56936126667944
0.9 9.85643904157995
0.925 10.1441873950806
0.95 10.433077466083
0.975 10.7224656480805
1 11.0127167555086
1.025 11.3034904188492
1.05 11.5947387518237
1.075 11.887385641588
1.1 12.1804610192359
1.125 12.4742121396362
1.15 12.7683488882898
1.175 13.062824362741
1.2 13.3588254919171
1.225 13.6557080474936
1.25 13.9534591832661
1.275 14.251671978615
1.3 14.5501558413824
1.325 14.8491569822015
1.35 15.1489230839679
1.375 15.4490843743412
1.4 15.750335523495
1.425 16.0519649092768
1.45 16.3539660612195
1.475 16.6564265837036
1.5 16.9594142258859
1.525 17.2634153854436
1.55 17.5683661644276
1.575 17.8736798231802
1.6 18.1794280739464
1.625 18.4860874386289
1.65 18.7935886118986
1.675 19.1024790975767
1.7 19.411710603426
1.725 19.7214588307272
1.75 20.032179493786
1.775 20.3432741351702
1.8 20.6550223515338
1.825 20.967147633385
1.85 21.2796458454345
1.875 21.5930742857301
1.9 21.907225658826
1.925 22.2227238535409
1.95 22.5384170551994
1.975 22.8543595121309
2 23.1716107222684
2.025 23.4891105077172
2.05 23.8068586894898
2.075 24.1253027699662
2.1 24.4442072362422
2.125 24.7639584560957
2.15 25.0850390785688
2.175 25.4066325969878
2.2 25.7288838265255
2.225 26.0513952825152
2.25 26.3741465136809
2.275 26.6975631761292
2.3 27.0214125037671
2.325 27.346163742249
2.35 27.6719988155741
2.375 27.9982945634719
2.4 28.3250458094585
2.425 28.6521380420398
2.45 28.9796392913347
2.475 29.3080281298727
2.5 29.636983115052
2.525 29.9672343747508
2.55 30.2978388774031
2.575 30.628672880173
2.6 30.960420981826
2.625 31.2925314484764
2.65 31.6250327168406
2.675 31.9584104643316
2.7 32.2922597445911
2.725 32.6270078289877
2.75 32.9619403684972
2.775 33.2970483051736
2.8 33.6330463825377
2.825 33.9698240744067
2.85 34.3073221433631
2.875 34.6459242214461
2.9 34.984675513407
2.925 35.3234235750002
2.95 35.6636355529315
2.975 36.0044572101355
3 36.3459585129094
3.025 36.6874598155729
3.05 37.0289578616157
3.075 37.3720140229317
3.1 37.7154711638008
3.125 38.059472476369
3.15 38.4036550366132
3.175 38.7480886040466
3.2 39.0938399231009
3.225 39.4402095496466
3.25 39.7872108732282
3.275 40.13421219692
3.3 40.4816329847557
3.325 40.8301343138831
3.35 41.1791705066855
3.375 41.5284692032267
3.4 41.8782205371018
3.425 42.2287849764227
3.45 42.5802863175238
3.475 42.9336112026179
3.5 43.2858029294465
3.525 43.6385542750956
3.55 43.9922081040969
3.575 44.3462880823875
3.6 44.701036052281
3.625 45.0561424279563
3.65 45.4119552105168
3.675 45.7683360375251
3.7 46.1249070036881
3.725 46.4816549625468
3.75 46.8396863596582
3.775 47.1982156151506
3.8 47.5569635625524
3.825 47.9162025256412
3.85 48.2759504670549
3.875 48.6365878503115
3.9 48.9976905244964
3.925 49.3596884530781
3.95 49.7222319975404
3.975 50.0850365463435
4 50.4482879320696
4.025 50.8120433178526
4.05 51.1770412297293
4.075 51.54204262213
4.1 51.9070405336759
4.125 52.2727862594068
4.15 52.6387929412868
4.175 53.0050408451706
4.2 53.3715623706885
4.225 53.7388102691035
4.25 54.1065798587483
4.275 54.474831263213
4.3 54.8440537278405
4.325 55.2135551376505
4.35 55.5842861957257
4.375 55.9550376099848
4.4 56.3257890243221
4.425 56.6966088589746
4.45 57.0674286936271
4.475 57.4407575185742
4.5 57.8145540007273
4.525 58.1890554293994
4.55 58.5640556914581
4.575 58.9395571237201
4.6 59.3154535252874
4.625 59.691454959554
4.65 60.0686987214836
4.675 60.4460048255133
4.7 60.8235062653745
4.725 61.2017689854946
4.75 61.5807704316091
4.775 61.9604846440389
4.8 62.3404824696438
4.825 62.7209637937424
4.85 63.1016579495299
4.875 63.4831594050287
4.9 63.8648756468233
4.925 64.2468734608421
4.95 64.6301388656321
4.975 65.0139943992676
5 65.39899586798
5.025 65.7841196712896
5.05 66.1693674671298
5.075 66.5558290946178
5.1 66.9428268798009
5.125 67.33030752053
5.15 67.7178473172379
5.175 68.1055950981761
5.2 68.4938613822891
5.225 68.8828591567446
5.25 69.2720298009578
5.275 69.6612812859159
5.3 70.0516719179699
5.325 70.4422479561029
5.35 70.8332457187527
5.375 71.2245342306796
5.4 71.6167857268467
5.425 72.0095210869472
5.45 72.4027725868141
5.475 72.7971195298495
5.5 73.1915517481346
5.525 73.585803251901
5.55 73.9806952269217
5.575 74.3761967355449
5.6 74.7728825553026
5.625 75.1698840703789
5.65 75.5668855849261
5.675 75.9640544594594
5.7 76.3625559796012
5.725 76.7612101736335
5.75 77.1602116954704
5.775 77.5596247041681
5.8 77.9593724166875
5.825 78.3599166357962
5.85 78.7606834115063
5.875 79.1621849428373
5.9 79.5641831648777
5.925 79.967430857341
5.95 80.3712314543305
5.975 80.7752961626738
6 81.1800477061512
6.025 81.5850285469866
6.05 81.9905262263156
6.075 82.3965775473785
6.1 82.8030752213536
6.125 83.2097746046107
6.15 83.6165423423708
6.175 84.0240400105529
6.2 84.4318817407822
6.225 84.8406294018228
6.25 85.249508086296
6.275 85.6585096461487
6.3 86.0681010755756
6.325 86.4778526385001
6.35 86.8890964228488
6.375 87.3005251390336
6.4 87.7125267107144
6.425 88.1252467188818
6.45 88.5387482959575
6.475 88.9531272492653
6.5 89.3678466684502
6.525 89.7830982524971
6.55 90.1985995157531
6.575 90.6143511019144
6.6 91.0308779320161
6.625 91.4476295218772
6.65 91.8653638728332
6.675 92.2830514032246
6.7 92.7005529959379
6.725 93.1186889491605
6.75 93.537440546511
6.775 93.9567208325666
6.8 94.37621843222
6.825 94.7964435278656
6.85 95.2168843375002
6.875 95.6381359445644
6.9 96.0598146544741
6.925 96.4820622384063
6.95 96.9052138794563
6.975 97.328762229537
7 97.7527638469627
7.025 98.1770068156296
7.05 98.6015043864275
7.075 99.0265362890509
7.1 99.4517838555139
7.125 99.8779951124609
7.15 100.304417026377
7.175 100.731664582003
7.2 101.159240086809
7.225 101.58723763788
7.25 102.015595439329
7.275 102.444097073538
7.3 102.873049233768
7.325 103.3021273212
7.35 103.731624863692
7.375 104.161615641678
7.4 104.592367285192
7.425 105.023591101841
7.45 105.455092747878
7.475 105.887700355722
7.5 106.320807560113
7.525 106.754809215884
7.55 107.189358824381
7.575 107.624129432265
7.6 108.059126943114
7.625 108.494445670542
7.65 108.930447333273
7.675 109.366991351278
7.7 109.803993018129
7.725 110.240990517721
7.75 110.678033622035
7.775 111.116285293962
7.8 111.554617729693
7.825 111.993619404252
7.85 112.432812105002
7.875 112.872313781257
7.9 113.31240915853
7.925 113.752906637537
7.95 114.193825536232
7.975 114.634980443086
8 115.077732131906
8.025 115.52096922236
8.05 115.964716683189
8.075 116.408690937301
8.1 116.852782133547
8.125 117.297783831343
8.15 117.743053764279
8.175 118.188801213522
8.2 118.634703243948
8.225 119.080700691842
8.25 119.527120603037
8.275 119.97375911134
8.3 120.421506549216
8.325 120.869413109917
8.35 121.317660544896
8.375 121.766185930525
8.4 122.214937642273
8.425 122.664475870776
8.45 123.114227586532
8.475 123.565081906894
8.5 124.016307824374
8.525 124.46805954753
8.55 124.920357738167
8.575 125.37295402111
8.6 125.826451426124
8.625 126.280106707357
8.65 126.73410843927
8.675 127.188245049308
8.7 127.642496782287
8.725 128.097728792242
8.75 128.553251205219
8.775 129.009752946536
8.8 129.466182003885
8.825 129.922433744573
8.85 130.37906483983
8.875 130.836066583371
8.9 131.293635410111
8.925 131.751382790715
8.95 132.208394382805
8.975 132.665099037597
9 133.120850776062
9.025 133.576480200871
9.05 134.031977594609
9.075 134.486120975529
9.1 134.939689925301
9.125 135.392191651619
9.15 135.84507789189
9.175 136.298325298343
9.2 136.751001317646
9.225 137.203498728434
9.25 137.655257361142
9.275 138.107068648329
9.3 138.559066062343
9.325 139.010433247241
9.35 139.461180668287
9.375 139.911404185718
9.4 140.361405902852
9.425 140.810222106983
9.45 141.258867932484
9.475 141.706115373566
9.5 142.152918370915
9.525 142.598670071363
9.55 143.043108712263
9.575 143.486947428113
9.6 143.929444896283
9.625 144.371555339458
9.65 144.812807022464
9.675 145.253506182425
9.7 145.693757861841
9.725 146.133807958207
9.75 146.573769281605
9.775 147.012270954379
9.8 147.450701337496
9.825 147.888953009313
9.85 148.326696484506
9.875 148.763948152843
9.9 149.200438664557
9.925 149.636686167955
9.95 150.071954107829
9.975 150.507259431324
10 150.942761092999
10.025 151.378262754674
10.05 151.811978307817
10.075 152.245951401073
10.1 152.679641176494
10.125 153.112642827968
10.15 153.545393830376
10.175 153.977891355234
10.2 154.410007721058
10.225 154.842005249105
10.25 155.273265155263
10.275 155.704643305329
10.3 156.13639083465
10.325 156.567806504821
10.35 156.998804038714
10.375 157.429805683125
10.4 157.860807327107
10.425 158.291804860891
10.45 158.722661057976
10.475 159.152908596333
10.5 159.582866539393
10.525 160.012368177916
10.55 160.441844981079
10.575 160.87200847697
10.6 161.3012560209
10.625 161.730259006635
10.65 162.159010642111
10.675 162.588113094056
10.7 163.017364731337
10.725 163.446140297411
10.75 163.874891932577
10.775 164.303643568163
10.8 164.732316015158
10.825 165.160817649677
10.85 165.589186828193
10.875 166.017438461994
10.9 166.445690095685
10.925 166.873937645188
10.95 167.301941318763
10.975 167.73005421192
11 168.158555846439
11.025 168.586448679242
11.05 169.013696234648
11.075 169.441174526571
11.1 169.868525643825
11.125 170.29552727255
11.15 170.72278352312
11.175 171.150281076648
11.2 171.576845101816
11.225 172.003409126985
11.25 172.429915406763
11.275 172.856417034037
11.3 173.282914593959
11.325 173.709159177098
11.35 174.135156739459
11.375 174.561510113147
11.4 174.9880117401
11.425 175.414268773403
11.45 175.840472889989
11.475 176.266470452343
11.5 176.692472077261
11.525 177.118473702499
11.55 177.544102669792
11.575 177.969634438277
11.6 178.395382002232
11.625 178.820753890066
11.65 179.245755511233
11.675 179.670395060502
11.7 180.094896679954
11.725 180.520592530844
11.75 180.946299615823
11.775 181.371801239026
11.8 181.797298804357
11.825 182.22280042756
11.85 182.648129332983
11.875 183.073380955383
11.9 183.49840677597
11.925 183.923404343952
11.95 184.348405965549
11.975 184.773306053595
12 185.198057674288
12.025 185.622809294543
12.05 186.047556864119
12.075 186.472751909592
12.1 186.897923754475
12.125 187.322925375534
12.15 187.747742937783
12.175 188.1724592125
12.2 188.596710830939
12.225 189.020958403173
12.25 189.445210021932
12.275 189.869739307322
12.3 190.294486876788
12.325 190.719238497262
12.35 191.143971976239
12.375 191.568223594569
12.4 191.992547865983
12.425 192.417045437262
12.45 192.841547056613
12.475 193.266044627354
12.5 193.690330578226
12.525 194.114332195752
12.55 194.538878957259
12.575 194.963681866214
12.6 195.388929432929
12.625 195.814181055219
12.65 196.23943267719
12.675 196.664856702347
12.7 197.09035832544
12.725 197.516350437802
12.75 197.942477790568
12.775 198.369229418324
12.8 198.795655798272
12.825 199.22165742319
12.85 199.647654985653
12.875 200.073656610891
12.9 200.499202083877
12.925 200.924851617297
12.95 201.351353243922
12.975 201.77755484508
13 202.203306468976
13.025 202.629058093411
13.05 203.054805657476
13.075 203.480108882849
13.1 203.905356449556
13.125 204.330608071417
13.15 204.755859693598
13.175 205.181107260094
13.2 205.607042845777
13.225 206.03307804684
13.25 206.459329673209
13.275 206.885581299031
13.3 207.311828860548
13.325 207.73841541187
13.35 208.16516296976
13.375 208.591702443398
13.4 209.018176413901
13.425 209.444423974769
13.45 209.87076240205
13.475 210.297259961754
13.5 210.724060036412
13.525 211.151061665181
13.55 211.57764133645
13.575 212.004108609497
13.6 212.430356170584
13.625 212.85686953294
13.65 213.283621160696
13.675 213.70946438756
13.7 214.135187680366
13.725 214.560685245916
13.75 214.986268217642
13.775 215.412019841866
13.8 215.837916066759
13.825 216.263917691997
13.85 216.68991525414
13.875 217.115916879159
13.9 217.541918503858
13.925 217.968067271797
13.95 218.394568898531
13.975 218.821066458125
14 219.247568085297
14.025 219.674468281122
14.05 220.101474246929
14.075 220.528725876931
14.1 220.955973431953
14.125 221.383225061635
14.15 221.810846597833
14.175 222.238482470437
14.2 222.665484099425
14.225 223.09248165626
14.25 223.519483285248
14.275 223.946484914237
14.3 224.373482470314
14.325 224.800961795139
14.35 225.228415131325
14.375 225.655666761226
14.4 226.082810767044
14.425 226.509808323668
14.45 226.937152917695
14.475 227.36465047145
14.5 227.792152102255
14.525 228.219591768274
14.55 228.646589325
14.575 229.073788183384
14.6 229.501285737132
14.625 229.928787367944
14.65 230.356288999077
14.675 230.783786552934
14.7 231.211366230351
14.725 231.639113782832
14.75 232.066125544294
14.775 232.492942661222
14.8 232.920690213594
14.825 233.348217455343
14.85 233.775215011858
14.875 234.202216640846
14.9 234.629218270482
14.925 235.056612455335
14.95 235.484118983148
14.975 235.911866536168
15 236.339545867117
15.025 236.767047497922
15.05 237.194415876709
15.075 237.621667506172
15.1 238.048731286134
15.125 238.475732915013
15.15 238.902982709132
15.175 239.330287997638
15.2 239.757789627904
15.225 240.185154847465
15.25 240.612406477366
15.275 241.039428683709
15.3 241.466501940901
15.325 241.893753570794
15.35 242.321128903972
15.375 242.748630535104
15.4 243.176132165698
15.425 243.603629719446
15.45 244.031131349821
15.475 244.458564797895
15.5 244.885816427359
15.525 245.313068056931
15.55 245.740315612281
15.575 246.167747379371
15.6 246.595244933337
15.625 247.02299135346
15.65 247.450697452231
15.675 247.878195005978
15.7 248.305696637002
15.725 248.733194190429
15.75 249.160881494427
15.775 249.588553856922
15.8 250.015801411725
15.825 250.442932738141
15.85 250.869930295194
15.875 251.296753975513
15.9 251.723505603589
15.925 252.150968188095
15.95 252.578377355329
15.975 253.005374912054
16 253.432376540933
16.025 253.859378169703
16.05 254.286554273222
16.075 254.713805902904
16.1 255.141281910887
16.125 255.568706766534
16.15 255.995708395522
16.175 256.422608422275
16.2 256.84936005057
16.225 257.276423374993
16.25 257.703675004784
16.275 258.130926635114
16.3 258.558174190245
16.325 258.985425820146
16.35 259.412486157142
16.375 259.83948778613
16.4 260.266489415009
16.425 260.6934039952
16.45 261.120155623175
16.475 261.547047315294
16.5 261.974048943743
16.525 262.401246128546
16.55 262.828489730996
16.575 263.255491359765
16.6 263.682345441091
16.625 264.108847068044
16.65 264.53560472341
16.675 264.962602279924
16.7 265.389053197757
16.725 265.815300758516
16.75 266.241800282625
16.775 266.668242860709
16.8 267.094490421897
16.825 267.520487436172
16.85 267.94623499958
16.875 268.371986624234
16.9 268.797588392475
16.925 269.222835958971
16.95 269.648087581481
16.975 270.073335147758
17 270.498586770267
17.025 270.923838392129
17.05 271.34883737425
17.075 271.773838995308
17.1 272.198836564266
17.125 272.623956631443
17.15 273.049208254062
17.175 273.474613143475
17.2 273.900114766568
17.225 274.324626346815
17.25 274.749224149144
17.275 275.174225770413
17.3 275.599330767675
17.325 276.024582389646
17.35 276.449452811711
17.375 276.874051625069
17.4 277.297803241582
17.425 277.721796369145
17.45 278.146047987365
17.475 278.570825156752
17.5 278.995826778029
17.525 279.420346202154
17.55 279.844938816541
17.575 280.269940437599
17.6 280.695046641339
17.625 281.120298263529
17.65 281.546730870581
17.675 281.973728427088
17.7 282.400272960044
17.725 282.826752023435
17.75 283.253003649148
17.775 283.679642562229
17.8 284.10689011736
17.825 284.534457737581
17.85 284.962205290171
17.875 285.390635642288
17.9 285.819218262485
17.925 286.2479658089
17.95 286.677280901634
17.975 287.106917768677
18 287.538669416145
18.025 287.97010455726
18.05 288.400852092748
18.075 288.832435120209
18.1 289.264682647092
18.125 289.698321627717
18.15 290.13232328305
18.175 290.56653613487
18.2 291.000816101154
18.225 291.435313614975
18.25 291.870401939797
18.275 292.306403602842
18.3 292.742698292375
18.325 293.179199957346
18.35 293.616964664141
18.375 294.055180239324
18.4 294.494681915757
18.425 294.93468702269
18.45 295.375188702873
18.475 295.81682120972
18.5 296.25854461623
18.525 296.700046300272
18.55 297.141808504656
18.575 297.584060191099
18.6 298.026904889814
18.625 298.470156580874
18.65 298.914239446266
18.675 299.358560290139
18.7 299.803811988801
18.725 300.249249191475
18.75 300.695250892327
18.775 301.141670617722
18.8 301.588418061
18.825 302.035959693616
18.85 302.483779506271
18.875 302.93253121789
18.9 303.381831106474
18.925 303.832328528624
18.95 304.283393088866
18.975 304.734678897397
19 305.186180619676
19.025 305.637801253522
19.05 306.089548668341
19.075 306.542537736181
19.1 306.996035141448
19.125 307.45001111418
19.15 307.904107045811
19.175 308.358604444802
19.2 308.813750108901
19.225 309.269747499286
19.25 309.72608476089
19.275 310.182586502232
19.3 310.639536785113
19.325 311.096612577068
19.35 311.554109958957
19.375 312.011901533058
19.4 312.470153281309
19.425 312.929188060354
19.45 313.388689813369
19.475 313.849323165367
19.5 314.310230467088
19.525 314.771482226691
19.55 315.233485211859
19.575 315.695791902328
19.6 316.158539254482
19.625 316.621521117311
19.65 317.08502288521
19.675 317.549214996646
19.7 318.013966769746
19.725 318.479111458757
19.75 318.944399940488
19.775 319.410651718834
19.8 319.877100374041
19.825 320.344102155545
19.85 320.81084486202
19.875 321.277346641598
19.9 321.744210030775
19.925 322.211207358394
19.95 322.679450871303
19.975 323.14758357911
20 323.615335363671
20.025 324.084088353666
20.05 324.553835665347
20.075 325.02335887363
20.1 325.492933053089
20.125 325.962684845143
20.15 326.432691768712
20.175 326.902939077963
20.2 327.373586933016
20.225 327.84433423896
20.25 328.315584377071
20.275 328.786836174808
20.3 329.258083478423
20.325 329.729689048315
20.35 330.201686347794
20.375 330.67386486908
20.4 331.146116670779
20.425 331.618841005884
20.45 332.091687075183
20.475 332.564934367222
20.5 333.038513530827
20.525 333.512515339114
20.55 333.98651262669
20.575 334.460594805596
20.6 334.935092090522
20.625 335.409593900626
20.65 335.88409571073
20.675 336.358768033195
20.7 336.833519844101
20.725 337.308766142303
20.75 337.784105860061
20.775 338.259857674718
20.8 338.735604952632
20.825 339.211356767609
20.85 339.687606519609
20.875 340.164108337104
20.9 340.641062673229
20.925 341.118022878918
20.95 341.594774697755
20.975 342.071815818222
21 342.549317639678
21.025 343.026819461462
21.05 343.504316729176
21.075 343.982276228547
21.1 344.459952010126
21.125 344.936953829437
21.15 345.414508827911
21.175 345.892238644072
21.2 346.369740465528
21.225 346.847472816194
21.25 347.325724640598
21.275 347.804125023802
21.3 348.282622285239
21.325 348.761124110875
21.35 349.23963617854
21.375 349.718638005672
21.4 350.197795506018
21.425 350.677292762148
21.45 351.156666064836
21.475 351.635913322459
21.5 352.115165150605
21.525 352.59441697865
21.55 353.074157982108
21.575 353.553909812399
21.6 354.033657067262
21.625 354.513534050406
21.65 354.993535881288
21.675 355.47376419822
21.7 355.95393826278
21.725 356.433935516268
21.75 356.9143135648
21.775 357.395065398739
21.8 357.876207392469
21.825 358.357459228544
21.85 358.838209410165
21.875 359.319168067469
21.9 359.800669904246
21.925 360.281933270475
21.95 360.762935105537
21.975 361.244109827124
22 361.725361663098
22.025 362.206377631057
22.05 362.687463980732
22.075 363.16896581729
22.1 363.64993314104
22.125 364.130184972936
22.15 364.611000891504
22.175 365.092039805203
22.2 365.573291641388
22.225 366.054538887435
22.25 366.535790723292
22.275 367.017580231293
22.3 367.499577473306
22.325 367.981579312118
22.35 368.46367008575
22.375 368.94617192616
22.4 369.42889054659
22.425 369.911887783187
22.45 370.39488962575
22.475 370.877886861917
22.5 371.360888704799
22.525 371.843890547463
22.55 372.32688778374
22.575 372.809889626084
22.6 373.292886862469
22.625 373.776047205542
22.65 374.259299049118
22.675 374.742087581005
22.7 375.224839422218
22.725 375.707586660089
22.75 376.190528873868
22.775 376.673547787046
22.8 377.156795021939
22.825 377.640197479309
22.85 378.123944711341
22.875 378.607843012479
22.9 379.091844858893
22.925 379.575842089433
22.95 380.059846978107
22.975 380.544094206833
23 381.02834605405
23.025 381.512597901696
23.05 381.996596440967
23.075 382.480348286149
23.1 382.96465316726
23.125 383.449155015819
23.15 383.933162507892
23.175 384.417159738213
23.2 384.901161584627
23.225 385.385285422702
23.25 385.86953726981
23.275 386.353559527341
23.3 386.837474028685
23.325 387.321225873976
23.35 387.804973106219
23.375 388.288724951721
23.4 388.772673938109
23.425 389.256671168538
23.45 389.740673015171
23.475 390.224803970056
23.5 390.709305818397
23.525 391.193577325785
23.55 391.677574556433
23.575 392.161928694808
23.6 392.646425922479
23.625 393.130927770928
23.65 393.615429619057
23.675 394.099926846837
23.7 394.5846322089
23.725 395.069629433811
23.75 395.554444759878
23.775 396.039196608912
23.8 396.523943834988
23.825 397.008695683912
23.85 397.493442910207
23.875 397.978194759459
23.9 398.462946608712
23.925 398.947693835217
23.95 399.432397971616
23.975 399.916895199178
24 400.401503399054
24.025 400.886255248087
24.05 401.370670343046
24.075 401.854922189936
24.1 402.339855220335
24.125 402.824786007382
24.15 403.309287855941
24.175 403.793982746635
24.2 404.2789845968
24.225 404.763981821711
24.25 405.248983671657
24.275 405.733985521923
24.3 406.219063084322
24.325 406.70431493518
24.35 407.189562158387
24.375 407.674814009675
24.4 408.160065861072
24.425 408.645275087599
24.45 409.130276937865
24.475 409.615369887717
24.5 410.100621738786
24.525 410.585873589644
24.55 411.07112081307
24.575 411.555943145846
24.6 412.040765478622
24.625 412.526017329377
24.65 413.01107785937
24.675 413.495825085669
24.7 413.980721740563
24.725 414.4657189655
24.75 414.95030565699
24.775 415.434923584181
24.8 415.920170807755
24.825 416.40555010304
24.85 416.891047325471
24.875 417.37677869357
24.9 417.862501207023
24.925 418.347998429236
24.95 418.833587902377
24.975 419.3193351229
25 419.805381470179
25.025 420.291633325405
25.05 420.77747134995
25.075 421.263223203213
25.1 421.74897042341
25.125 422.234644699845
25.15 422.720146551745
25.175 423.205779822427
25.2 423.691531675362
25.225 424.177079847352
25.25 424.662574861816
25.275 425.14782671268
25.3 425.633073936146
25.325 426.118325786901
25.35 426.603573010257
25.375 427.088865502151
25.4 427.57436735416
25.425 428.059864576045
25.45 428.545366427945
25.475 429.031037268713
25.5 429.516789121649
25.525 430.002299764101
25.55 430.487841172014
25.575 430.973593025059
25.6 431.459340245473
25.625 431.945092098736
25.65 432.430675663911
25.675 432.916172886015
25.7 433.402129693271
25.725 433.88809059163
25.75 434.373842444348
25.775 434.859594297501
25.8 435.345341518352
25.825 435.831249716345
25.85 436.317246935507
25.875 436.802786233665
25.9 437.288370864673
25.925 437.774118084978
25.95 438.259681899779
25.975 438.745196850577
26 439.230948703513
26.025 439.716471381429
26.05 440.201468606366
26.075 440.686894166694
26.1 441.17264138689
26.125 441.657804448545
26.15 442.142823574596
26.175 442.628320795936
26.2 443.113890948175
26.225 443.599638168481
26.25 444.085390021634
26.275 444.571141874897
26.3 445.056889095203
26.325 445.542640948356
26.35 446.028388168553
26.375 446.514199573896
26.4 447.000201427759
26.425 447.485802666909
26.45 447.9710545181
26.475 448.456764044709
26.5 448.942436344593
26.525 449.427938196928
26.55 449.913562929618
26.575 450.399314782881
26.6 450.885062003404
26.625 451.370813856667
26.65 451.856318398108
26.675 452.341752605867
26.7 452.827004457167
26.725 453.312132793279
26.75 453.797134643216
26.775 454.282494559936
26.8 454.767991782149
26.825 455.253741472791
26.85 455.739538461171
26.875 456.225540314925
26.9 456.711542169005
26.925 457.197539388167
26.95 457.683541242029
26.975 458.169538461409
27 458.655540315271
27.025 459.141418384709
27.05 459.62716560545
27.075 460.113266703551
27.1 460.599513921132
27.125 461.085294141244
27.15 461.571045994616
27.175 462.056793214921
27.2 462.542440149573
27.225 463.027937371241
27.25 463.513617750976
27.275 463.999369604348
27.3 464.48511682498
27.325 464.970868678134
27.35 465.456615898766
27.375 465.942367751811
27.4 466.428119604856
27.425 466.913866825161
27.45 467.399618678424
27.475 467.885133421571
27.5 468.370716662463
27.525 468.856468515834
27.55 469.342405495978
27.575 469.828389442599
27.6 470.314136663014
27.625 470.79996604972
27.65 471.285967903692
27.675 471.77210224424
27.7 472.258354099356
27.725 472.744201889377
27.75 473.229953742313
27.775 473.715705595358
27.8 474.201519184578
27.825 474.68752103844
27.85 475.173390563018
27.875 475.659142416063
27.9 476.145092418694
27.925 476.631089637964
27.95 477.116844273094
27.975 477.602591493291
28 478.088343346445
28.025 478.57397076388
28.05 479.059467986092
28.075 479.545201545044
28.1 480.030869043233
28.125 480.516370895351
28.15 481.002001727089
28.175 481.487748947067
28.2 481.973693542492
28.225 482.459690761762
28.25 482.945445082528
28.275 483.431196935354
28.3 483.916944156096
28.325 484.40257982757
28.35 484.888077049238
28.375 485.373752470205
28.4 485.859504323032
28.425 486.345251543992
28.45 486.831003397254
28.475 487.31675061756
28.5 487.802963464731
28.525 488.289965322519
28.55 488.776201666173
28.575 489.262168260051
28.6 489.747915480575
28.625 490.233411768195
28.65 490.718663619168
28.675 491.204088314229
28.7 491.68959016602
28.725 492.174862645169
28.75 492.660209796867
28.775 493.145961649912
28.8 493.631708870545
28.825 494.117460723589
28.85 494.603041275048
28.875 495.088543126948
28.9 495.574271598508
28.925 496.059981675583
28.95 496.545483527374
28.975 497.031073198792
29 497.516825052055
29.025 498.002576904663
29.05 498.488324125732
29.075 498.974306760501
29.1 499.460303979881
29.125 499.946305833743
29.15 500.432307687715
29.175 500.918286412762
29.2 501.404038265807
29.225 501.889785486549
29.25 502.375537339484
29.275 502.861427243821
29.3 503.347424462873
29.325 503.833026798812
29.35 504.318532353546
29.375 504.804284206591
29.4 505.290104685894
29.425 505.776101905164
29.45 506.261725674485
29.475 506.746972898169
29.5 507.232619818214
29.525 507.71837167115
29.55 508.203625107463
29.575 508.688988390852
29.6 509.174735611158
29.625 509.660617390514
29.65 510.146619244377
29.675 510.632616463429
29.7 511.118534308085
29.725 511.604281528172
29.75 512.090366089418
29.775 512.576617944317
29.8 513.062467794387
29.825 513.548219647432
29.85 514.033717930135
29.875 514.51933723687
29.9 515.005339090515
29.925 515.491336309895
29.95 515.977338163648
29.975 516.46333538281
30 516.949337236454
30.025 517.435339090644
30.05 517.921289341475
30.075 518.407041194083
30.1 518.892788414825
30.125 519.37854026787
30.15 519.864292120478
30.175 520.350039341438
30.2 520.835791194374
30.225 521.321672854263
30.25 521.807674708126
30.275 522.293458046597
30.3 522.779179873548
30.325 523.264681725448
30.35 523.750263626692
30.375 524.236015479627
30.4 524.721915319157
30.425 525.2079125381
30.45 525.693291604806
30.475 526.178577881236
30.5 526.664579735098
30.525 527.150441065044
30.55 527.635938287039
30.575 528.121635098768
30.6 528.607632317711
30.625 529.093326716875
30.65 529.578828569102
30.675 530.064791026641
30.7 530.55071253598
30.725 531.036459756504
30.75 531.522403841794
30.775 532.008405695657
30.8 532.4944029146
30.825 532.980404768899
30.85 533.466401988169
30.875 533.952266226644
30.9 534.438018079798
30.925 534.923366793074
30.95 535.408660930587
30.975 535.894658149421
31 536.380590387604
31.025 536.866342240758
31.05 537.352218810641
31.075 537.838220664721
31.1 538.32441550197
31.125 538.810667356651
31.15 539.29642272875
31.175 539.782284928977
31.2 540.268536783657
31.225 540.754652323029
31.25 541.240654177219
31.275 541.727121748202
31.3 542.213367716717
31.325 542.699119569761
31.35 543.185034712994
31.375 543.671036566966
31.4 544.157038421046
31.425 544.642996325418
31.45 545.128748178463
31.475 545.6145943475
31.5 546.100596201253
31.525 546.586436319203
31.55 547.072183539726
31.575 547.558153669446
31.6 548.044150888716
31.625 548.53015274247
31.65 549.016154596115
31.675 549.502151815385
31.7 549.988153669248
31.725 550.474150888409
31.75 550.960332066593
31.775 551.446525308145
31.8 551.932272529214
31.825 552.418163229817
31.85 552.904160448979
31.875 553.390162302842
31.9 553.87613633607
31.925 554.361883556594
31.95 554.847722692098
31.975 555.333719911151
32 555.819572941712
32.025 556.305324794757
32.05 556.79128128961
32.075 557.2773191401
32.1 557.763816356101
32.125 558.25008400162
32.15 558.735835854447
32.175 559.22158307497
32.2 559.707334928233
32.225 560.193280252779
32.25 560.67928210686
32.275 561.165283961377
32.3 561.651281180211
32.325 562.137283034182
32.35 562.623085992333
32.375 563.108885486366
32.4 563.594887340556
32.425 564.080766153721
32.45 564.566518006766
32.475 565.052616317659
32.5 565.538868172666
32.525 566.024646547683
32.55 566.510488181756
32.575 566.996740036764
32.6 567.482772806402
32.625 567.968524658902
32.65 568.454441750565
32.675 568.940438969181
32.7 569.426477930739
32.725 569.912516892297
32.75 570.398518746396
32.775 570.884520600174
32.8 571.370517819371
32.825 571.856675348607
32.85 572.342922566408
32.875 572.829174421078
32.9 573.315164270645
32.925 573.800661491979
32.95 574.286553215537
32.975 574.772550434734
33 575.258552288943
33.025 575.74455414283
33.05 576.230551361808
33.075 576.716553216235
33.1 577.202550435214
33.125 577.688547654192
33.15 578.17432342676
33.175 578.660570644341
33.2 579.146754361642
33.225 579.632751580259
33.25 580.118753434121
33.275 580.604755288093
33.3 581.090752507581
33.325 581.576754361553
33.35 582.062751580606
33.375 582.548810109195
33.4 583.035061964093
33.425 583.521174611178
33.45 584.007176465259
33.475 584.493640899242
33.5 584.979973258681
33.525 585.465975112326
33.55 585.951848279568
33.575 586.437600132613
33.6 586.923544730654
33.625 587.409546584408
33.65 587.895548438598
33.675 588.381545657977
33.7 588.86754751184
33.725 589.353473968392
33.75 589.838975820618
33.775 590.324574883304
33.8 590.81032210361
33.825 591.296073956764
33.85 591.781821177178
33.875 592.267573030005
33.9 592.753410484791
33.925 593.239407703953
33.95 593.725337603802
33.975 594.211267503652
34 594.69726935797
34.025 595.183271211639
34.05 595.669268430617
34.075 596.155418254594
34.1 596.641665471964
34.125 597.127718631145
34.15 597.613737079231
34.175 598.099984297032
34.2 598.58623615203
34.225 599.072483369509
34.25 599.558483406707
34.275 600.044235260023
34.3 600.530358838256
34.325 601.016610693145
34.35 601.502108029367
34.375 601.987724532024
34.4 602.473726386014
34.425 602.95985619695
34.45 603.44610805162
34.475 603.9323552691
34.5 604.418520656959
34.525 604.904522510846
34.55 605.390686487414
34.575 605.876938342302
34.6 606.362957085577
34.625 606.848958939676
34.65 607.334960793564
34.675 607.821153344561
34.7 608.307655200669
34.725 608.793817539111
34.75 609.279819393211
34.775 609.766038174371
34.8 610.252257263254
34.825 610.738259117681
34.85 611.224353639903
34.875 611.710605495004
34.9 612.196708674208
34.925 612.682705893296
34.95 613.168925449807
34.975 613.655215313207
35 614.141717168987
35.025 614.628084191774
35.05 615.114331409575
35.075 615.600583264135
35.1 616.086782020808
35.125 616.572533873906
35.15 617.058547989574
35.175 617.545045205767
35.2 618.031400670175
35.225 618.517647887764
35.25 619.003899742543
35.275 619.490132793472
35.3 619.976130012996
35.325 620.462210556924
35.35 620.948457774944
35.375 621.434846994803
35.4 621.921348850803
35.425 622.407446796989
35.45 622.893448650766
35.475 623.379445869963
35.5 623.865447724063
35.525 624.351449577841
35.55 624.837446797037
35.575 625.323495471742
35.6 625.809742689324
35.625 626.295994544213
35.65 626.782246398992
35.675 627.268671600518
35.7 627.755173456408
35.725 628.241177913994
35.75 628.727276479982
35.775 629.213778335654
35.8 629.699954437646
35.825 630.185706290744
35.85 630.671622521306
35.875 631.157624375515
35.9 631.643626229074
35.925 632.129593758067
35.95 632.615345611274
35.975 633.101289002571
36 633.587540857569
36.025 634.073635718795
36.05 634.559632937882
36.075 635.045634792091
36.1 635.531547792062
36.125 636.017299644942
36.15 636.50324771137
36.175 636.989226274314
36.2 637.474978127303
36.225 637.960800327833
36.25 638.446802182041
36.275 638.93294662465
36.3 639.419193842451
36.325 639.905050918898
36.35 640.390813765398
36.375 640.877065620286
36.4 641.363182234712
36.425 641.848929454977
36.45 642.334681308184
36.475 642.820428528668
36.5 643.306359930145
36.525 643.792361784032
36.55 644.27811098394
36.575 644.763921156285
36.6 645.249918375045
36.625 645.735780027097
36.65 646.221531880194
36.675 646.70798601673
36.7 647.194314275387
36.725 647.680311494257
36.75 648.166313348465
36.775 648.652315202243
36.8 649.13831242144
36.825 649.624236026209
36.85 650.109733248519
36.875 650.59543175258
36.9 651.081433605921
36.925 651.567261772201
36.95 652.053013624971
36.975 652.538760845454
37 653.02448560256
37.025 653.509987454438
37.05 653.995649232874
37.075 654.481651087083
37.1 654.967648306171
37.125 655.453650159948
37.15 655.939033224624
37.175 656.424366998851
37.2 656.910118851409
37.225 657.395866072214
37.25 657.881617924772
37.275 658.367135014492
37.3 658.852647183396
37.325 659.339886878733
37.35 659.825677237184
37.375 660.311679091251
37.4 660.797680945545
37.425 661.283678164462
37.45 661.769473495222
37.475 662.255239197703
37.5 662.741241051444
37.525 663.227126834174
37.55 663.712624055662
37.575 664.198125907829
37.6 664.683623129534
37.625 665.168927211127
37.65 665.654198264138
37.675 666.13969548606
37.7 666.62511190306
37.725 667.110359126605
37.75 667.595808266529
37.775 668.081386698343
37.8 668.56738391726
37.825 669.053385771662
37.85 669.53938299047
37.875 670.025207038749
37.9 670.510958892473
37.925 670.996215514062
37.95 671.481562002892
37.975 671.967309223321
38 672.453061076167
38.025 672.938812929229
38.05 673.424560149432
38.075 673.910312002277
38.1 674.395830739206
38.125 674.881332591047
38.15 675.366834443105
38.175 675.852331665364
38.2 676.337833517531
38.225 676.823330739345
38.25 677.308832591739
38.275 677.794334443797
38.3 678.279744797741
38.325 678.764996648913
38.35 679.250634683398
38.375 679.736350682918
38.4 680.221602534091
38.425 680.707008386857
38.45 681.192760239702
38.475 681.678369170348
38.5 682.163871021972
38.525 682.649175208652
38.55 683.134467643121
38.575 683.620469497079
38.6 684.106398798061
38.625 684.592150651123
38.65 685.077902504077
38.675 685.563649724171
38.7 686.049401577017
38.725 686.535148797772
38.75 687.020900650726
38.775 687.506708963521
38.8 687.992706182773
38.825 688.478570345083
38.85 688.964317565611
38.875 689.450069418457
38.9 689.935649815046
38.925 690.420897038254
38.95 690.906314068207
38.975 691.391811290465
39 691.877084797033
39.025 692.362413032832
39.05 692.848160253143
39.075 693.333814943913
39.1 693.81931216551
39.125 694.304479091878
39.15 694.789480942371
39.175 695.274694653656
39.2 695.759918774046
39.225 696.244915998649
39.25 696.7300040393
39.275 697.215255890245
39.3 697.700503114007
39.325 698.185754964952
39.35 698.671002188269
39.375 699.156299718457
39.4 699.641801570298
39.425 700.127435077228
39.45 700.613186929856
39.475 701.098693221392
39.5 701.584195074003
39.525 702.069696925844
39.55 702.555322264553
39.575 703.04107411729
39.6 703.526228111882
39.625 704.011229961823
39.65 704.496478568464
39.675 704.981725791564
39.7 705.466977642509
39.725 705.951992922515
39.75 706.436744771452
39.775 706.921853557508
39.8 707.407100780825
39.825 707.892352631553
39.85 708.377692843025
39.875 708.863444696196
39.9 709.349089852431
39.925 709.834587074245
39.95 710.320473261608
39.975 710.806382490089
40 711.291884342039
40.025 711.777521096225
40.05 712.263268316753
40.075 712.749020169707
40.1 713.234767390571
40.125 713.720519243633
40.15 714.206357718862
40.175 714.692354937887
40.2 715.178064762304
40.225 715.663561984226
40.25 716.149063836176
40.275 716.634581974536
40.3 717.120329194965
40.325 717.605928274511
40.35 718.091175497502
40.375 718.576563648174
40.4 719.062065500232
40.425 719.547364780114
40.45 720.032626780294
40.475 720.518124002325
40.5 721.003453578168
40.525 721.488455428218
40.55 721.973848732912
40.575 722.459350584862
40.6 722.944847807012
40.625 723.430463408533
40.65 723.916215261379
40.675 724.401962482134
40.7 724.887714334979
40.725 725.37346155529
40.75 725.859213407919
40.775 726.344965260873
40.8 726.830712481737
40.825 727.316464334908
40.85 727.802379747654
40.875 728.288381601395
40.9 728.774383455462
40.925 729.260342193739
40.95 729.746094046585
40.975 730.231743916344
41 730.717245768511
41.025 731.202747620352
41.05 731.688244842719
41.075 732.173746694669
41.1 732.659159158584
41.125 733.144411009421
41.15 733.629857782119
41.175 734.115318024884
41.2 734.60031987516
41.225 735.085472563065
41.25 735.570974415231
41.275 736.056189688167
41.3 736.541186913215
41.325 737.026188763373
41.35 737.511236321116
41.375 737.996988173744
41.4 738.482607971438
41.425 738.967855194864
41.45 739.453107046253
41.475 739.938354269462
41.5 740.424002820213
41.525 740.909754673059
41.55 741.395004977072
41.575 741.880315037996
41.6 742.365812260136
41.625 742.851451467227
41.65 743.337203320073
41.675 743.82246983146
41.7 744.30763588773
41.725 744.792633112778
41.75 745.277803105339
41.775 745.763054956176
41.8 746.248302179276
41.825 746.733590835296
41.85 747.219088057446
41.875 747.70439540744
41.9 748.189397257381
41.925 748.674730910025
41.95 749.16023276231
41.975 749.646152314109
42 750.132056473694
42.025 750.61705832397
42.05 751.102153028853
42.075 751.587654880803
42.1 752.073260118679
42.125 752.559011971525
42.15 753.044959930822
42.175 753.530918850378
42.2 754.016670703115
42.225 754.50214542068
42.25 754.987397271734
42.275 755.472866601013
42.3 755.95833468334
42.325 756.443586534176
};
\addplot [line width=2.4000000000000004pt, color1, dashed, forget plot]
table {%
0 4.96315940142527e-14
0.025 0.261539436072527
0.05 0.523822961657963
0.075 0.786849277751799
0.1 1.05061708534948
0.125 1.31512508544644
0.15 1.58037197903812
0.175 1.84635646711996
0.2 2.11307725068742
0.225 2.38053303073591
0.25 2.6487225082609
0.275 2.91764438425781
0.3 3.1872973597221
0.325 3.45768013564919
0.35 3.72879141303454
0.375 4.00062989287357
0.4 4.27319427616174
0.425 4.54648326389449
0.45 4.82049555706725
0.475 5.09522985667547
0.5 5.37068486371459
0.525 5.64685927918004
0.55 5.92375180406728
0.575 6.20136113937173
0.6 6.47968598608885
0.625 6.75872504521407
0.65 7.03847701774284
0.675 7.31894060467059
0.7 7.60011450699276
0.725 7.88199742570481
0.75 8.16458806180216
0.775 8.44788511628026
0.8 8.73188729013455
0.825 9.01659328436047
0.85 9.30200179995347
0.875 9.58811153790897
0.9 9.87492119922244
0.925 10.1624294848893
0.95 10.450635095905
0.975 10.739536733265
1 11.0291330979646
1.025 11.3194228909995
1.05 11.610404813365
1.075 11.9020775660564
1.1 12.1944398500694
1.125 12.4874903663993
1.15 12.7812278160416
1.175 13.0756508999917
1.2 13.370758319245
1.225 13.666548774797
1.25 13.9630209676431
1.275 14.2601735987788
1.3 14.5580053691995
1.325 14.8565149799007
1.35 15.1557011318778
1.375 15.4555625261262
1.4 15.7560978636413
1.425 16.0573058454187
1.45 16.3591851724538
1.475 16.6617345457419
1.5 16.9649526662786
1.525 17.2688382350593
1.55 17.5733899530793
1.575 17.8786065213343
1.6 18.1844866408195
1.625 18.4910290125305
1.65 18.7982323374627
1.675 19.1060953166115
1.7 19.4146166509724
1.725 19.7237950415408
1.75 20.0336291893121
1.775 20.3441177952818
1.8 20.6552595604454
1.825 20.9670531857982
1.85 21.2794973723358
1.875 21.5925908210535
1.9 21.9063322329468
1.925 22.2207203090111
1.95 22.5357537502419
1.975 22.8514312576347
2 23.1677515321848
2.025 23.4847132748876
2.05 23.8023151867388
2.075 24.1205559687336
2.1 24.4394343218675
2.125 24.7589489471361
2.15 25.0790985455346
2.175 25.3998818180585
2.2 25.7212974657034
2.225 26.0433441894646
2.25 26.3660206903375
2.275 26.6893256693177
2.3 27.0132578274005
2.325 27.3378158655814
2.35 27.6629984848559
2.375 27.9888043862193
2.4 28.3152322706671
2.425 28.6422808391948
2.45 28.9699487927978
2.475 29.2982348324716
2.5 29.6271376592115
2.525 29.956655974013
2.55 30.2867884778716
2.575 30.6175338717827
2.6 30.9488908567418
2.625 31.2808581337442
2.65 31.6134344037855
2.675 31.946618367861
2.7 32.2804087269663
2.725 32.6148041820967
2.75 32.9498034342477
2.775 33.2854051844147
2.8 33.6216081335932
2.825 33.9584109827786
2.85 34.2958124329664
2.875 34.6338111851519
2.9 34.9724059403307
2.925 35.3115953994982
2.95 35.6513782636498
2.975 35.991753233781
3 36.3327190108871
3.025 36.6742742959637
3.05 37.0164177900062
3.075 37.3591481940101
3.1 37.7024642089706
3.125 38.0463645358834
3.15 38.3908478757439
3.175 38.7359129295474
3.2 39.0815583982895
3.225 39.4277829829655
3.25 39.7745853845709
3.275 40.1219643041012
3.3 40.4699184425518
3.325 40.8184465009181
3.35 41.1675471801956
3.375 41.5172191813798
3.4 41.8674612054659
3.425 42.2182719534496
3.45 42.5696501263262
3.475 42.9215944250912
3.5 43.2741035507401
3.525 43.6271762042682
3.55 43.9808110866709
3.575 44.3350068989439
3.6 44.6897623420824
3.625 45.045076117082
3.65 45.400946924938
3.675 45.7573734666459
3.7 46.1143544432012
3.725 46.4718885555993
3.75 46.8299745048356
3.775 47.1886109919056
3.8 47.5477967178047
3.825 47.9075303835284
3.85 48.267810690072
3.875 48.6286363384311
3.9 48.9900060296011
3.925 49.3519184645774
3.95 49.7143723443554
3.975 50.0773663699307
4 50.4408992422986
4.025 50.8049696624546
4.05 51.1695763313941
4.075 51.5347179501126
4.1 51.9003932196055
4.125 52.2666008408682
4.15 52.6333395148962
4.175 53.000607942685
4.2 53.3684048252299
4.225 53.7367288635265
4.25 54.1055787585701
4.275 54.4749532113562
4.3 54.8448509228802
4.325 55.2152705941376
4.35 55.5862109261239
4.375 55.9576706198343
4.4 56.3296483762645
4.425 56.7021428964099
4.45 57.0751528812658
4.475 57.4486770318277
4.5 57.8227140683282
4.525 58.1972631534544
4.55 58.5723238923477
4.575 58.9478959093868
4.6 59.3239788289501
4.625 59.7005722754164
4.65 60.0776758731643
4.675 60.4552892465724
4.7 60.8334120200193
4.725 61.2120438178836
4.75 61.5911842645439
4.775 61.9708329843789
4.8 62.3509896017672
4.825 62.7316537410873
4.85 63.112825026718
4.875 63.4945030830377
4.9 63.8766875344252
4.925 64.259378005259
4.95 64.6425741199178
4.975 65.0262755027801
5 65.4104817782247
5.025 65.79519257063
5.05 66.1804075043748
5.075 66.5661262038376
5.1 66.9523482933971
5.125 67.3390733974318
5.15 67.7263011403204
5.175 68.1140311464415
5.2 68.5022630401738
5.225 68.8909964458958
5.25 69.2802309879861
5.275 69.6699662908234
5.3 70.0602019787862
5.325 70.4509376762533
5.35 70.8421730076032
5.375 71.2339075972145
5.4 71.6261410694658
5.425 72.0188730487358
5.45 72.4121031594031
5.475 72.8058310258462
5.5 73.2000562724439
5.525 73.5947785235747
5.55 73.9899974036172
5.575 74.3857125369501
5.6 74.7819235479519
5.625 75.1786300610013
5.65 75.575831700477
5.675 75.9735280907574
5.7 76.3717188562213
5.725 76.7704036212472
5.75 77.1695820102138
5.775 77.5692536474996
5.8 77.9694181574834
5.825 78.3700751645436
5.85 78.771224293059
5.875 79.1728651674081
5.9 79.5749974119696
5.925 79.9776206511221
5.95 80.3807345092441
5.975 80.7843386107144
6 81.1884325799114
6.025 81.5930160412139
6.05 81.9980886190005
6.075 82.4036499376497
6.1 82.8096996215402
6.125 83.2162372950506
6.15 83.6232625825595
6.175 84.0307751084456
6.2 84.4387744970873
6.225 84.8472603728635
6.25 85.2562323601526
6.275 85.6656900833333
6.3 86.0756331667842
6.325 86.4860612348839
6.35 86.8969739120111
6.375 87.3083708225443
6.4 87.7202515908623
6.425 88.1326158413435
6.45 88.5454631983665
6.475 88.9587932863101
6.5 89.3726057295529
6.525 89.7869001524734
6.55 90.2016761794502
6.575 90.616933434862
6.6 91.0326715430875
6.625 91.4488901285051
6.65 91.8655888154936
6.675 92.2827672284315
6.7 92.7004249916975
6.725 93.1185617296701
6.75 93.537177066728
6.775 93.9562706272499
6.8 94.3758420356142
6.825 94.7958909161997
6.85 95.216416893385
6.875 95.6374195915486
6.9 96.0588986350692
6.925 96.4808536483254
6.95 96.9032842556958
6.975 97.326190081559
7 97.7495707502937
7.025 98.1734258862785
7.05 98.5977551138919
7.075 99.0225580575127
7.1 99.4478343415193
7.125 99.8735835902904
7.15 100.299805428205
7.175 100.726499479641
7.2 101.153665368977
7.225 101.581302720593
7.25 102.009411158866
7.275 102.437990308175
7.3 102.867039792899
7.325 103.296559237417
7.35 103.726548266106
7.375 104.157006503346
7.4 104.587933573516
7.425 105.019329100993
7.45 105.451192710157
7.475 105.883524025387
7.5 106.31632267106
7.525 106.749588271555
7.55 107.183320451251
7.575 107.617518834527
7.6 108.052183045761
7.625 108.487312709332
7.65 108.922907449619
7.675 109.358966890999
7.7 109.795490657852
7.725 110.232478374557
7.75 110.669929665491
7.775 111.107844155034
7.8 111.546221467564
7.825 111.98506122746
7.85 112.4243630591
7.875 112.864126586863
7.9 113.304351435128
7.925 113.745037228273
7.95 114.186183590677
7.975 114.627790146718
8 115.069856520776
8.025 115.512382337228
8.05 115.955367220453
8.075 116.398810794831
8.1 116.842712684739
8.125 117.287072514556
8.15 117.731889908661
8.175 118.177164491432
8.2 118.622895887248
8.225 119.069083720488
8.25 119.51572761553
8.275 119.962827196753
8.3 120.410382088535
8.325 120.858391915255
8.35 121.306856301292
8.375 121.755774871025
8.4 122.205147248831
8.425 122.65497305909
8.45 123.10525192618
8.475 123.55598347448
8.5 124.007167328368
8.525 124.458803112224
8.55 124.910890450424
8.575 125.363428967349
8.6 125.816418287377
8.625 126.269858034887
8.65 126.723747834256
8.675 127.178087309864
8.7 127.63287608609
8.725 128.088113787311
8.75 128.543800037907
8.775 128.999934462256
8.8 129.456516684737
8.825 129.913546329729
8.85 130.371023021609
8.875 130.828946384757
8.9 131.287316043551
8.925 131.74613162237
8.95 132.205392745592
8.975 132.665099037597
9 133.120850776062
9.025 133.578256789612
9.05 134.03497431605
9.075 134.49100992786
9.1 134.946370197529
9.125 135.401061697542
9.15 135.855091000385
9.175 136.308464678545
9.2 136.761189304507
9.225 137.213271450756
9.25 137.66471768978
9.275 138.115534594063
9.3 138.565728736092
9.325 139.015306688352
9.35 139.464275023329
9.375 139.91264031351
9.4 140.36040913138
9.425 140.807588049425
9.45 141.254183640131
9.475 141.700202475984
9.5 142.145651129469
9.525 142.590536173073
9.55 143.034864179281
9.575 143.47864172058
9.6 143.921875369455
9.625 144.364571698392
9.65 144.806737279877
9.675 145.248378686395
9.7 145.689502490434
9.725 146.130115264478
9.75 146.570223581013
9.775 147.009834012526
9.8 147.448953131502
9.825 147.887587510427
9.85 148.325743721787
9.875 148.763428338069
9.9 149.200647931756
9.925 149.637409075337
9.95 150.073718341296
9.975 150.50958230212
10 150.945007683652
10.025 151.380004738974
10.05 151.814587248405
10.075 152.248769145622
10.1 152.682564364303
10.125 153.115986838124
10.15 153.549050500764
10.175 153.981769285898
10.2 154.414157127204
10.225 154.84622795836
10.25 155.277995713043
10.275 155.709474324929
10.3 156.140677727697
10.325 156.571619855022
10.35 157.002314640583
10.375 157.432776018057
10.4 157.86301792112
10.425 158.293054283451
10.45 158.722899038725
10.475 159.152566120621
10.5 159.582069462816
10.525 160.011422998987
10.55 160.44064066281
10.575 160.869736387964
10.6 161.298724108125
10.625 161.72761775697
10.65 162.156431268178
10.675 162.585178575424
10.7 163.013873612386
10.725 163.442530312742
10.75 163.871162610169
10.775 164.299784438343
10.8 164.728409730942
10.825 165.157052421643
10.85 165.585726444124
10.875 166.014445732061
10.9 166.443224219132
10.925 166.872075839014
10.95 167.301014525384
10.975 167.73005421192
11 168.158555846439
11.025 168.584894061214
11.05 169.011232275989
11.075 169.437570490764
11.1 169.863908705539
11.125 170.290246920314
11.15 170.716585135089
11.175 171.142923349864
11.2 171.569261564639
11.225 171.995599779414
11.25 172.421937994189
11.275 172.848276208964
11.3 173.274614423739
11.325 173.700952638514
11.35 174.12729085329
11.375 174.553629068065
11.4 174.97996728284
11.425 175.406305497615
11.45 175.83264371239
11.475 176.258981927165
11.5 176.68532014194
11.525 177.111658356715
11.55 177.53799657149
11.575 177.964334786265
11.6 178.39067300104
11.625 178.817011215815
11.65 179.24334943059
11.675 179.669687645365
11.7 180.09602586014
11.725 180.522364074915
11.75 180.94870228969
11.775 181.375040504465
11.8 181.80137871924
11.825 182.227716934015
11.85 182.65405514879
11.875 183.080393363566
11.9 183.506731578341
11.925 183.933069793116
11.95 184.359408007891
11.975 184.785746222666
12 185.212084437441
12.025 185.638422652216
12.05 186.064760866991
12.075 186.491099081766
12.1 186.917437296541
12.125 187.343775511316
12.15 187.770113726091
12.175 188.196451940866
12.2 188.622790155641
12.225 189.049128370416
12.25 189.475466585191
12.275 189.901804799966
12.3 190.328143014741
12.325 190.754481229516
12.35 191.180819444291
12.375 191.607157659067
12.4 192.033495873842
12.425 192.459834088617
12.45 192.886172303392
12.475 193.312510518167
12.5 193.738848732942
12.525 194.165186947717
12.55 194.591525162492
12.575 195.017863377267
12.6 195.444201592042
12.625 195.870539806817
12.65 196.296878021592
12.675 196.723216236367
12.7 197.149554451142
12.725 197.575892665917
12.75 198.002230880692
12.775 198.428569095467
12.8 198.854907310242
12.825 199.281245525017
12.85 199.707583739792
12.875 200.133921954567
12.9 200.560260169343
12.925 200.986598384118
12.95 201.412936598893
12.975 201.839274813668
13 202.265613028443
13.025 202.691951243218
13.05 203.118289457993
13.075 203.544627672768
13.1 203.970965887543
13.125 204.397304102318
13.15 204.823642317093
13.175 205.249980531868
13.2 205.676318746643
13.225 206.102656961418
13.25 206.528995176193
13.275 206.955333390968
13.3 207.381671605743
13.325 207.808009820518
13.35 208.234348035293
13.375 208.660686250068
13.4 209.087024464843
13.425 209.513362679619
13.45 209.939700894394
13.475 210.366039109169
13.5 210.792377323944
13.525 211.218715538719
13.55 211.645053753494
13.575 212.071391968269
13.6 212.497730183044
13.625 212.924068397819
13.65 213.350406612594
13.675 213.776744827369
13.7 214.203083042144
13.725 214.629421256919
13.75 215.055759471694
13.775 215.482097686469
13.8 215.908435901244
13.825 216.334774116019
13.85 216.761112330794
13.875 217.187450545569
13.9 217.613788760344
13.925 218.04012697512
13.95 218.466465189895
13.975 218.89280340467
14 219.319141619445
14.025 219.74547983422
14.05 220.171818048995
14.075 220.59815626377
14.1 221.024494478545
14.125 221.45083269332
14.15 221.877170908095
14.175 222.30350912287
14.2 222.729847337645
14.225 223.15618555242
14.25 223.582523767195
14.275 224.00886198197
14.3 224.435200196745
14.325 224.86153841152
14.35 225.287876626295
14.375 225.71421484107
14.4 226.140553055845
14.425 226.56689127062
14.45 226.993229485396
14.475 227.419567700171
14.5 227.845905914946
14.525 228.272244129721
14.55 228.698582344496
14.575 229.124920559271
14.6 229.551258774046
14.625 229.977596988821
14.65 230.403935203596
14.675 230.830273418371
14.7 231.256611633146
14.725 231.682949847921
14.75 232.109288062696
14.775 232.535626277471
14.8 232.961964492246
14.825 233.388302707021
14.85 233.814640921796
14.875 234.240979136571
14.9 234.667317351346
14.925 235.093655566121
14.95 235.519993780897
14.975 235.946331995672
15 236.372670210447
15.025 236.799008425222
15.05 237.225346639997
15.075 237.651684854772
15.1 238.078023069547
15.125 238.504361284322
15.15 238.930699499097
15.175 239.357037713872
15.2 239.783375928647
15.225 240.209714143422
15.25 240.636052358197
15.275 241.062390572972
15.3 241.488728787747
15.325 241.915067002522
15.35 242.341405217297
15.375 242.767743432072
15.4 243.194081646847
15.425 243.620419861622
15.45 244.046758076398
15.475 244.473096291173
15.5 244.899434505948
15.525 245.325772720723
15.55 245.752110935498
15.575 246.178449150273
15.6 246.604787365048
15.625 247.031125579823
15.65 247.457463794598
15.675 247.883802009373
15.7 248.310140224148
15.725 248.736478438923
15.75 249.162816653698
15.775 249.589154868473
15.8 250.015493083248
15.825 250.441831298023
15.85 250.868169512798
15.875 251.294507727573
15.9 251.720845942348
15.925 252.147184157123
15.95 252.573522371898
15.975 252.999860586674
16 253.426198801449
16.025 253.852537016224
16.05 254.278875230999
16.075 254.705213445774
16.1 255.131551660549
16.125 255.557889875324
16.15 255.984228090099
16.175 256.410566304874
16.2 256.836904519649
16.225 257.263242734424
16.25 257.689580949199
16.275 258.115919163974
16.3 258.542257378749
16.325 258.968595593524
16.35 259.394933808299
16.375 259.821272023074
16.4 260.247610237849
16.425 260.673948452624
16.45 261.100286667399
16.475 261.526624882174
16.5 261.95296309695
16.525 262.379301311725
16.55 262.8056395265
16.575 263.231977741275
16.6 263.65831595605
16.625 264.084654170825
16.65 264.5109923856
16.675 264.937330600375
16.7 265.36366881515
16.725 265.790007029925
16.75 266.2163452447
16.775 266.642683459475
16.8 267.06902167425
16.825 267.495359889025
16.85 267.9216981038
16.875 268.348036318575
16.9 268.77437453335
16.925 269.200712748125
16.95 269.6270509629
16.975 270.053389177675
17 270.479727392451
17.025 270.906065607226
17.05 271.332403822001
17.075 271.758742036776
17.1 272.185080251551
17.125 272.611418466326
17.15 273.037756681101
17.175 273.464094895876
17.2 273.890433110651
17.225 274.316771325426
17.25 274.743109540201
17.275 275.169447754976
17.3 275.595785969751
17.325 276.022124184526
17.35 276.448462399301
17.375 276.874800614076
17.4 277.301138828851
17.425 277.727477043626
17.45 278.153815258401
17.475 278.580153473176
17.5 279.006491687951
17.525 279.432829902727
17.55 279.859168117502
17.575 280.285506332277
17.6 280.711844547052
17.625 281.138182761827
17.65 281.564520976602
17.675 281.990859191377
17.7 282.417197406152
17.725 282.843535620927
17.75 283.269873835702
17.775 283.696212050477
17.8 284.122550265252
17.825 284.548888480027
17.85 284.975226694802
17.875 285.401564909577
17.9 285.827903124352
17.925 286.254241339127
17.95 286.680579553902
17.975 287.106917768677
18 287.538669416145
18.025 287.967710856429
18.05 288.397433837789
18.075 288.827833545514
18.1 289.258905164891
18.125 289.69064388121
18.15 290.123044879757
18.175 290.556103345822
18.2 290.989814464693
18.225 291.424173421658
18.25 291.859175402004
18.275 292.294815591022
18.3 292.731089173997
18.325 293.16799133622
18.35 293.605517262977
18.375 294.043662139558
18.4 294.482421151251
18.425 294.921789483343
18.45 295.361762321123
18.475 295.802334849879
18.5 296.2435022549
18.525 296.685259721473
18.55 297.127602434887
18.575 297.570525580431
18.6 298.014024343391
18.625 298.458093909058
18.65 298.902729462718
18.675 299.34792618966
18.7 299.793679275172
18.725 300.239983904543
18.75 300.686835263061
18.775 301.134228536013
18.8 301.582158908689
18.825 302.030621566376
18.85 302.479611694363
18.875 302.929124477938
18.9 303.379155102388
18.925 303.829698753004
18.95 304.280750615071
18.975 304.73230587388
19 305.184359714717
19.025 305.636907322872
19.05 306.089943883632
19.075 306.543464582287
19.1 306.997464604123
19.125 307.451939134429
19.15 307.906883358493
19.175 308.362292461604
19.2 308.818161629051
19.225 309.27448604612
19.25 309.7312608981
19.275 310.18848137028
19.3 310.646142647948
19.325 311.104239916392
19.35 311.5627683609
19.375 312.021723166761
19.4 312.481099519262
19.425 312.940892603692
19.45 313.40109760534
19.475 313.861709709493
19.5 314.322724101439
19.525 314.784135966467
19.55 315.245940489866
19.575 315.708132856922
19.6 316.170708252926
19.625 316.633661863164
19.65 317.096988872925
19.675 317.560684467497
19.7 318.024743832169
19.725 318.489162152229
19.75 318.953934612964
19.775 319.419056399664
19.8 319.884522697616
19.825 320.350328692108
19.85 320.81646956843
19.875 321.282940511868
19.9 321.749736707712
19.925 322.216853341249
19.95 322.684285597769
19.975 323.152028662558
20 323.620077720905
20.025 324.088427958099
20.05 324.557074559428
20.075 325.026012710179
20.1 325.495237595642
20.125 325.964744401104
20.15 326.434528311853
20.175 326.904584513179
20.2 327.374908190369
20.225 327.84549452871
20.25 328.316338713493
20.275 328.787435930004
20.3 329.258781363532
20.325 329.730370199365
20.35 330.202197622792
20.375 330.674258819101
20.4 331.146548973579
20.425 331.619063271516
20.45 332.091796898199
20.475 332.564745038916
20.5 333.037902933258
20.525 333.511267069738
20.55 333.984835185797
20.575 334.458605073176
20.6 334.932574523617
20.625 335.40674132886
20.65 335.881103280647
20.675 336.355658170719
20.7 336.830403790818
20.725 337.305337932684
20.75 337.780458388059
20.775 338.255762948684
20.8 338.731249406301
20.825 339.206915552651
20.85 339.682759179474
20.875 340.158778078513
20.9 340.634970041508
20.925 341.111332860201
20.95 341.587864326333
20.975 342.064562231646
21 342.541424367879
21.025 343.018448526776
21.05 343.495632500077
21.075 343.972974079523
21.1 344.450471056855
21.125 344.928121223816
21.15 345.405922372146
21.175 345.883872293585
21.2 346.361968779877
21.225 346.840209622762
21.25 347.31859261398
21.275 347.797115545274
21.3 348.275776208385
21.325 348.754572395054
21.35 349.233501897022
21.375 349.71256250603
21.4 350.19175201382
21.425 350.671068212133
21.45 351.15050889271
21.475 351.630071847293
21.5 352.109754867623
21.525 352.58955574544
21.55 353.069472272487
21.575 353.549502240504
21.6 354.029643441234
21.625 354.509893666416
21.65 354.990250707793
21.675 355.470712357105
21.7 355.951276406094
21.725 356.431940646501
21.75 356.912702870068
21.775 357.393560868535
21.8 357.874512433644
21.825 358.355555357136
21.85 358.836687430753
21.875 359.317906446235
21.9 359.799210195325
21.925 360.280596469762
21.95 360.762063061289
21.975 361.243607761647
22 361.725228362577
22.025 362.20692265582
22.05 362.688688433117
22.075 363.170523486211
22.1 363.652425606841
22.125 364.134392586749
22.15 364.616422217678
22.175 365.098512291367
22.2 365.580660599558
22.225 366.062864933992
22.25 366.545123086411
22.275 367.027432848556
22.3 367.509792012168
22.325 367.992198368988
22.35 368.474649710758
22.375 368.957143829219
22.4 369.439678516112
22.425 369.922251563179
22.45 370.40486076216
22.475 370.887503904797
22.5 371.370178782832
22.525 371.852883188004
22.55 372.335614912057
22.575 372.818371746731
22.6 373.301151483766
22.625 373.783951914906
22.65 374.26677083189
22.675 374.74960602646
22.7 375.232455290358
22.725 375.715316415324
22.75 376.198187193099
22.775 376.681065415426
22.8 377.163948874046
22.825 377.646835360699
22.85 378.129722667126
22.875 378.61260858507
22.9 379.095490906272
22.925 379.578367422472
22.95 380.061235925412
22.975 380.544094206833
23 381.02834605405
23.025 381.514006132291
23.05 381.999666210531
23.075 382.485326288772
23.1 382.970986367013
23.125 383.456646445254
23.15 383.942306523495
23.175 384.427966601735
23.2 384.913626679976
23.225 385.399286758217
23.25 385.884946836458
23.275 386.370606914698
23.3 386.856266992939
23.325 387.34192707118
23.35 387.827587149421
23.375 388.313247227662
23.4 388.798907305902
23.425 389.284567384143
23.45 389.770227462384
23.475 390.255887540625
23.5 390.741547618865
23.525 391.227207697106
23.55 391.712867775347
23.575 392.198527853588
23.6 392.684187931829
23.625 393.169848010069
23.65 393.65550808831
23.675 394.141168166551
23.7 394.626828244792
23.725 395.112488323033
23.75 395.598148401273
23.775 396.083808479514
23.8 396.569468557755
23.825 397.055128635996
23.85 397.540788714236
23.875 398.026448792477
23.9 398.512108870718
23.925 398.997768948959
23.95 399.4834290272
23.975 399.96908910544
24 400.454749183681
24.025 400.940409261922
24.05 401.426069340163
24.075 401.911729418404
24.1 402.397389496644
24.125 402.883049574885
24.15 403.368709653126
24.175 403.854369731367
24.2 404.340029809608
24.225 404.825689887848
24.25 405.311349966089
24.275 405.79701004433
24.3 406.282670122571
24.325 406.768330200811
24.35 407.253990279052
24.375 407.739650357293
24.4 408.225310435534
24.425 408.710970513775
24.45 409.196630592015
24.475 409.682290670256
24.5 410.167950748497
24.525 410.653610826738
24.55 411.139270904978
24.575 411.624930983219
24.6 412.11059106146
24.625 412.596251139701
24.65 413.081911217942
24.675 413.567571296182
24.7 414.053231374423
24.725 414.538891452664
24.75 415.024551530905
24.775 415.510211609146
24.8 415.995871687386
24.825 416.481531765627
24.85 416.967191843868
24.875 417.452851922109
24.9 417.938512000349
24.925 418.42417207859
24.95 418.909832156831
24.975 419.395492235072
25 419.881152313313
25.025 420.366812391553
25.05 420.852472469794
25.075 421.338132548035
25.1 421.823792626276
25.125 422.309452704517
25.15 422.795112782757
25.175 423.280772860998
25.2 423.766432939239
25.225 424.25209301748
25.25 424.73775309572
25.275 425.223413173961
25.3 425.709073252202
25.325 426.194733330443
25.35 426.680393408684
25.375 427.166053486924
25.4 427.651713565165
25.425 428.137373643406
25.45 428.623033721647
25.475 429.108693799887
25.5 429.594353878128
25.525 430.080013956369
25.55 430.56567403461
25.575 431.051334112851
25.6 431.536994191091
25.625 432.022654269332
25.65 432.508314347573
25.675 432.993974425814
25.7 433.479634504055
25.725 433.965294582295
25.75 434.450954660536
25.775 434.936614738777
25.8 435.422274817018
25.825 435.907934895259
25.85 436.393594973499
25.875 436.87925505174
25.9 437.364915129981
25.925 437.850575208222
25.95 438.336235286462
25.975 438.821895364703
26 439.307555442944
26.025 439.793215521185
26.05 440.278875599426
26.075 440.764535677666
26.1 441.250195755907
26.125 441.735855834148
26.15 442.221515912389
26.175 442.707175990629
26.2 443.19283606887
26.225 443.678496147111
26.25 444.164156225352
26.275 444.649816303593
26.3 445.135476381833
26.325 445.621136460074
26.35 446.106796538315
26.375 446.592456616556
26.4 447.078116694797
26.425 447.563776773037
26.45 448.049436851278
26.475 448.535096929519
26.5 449.02075700776
26.525 449.506417086
26.55 449.992077164241
26.575 450.477737242482
26.6 450.963397320723
26.625 451.449057398964
26.65 451.934717477204
26.675 452.420377555445
26.7 452.906037633686
26.725 453.391697711927
26.75 453.877357790167
26.775 454.363017868408
26.8 454.848677946649
26.825 455.33433802489
26.85 455.819998103131
26.875 456.305658181371
26.9 456.791318259612
26.925 457.276978337853
26.95 457.762638416094
26.975 458.248298494335
27 458.733958572575
27.025 459.219618650816
27.05 459.705278729057
27.075 460.190938807298
27.1 460.676598885539
27.125 461.162258963779
27.15 461.64791904202
27.175 462.133579120261
27.2 462.619239198502
27.225 463.104899276742
27.25 463.590559354983
27.275 464.076219433224
27.3 464.561879511465
27.325 465.047539589706
27.35 465.533199667946
27.375 466.018859746187
27.4 466.504519824428
27.425 466.990179902669
27.45 467.47583998091
27.475 467.96150005915
27.5 468.447160137391
27.525 468.932820215632
27.55 469.418480293873
27.575 469.904140372113
27.6 470.389800450354
27.625 470.875460528595
27.65 471.361120606836
27.675 471.846780685077
27.7 472.332440763317
27.725 472.818100841558
27.75 473.303760919799
27.775 473.78942099804
27.8 474.27508107628
27.825 474.760741154521
27.85 475.246401232762
27.875 475.732061311003
27.9 476.217721389244
27.925 476.703381467484
27.95 477.189041545725
27.975 477.674701623966
28 478.160361702207
28.025 478.646021780448
28.05 479.131681858688
28.075 479.617341936929
28.1 480.10300201517
28.125 480.588662093411
28.15 481.074322171651
28.175 481.559982249892
28.2 482.045642328133
28.225 482.531302406374
28.25 483.016962484615
28.275 483.502622562855
28.3 483.988282641096
28.325 484.473942719337
28.35 484.959602797578
28.375 485.445262875818
28.4 485.930922954059
28.425 486.4165830323
28.45 486.902243110541
28.475 487.387903188782
28.5 487.873563267022
28.525 488.359223345263
28.55 488.844883423504
28.575 489.330543501745
28.6 489.816203579986
28.625 490.301863658226
28.65 490.787523736467
28.675 491.273183814708
28.7 491.758843892949
28.725 492.244503971189
28.75 492.73016404943
28.775 493.215824127671
28.8 493.701484205912
28.825 494.187144284153
28.85 494.672804362393
28.875 495.158464440634
28.9 495.644124518875
28.925 496.129784597116
28.95 496.615444675357
28.975 497.101104753597
29 497.586764831838
29.025 498.072424910079
29.05 498.55808498832
29.075 499.043745066561
29.1 499.529405144801
29.125 500.015065223042
29.15 500.500725301283
29.175 500.986385379524
29.2 501.472045457764
29.225 501.957705536005
29.25 502.443365614246
29.275 502.929025692487
29.3 503.414685770728
29.325 503.900345848968
29.35 504.386005927209
29.375 504.87166600545
29.4 505.357326083691
29.425 505.842986161931
29.45 506.328646240172
29.475 506.814306318413
29.5 507.299966396654
29.525 507.785626474895
29.55 508.271286553135
29.575 508.756946631376
29.6 509.242606709617
29.625 509.728266787858
29.65 510.213926866099
29.675 510.699586944339
29.7 511.18524702258
29.725 511.670907100821
29.75 512.156567179062
29.775 512.642227257302
29.8 513.127887335543
29.825 513.613547413784
29.85 514.099207492025
29.875 514.584867570266
29.9 515.070527648506
29.925 515.556187726747
29.95 516.041847804988
29.975 516.527507883229
30 517.013167961469
30.025 517.49882803971
30.05 517.984488117951
30.075 518.470148196192
30.1 518.955808274433
30.125 519.441468352673
30.15 519.927128430914
30.175 520.412788509155
30.2 520.898448587396
30.225 521.384108665637
30.25 521.869768743877
30.275 522.355428822118
30.3 522.841088900359
30.325 523.3267489786
30.35 523.81240905684
30.375 524.298069135081
30.4 524.783729213322
30.425 525.269389291563
30.45 525.755049369804
30.475 526.240709448044
30.5 526.726369526285
30.525 527.212029604526
30.55 527.697689682767
30.575 528.183349761008
30.6 528.669009839248
30.625 529.154669917489
30.65 529.64032999573
30.675 530.125990073971
30.7 530.611650152211
30.725 531.097310230452
30.75 531.582970308693
30.775 532.068630386934
30.8 532.554290465175
30.825 533.039950543415
30.85 533.525610621656
30.875 534.011270699897
30.9 534.496930778138
30.925 534.982590856378
30.95 535.468250934619
30.975 535.95391101286
31 536.439571091101
31.025 536.925231169342
31.05 537.410891247582
31.075 537.896551325823
31.1 538.382211404064
31.125 538.867871482305
31.15 539.353531560546
31.175 539.839191638786
31.2 540.324851717027
31.225 540.810511795268
31.25 541.296171873509
31.275 541.78183195175
31.3 542.26749202999
31.325 542.753152108231
31.35 543.238812186472
31.375 543.724472264713
31.4 544.210132342953
31.425 544.695792421194
31.45 545.181452499435
31.475 545.667112577676
31.5 546.152772655917
31.525 546.638432734157
31.55 547.124092812398
31.575 547.609752890639
31.6 548.09541296888
31.625 548.581073047121
31.65 549.066733125361
31.675 549.552393203602
31.7 550.038053281843
31.725 550.523713360084
31.75 551.009373438324
31.775 551.495033516565
31.8 551.980693594806
31.825 552.466353673047
31.85 552.952013751288
31.875 553.437673829528
31.9 553.923333907769
31.925 554.40899398601
31.95 554.894654064251
31.975 555.380314142491
32 555.865974220732
32.025 556.351634298973
32.05 556.837294377214
32.075 557.322954455455
32.1 557.808614533695
32.125 558.294274611936
32.15 558.779934690177
32.175 559.265594768418
32.2 559.751254846659
32.225 560.236914924899
32.25 560.72257500314
32.275 561.208235081381
32.3 561.693895159622
32.325 562.179555237862
32.35 562.665215316103
32.375 563.150875394344
32.4 563.636535472585
32.425 564.122195550826
32.45 564.607855629066
32.475 565.093515707307
32.5 565.579175785548
32.525 566.064835863789
32.55 566.55049594203
32.575 567.03615602027
32.6 567.521816098511
32.625 568.007476176752
32.65 568.493136254993
32.675 568.978796333234
32.7 569.464456411474
32.725 569.950116489715
32.75 570.435776567956
32.775 570.921436646197
32.8 571.407096724437
32.825 571.892756802678
32.85 572.378416880919
32.875 572.86407695916
32.9 573.3497370374
32.925 573.835397115641
32.95 574.321057193882
32.975 574.806717272123
33 575.292377350364
33.025 575.778037428604
33.05 576.263697506845
33.075 576.749357585086
33.1 577.235017663327
33.125 577.720677741568
33.15 578.206337819808
33.175 578.691997898049
33.2 579.17765797629
33.225 579.663318054531
33.25 580.148978132772
33.275 580.634638211012
33.3 581.120298289253
33.325 581.605958367494
33.35 582.091618445735
33.375 582.577278523975
33.4 583.062938602216
33.425 583.548598680457
33.45 584.034258758698
33.475 584.519918836939
33.5 585.005578915179
33.525 585.49123899342
33.55 585.976899071661
33.575 586.462559149902
33.6 586.948219228143
33.625 587.433879306383
33.65 587.919539384624
33.675 588.405199462865
33.7 588.890859541106
33.725 589.376519619346
33.75 589.862179697587
33.775 590.347839775828
33.8 590.833499854069
33.825 591.31915993231
33.85 591.80482001055
33.875 592.290480088791
33.9 592.776140167032
33.925 593.261800245273
33.95 593.747460323513
33.975 594.233120401754
34 594.718780479995
34.025 595.204440558236
34.05 595.690100636477
34.075 596.175760714717
34.1 596.661420792958
34.125 597.147080871199
34.15 597.63274094944
34.175 598.118401027681
34.2 598.604061105921
34.225 599.089721184162
34.25 599.575381262403
34.275 600.061041340644
34.3 600.546701418885
34.325 601.032361497125
34.35 601.518021575366
34.375 602.003681653607
34.4 602.489341731848
34.425 602.975001810088
34.45 603.460661888329
34.475 603.94632196657
34.5 604.431982044811
34.525 604.917642123051
34.55 605.403302201292
34.575 605.888962279533
34.6 606.374622357774
34.625 606.860282436015
34.65 607.345942514255
34.675 607.831602592496
34.7 608.317262670737
34.725 608.802922748978
34.75 609.288582827219
34.775 609.774242905459
34.8 610.2599029837
34.825 610.745563061941
34.85 611.231223140182
34.875 611.716883218422
34.9 612.202543296663
34.925 612.688203374904
34.95 613.173863453145
34.975 613.659523531386
35 614.145183609626
35.025 614.630843687867
35.05 615.116503766108
35.075 615.602163844349
35.1 616.08782392259
35.125 616.57348400083
35.15 617.059144079071
35.175 617.544804157312
35.2 618.030464235553
35.225 618.516124313794
35.25 619.001784392034
35.275 619.487444470275
35.3 619.973104548516
35.325 620.458764626757
35.35 620.944424704997
35.375 621.430084783238
35.4 621.915744861479
35.425 622.40140493972
35.45 622.887065017961
35.475 623.372725096201
35.5 623.858385174442
35.525 624.344045252683
35.55 624.829705330924
35.575 625.315365409165
35.6 625.801025487405
35.625 626.286685565646
35.65 626.772345643887
35.675 627.258005722128
35.7 627.743665800368
35.725 628.229325878609
35.75 628.71498595685
35.775 629.200646035091
35.8 629.686306113332
35.825 630.171966191572
35.85 630.657626269813
35.875 631.143286348054
35.9 631.628946426295
35.925 632.114606504535
35.95 632.600266582776
35.975 633.085926661017
36 633.571586739258
36.025 634.057246817499
36.05 634.542906895739
36.075 635.02856697398
36.1 635.514227052221
36.125 635.999887130462
36.15 636.485547208702
36.175 636.971207286943
36.2 637.456867365184
36.225 637.942527443425
36.25 638.428187521666
36.275 638.913847599906
36.3 639.399507678147
36.325 639.885167756388
36.35 640.370827834629
36.375 640.85648791287
36.4 641.34214799111
36.425 641.827808069351
36.45 642.313468147592
36.475 642.799128225833
36.5 643.284788304074
36.525 643.770448382314
36.55 644.256108460555
36.575 644.741768538796
36.6 645.227428617037
36.625 645.713088695277
36.65 646.198748773518
36.675 646.684408851759
36.7 647.17006893
36.725 647.655729008241
36.75 648.141389086481
36.775 648.627049164722
36.8 649.112709242963
36.825 649.598369321204
36.85 650.084029399444
36.875 650.569689477685
36.9 651.055349555926
36.925 651.541009634167
36.95 652.026669712408
36.975 652.512329790648
37 652.997989868889
37.025 653.48364994713
37.05 653.969310025371
37.075 654.454970103612
37.1 654.940630181852
37.125 655.426290260093
37.15 655.911950338334
37.175 656.397610416575
37.2 656.883270494816
37.225 657.368930573056
37.25 657.854590651297
37.275 658.340250729538
37.3 658.825910807779
37.325 659.311570886019
37.35 659.79723096426
37.375 660.282891042501
37.4 660.768551120742
37.425 661.254211198983
37.45 661.739871277223
37.475 662.225531355464
37.5 662.711191433705
37.525 663.196851511946
37.55 663.682511590187
37.575 664.168171668427
37.6 664.653831746668
37.625 665.139491824909
37.65 665.62515190315
37.675 666.11081198139
37.7 666.596472059631
37.725 667.082132137872
37.75 667.567792216113
37.775 668.053452294354
37.8 668.539112372594
37.825 669.024772450835
37.85 669.510432529076
37.875 669.996092607317
37.9 670.481752685557
37.925 670.967412763798
37.95 671.453072842039
37.975 671.93873292028
38 672.42439299852
38.025 672.910053076761
38.05 673.395713155002
38.075 673.881373233243
38.1 674.367033311484
38.125 674.852693389724
38.15 675.338353467965
38.175 675.824013546206
38.2 676.309673624447
38.225 676.795333702688
38.25 677.280993780928
38.275 677.766653859169
38.3 678.25231393741
38.325 678.737974015651
38.35 679.223634093892
38.375 679.709294172132
38.4 680.194954250373
38.425 680.680614328614
38.45 681.166274406855
38.475 681.651934485096
38.5 682.137594563336
38.525 682.623254641577
38.55 683.108914719818
38.575 683.594574798059
38.6 684.080234876299
38.625 684.56589495454
38.65 685.051555032781
38.675 685.537215111022
38.7 686.022875189263
38.725 686.508535267503
38.75 686.994195345744
38.775 687.479855423985
38.8 687.965515502226
38.825 688.451175580466
38.85 688.936835658707
38.875 689.422495736948
38.9 689.908155815189
38.925 690.39381589343
38.95 690.87947597167
38.975 691.365136049911
39 691.850796128152
39.025 692.336456206393
39.05 692.822116284634
39.075 693.307776362874
39.1 693.793436441115
39.125 694.279096519356
39.15 694.764756597597
39.175 695.250416675838
39.2 695.736076754078
39.225 696.221736832319
39.25 696.70739691056
39.275 697.193056988801
39.3 697.678717067041
39.325 698.164377145282
39.35 698.650037223523
39.375 699.135697301764
39.4 699.621357380005
39.425 700.107017458245
39.45 700.592677536486
39.475 701.078337614727
39.5 701.563997692968
39.525 702.049657771209
39.55 702.535317849449
39.575 703.02097792769
39.6 703.506638005931
39.625 703.992298084172
39.65 704.477958162413
39.675 704.963618240653
39.7 705.449278318894
39.725 705.934938397135
39.75 706.420598475376
39.775 706.906258553616
39.8 707.391918631857
39.825 707.877578710098
39.85 708.363238788339
39.875 708.848898866579
39.9 709.33455894482
39.925 709.820219023061
39.95 710.305879101302
39.975 710.791539179543
40 711.277199257783
40.025 711.762859336024
40.05 712.248519414265
40.075 712.734179492506
40.1 713.219839570746
40.125 713.705499648987
40.15 714.191159727228
40.175 714.676819805469
40.2 715.16247988371
40.225 715.64813996195
40.25 716.133800040191
40.275 716.619460118432
40.3 717.105120196673
40.325 717.590780274914
40.35 718.076440353154
40.375 718.562100431395
40.4 719.047760509636
40.425 719.533420587877
40.45 720.019080666118
40.475 720.504740744358
40.5 720.990400822599
40.525 721.47606090084
40.55 721.961720979081
40.575 722.447381057321
40.6 722.933041135562
40.625 723.418701213803
40.65 723.904361292044
40.675 724.390021370285
40.7 724.875681448525
40.725 725.361341526766
40.75 725.847001605007
40.775 726.332661683248
40.8 726.818321761489
40.825 727.303981839729
40.85 727.78964191797
40.875 728.275301996211
40.9 728.760962074452
40.925 729.246622152692
40.95 729.732282230933
40.975 730.217942309174
41 730.703602387415
41.025 731.189262465656
41.05 731.674922543896
41.075 732.160582622137
41.1 732.646242700378
41.125 733.131902778619
41.15 733.61756285686
41.175 734.1032229351
41.2 734.588883013341
41.225 735.074543091582
41.25 735.560203169823
41.275 736.045863248063
41.3 736.531523326304
41.325 737.017183404545
41.35 737.502843482786
41.375 737.988503561027
41.4 738.474163639267
41.425 738.959823717508
41.45 739.445483795749
41.475 739.93114387399
41.5 740.41680395223
41.525 740.902464030471
41.55 741.388124108712
41.575 741.873784186953
41.6 742.359444265194
41.625 742.845104343434
41.65 743.330764421675
41.675 743.816424499916
41.7 744.302084578157
41.725 744.787744656398
41.75 745.273404734638
41.775 745.759064812879
41.8 746.24472489112
41.825 746.730384969361
41.85 747.216045047601
41.875 747.701705125842
41.9 748.187365204083
41.925 748.673025282324
41.95 749.158685360565
41.975 749.644345438805
42 750.130005517046
42.025 750.615665595287
42.05 751.101325673528
42.075 751.586985751768
42.1 752.072645830009
42.125 752.55830590825
42.15 753.043965986491
42.175 753.529626064732
42.2 754.015286142972
42.225 754.500946221213
42.25 754.986606299454
42.275 755.472266377695
42.3 755.957926455936
42.325 756.443586534176
};
\addplot [ultra thick, black, forget plot]
table {%
9 -37.8221793267088
9 794.265765860885
};
\addplot [ultra thick, black, forget plot]
table {%
11 -37.8221793267088
11 794.265765860885
};
\addplot [ultra thick, black, forget plot]
table {%
18 -37.8221793267088
18 794.265765860885
};
\addplot [ultra thick, black, forget plot]
table {%
23 -37.8221793267088
23 794.265765860885
};
\node at (axis cs:4.5,40)[
  anchor=base,
  text=black,
  rotate=0.0
]{ A};
\node at (axis cs:10,40)[
  anchor=base,
  text=black,
  rotate=0.0
]{ B};
\node at (axis cs:14.5,40)[
  anchor=base,
  text=black,
  rotate=0.0
]{ C};
\node at (axis cs:20.5,40)[
  anchor=base,
  text=black,
  rotate=0.0
]{ D};
\node at (axis cs:31.5,40)[
  anchor=base,
  text=black,
  rotate=0.0
]{ E};
\end{axis}

\end{tikzpicture}%
		\subcaption{Longitudinal position.}
		\label{fig:example ego states lon pos}
	\end{minipage}%
	\hfill
	\begin{minipage}[b]{0.49\linewidth}
		\centering
		% This file was created by matplotlib2tikz v0.6.14.
\begin{tikzpicture}

\definecolor{color1}{rgb}{1,0.498039215686275,0.0549019607843137}
\definecolor{color0}{rgb}{0.12156862745098,0.466666666666667,0.705882352941177}

\begin{axis}[
xlabel={$t$ [s]},
ylabel={Acceleration [m/s$^2$]},
xmin=0, xmax=40,
ymin=-1.90619599582758, ymax=2.38951469489845,
width=\figurewidth,
height=\figureheight,
tick align=outside,
tick pos=left,
xmajorgrids,
x grid style={white!69.019607843137251!black},
ymajorgrids,
y grid style={white!69.019607843137251!black},
clip marker paths
]
\addplot [semithick, color0, mark=*, mark size=1, mark options={solid}, only marks, forget plot]
table {%
0 -0.0019836657013883
0.025 0.558029076401425
0.05 1.32171431854022
0.075 1.30367666158453
0.1 1.496801347221
0.125 1.79441637298755
0.15 1.32108684158922
0.175 0.998795410533973
0.2 1.01630904205699
0.225 1.3888361775544
0.25 1.6860511352785
0.275 1.15348789128752
0.3 0.940072046620024
0.325 1.16410062944787
0.35 1.29928871923219
0.375 1.54827208012001
0.4 1.42928491335592
0.425 1.2651082722796
0.45 1.3303663468573
0.475 1.49058351856866
0.5 1.59063832293757
0.525 1.14079007615945
0.55 1.04384417046077
0.575 1.39915423406208
0.6 1.1318500118216
0.625 0.872492238372972
0.65 0.961560537265384
0.675 1.11137853934728
0.7 1.49427492146131
0.725 1.52615916145376
0.75 1.18888840519013
0.775 1.08409131865468
0.8 0.94832997244083
0.825 1.05563045346368
0.85 1.28739654045145
0.875 1.07465569072591
0.9 1.16781048382055
0.925 1.38084983949085
0.95 1.20034596902112
0.975 1.09860710750596
1 0.953082755840916
1.025 1.14818078808799
1.05 1.48010843879476
1.075 1.17251012372431
1.1 0.866240656691274
1.125 0.714290022613539
1.15 1.03549382933323
1.175 1.7085846591705
1.2 1.66283508872418
1.225 1.2320985474851
1.25 0.825186706917335
1.275 0.608428986063813
1.3 0.828233787655818
1.325 0.976955421277025
1.35 1.05807877665001
1.375 1.18125711833521
1.4 0.887239278944207
1.425 0.632455796528574
1.45 0.727050776802329
1.475 1.01085092524045
1.5 1.4015095500671
1.525 1.31025439865908
1.55 0.843988390876404
1.575 0.857271084796806
1.6 1.23945137335824
1.625 1.59361739967245
1.65 1.58458143005475
1.675 1.03522968108649
1.7 0.938759533053392
1.725 1.13422851700023
1.75 0.949586955146095
1.775 0.823277508708031
1.8 0.712254461132033
1.825 0.821261652151541
1.85 1.18252779631405
1.875 1.48916618624387
1.9 1.4446331927715
1.925 0.794436311707613
1.95 0.800908278247582
1.975 1.24613479848503
2 0.821720060942862
2.025 0.576506665081557
2.05 0.840231812392136
2.075 0.98536955224461
2.1 1.39331822969812
2.125 1.60738190501118
2.15 1.20516225204507
2.175 0.835417854107763
2.2 0.567175679494804
2.225 0.562083234640911
2.25 0.801321172263239
2.275 0.973069002333732
2.3 1.32812868837905
2.325 1.41210204125973
2.35 0.984272830942782
2.375 0.685062937988864
2.4 0.618595196697843
2.425 0.818643705970071
2.45 1.100136736477
2.475 1.32646281808399
2.5 1.40477545348481
2.525 0.892904217657029
2.55 0.690536828713102
2.575 0.968025152501223
2.6 0.811852236735898
2.625 0.80817902069771
2.65 1.04611709431832
2.675 1.08733952037266
2.7 0.981438462234511
2.725 0.577244611875935
2.75 0.570156053772057
2.775 1.09411721891444
2.8 1.26789871397318
2.825 1.32975112252325
2.85 1.23104368736378
2.875 0.559682605865387
2.9 0.642667792288446
2.925 1.41371263252055
2.95 1.34516818130806
2.975 0.787588120798406
3 0.270555491400584
3.025 0.620640768542128
3.05 1.40558139149463
3.075 1.16169843142587
3.1 0.668228250944909
3.125 0.463069696121889
3.15 0.800405470087244
3.175 1.40192716898998
3.2 1.27442545586974
3.225 0.752680669319619
3.25 0.420464560068365
3.275 0.767787875880895
3.3 1.24616416093453
3.325 0.965108952203767
3.35 0.605003394594308
3.375 0.792353540964892
3.4 1.2063000022632
3.425 1.80418119969055
3.45 1.38033260026305
3.475 0.0467385129727662
3.5 0.355425091121118
3.525 1.11629392571863
3.55 0.969109413463446
3.575 0.848215310730896
3.6 0.836484020672724
3.625 0.935705600020102
3.65 0.813053974237898
3.675 0.450126181155214
3.7 0.731025119466724
3.725 1.29669103280918
3.75 0.999138769685993
3.775 0.570503154739012
3.8 0.683880643302077
3.825 0.959365671880622
3.85 1.10126117556348
3.875 1.0861112385416
3.9 1.12056624104724
3.925 0.898996199498189
3.95 0.605784594048944
3.975 0.663471297488878
4 1.07894525222604
4.025 1.19701310731248
4.05 0.49840251475125
4.075 0.297733199752766
4.1 0.701241465722051
4.125 0.604379394823697
4.15 0.406808716297748
4.175 0.605935267617996
4.2 0.899223463298414
4.225 0.900628070652448
4.25 0.982552412884278
4.275 1.0811521311894
4.3 1.10343951726624
4.325 1.10343915871168
4.35 0.508144284423937
4.375 0.0355106622265568
4.4 0.0547362834907974
4.425 1.03096424391861
4.45 2.19425511804726
4.475 1.65970049019923
4.5 0.950553452261719
4.525 0.881513398257997
4.55 0.758457239371427
4.575 0.558456605259031
4.6 0.738944946760614
4.625 1.06081205012174
4.65 0.624939077889055
4.675 0.48571760883963
4.7 0.982648937522854
4.725 1.18059942517985
4.75 0.979148719960964
4.775 0.705396463726505
4.8 0.585376740604888
4.825 0.686584633061218
4.85 0.816886962940127
4.875 0.607377810854646
4.9 0.818208606128212
4.925 1.36275304484717
4.95 1.43751341565519
4.975 1.20173343862007
5 0.605838720736289
5.025 0.633660522453274
5.05 1.23512540853291
5.075 1.10760103362395
5.1 0.624409906396297
5.125 0.32366069353902
5.15 0.397451045637348
5.175 0.790592368963242
5.2 0.861741446988165
5.225 0.46322824113645
5.25 0.589479337412657
5.275 1.01781640625516
5.3 0.772673508225668
5.325 0.527841755871352
5.35 0.786482924519308
5.375 1.08023267640658
5.4 0.978740749349058
5.425 1.04463465385145
5.45 1.11692054124437
5.475 0.434111659683369
5.5 0.145726986613753
5.525 0.683904636957767
5.55 1.21753983752335
5.575 1.31754047603181
5.6 0.726280497053724
5.625 0.193221698612015
5.65 0.666946020641025
5.675 1.19413003744736
5.7 0.79412847764502
5.725 0.503526544179069
5.75 0.602002139152093
5.775 0.750960437409265
5.8 0.86010944070836
5.825 0.790550165180551
5.85 0.875503421013949
5.875 1.19104298506727
5.9 1.41941443254721
5.925 1.04775633163286
5.95 0.707184947157205
5.975 0.746831592010724
6 0.6649073374944
6.025 0.726646431621692
6.05 0.828189949442123
6.075 0.659222736140921
6.1 0.367250391701647
6.125 0.427339483962827
6.15 0.748910957651095
6.175 0.929594131116648
6.2 0.914778840945022
6.225 0.516341222396479
6.25 0.386657506351115
6.275 0.58509921020117
6.3 0.96094319744687
6.325 1.33180327280229
6.35 0.973976236934963
6.375 0.819631725926089
6.4 1.11651895100527
6.425 1.26358021408941
6.45 1.15071489984189
6.475 0.836189139351976
6.5 0.66178992406073
6.525 0.512738474179734
6.55 0.610227584087613
6.575 0.81022821819943
6.6 0.883009821598083
6.625 0.857384553836201
6.65 0.281272915071891
6.675 0.0862658353756277
6.7 0.679370987325003
6.725 0.957734988111273
6.75 0.756134054381337
6.775 0.676324757159809
6.8 0.755207828467806
6.825 0.787888559915476
6.85 0.905764677514756
6.875 0.893550857267303
6.9 0.987563203284467
6.925 1.10947891549849
6.95 0.860297009666059
6.975 0.617837984748351
7 0.476228783412226
7.025 0.513954931483333
7.05 0.615571848613075
7.075 0.771739995485635
7.1 0.94148071077598
7.125 0.88425845281073
7.15 0.875955827331142
7.175 0.761434534070207
7.2 0.612916835393664
7.225 0.514551912448271
7.25 0.439376767308916
7.275 0.468324801369129
7.3 0.448734194248319
7.325 0.583229126300468
7.35 0.866716630679324
7.375 0.994855874148048
7.4 0.793216473971359
7.425 0.853517487366136
7.45 1.19573981978647
7.475 1.19984251237498
7.5 1.13458081324851
7.525 0.884542487420946
7.55 0.486741785716022
7.575 0.39840875865025
7.6 0.620908570766689
7.625 0.891776983427878
7.65 0.890117878236639
7.675 0.581394282682412
7.7 0.197967620721329
7.725 0.518243919401657
7.75 1.01740150017755
7.775 0.815733619492664
7.8 0.64410706089632
7.825 0.544106685890713
7.85 0.561071287319239
7.875 0.759391710374191
7.9 0.727729669659983
7.925 0.592379707132411
7.95 0.996087188701726
7.975 1.56598949039904
8 1.23118224357704
8.025 0.693174266648872
8.05 0.432359630396058
8.075 0.548471640514817
8.1 0.882472149429408
8.125 0.769795254488901
8.15 0.551139574304208
8.175 0.352838456451673
8.2 0.307151768534062
8.225 0.463576471678247
8.25 0.787434836023297
8.275 1.03823163094026
8.3 0.707219800115401
8.325 0.447528812937046
8.35 0.449240678794648
8.375 0.606847856759876
8.4 0.805138753219114
8.425 0.926438347012479
8.45 1.11611743274838
8.475 0.948641807467274
8.5 0.787870380723621
8.525 0.766733177090362
8.55 0.817509665468267
8.575 0.903285067251237
8.6 0.625330075238977
8.625 0.394262281188981
8.65 0.292531948969095
8.675 0.53816039346244
8.7 0.946431965651746
8.725 1.01616454392529
8.75 0.870550293494278
8.775 0.26265749743288
8.8 -0.0191850879998154
8.825 0.380816304051956
8.85 0.675093734171895
8.875 0.673347418545011
8.9 0.0753609652861076
8.925 -0.639984184713915
8.95 -0.921031775008762
8.975 -0.93403344334485
9 -0.531829883563396
9.025 -0.696155446291859
9.05 -1.36579514239498
9.075 -1.42803942733127
9.1 -0.929745641214288
9.125 0.0251882536289113
9.15 0.214183667071663
9.175 -0.384086653058015
9.2 -0.666952904235472
9.225 -0.641404078375359
9.25 -0.178936917633689
9.275 -0.0821283935351846
9.3 -0.677638103410051
9.325 -0.957464173573917
9.35 -0.755748551127072
9.375 -0.861206885134879
9.4 -1.10528197329812
9.425 -1.16986187274362
9.45 -1.36463648055819
9.475 -1.33542751422193
9.5 -1.5440388330161
9.525 -1.71093641897639
9.55 -1.54166293126534
9.575 -1.46777816182625
9.6 -1.18962313534776
9.625 -1.06282735158752
9.65 -0.964514722103318
9.675 -0.659626873652916
9.7 -0.375767844831211
9.725 -0.73551184388984
9.75 -1.23174554949062
9.775 -0.71237649568835
9.8 -0.374763552053986
9.825 -0.6747645611643
9.85 -0.901186783709226
9.875 -0.902851367142096
9.9 -0.890694711621252
9.925 -0.865900697374388
9.95 -0.283383240774949
9.975 0.172023992457685
10 -0.635908141181161
10.025 -1.32587078075046
10.05 -0.621738455845309
10.075 -0.398887801110484
10.1 -0.76408591825066
10.125 -0.577159851753576
10.15 -0.455505280081987
10.175 -0.453853358226291
10.2 -0.542582590639064
10.225 -0.590335058676827
10.25 -0.0527019272121987
10.275 0.210057307162043
10.3 -0.284990129205194
10.325 -0.465608474882224
10.35 -0.163966268711707
10.375 -2.15539728287699e-07
10.4 -0.0598230089858012
10.425 -0.358176929421456
10.45 -0.659299781511029
10.475 -0.657661543527027
10.5 -0.490815891897409
10.525 0.0722869882216049
10.55 0.173039253934346
10.575 -0.555907755766888
10.6 -0.662567444077538
10.625 -0.158576897207187
10.65 0.239787206157015
10.675 0.0692463730956651
10.7 -0.33075519463551
10.725 -0.209573040591811
10.75 -0.0412474635604108
10.775 -0.131675695582485
10.8 -0.221307818594593
10.825 -0.221307678906723
10.85 -0.147018216886468
10.875 -0.0486516488544463
10.9 -0.100817765565749
10.925 -0.153046584978183
10.95 0.145321839000585
10.975 0.111160236202397
11 -0.589655786485194
11.025 -0.667447997467069
11.05 -0.124391613201167
11.075 -0.14924054012148
11.1 -0.228611952774505
11.125 0.160423248029605
11.15 -0.0785202391065087
11.175 -0.650301504742847
11.2 -0.396509499966407
11.225 -0.0480573139377337
11.25 -0.0284471006928306
11.275 -0.106305596364818
11.3 -0.3028166780723
11.325 -0.156482804754174
11.35 0.24514205638369
11.375 0.163089682995835
11.4 -0.1575403004729
11.425 -0.222792526756163
11.45 -0.184785047451183
11.475 -0.0793715135039025
11.5 -0.147437895839889
11.525 -0.337005751044468
11.55 -0.14050403655119
11.575 -0.016513595834766
11.6 -0.362329375763445
11.625 -0.591312541076418
11.65 -0.492936111937318
11.675 0.222519962392195
11.7 0.904706859137848
11.725 0.404495136093183
11.75 -0.161498933892549
11.775 -0.0838078592558122
11.8 -0.0674639633302822
11.825 -0.167464284458416
11.85 -0.221234255240437
11.875 -0.222855701623108
11.9 -0.111301363723014
11.925 -0.0486715700617417
11.95 -0.138992335597974
11.975 -0.159387478219489
12 -0.0610075633183284
12.025 0.175749640060727
12.05 0.347080210156032
12.075 0.092340357492303
12.1 -0.21908281946637
12.125 -0.255851590134887
12.15 -0.340516060669666
12.175 -0.413858517799775
12.2 -0.187480865462672
12.225 0.112685390741521
12.25 0.311065545542284
12.275 0.287314316528149
12.3 0.0832978380117299
12.325 -0.205637053419846
12.35 -0.363683883085039
12.375 -0.0653018453824941
12.4 0.169320354699918
12.425 0.0709389592657307
12.45 -0.0862676072181033
12.475 -0.284648677222072
12.5 -0.0940570314703868
12.525 0.424840825211697
12.55 0.600838654281679
12.575 0.459807417565366
12.6 0.181107436709031
12.625 0.0705832487589945
12.65 0.168961595692705
12.675 0.327235331235438
12.7 0.477526751024719
12.725 0.554098026578498
12.75 0.423417030503268
12.775 -0.180390262301697
12.8 -0.471528129480696
12.825 -0.171526866164837
12.85 -0.180835662524008
12.875 -0.321672517885226
12.9 0.241624728322872
12.925 0.603288550928474
12.95 -0.0791739961520921
12.975 -0.479991780706978
13 -0.181614621840968
13.025 -0.180983557743983
13.05 -0.37935856823878
13.075 -0.220640347605891
13.1 -0.019019214766303
13.125 4.37232472449987e-08
13.15 0.271963254726156
13.175 0.588639227453305
13.2 0.441470101021218
13.225 0.212986177997014
13.25 0.08494396271246
13.275 0.13234425905395
13.3 0.333968749612263
13.325 0.181167475592045
13.35 -0.128266028865912
13.375 -0.226200062559916
13.4 -0.170982396492289
13.425 0.0457822452290195
13.45 0.284658484884655
13.475 0.386286577327874
13.5 0.113466270908449
13.525 -0.301903644845676
13.55 -0.346586362161356
13.575 -0.114408349030979
13.6 0.2200623914753
13.625 -0.0664275290439789
13.65 -0.679388176798303
13.675 -0.549598505995156
13.7 -0.194392957495634
13.725 0.0454950381033835
13.75 0.226924736329442
13.775 0.225301672070796
13.8 0.140535305342835
13.825 0.0405348123997834
13.85 0.00162493524911159
13.875 0.0604821903600352
13.9 0.258857981747692
13.925 0.339517476265878
13.95 0.13951683711255
13.975 0.161054667456287
14 0.362789946620275
14.025 0.342309124743418
14.05 0.237209356555464
14.075 0.0966355578506026
14.1 0.149592342598556
14.125 0.303289639565853
14.15 -0.0942657151320945
14.175 -0.503289191783544
14.2 -0.255326307728151
14.225 0.00162886141268363
14.25 -3.02998870438393e-07
14.275 0.189449170488129
14.3 0.373390378456401
14.325 0.0912340735575144
14.35 -0.214810117301245
14.375 -0.225361457922522
14.4 -0.0213940269198784
14.425 0.280234136278064
14.45 0.262813563598097
14.475 0.0396596166532257
14.5 -0.224784726140901
14.525 -0.297952685731029
14.55 0.103675754849063
14.575 0.32110778009915
14.6 0.122739925394273
14.625 1.718027675679e-07
14.65 0.0295877312274935
14.675 0.131217963235029
14.7 -0.127166932030958
14.725 -0.599340602786226
14.75 -0.0778578572870731
14.775 0.578366292051405
14.8 -0.0159484145910935
14.825 -0.510243447205731
14.85 -0.208615855591892
14.875 0.158651594392015
14.9 0.358981616488302
14.925 0.343306537831717
14.95 0.210468521334729
14.975 -0.0292476316864316
15 -0.222749751186133
15.025 -0.224381401825081
15.05 -0.221840066330966
15.075 -0.221839763332096
15.1 -0.025594571172079
15.125 0.195869513435558
15.15 0.222198309915381
15.175 0.124706881058499
15.2 -0.0760277241389662
15.225 -0.23720543349477
15.25 -0.208554370613787
15.275 0.0204203363856689
15.3 0.212599813983158
15.325 0.220830890317171
15.35 0.150519462545947
15.375 0.0488880123604229
15.4 -0.00163104158524161
15.425 -0.0256423567975613
15.45 -0.125642634190298
15.475 -0.172727765526947
15.5 -0.0743570460834791
15.525 0.0704253616951424
15.55 0.172054453696544
15.575 0.225860659566802
15.6 0.209279134981775
15.625 -0.016128628056649
15.65 -0.181333649141704
15.675 -0.0817872271454678
15.7 0.0742690617024511
15.725 0.144192816743146
15.75 -0.105976050952779
15.775 -0.392314109468183
15.8 -0.316413531413668
15.825 -0.223057538652679
15.85 -0.22143002970481
15.875 0.157190084075157
15.9 0.518577338004171
15.925 0.0770045506101258
15.95 -0.349026454249918
15.975 -0.161386524200683
16 0.0714186736779244
16.025 0.169790220924495
16.05 0.219962150720233
16.075 0.189252171638469
16.1 -0.120461212316057
16.125 -0.399683155376346
16.15 -0.309931834645663
16.175 -0.034681209365319
16.2 0.265319666594546
16.225 0.275322961169873
16.25 0.0736924985858412
16.275 -0.00163003580688323
16.3 -0.0748874259215924
16.325 -0.174887619006438
16.35 -0.12348361158274
16.375 -0.0583027653647142
16.4 -0.134819880393804
16.425 -0.109155590803318
16.45 0.0908449609126892
16.475 0.222197263587983
16.5 0.218986673576179
16.525 0.0185671865907719
16.55 -0.23403086306645
16.575 -0.355809176336729
16.6 -0.23857111023062
16.625 0.15980144100922
16.65 0.0756768112766071
16.675 -0.422693315590574
16.7 -0.280555791870607
16.725 0.0974494403635603
16.75 -0.0227782381443827
16.775 -0.279010691859507
16.8 -0.378224635596801
16.825 -0.298154960116004
16.85 -0.156473915421882
16.875 -0.259941330227772
16.9 -0.34168155573155
16.925 -0.140058380134186
16.95 -8.74720740284816e-08
16.975 0.00162223345796517
17 -0.0994339217186266
17.025 -0.201056476292646
17.05 -0.100565586740237
17.075 0.0468131822330513
17.1 0.148999912357792
17.125 0.215550358348079
17.15 0.213929083840011
17.175 -0.257323476512298
17.2 -0.718851972148105
17.225 -0.165511896534551
17.25 0.398894382237808
17.275 0.302878253684464
17.3 0.00617020175468532
17.325 -0.354953523925867
17.35 -0.708645666054498
17.375 -0.689796538586042
17.4 -0.0422736347204727
17.425 0.513617412207197
17.45 0.613617951943075
17.475 0.196903117875988
17.5 -0.266700860629495
17.525 0.0292760176094471
17.55 0.398314514495746
17.575 0.305436193364979
17.6 0.630561777484209
17.625 1.22893505224509
17.65 0.743144088279795
17.675 -0.162626885162363
17.7 -0.324560143577979
17.725 -0.0532230210410482
17.75 0.462311643423157
17.775 0.769854622969166
17.8 0.571481839415388
17.825 0.5450917421497
17.85 0.679119800974978
17.875 0.460904762212444
17.9 0.41986673457074
17.925 0.648717265789892
17.95 1.3303501444153
17.975 1.6939315222771
18 0.317664836938998
18.025 -0.342490253340694
18.05 0.659151096147108
18.075 1.42237782315353
18.1 1.52403264605707
18.125 0.931199858291607
18.15 0.340872858170864
18.175 0.22518918115105
18.2 0.437208215292054
18.225 0.925003104111952
18.25 1.08420557442058
18.275 0.682546655277524
18.3 0.788007674941627
18.325 1.27357098965396
18.35 1.38035194013353
18.375 1.41060055502747
18.4 1.11581420010451
18.425 1.05096146560298
18.45 1.13965049613162
18.475 0.436361408560509
18.5 -0.0368099720390086
18.525 0.315520110373337
18.55 0.732998692524275
18.575 0.832999579003927
18.6 0.895268517842838
18.625 0.922927796091231
18.65 0.895194432885363
18.675 0.914076828189588
18.7 0.746544396335754
18.725 0.693009964652437
18.75 0.691306059343323
18.775 0.747059858670472
18.8 0.877710638474106
18.825 0.912979352344352
18.85 1.07606197327639
18.875 1.29031458363896
18.9 1.40415287569622
18.925 1.02122321552548
18.95 0.490219366767512
18.975 0.308794940974617
19 0.232207142380499
19.025 0.64565061352666
19.05 1.24736977656994
19.075 1.09475813626125
19.1 0.63417250278234
19.125 0.447981048887407
19.15 0.628463490306777
19.175 1.01988954475017
19.2 1.07663555980116
19.225 0.678379385021837
19.25 0.446948893656079
19.275 0.474828755932315
19.3 0.448459848166743
19.325 0.505152461577005
19.35 0.588059403435215
19.375 0.79902852255272
19.4 0.997283883248187
19.425 1.13943108656249
19.45 1.20164866348544
19.475 0.809582524357257
19.5 0.685636420981756
19.525 0.860245725334595
19.55 0.719719140647602
19.575 0.567815738580882
19.6 0.571835241680706
19.625 0.785865740655101
19.65 0.984101523636113
19.675 0.881033110340455
19.7 0.595714482077759
19.725 0.657519186233879
19.75 0.906905124675946
19.775 0.764070653963234
19.8 0.417621770293402
19.825 -0.0823802632983461
19.85 -0.151727689444678
19.875 0.24649229701097
19.9 0.75026870883902
19.925 1.00620156795229
19.95 0.257460736410589
19.975 0.0514215361818089
20 1.04632372317225
20.025 1.10629816306299
20.05 0.238834426068021
20.075 0.0221806198891269
20.1 0.264531152720353
20.125 0.371304522532791
20.15 0.454579471966092
20.175 0.456371271047822
20.2 0.440911900640793
20.225 0.442709940352895
20.25 0.200662918992404
20.275 0.140375063210738
20.3 0.441507207551688
20.325 0.529178903688666
20.35 0.330981445426862
20.375 0.320126415268045
20.4 0.456032567694813
20.425 0.446889813679257
20.45 0.502420495926891
20.475 0.595044221813657
20.5 0.469056087577322
20.525 0.199397835967827
20.55 0.232147187807638
20.575 0.367851418836546
20.6 0.169662550047188
20.625 0.0700150158721148
20.65 0.168205265799768
20.675 0.329590615547199
20.7 0.464753034957248
20.725 0.437369322708062
20.75 0.365230644679565
20.775 0.163024190578651
20.8 0.20098976262943
20.825 0.500990641766982
20.85 0.581834656964053
20.875 0.465188927382976
20.9 0.102348362202063
20.925 -0.0486410040139156
20.95 0.332366959177998
20.975 0.484281574358675
21 0.182459029788544
21.025 0.18124953787719
21.05 0.254476580539627
21.075 -0.311666477769204
21.1 -0.431385265778797
21.125 0.242885497887283
21.15 0.269927933049981
21.175 -0.0202570051669679
21.2 0.301014980804975
21.225 0.56721419416586
21.25 0.365387828987943
21.275 0.138751786516876
21.3 0.0464994633603049
21.325 0.205923089224524
21.35 0.458173671063404
21.375 0.456344671756661
21.4 0.284492536206922
21.425 -0.0136784659162002
21.45 -0.148589219834319
21.475 -0.046761648513896
21.5 0.197498293568967
21.525 0.395671023557037
21.55 0.202501460512391
21.575 0.0525617032553775
21.6 0.151891548557623
21.625 0.24236392300281
21.65 0.209426986452851
21.675 -0.0235319068906389
21.7 -0.0108317885587893
21.725 0.383425769541667
21.75 0.607410259426606
21.775 0.505578933734789
21.8 0.04327601100158
21.825 -0.273996352457214
21.85 0.183390553820583
21.875 0.422545632813958
21.9 -0.0781171157768756
21.925 -0.235458543306773
21.95 0.0645425075891372
21.975 0.036498914278129
22 -0.129695970581521
22.025 0.128152919546663
22.05 0.146737069553635
22.075 -0.547612234886322
22.1 -0.560563937438051
22.125 0.254270648606433
22.15 0.489199767866914
22.175 0.257699986036638
22.2 0.083332808144475
22.225 0.216904650164919
22.25 0.515067243868465
22.275 0.383094787321164
22.3 0.122344971941855
22.325 0.237413286372569
22.35 0.451139358642675
22.375 0.449297193990639
22.4 0.31144732747634
22.425 0.113288681109223
22.45 -4.37192682056769e-08
22.475 0.00184272689693909
22.5 -4.37513847373332e-08
22.525 -0.00184277059815938
22.55 -8.43328962218936e-08
22.575 0.0634003352890024
22.6 0.165243167771232
22.625 -0.0468815982277697
22.65 -0.348725419712532
22.675 -0.216518551633268
22.7 0.059631419963039
22.725 0.18481914921723
22.75 0.230678568500334
22.775 0.275426122357199
22.8 0.353416532425394
22.825 0.398336363018851
22.85 0.30018326024221
22.875 0.141417513232012
22.9 0.0407886642278748
22.925 0.101216178677531
22.95 0.198782692068772
22.975 0.100630985480171
23 -0.0994756103457192
23.025 -0.301324164456669
23.05 -0.0774642500346079
23.075 0.422538086677585
23.1 0.181045736167117
23.125 -0.320802909964395
23.15 -0.2041055591917
23.175 0.0483848377371025
23.2 0.150643379042987
23.225 0.0593680603287828
23.25 -0.175570522974551
23.275 -0.243103202069502
23.3 -0.175072536330842
23.325 -0.0669075557605216
23.35 0.0807017422403124
23.375 0.178855628969143
23.4 0.119298069216709
23.425 0.0746338801718583
23.45 0.253490464961242
23.475 0.256221684116952
23.5 -0.1456260754869
23.525 -0.16879468189465
23.55 0.233051203622306
23.575 0.25908283865213
23.6 0.0609322133186686
23.625 -8.43604652800411e-08
23.65 0.0795573054347176
23.675 0.281404425845718
23.7 0.24398445425987
23.725 -0.0541647493264494
23.75 -0.125390347289168
23.775 -0.0272400399808244
23.8 4.37239577877329e-08
23.825 2.18647357996812e-07
23.85 0.00184931420790246
23.875 8.43280645312904e-08
23.9 -0.0209342400239617
23.925 -0.119084719108285
23.95 -0.138374186867336
23.975 0.0634739798670161
24 0.0246066221231441
24.025 -0.277242824654347
24.05 0.00717331827217994
24.075 0.478750239112948
24.1 0.0991033268519515
24.125 -0.266831276990441
24.15 0.105642101470947
24.175 0.320934328897233
24.2 0.120933599173723
24.225 0.0018500544683776
24.25 0.0321351230874707
24.275 0.130285217960589
24.3 0.167864560262385
24.325 0.0678644953089948
24.35 0.00185144756414967
24.375 -0.0151986282893544
24.4 -0.117050356357138
24.425 -0.146511122602531
24.45 0.0534896509001462
24.475 0.163560723626697
24.5 0.0617093453634956
24.525 -0.173658290009368
24.55 -0.341763493111098
24.575 0.00185093145347534
24.6 0.267086078451954
24.625 -0.106570895518699
24.65 -0.267399821768773
24.675 0.0344494157388903
24.7 -0.0239859063137257
24.725 -0.275704460000286
24.75 0.112493735213944
24.775 0.568760071382704
24.8 0.404546780146114
24.825 0.240828667988705
24.85 0.230945534349316
24.875 -0.00354194558475735
24.9 -0.146874479128199
24.925 0.0467831996957095
24.95 0.282748979366612
24.975 0.384603535781047
25 0.118524787598842
25.025 -0.283329878402579
25.05 -0.236322524801832
25.075 -0.0673524704276209
25.1 -0.129178049542141
25.125 -0.114549620302922
25.15 0.0835981128602713
25.175 0.0659609373542471
25.2 -0.136774865558989
25.225 -0.221263838602694
25.25 -0.217644849595115
25.275 -0.0991164430843114
25.3 -8.71978045324783e-08
25.325 0.016256411873101
25.35 0.118107916250672
25.375 0.183743457677889
25.4 0.0818919532093775
25.425 0.0694475093041537
25.45 0.169447967441909
25.475 0.035921087942372
25.5 -0.148257335631923
25.525 0.0123062280636075
25.55 0.178809237763741
25.575 0.0823250878105585
25.6 -0.0654620085015304
25.625 -0.167314559398972
25.65 0.0474443687380692
25.675 0.334767334652492
25.7 0.103488686823211
25.725 -0.165599897627554
25.75 -0.0854708290083993
25.775 0.0606851890324825
25.8 0.162537260166571
25.825 -0.0475606095229601
25.85 -0.312595195600824
25.875 -0.0818664030297344
25.9 0.0748423762388473
25.925 -0.101234285375538
25.95 -0.0176925490173827
25.975 0.0783061011588915
26 -0.298760352239138
26.025 -0.340698234981569
26.05 0.26115106838915
26.075 0.194998634611136
26.1 -0.396237127574773
26.125 -0.157573783990654
26.15 0.354074349752404
26.175 0.320410061408434
26.2 0.172679951838643
26.225 0.0745335487636154
26.25 4.35826308375902e-08
26.275 -0.00185318294825263
26.3 -8.71978045324783e-08
26.325 0.0238208321710683
26.35 0.125674342211831
26.375 0.0177869892873869
26.4 -0.384067545905958
26.425 -0.256686084890131
26.45 0.211494460773235
26.475 0.0851097673148664
26.5 -0.102096586970573
26.525 0.0809734934136929
26.55 0.148995504287084
26.575 0.0489951335388376
26.6 -0.0970716328083709
26.625 -0.224129834378243
26.65 -0.228134258190945
26.675 -0.223424715318285
26.7 -0.222349203790767
26.725 -0.00736830199912261
26.75 0.290781153416613
26.775 0.354058479001646
26.8 0.275816035582462
26.825 0.220771711615623
26.85 0.182811525041881
26.875 0.0800924434039263
26.9 -0.0018539240954851
26.925 0
26.95 8.72254446448993e-08
26.975 -0.0476599768717306
27 -0.149513225306135
27.025 -0.0126417829505952
27.05 0.28921020129928
27.075 0.0716475400898976
27.1 -0.326496879574752
27.125 -0.211345606231032
27.15 -0.0551674103462566
27.175 -0.141966943056102
27.2 -0.0866214217645466
27.225 0.115230714982324
27.25 0.128589040509366
27.275 0.0267362716840576
27.3 -8.71978045324783e-08
27.325 -4.35598934700459e-08
27.35 0.00185292135483916
27.375 -1.30835715594911e-07
27.4 -0.00185300855264356
27.425 -0.0929907761019409
27.45 -0.160435811921713
27.475 0.0273991411859258
27.5 0.236339790848984
27.525 0.234333000798866
27.55 0.0169334080283077
27.575 -0.13772785748152
27.6 0.0400294571136351
27.625 0.223834959224046
27.65 0.221981994377526
27.675 -0.0146197526419199
27.7 -0.314621083016031
27.725 -0.238375662843296
27.75 -0.0136802769751652
27.775 0.124694840656046
27.8 0.122374470561224
27.825 -0.0776261836968917
27.85 -0.06780910603084
27.875 0.1303377113296
27.9 0.0199994898594011
27.925 -0.178146629968976
27.95 -0.101112420015495
27.975 -0.0490338949703101
28 -0.14977348098945
28.025 -0.0593957699720696
28.05 0.110566997045679
28.075 -0.024572342881228
28.1 -0.107349313816491
28.125 0.0834805639393454
28.15 0.223652618933553
28.175 0.225505191406225
28.2 0.0238898529736528
28.225 -0.174256441277986
28.25 -0.100986586976786
28.275 -0.0493125505727214
28.3 -0.146472170581544
28.325 -0.0840998319696228
28.35 0.117752261098765
28.375 0.130572460211267
28.4 0.0287201711444136
28.425 -8.71750671649339e-08
28.45 0.184397301614252
28.475 0.686252556442962
28.5 0.511253586268687
28.525 -0.404706970971341
28.55 -0.609754816173327
28.575 -0.383771755191086
28.6 -0.386270323360826
28.625 -0.226784843817211
28.65 0.0713633036184537
28.675 0.0391139621115139
28.7 -0.12276640201854
28.725 0.129869521356348
28.75 0.351777131999919
28.775 0.160027573679142
28.8 -0.0666676696027925
28.825 -0.166668127526106
28.85 -0.0408324171758091
28.875 0.142458110821833
28.9 -0.00735783778509347
28.925 -0.138810170687973
28.95 0.051838326026327
28.975 0.164873065045877
29 0.0630195985146278
29.025 0.0904599865862097
29.05 0.192312188596586
29.075 0.107686961728531
29.1 0.0095414742645522
29.125 -0.00739768918755591
29.15 -0.109251896901695
29.175 -0.19260209319327
29.2 -0.0926017658408527
29.225 0.0570733942822731
29.25 0.155219884363402
29.275 -0.0168809122442324
29.3 -0.311693086200644
29.325 -0.136858885043336
29.35 0.185776670082315
29.375 0.224116318003876
29.4 0.0194624971072696
29.425 -0.378682227778876
29.45 -0.290737944642028
29.475 0.211111990756775
29.5 0.0444582079531835
29.525 -0.312821311513285
29.55 0.0420857781136868
29.575 0.404911983928636
29.6 0.309251809972029
29.625 0.148029301153299
29.65 0.0125721954645995
29.675 -0.133603268818659
29.7 -0.0333129497627022
29.725 0.268540560659873
29.75 0.107969454001802
29.775 -0.293885211931406
29.8 -0.340627688137824
29.825 -0.193645470489724
29.85 0.148409852590711
29.875 0.352593434190207
29.9 0.151165101351578
29.925 -4.36102709500113e-08
29.95 -1.30803172737613e-07
29.975 0.00185396770575618
30 -0.018787114213765
30.025 -0.120641757939664
30.05 -0.181212668167061
30.075 -0.08121186121535
30.1 0.00185292128662706
30.125 -8.72936567475323e-08
30.15 -0.00185270318979744
30.175 0.0537757027814215
30.2 0.15377594270866
30.225 0.0606718033395771
30.25 -0.151339331914995
30.275 -0.224605392895679
30.3 -0.168564911655264
30.325 0.0440301311576968
30.35 0.227175728741726
30.375 0.225321717534467
30.4 -0.110162726423367
30.425 -0.492686135066826
30.45 -0.0352621426096528
30.475 0.479136269223118
30.5 0.0281686591002735
30.525 -0.267660033887292
30.55 0.134191492907119
30.575 0.19903375367029
30.6 -0.199111712527167
30.625 -0.0909233367841709
30.65 0.275086194642423
30.675 0.0817680385353003
30.7 -0.0770644255404551
30.725 0.11088371638835
30.75 0.123106796872532
30.775 0.0212536355775225
30.8 3.05180947179902e-07
30.825 -0.0550461988674292
30.85 -0.153192776315345
30.875 -0.304356526003526
30.9 -0.389296335697503
30.925 0.0763159668622382
30.95 0.514642491434991
30.975 0.157093996040212
31 -0.120413592622198
31.025 0.0777330502842222
31.05 0.227307317300071
31.075 0.227307186587847
31.1 -0.0757858196266213
31.125 -0.331647841385698
31.15 0.0427312506417359
31.175 0.299928690674274
31.2 0.00133546193900711
31.225 0.040812447878127
31.25 0.238458374423868
31.275 -0.188641444947277
31.3 -0.418617288141826
31.325 -0.0323297422433906
31.35 0.134684709928486
31.375 0.0171044993472689
31.4 -0.117580254163839
31.425 -0.144694548589539
31.45 0.0553061493341281
31.475 0.0975798487616031
31.5 -0.104273726929377
31.525 -0.0498485840125788
31.55 0.15200420689325
31.575 0.112689112420766
31.6 0.0145433632951616
31.625 -4.36379110624324e-08
31.65 -0.00185370608960509
31.675 4.36102709500113e-08
31.7 0.071729684741867
31.725 0.150138684784915
31.75 -0.095173889558636
31.775 -0.294599225617272
31.8 -0.0210171424896541
31.825 0.144460540860008
31.85 0.0351869305632155
31.875 -0.111127708878838
31.9 -0.155812425236943
31.925 0.0460403221040899
31.95 0.105557434272113
31.975 -0.092588380420798
32 -0.0567606862757263
32.025 0.155784694824619
32.05 0.330687437108893
32.075 0.308206470582082
32.1 -0.206227258195225
32.125 -0.506315267840123
32.15 -0.208169823745887
32.175 0.0792417835879821
32.2 0.179241936358007
32.225 0.122612315408339
32.25 0.0207578896053929
32.275 -0.00185431673415321
32.3 -0.0777044914586545
32.325 -0.158648248715716
32.35 -0.00138555974444906
32.375 0.111286068376308
32.4 -0.0682728051840029
32.425 -0.0122013669526656
32.45 0.287799876114806
32.475 0.0720264346910012
32.5 -0.292062724348144
32.525 0.0253036225232939
32.55 0.265846222883445
32.575 -0.123546776821264
32.6 -0.246272193726341
32.625 0.0518752561576748
32.65 0.146894404760319
32.675 0.0654451349062357
32.7 0.00185419327024771
32.725 -0.029686095558647
32.75 -0.0166970730886372
32.775 0.0604162223582653
32.8 0.162269625034938
32.825 0.137729615117487
32.85 -0.0652171200727062
32.875 -0.404800628312595
32.9 -0.341103738138073
32.925 0.160748742119438
32.95 0.244051405489443
32.975 0.0459061361198136
33 -2.16014797160836e-07
33.025 -0.00185387617491967
33.05 2.16014797160806e-07
33.075 -0.00185417942631726
33.1 -0.0904327434909698
33.125 0.0114208770712316
33.15 0.263177334126516
33.175 0.0631783071283056
33.2 -0.1727449612092
33.225 -0.0708912331447249
33.25 3.92391683590199e-07
33.275 -0.00185374975515613
33.3 -1.74340328840102e-07
33.325 0.0226696726811339
33.35 0.124524185457638
33.375 0.123501736597269
33.4 -0.0783529290045171
33.425 0.0407144324169871
33.45 0.2729169029506
33.475 -0.0528299921825267
33.5 -0.368709014229175
33.525 -0.283677118623633
33.55 -0.071427920546932
33.575 0.128572602953419
33.6 0.122902743078164
33.625 0.021048709767656
33.65 -0.00185388054049424
33.675 -0.0283052617662325
33.7 -0.228305785312144
33.725 -0.330958200675032
33.75 -0.0328103149610844
33.775 0.159263418517441
33.8 0.0611162307547578
33.825 -8.72108074645489e-08
33.85 0.036093618018642
33.875 0.134240282649721
33.9 0.135124559064212
33.925 0.0100503006348865
33.95 0.0018540626213338
33.975 0.0575633148400844
34 0.0269273917540162
34.025 0.0573339874237839
34.05 0.159187480049483
34.075 0.061333438507777
34.1 -0.130173631851001
34.125 -0.0138842657277402
34.15 0.17098621246005
34.175 0.0933626361172833
34.2 -0.100727248754288
34.225 -0.298872785469956
34.25 -0.14872925081435
34.275 0.249417043268473
34.3 -0.0504961755625245
34.325 -0.504637697021508
34.35 -0.0523337854247104
34.375 0.407130418706624
34.4 0.305323583542477
34.425 0.14696288941265
34.45 0.0123758930028835
34.475 -0.132732161318935
34.5 -0.0987099535548508
34.525 0.0994358839208957
34.55 0.0419070836058175
34.575 -0.158093632662011
34.6 -0.106756070863418
34.625 0.0695230039484369
34.65 0.276279647631957
34.675 0.187915866545849
34.7 -0.212085825650531
34.725 -0.17742371624082
34.75 0.109471000804866
34.775 0.000123219924006424
34.8 -0.136597357858778
34.825 0.0501736050507901
34.85 0.103463062623638
34.875 -0.0983916121283868
34.9 -0.0553034827232987
34.925 0.163608647390689
34.95 0.229977432682786
34.975 0.143783462681321
35 -0.0709914371880456
35.025 -0.147922482311303
35.05 -0.0654517414561384
35.075 -0.219385036315755
35.1 -0.273856987178094
35.125 0.224288835888031
35.15 0.434684734217788
35.175 0.0365400547350672
35.2 -0.141443293286309
35.225 -0.0471105157885203
35.25 -0.107520766014921
35.275 -0.162856902180252
35.3 0.0389965977417683
35.325 0.22346977067245
35.35 0.225325564588287
35.375 -0.0146542775935651
35.4 -0.316510358025654
35.425 -0.240291684112159
35.45 -0.0402906670289861
35.475 0.00185396097549539
35.5 -1.28612356320453e-07
35.525 0.0168744095746817
35.55 0.118727724848914
35.575 0.181271427901077
35.6 0.0831269524496794
35.625 0.0711935335304799
35.65 0.169339099001448
35.675 -0.0682971311918776
35.7 -0.329613536213619
35.725 0.0376432731172848
35.75 0.229973635340896
35.775 -0.268986627962255
35.8 -0.403949601003825
35.825 -0.00394812731357797
35.85 0.134249642860603
35.875 0.0205191120487325
35.9 -0.113730227404351
35.925 -0.109655218865745
35.95 0.190345638257636
35.975 0.260588687751821
36 -0.0412663926552194
36.025 -0.13905717106424
36.05 -0.0708904531977279
36.075 -0.13368817793819
36.1 -0.0859739490806581
36.125 0.104710608265327
36.15 0.0121986498373151
36.175 -0.141030341218725
36.2 0.0374555225887029
36.225 0.22889731931174
36.25 0.22704226820609
36.275 -0.016801027378078
36.3 -0.308694984905005
36.325 -0.0358371444137638
36.35 0.299418547137992
36.375 -0.0603466785930385
36.4 -0.347758337111982
36.425 -0.1459044001779
36.45 0.0718193954168852
36.475 0.173672669416202
36.5 0.0289727336524253
36.525 -0.149553244601238
36.55 0.0225349241288323
36.575 0.119799423519644
36.6 -0.077554382187377
36.625 0.138847528877911
36.65 0.467556017328366
36.675 0.0477951569889927
36.7 -0.313328846086041
36.725 -0.128707815883907
36.75 -4.12099154800695e-08
36.775 -0.0331536082838085
36.8 -0.231298358120691
36.825 -0.290039037681015
36.85 0.111812129080142
36.875 0.2537172993828
36.9 -0.0481353411129918
36.925 -0.1323785465859
36.95 -0.0432165837894896
36.975 -0.108985707765902
37 -0.123338910371522
37.025 0.174809464160062
37.05 0.334177192926901
37.075 0.134176088395826
37.1 -0.245661937015043
37.125 -0.512893585174184
37.15 -0.119716667550794
37.175 0.312893784130823
37.2 0.165378631264242
37.225 -0.0920524339562742
37.25 -0.187925895771254
37.275 0.593168784952616
37.3 0.800318065976742
37.325 -0.383860688761381
37.35 -0.410538170854148
37.375 0.0827442768142864
37.4 -0.0844634737956723
37.425 -0.175215987646154
37.45 -0.00999738180262057
37.475 0.130641291661967
37.5 -0.15382080197079
37.525 -0.355425126245491
37.55 -0.153572138005331
37.575 -0.079108142917268
37.6 -0.169575707650447
37.625 -0.0132113456499646
37.65 0.135401727756547
37.675 -0.0418537556827623
37.7 -0.0869101816436362
37.725 0.145572477731556
37.75 0.351714904890577
37.775 0.388600632176515
37.8 0.169368991974522
37.825 -0.0711224929447992
37.85 -0.169268483114147
37.875 -0.325116709528378
37.9 -0.389116633833206
37.925 0.0340935784249268
37.95 0.358385142690827
37.975 0.16399865967288
38 -4.05698585836967e-09
38.025 -0.0018531441680379
38.05 -0.0913933962140512
38.075 -0.191393711146983
38.1 -0.106754350008493
38.125 -0.0086057813502554
38.15 -0.00185178970305328
38.175 -1.34231683546193e-07
38.2 -8.68470095838397e-08
38.225 0.00185218850674584
38.25 -0.0347472824257286
38.275 -0.136599734524398
38.3 -0.00892813763066082
38.325 0.276731555878287
38.35 0.0311860137474458
38.375 -0.278532026969221
38.4 0.0759419677348851
38.425 0.28123182142801
38.45 -0.018769336566039
38.475 -0.221898075333881
38.5 -0.205664448695355
38.525 0.195300049618422
38.55 0.533813516550014
38.575 0.154746246562363
38.6 -0.170979569257464
38.625 -0.072832398437364
38.65 -0.00185323091436358
38.675 2.21198845906669e-07
38.7 3.0800038075542e-07
38.725 0.0244368593387648
38.75 0.122583335168461
38.775 0.120486325168433
38.8 -0.0776596837015875
38.825 -0.144923275320261
38.85 -0.111653361419221
38.875 -0.268581430519296
38.9 -0.267198509335932
38.925 0.0346529654149938
38.95 0.0425902661845636
38.975 -0.125003937783468
39 0.121890913068715
39.025 0.320067485755019
39.05 0.0305825030506685
39.075 -0.295105246323163
39.1 -0.393254202081814
39.125 -0.179434474840206
39.15 0.107621925108736
39.175 0.00231328600477822
39.2 -0.141026568053428
39.225 0.0474186413515581
39.25 0.165523780869407
39.275 0.0636732440950229
39.3 -1.77660055555872e-07
39.325 0.0182715192750048
39.35 0.120123106525447
39.375 0.236242106383793
39.4 0.234391011836743
39.425 0.0491141574099913
39.45 -0.1508861644652
39.475 -0.101775884887238
39.5 0.04761856115628
39.525 0.14939479757294
39.55 -0.0885372884925368
39.575 -0.488538765375878
39.6 -0.262156298832196
39.625 0.135994083292974
39.65 0.0994469854146644
39.675 -0.0914795156307998
39.7 -0.292778040920981
39.725 -0.162598383722648
39.75 0.23555217207045
39.775 0.255375621109742
39.8 0.09442113074393
39.825 0.237196239140771
39.85 0.321978882438377
39.875 0.0201253623247749
39.9 -0.00544009166269131
39.925 0.261215118153559
39.95 0.0110685016807117
39.975 -0.262723883552753
40 -0.0108422868476756
40.025 0.144186938778775
40.05 0.046039641584417
40.075 1.77669789991352e-07
40.1 0.0365017893630458
40.125 0.134648130905847
40.15 0.0466860601233736
40.175 -0.251459165828521
40.2 -0.283187827744785
40.225 -0.0748224119076511
40.25 0.106513966820136
40.275 0.130523866310515
40.3 -0.167622500920795
40.325 -0.284370524877902
40.35 0.0174800771122118
40.375 0.0663033105768506
40.4 -0.131489067361403
40.425 -0.0167638912346755
40.45 0.106207125327279
40.475 -0.171118527405919
40.5 -0.172657252333635
40.525 0.225492300533626
40.55 0.241567742214755
40.575 0.0870668108746081
40.6 0.147352107014109
40.625 0.154499971712796
40.65 0.0526476935289776
40.675 -1.77660055555872e-07
40.7 -2.6457200874576e-07
40.725 0.00185297034413178
40.75 3.51373543594461e-07
40.775 -0.00185261891537875
40.8 0.0672768398641959
40.825 0.167276980982943
40.85 0.13457675635366
40.875 0.0173303425428613
40.9 -0.117246674371206
40.925 -0.223547895611075
40.95 -0.223547678479021
40.975 -0.159207438819715
41 -0.0610590873390696
41.025 -0.0018518764948537
41.05 -0.0339033369430572
41.075 -0.133903825599199
41.1 -0.0862769324255246
41.125 0.0970802840031754
41.15 -0.094612197234972
41.175 -0.300990912718947
41.2 0.0769788119214398
41.225 0.225034768836294
41.25 -0.176816835128817
41.275 -0.287219958990452
41.3 -0.0652360329144168
41.325 0.320134066396438
41.35 0.528896968480156
41.375 0.027044299294443
41.4 -0.349030202768758
41.425 -0.147178608621567
41.45 0.158679657897674
41.475 0.360531599553105
41.5 0.0425531592457381
41.525 -0.336015464007928
41.55 -0.0779495176987455
41.575 0.230425717745688
41.6 0.233510749184744
41.625 -0.0472259992518036
41.65 -0.383396911674264
41.675 -0.3420331663523
41.7 -0.106140019397649
41.725 0.103424832458933
41.75 0.132742530913095
41.775 0.0456142887384875
41.8 0.114721693453894
41.825 0.107477209439827
41.85 -0.190671294158815
41.875 -0.18762782384391
41.9 0.210521997688087
41.925 0.434360600102746
41.95 0.395282581985179
41.975 -0.206157689363605
42 -0.69086248991938
42.025 -0.123781211495302
42.05 0.404213866896251
42.075 0.304213555955374
42.1 0.237088926532407
42.125 0.219915252536519
42.15 0.00438406028543133
42.175 -0.272123420872263
42.2 -0.393681470040996
42.225 -0.202155987481743
42.25 0.0843371950298105
42.275 -0.000498867611256017
42.3 -0.173483972993661
42.325 -0.172985192239149
};
\addplot [line width=2.4000000000000004pt, color1, dashed, forget plot]
table {%
0 1.19262162802871
0.025 1.19054322073438
0.05 1.18846481344005
0.075 1.18638640614572
0.1 1.18430799885139
0.125 1.18222959155706
0.15 1.18015118426273
0.175 1.1780727769684
0.2 1.17599436967407
0.225 1.17391596237974
0.25 1.17183755508541
0.275 1.16975914779108
0.3 1.16768074049675
0.325 1.16560233320242
0.35 1.16352392590809
0.375 1.16144551861376
0.4 1.15936711131943
0.425 1.1572887040251
0.45 1.15521029673077
0.475 1.15313188943644
0.5 1.15105348214211
0.525 1.14897507484777
0.55 1.14689666755344
0.575 1.14481826025911
0.6 1.14273985296478
0.625 1.14066144567045
0.65 1.13858303837612
0.675 1.13650463108179
0.7 1.13442622378746
0.725 1.13234781649313
0.75 1.1302694091988
0.775 1.12819100190447
0.8 1.12611259461014
0.825 1.12403418731581
0.85 1.12195578002148
0.875 1.11987737272715
0.9 1.11779896543282
0.925 1.11572055813849
0.95 1.11364215084416
0.975 1.11156374354983
1 1.1094853362555
1.025 1.10740692896117
1.05 1.10532852166684
1.075 1.10325011437251
1.1 1.10117170707818
1.125 1.09909329978385
1.15 1.09701489248952
1.175 1.09493648519519
1.2 1.09285807790086
1.225 1.09077967060653
1.25 1.0887012633122
1.275 1.08662285601787
1.3 1.08454444872354
1.325 1.08246604142921
1.35 1.08038763413488
1.375 1.07830922684054
1.4 1.07623081954621
1.425 1.07415241225188
1.45 1.07207400495755
1.475 1.06999559766322
1.5 1.06791719036889
1.525 1.06583878307456
1.55 1.06376037578023
1.575 1.0616819684859
1.6 1.05960356119157
1.625 1.05752515389724
1.65 1.05544674660291
1.675 1.05336833930858
1.7 1.05128993201425
1.725 1.04921152471992
1.75 1.04713311742559
1.775 1.04505471013126
1.8 1.04297630283693
1.825 1.0408978955426
1.85 1.03881948824827
1.875 1.03674108095394
1.9 1.03466267365961
1.925 1.03258426636528
1.95 1.03050585907095
1.975 1.02842745177662
2 1.02634904448229
2.025 1.02427063718796
2.05 1.02219222989363
2.075 1.0201138225993
2.1 1.01803541530497
2.125 1.01595700801064
2.15 1.0138786007163
2.175 1.01180019342197
2.2 1.00972178612764
2.225 1.00764337883331
2.25 1.00556497153898
2.275 1.00348656424465
2.3 1.00140815695032
2.325 0.999329749655993
2.35 0.997251342361663
2.375 0.995172935067332
2.4 0.993094527773002
2.425 0.991016120478672
2.45 0.988937713184341
2.475 0.986859305890011
2.5 0.984780898595681
2.525 0.98270249130135
2.55 0.98062408400702
2.575 0.97854567671269
2.6 0.976467269418359
2.625 0.974388862124029
2.65 0.972310454829699
2.675 0.970232047535369
2.7 0.968153640241038
2.725 0.966075232946708
2.75 0.963996825652378
2.775 0.961918418358047
2.8 0.959840011063717
2.825 0.957761603769387
2.85 0.955683196475056
2.875 0.953604789180726
2.9 0.951526381886396
2.925 0.949447974592065
2.95 0.947369567297735
2.975 0.945291160003405
3 0.943212752709075
3.025 0.941134345414744
3.05 0.939055938120414
3.075 0.936977530826084
3.1 0.934899123531753
3.125 0.932820716237423
3.15 0.930742308943093
3.175 0.928663901648762
3.2 0.926585494354432
3.225 0.924507087060102
3.25 0.922428679765772
3.275 0.920350272471441
3.3 0.918271865177111
3.325 0.916193457882781
3.35 0.91411505058845
3.375 0.91203664329412
3.4 0.90995823599979
3.425 0.907879828705459
3.45 0.905801421411129
3.475 0.903723014116799
3.5 0.901644606822468
3.525 0.899566199528138
3.55 0.897487792233808
3.575 0.895409384939478
3.6 0.893330977645147
3.625 0.891252570350817
3.65 0.889174163056487
3.675 0.887095755762156
3.7 0.885017348467826
3.725 0.882938941173496
3.75 0.880860533879165
3.775 0.878782126584835
3.8 0.876703719290505
3.825 0.874625311996174
3.85 0.872546904701844
3.875 0.870468497407514
3.9 0.868390090113183
3.925 0.866311682818853
3.95 0.864233275524523
3.975 0.862154868230193
4 0.860076460935862
4.025 0.857998053641532
4.05 0.855919646347202
4.075 0.853841239052871
4.1 0.851762831758541
4.125 0.849684424464211
4.15 0.84760601716988
4.175 0.84552760987555
4.2 0.84344920258122
4.225 0.841370795286889
4.25 0.839292387992559
4.275 0.837213980698229
4.3 0.835135573403899
4.325 0.833057166109568
4.35 0.830978758815238
4.375 0.828900351520908
4.4 0.826821944226577
4.425 0.824743536932247
4.45 0.822665129637917
4.475 0.820586722343586
4.5 0.819247021588974
4.525 0.818646027374083
4.55 0.818045033159192
4.575 0.817444038944301
4.6 0.81684304472941
4.625 0.816242050514519
4.65 0.815641056299628
4.675 0.815040062084737
4.7 0.814439067869846
4.725 0.813838073654955
4.75 0.813237079440064
4.775 0.812636085225172
4.8 0.812035091010281
4.825 0.81143409679539
4.85 0.810833102580499
4.875 0.810232108365608
4.9 0.809631114150717
4.925 0.809030119935826
4.95 0.808429125720935
4.975 0.807828131506044
5 0.807227137291153
5.025 0.806626143076262
5.05 0.806025148861371
5.075 0.80542415464648
5.1 0.804823160431589
5.125 0.804222166216698
5.15 0.803621172001806
5.175 0.803020177786915
5.2 0.802419183572024
5.225 0.801818189357133
5.25 0.801217195142242
5.275 0.800616200927351
5.3 0.80001520671246
5.325 0.799414212497569
5.35 0.798813218282678
5.375 0.798212224067787
5.4 0.797611229852896
5.425 0.797010235638005
5.45 0.796409241423114
5.475 0.795808247208223
5.5 0.795207252993332
5.525 0.794606258778441
5.55 0.79400526456355
5.575 0.793404270348658
5.6 0.792803276133768
5.625 0.792202281918877
5.65 0.791601287703985
5.675 0.791000293489094
5.7 0.790399299274203
5.725 0.789798305059312
5.75 0.789197310844421
5.775 0.78859631662953
5.8 0.787995322414639
5.825 0.787394328199748
5.85 0.786793333984857
5.875 0.786192339769966
5.9 0.785591345555075
5.925 0.784990351340184
5.95 0.784389357125293
5.975 0.783788362910402
6 0.783187368695511
6.025 0.78258637448062
6.05 0.781985380265729
6.075 0.781384386050837
6.1 0.780783391835946
6.125 0.780182397621055
6.15 0.779581403406164
6.175 0.778980409191273
6.2 0.778379414976382
6.225 0.777778420761491
6.25 0.7771774265466
6.275 0.776576432331709
6.3 0.775975438116818
6.325 0.775374443901927
6.35 0.774773449687036
6.375 0.774172455472145
6.4 0.773571461257254
6.425 0.772970467042363
6.45 0.772369472827472
6.475 0.771768478612581
6.5 0.771167484397689
6.525 0.770566490182798
6.55 0.769965495967907
6.575 0.769364501753016
6.6 0.768763507538125
6.625 0.768162513323234
6.65 0.767561519108343
6.675 0.766960524893452
6.7 0.766359530678561
6.725 0.76575853646367
6.75 0.765157542248779
6.775 0.764556548033888
6.8 0.763955553818997
6.825 0.763354559604106
6.85 0.762753565389215
6.875 0.762152571174324
6.9 0.761551576959433
6.925 0.760950582744542
6.95 0.76034958852965
6.975 0.759748594314759
7 0.759147600099868
7.025 0.758546605884977
7.05 0.757945611670086
7.075 0.757344617455195
7.1 0.756743623240304
7.125 0.756142629025413
7.15 0.755541634810522
7.175 0.754940640595631
7.2 0.75433964638074
7.225 0.753738652165849
7.25 0.753137657950958
7.275 0.752536663736067
7.3 0.751935669521176
7.325 0.751334675306285
7.35 0.750733681091394
7.375 0.750132686876502
7.4 0.749531692661612
7.425 0.74893069844672
7.45 0.748329704231829
7.475 0.747728710016938
7.5 0.747127715802047
7.525 0.746526721587156
7.55 0.745925727372265
7.575 0.745324733157374
7.6 0.744723738942483
7.625 0.744122744727592
7.65 0.743521750512701
7.675 0.74292075629781
7.7 0.742319762082919
7.725 0.741718767868028
7.75 0.741117773653137
7.775 0.740516779438246
7.8 0.739915785223355
7.825 0.739314791008463
7.85 0.738713796793573
7.875 0.738112802578681
7.9 0.73751180836379
7.925 0.736910814148899
7.95 0.736309819934008
7.975 0.735708825719117
8 0.735107831504226
8.025 0.734506837289335
8.05 0.733905843074444
8.075 0.733304848859553
8.1 0.732703854644662
8.125 0.732102860429771
8.15 0.73150186621488
8.175 0.730900871999989
8.2 0.730299877785098
8.225 0.729698883570207
8.25 0.729097889355316
8.275 0.728496895140425
8.3 0.727895900925533
8.325 0.727294906710642
8.35 0.726693912495751
8.375 0.72609291828086
8.4 0.725491924065969
8.425 0.724890929851078
8.45 0.724289935636187
8.475 0.723688941421296
8.5 0.723087947206405
8.525 0.722486952991514
8.55 0.721885958776623
8.575 0.721284964561732
8.6 0.720683970346841
8.625 0.72008297613195
8.65 0.719481981917059
8.675 0.718880987702168
8.7 0.718279993487277
8.725 0.717678999272386
8.75 0.717078005057494
8.775 0.716477010842603
8.8 0.715876016627712
8.825 0.715275022412821
8.85 0.71467402819793
8.875 0.714073033983039
8.9 0.713472039768148
8.925 0.712871045553257
8.95 0.712270051338366
8.975 0.711669057123475
9 -1.11209535857595
9.025 -1.10157938116379
9.05 -1.09106340375164
9.075 -1.08054742633949
9.1 -1.07003144892734
9.125 -1.05951547151519
9.15 -1.04899949410304
9.175 -1.03848351669089
9.2 -1.02796753927874
9.225 -1.01745156186659
9.25 -1.00693558445444
9.275 -0.996419607042287
9.3 -0.985903629630136
9.325 -0.975387652217985
9.35 -0.964871674805835
9.375 -0.954355697393684
9.4 -0.943839719981533
9.425 -0.933323742569382
9.45 -0.922807765157231
9.475 -0.912291787745081
9.5 -0.90177581033293
9.525 -0.891259832920779
9.55 -0.880743855508628
9.575 -0.870227878096477
9.6 -0.859711900684326
9.625 -0.849195923272176
9.65 -0.838679945860025
9.675 -0.828163968447874
9.7 -0.817647991035723
9.725 -0.807132013623572
9.75 -0.796616036211422
9.775 -0.786100058799271
9.8 -0.77558408138712
9.825 -0.76506810397497
9.85 -0.754552126562819
9.875 -0.744036149150669
9.9 -0.733520171738518
9.925 -0.723004194326367
9.95 -0.712488216914216
9.975 -0.701972239502065
10 -0.685567309038302
10.025 -0.663273425522912
10.05 -0.640979542007524
10.075 -0.618685658492135
10.1 -0.596391774976746
10.125 -0.574097891461358
10.15 -0.55180400794597
10.175 -0.529510124430581
10.2 -0.507216240915192
10.225 -0.484922357399803
10.25 -0.462628473884415
10.275 -0.440334590369027
10.3 -0.418040706853638
10.325 -0.395746823338249
10.35 -0.37345293982286
10.375 -0.351159056307472
10.4 -0.328865172792083
10.425 -0.306571289276694
10.45 -0.284277405761306
10.475 -0.261983522245917
10.5 -0.239689638730529
10.525 -0.21739575521514
10.55 -0.195101871699751
10.575 -0.172807988184362
10.6 -0.150514104668974
10.625 -0.128220221153586
10.65 -0.105926337638197
10.675 -0.0836324541228083
10.7 -0.0613385706074195
10.725 -0.0390446870920304
10.75 -0.0167508035766431
10.775 0.00554307993874602
10.8 0.0278369634541348
10.825 0.0501308469695237
10.85 0.0724247304849125
10.875 0.0947186140002999
10.9 0.117012497515689
10.925 0.139306381031078
10.95 0.161600264546467
10.975 0.183894148061856
11 0
11.025 0
11.05 0
11.075 0
11.1 0
11.125 0
11.15 0
11.175 0
11.2 0
11.225 0
11.25 0
11.275 0
11.3 0
11.325 0
11.35 0
11.375 0
11.4 0
11.425 0
11.45 0
11.475 0
11.5 0
11.525 0
11.55 0
11.575 0
11.6 0
11.625 0
11.65 0
11.675 0
11.7 0
11.725 0
11.75 0
11.775 0
11.8 0
11.825 0
11.85 0
11.875 0
11.9 0
11.925 0
11.95 0
11.975 0
12 0
12.025 0
12.05 0
12.075 0
12.1 0
12.125 0
12.15 0
12.175 0
12.2 0
12.225 0
12.25 0
12.275 0
12.3 0
12.325 0
12.35 0
12.375 0
12.4 0
12.425 0
12.45 0
12.475 0
12.5 0
12.525 0
12.55 0
12.575 0
12.6 0
12.625 0
12.65 0
12.675 0
12.7 0
12.725 0
12.75 0
12.775 0
12.8 0
12.825 0
12.85 0
12.875 0
12.9 0
12.925 0
12.95 0
12.975 0
13 0
13.025 0
13.05 0
13.075 0
13.1 0
13.125 0
13.15 0
13.175 0
13.2 0
13.225 0
13.25 0
13.275 0
13.3 0
13.325 0
13.35 0
13.375 0
13.4 0
13.425 0
13.45 0
13.475 0
13.5 0
13.525 0
13.55 0
13.575 0
13.6 0
13.625 0
13.65 0
13.675 0
13.7 0
13.725 0
13.75 0
13.775 0
13.8 0
13.825 0
13.85 0
13.875 0
13.9 0
13.925 0
13.95 0
13.975 0
14 0
14.025 0
14.05 0
14.075 0
14.1 0
14.125 0
14.15 0
14.175 0
14.2 0
14.225 0
14.25 0
14.275 0
14.3 0
14.325 0
14.35 0
14.375 0
14.4 0
14.425 0
14.45 0
14.475 0
14.5 0
14.525 0
14.55 0
14.575 0
14.6 0
14.625 0
14.65 0
14.675 0
14.7 0
14.725 0
14.75 0
14.775 0
14.8 0
14.825 0
14.85 0
14.875 0
14.9 0
14.925 0
14.95 0
14.975 0
15 0
15.025 0
15.05 0
15.075 0
15.1 0
15.125 0
15.15 0
15.175 0
15.2 0
15.225 0
15.25 0
15.275 0
15.3 0
15.325 0
15.35 0
15.375 0
15.4 0
15.425 0
15.45 0
15.475 0
15.5 0
15.525 0
15.55 0
15.575 0
15.6 0
15.625 0
15.65 0
15.675 0
15.7 0
15.725 0
15.75 0
15.775 0
15.8 0
15.825 0
15.85 0
15.875 0
15.9 0
15.925 0
15.95 0
15.975 0
16 0
16.025 0
16.05 0
16.075 0
16.1 0
16.125 0
16.15 0
16.175 0
16.2 0
16.225 0
16.25 0
16.275 0
16.3 0
16.325 0
16.35 0
16.375 0
16.4 0
16.425 0
16.45 0
16.475 0
16.5 0
16.525 0
16.55 0
16.575 0
16.6 0
16.625 0
16.65 0
16.675 0
16.7 0
16.725 0
16.75 0
16.775 0
16.8 0
16.825 0
16.85 0
16.875 0
16.9 0
16.925 0
16.95 0
16.975 0
17 0
17.025 0
17.05 0
17.075 0
17.1 0
17.125 0
17.15 0
17.175 0
17.2 0
17.225 0
17.25 0
17.275 0
17.3 0
17.325 0
17.35 0
17.375 0
17.4 0
17.425 0
17.45 0
17.475 0
17.5 0
17.525 0
17.55 0
17.575 0
17.6 0
17.625 0
17.65 0
17.675 0
17.7 0
17.725 0
17.75 0
17.775 0
17.8 0
17.825 0
17.85 0
17.875 0
17.9 0
17.925 0
17.95 0
17.975 0
18 1.09816926088942
18.025 1.09046572205757
18.05 1.08276218322573
18.075 1.07505864439388
18.1 1.06735510556203
18.125 1.05965156673018
18.15 1.05194802789833
18.175 1.04424448906648
18.2 1.03654095023463
18.225 1.02883741140278
18.25 1.02113387257093
18.275 1.01343033373908
18.3 1.00572679490723
18.325 0.998023256075386
18.35 0.990319717243537
18.375 0.982616178411688
18.4 0.974912639579838
18.425 0.96720910074799
18.45 0.959505561916141
18.475 0.951802023084291
18.5 0.944098484252442
18.525 0.936394945420593
18.55 0.928691406588744
18.575 0.920987867756895
18.6 0.913284328925046
18.625 0.905580790093197
18.65 0.897877251261347
18.675 0.890173712429499
18.7 0.88247017359765
18.725 0.874766634765801
18.75 0.867063095933952
18.775 0.859359557102102
18.8 0.851656018270254
18.825 0.843952479438405
18.85 0.836248940606555
18.875 0.828545401774706
18.9 0.820841862942857
18.925 0.813138324111008
18.95 0.805434785279159
18.975 0.79773124644731
19 0.790027707615461
19.025 0.782324168783611
19.05 0.774620629951763
19.075 0.766917091119914
19.1 0.759213552288064
19.125 0.751510013456216
19.15 0.743806474624366
19.175 0.736102935792517
19.2 0.728399396960667
19.225 0.720695858128819
19.25 0.71299231929697
19.275 0.705288780465121
19.3 0.697585241633272
19.325 0.689881702801422
19.35 0.682178163969574
19.375 0.674474625137725
19.4 0.666771086305875
19.425 0.659067547474027
19.45 0.651364008642177
19.475 0.643660469810328
19.5 0.63595693097848
19.525 0.62825339214663
19.55 0.620549853314781
19.575 0.612846314482931
19.6 0.605142775651083
19.625 0.597439236819234
19.65 0.589735697987385
19.675 0.582032159155536
19.7 0.574328620323686
19.725 0.566625081491838
19.75 0.558921542659989
19.775 0.551218003828139
19.8 0.543514464996291
19.825 0.535810926164441
19.85 0.528107387332592
19.875 0.520403848500744
19.9 0.512700309668894
19.925 0.504996770837045
19.95 0.497293232005195
19.975 0.489589693173347
20 0.481886154341498
20.025 0.474182615509648
20.05 0.4664790766778
20.075 0.45877553784595
20.1 0.451071999014101
20.125 0.443368460182253
20.15 0.435664921350403
20.175 0.427961382518554
20.2 0.420257843686705
20.225 0.412554304854856
20.25 0.404850766023008
20.275 0.397147227191158
20.3 0.389443688359309
20.325 0.381740149527459
20.35 0.374036610695611
20.375 0.366333071863762
20.4 0.358629533031912
20.425 0.350925994200064
20.45 0.343222455368214
20.475 0.335518916536365
20.5 0.329900540129375
20.525 0.326367326147246
20.55 0.322834112165118
20.575 0.31930089818299
20.6 0.315767684200862
20.625 0.312234470218734
20.65 0.308701256236605
20.675 0.305168042254477
20.7 0.301634828272349
20.725 0.298101614290221
20.75 0.294568400308093
20.775 0.291035186325964
20.8 0.287501972343836
20.825 0.283968758361708
20.85 0.28043554437958
20.875 0.276902330397452
20.9 0.273369116415323
20.925 0.269835902433195
20.95 0.266302688451067
20.975 0.262769474468939
21 0.259236260486811
21.025 0.255703046504683
21.05 0.252169832522555
21.075 0.248636618540426
21.1 0.245103404558298
21.125 0.24157019057617
21.15 0.238036976594042
21.175 0.234503762611914
21.2 0.230970548629785
21.225 0.227437334647657
21.25 0.223904120665529
21.275 0.220370906683401
21.3 0.216837692701273
21.325 0.213304478719144
21.35 0.209771264737016
21.375 0.206238050754888
21.4 0.20270483677276
21.425 0.199171622790632
21.45 0.195638408808503
21.475 0.192105194826375
21.5 0.188571980844247
21.525 0.185038766862119
21.55 0.181505552879991
21.575 0.177972338897863
21.6 0.174439124915735
21.625 0.170905910933607
21.65 0.167372696951478
21.675 0.16383948296935
21.7 0.160306268987222
21.725 0.156773055005094
21.75 0.153239841022966
21.775 0.149706627040837
21.8 0.146173413058709
21.825 0.142640199076581
21.85 0.139106985094453
21.875 0.135573771112325
21.9 0.132040557130196
21.925 0.128507343148068
21.95 0.12497412916594
21.975 0.121440915183812
22 0.117907701201684
22.025 0.114374487219555
22.05 0.110841273237427
22.075 0.107308059255299
22.1 0.103774845273171
22.125 0.100241631291043
22.15 0.0967084173089146
22.175 0.0931752033267866
22.2 0.0896419893446581
22.225 0.0861087753625302
22.25 0.0825755613804023
22.275 0.0790423473982737
22.3 0.0755091334161457
22.325 0.0719759194340173
22.35 0.0684427054518893
22.375 0.0649094914697612
22.4 0.0613762774876328
22.425 0.0578430635055048
22.45 0.0543098495233763
22.475 0.0507766355412484
22.5 0.0472434215591204
22.525 0.0437102075769919
22.55 0.0401769935948639
22.575 0.0366437796127355
22.6 0.0331105656306074
22.625 0.0295773516484794
22.65 0.026044137666351
22.675 0.022510923684223
22.7 0.0189777097020945
22.725 0.0154444957199666
22.75 0.0119112817378387
22.775 0.00837806775571018
22.8 0.00484485377358216
22.825 0.00131163979145373
22.85 -0.00222157419067429
22.875 -0.0057547881728023
22.9 -0.00928800215493073
22.925 -0.0128212161370588
22.95 -0.0163544301191873
22.975 -0.0198876441013152
23 0
23.025 0
23.05 0
23.075 0
23.1 0
23.125 0
23.15 0
23.175 0
23.2 0
23.225 0
23.25 0
23.275 0
23.3 0
23.325 0
23.35 0
23.375 0
23.4 0
23.425 0
23.45 0
23.475 0
23.5 0
23.525 0
23.55 0
23.575 0
23.6 0
23.625 0
23.65 0
23.675 0
23.7 0
23.725 0
23.75 0
23.775 0
23.8 0
23.825 0
23.85 0
23.875 0
23.9 0
23.925 0
23.95 0
23.975 0
24 0
24.025 0
24.05 0
24.075 0
24.1 0
24.125 0
24.15 0
24.175 0
24.2 0
24.225 0
24.25 0
24.275 0
24.3 0
24.325 0
24.35 0
24.375 0
24.4 0
24.425 0
24.45 0
24.475 0
24.5 0
24.525 0
24.55 0
24.575 0
24.6 0
24.625 0
24.65 0
24.675 0
24.7 0
24.725 0
24.75 0
24.775 0
24.8 0
24.825 0
24.85 0
24.875 0
24.9 0
24.925 0
24.95 0
24.975 0
25 0
25.025 0
25.05 0
25.075 0
25.1 0
25.125 0
25.15 0
25.175 0
25.2 0
25.225 0
25.25 0
25.275 0
25.3 0
25.325 0
25.35 0
25.375 0
25.4 0
25.425 0
25.45 0
25.475 0
25.5 0
25.525 0
25.55 0
25.575 0
25.6 0
25.625 0
25.65 0
25.675 0
25.7 0
25.725 0
25.75 0
25.775 0
25.8 0
25.825 0
25.85 0
25.875 0
25.9 0
25.925 0
25.95 0
25.975 0
26 0
26.025 0
26.05 0
26.075 0
26.1 0
26.125 0
26.15 0
26.175 0
26.2 0
26.225 0
26.25 0
26.275 0
26.3 0
26.325 0
26.35 0
26.375 0
26.4 0
26.425 0
26.45 0
26.475 0
26.5 0
26.525 0
26.55 0
26.575 0
26.6 0
26.625 0
26.65 0
26.675 0
26.7 0
26.725 0
26.75 0
26.775 0
26.8 0
26.825 0
26.85 0
26.875 0
26.9 0
26.925 0
26.95 0
26.975 0
27 0
27.025 0
27.05 0
27.075 0
27.1 0
27.125 0
27.15 0
27.175 0
27.2 0
27.225 0
27.25 0
27.275 0
27.3 0
27.325 0
27.35 0
27.375 0
27.4 0
27.425 0
27.45 0
27.475 0
27.5 0
27.525 0
27.55 0
27.575 0
27.6 0
27.625 0
27.65 0
27.675 0
27.7 0
27.725 0
27.75 0
27.775 0
27.8 0
27.825 0
27.85 0
27.875 0
27.9 0
27.925 0
27.95 0
27.975 0
28 0
28.025 0
28.05 0
28.075 0
28.1 0
28.125 0
28.15 0
28.175 0
28.2 0
28.225 0
28.25 0
28.275 0
28.3 0
28.325 0
28.35 0
28.375 0
28.4 0
28.425 0
28.45 0
28.475 0
28.5 0
28.525 0
28.55 0
28.575 0
28.6 0
28.625 0
28.65 0
28.675 0
28.7 0
28.725 0
28.75 0
28.775 0
28.8 0
28.825 0
28.85 0
28.875 0
28.9 0
28.925 0
28.95 0
28.975 0
29 0
29.025 0
29.05 0
29.075 0
29.1 0
29.125 0
29.15 0
29.175 0
29.2 0
29.225 0
29.25 0
29.275 0
29.3 0
29.325 0
29.35 0
29.375 0
29.4 0
29.425 0
29.45 0
29.475 0
29.5 0
29.525 0
29.55 0
29.575 0
29.6 0
29.625 0
29.65 0
29.675 0
29.7 0
29.725 0
29.75 0
29.775 0
29.8 0
29.825 0
29.85 0
29.875 0
29.9 0
29.925 0
29.95 0
29.975 0
30 0
30.025 0
30.05 0
30.075 0
30.1 0
30.125 0
30.15 0
30.175 0
30.2 0
30.225 0
30.25 0
30.275 0
30.3 0
30.325 0
30.35 0
30.375 0
30.4 0
30.425 0
30.45 0
30.475 0
30.5 0
30.525 0
30.55 0
30.575 0
30.6 0
30.625 0
30.65 0
30.675 0
30.7 0
30.725 0
30.75 0
30.775 0
30.8 0
30.825 0
30.85 0
30.875 0
30.9 0
30.925 0
30.95 0
30.975 0
31 0
31.025 0
31.05 0
31.075 0
31.1 0
31.125 0
31.15 0
31.175 0
31.2 0
31.225 0
31.25 0
31.275 0
31.3 0
31.325 0
31.35 0
31.375 0
31.4 0
31.425 0
31.45 0
31.475 0
31.5 0
31.525 0
31.55 0
31.575 0
31.6 0
31.625 0
31.65 0
31.675 0
31.7 0
31.725 0
31.75 0
31.775 0
31.8 0
31.825 0
31.85 0
31.875 0
31.9 0
31.925 0
31.95 0
31.975 0
32 0
32.025 0
32.05 0
32.075 0
32.1 0
32.125 0
32.15 0
32.175 0
32.2 0
32.225 0
32.25 0
32.275 0
32.3 0
32.325 0
32.35 0
32.375 0
32.4 0
32.425 0
32.45 0
32.475 0
32.5 0
32.525 0
32.55 0
32.575 0
32.6 0
32.625 0
32.65 0
32.675 0
32.7 0
32.725 0
32.75 0
32.775 0
32.8 0
32.825 0
32.85 0
32.875 0
32.9 0
32.925 0
32.95 0
32.975 0
33 0
33.025 0
33.05 0
33.075 0
33.1 0
33.125 0
33.15 0
33.175 0
33.2 0
33.225 0
33.25 0
33.275 0
33.3 0
33.325 0
33.35 0
33.375 0
33.4 0
33.425 0
33.45 0
33.475 0
33.5 0
33.525 0
33.55 0
33.575 0
33.6 0
33.625 0
33.65 0
33.675 0
33.7 0
33.725 0
33.75 0
33.775 0
33.8 0
33.825 0
33.85 0
33.875 0
33.9 0
33.925 0
33.95 0
33.975 0
34 0
34.025 0
34.05 0
34.075 0
34.1 0
34.125 0
34.15 0
34.175 0
34.2 0
34.225 0
34.25 0
34.275 0
34.3 0
34.325 0
34.35 0
34.375 0
34.4 0
34.425 0
34.45 0
34.475 0
34.5 0
34.525 0
34.55 0
34.575 0
34.6 0
34.625 0
34.65 0
34.675 0
34.7 0
34.725 0
34.75 0
34.775 0
34.8 0
34.825 0
34.85 0
34.875 0
34.9 0
34.925 0
34.95 0
34.975 0
35 0
35.025 0
35.05 0
35.075 0
35.1 0
35.125 0
35.15 0
35.175 0
35.2 0
35.225 0
35.25 0
35.275 0
35.3 0
35.325 0
35.35 0
35.375 0
35.4 0
35.425 0
35.45 0
35.475 0
35.5 0
35.525 0
35.55 0
35.575 0
35.6 0
35.625 0
35.65 0
35.675 0
35.7 0
35.725 0
35.75 0
35.775 0
35.8 0
35.825 0
35.85 0
35.875 0
35.9 0
35.925 0
35.95 0
35.975 0
36 0
36.025 0
36.05 0
36.075 0
36.1 0
36.125 0
36.15 0
36.175 0
36.2 0
36.225 0
36.25 0
36.275 0
36.3 0
36.325 0
36.35 0
36.375 0
36.4 0
36.425 0
36.45 0
36.475 0
36.5 0
36.525 0
36.55 0
36.575 0
36.6 0
36.625 0
36.65 0
36.675 0
36.7 0
36.725 0
36.75 0
36.775 0
36.8 0
36.825 0
36.85 0
36.875 0
36.9 0
36.925 0
36.95 0
36.975 0
37 0
37.025 0
37.05 0
37.075 0
37.1 0
37.125 0
37.15 0
37.175 0
37.2 0
37.225 0
37.25 0
37.275 0
37.3 0
37.325 0
37.35 0
37.375 0
37.4 0
37.425 0
37.45 0
37.475 0
37.5 0
37.525 0
37.55 0
37.575 0
37.6 0
37.625 0
37.65 0
37.675 0
37.7 0
37.725 0
37.75 0
37.775 0
37.8 0
37.825 0
37.85 0
37.875 0
37.9 0
37.925 0
37.95 0
37.975 0
38 0
38.025 0
38.05 0
38.075 0
38.1 0
38.125 0
38.15 0
38.175 0
38.2 0
38.225 0
38.25 0
38.275 0
38.3 0
38.325 0
38.35 0
38.375 0
38.4 0
38.425 0
38.45 0
38.475 0
38.5 0
38.525 0
38.55 0
38.575 0
38.6 0
38.625 0
38.65 0
38.675 0
38.7 0
38.725 0
38.75 0
38.775 0
38.8 0
38.825 0
38.85 0
38.875 0
38.9 0
38.925 0
38.95 0
38.975 0
39 0
39.025 0
39.05 0
39.075 0
39.1 0
39.125 0
39.15 0
39.175 0
39.2 0
39.225 0
39.25 0
39.275 0
39.3 0
39.325 0
39.35 0
39.375 0
39.4 0
39.425 0
39.45 0
39.475 0
39.5 0
39.525 0
39.55 0
39.575 0
39.6 0
39.625 0
39.65 0
39.675 0
39.7 0
39.725 0
39.75 0
39.775 0
39.8 0
39.825 0
39.85 0
39.875 0
39.9 0
39.925 0
39.95 0
39.975 0
40 0
40.025 0
40.05 0
40.075 0
40.1 0
40.125 0
40.15 0
40.175 0
40.2 0
40.225 0
40.25 0
40.275 0
40.3 0
40.325 0
40.35 0
40.375 0
40.4 0
40.425 0
40.45 0
40.475 0
40.5 0
40.525 0
40.55 0
40.575 0
40.6 0
40.625 0
40.65 0
40.675 0
40.7 0
40.725 0
40.75 0
40.775 0
40.8 0
40.825 0
40.85 0
40.875 0
40.9 0
40.925 0
40.95 0
40.975 0
41 0
41.025 0
41.05 0
41.075 0
41.1 0
41.125 0
41.15 0
41.175 0
41.2 0
41.225 0
41.25 0
41.275 0
41.3 0
41.325 0
41.35 0
41.375 0
41.4 0
41.425 0
41.45 0
41.475 0
41.5 0
41.525 0
41.55 0
41.575 0
41.6 0
41.625 0
41.65 0
41.675 0
41.7 0
41.725 0
41.75 0
41.775 0
41.8 0
41.825 0
41.85 0
41.875 0
41.9 0
41.925 0
41.95 0
41.975 0
42 0
42.025 0
42.05 0
42.075 0
42.1 0
42.125 0
42.15 0
42.175 0
42.2 0
42.225 0
42.25 0
42.275 0
42.3 0
42.325 0
};
\addplot [ultra thick, black, forget plot]
table {%
9 -1.90619599582758
9 2.38951469489845
};
\addplot [ultra thick, black, forget plot]
table {%
11 -1.90619599582758
11 2.38951469489845
};
\addplot [ultra thick, black, forget plot]
table {%
18 -1.90619599582758
18 2.38951469489845
};
\addplot [ultra thick, black, forget plot]
table {%
23 -1.90619599582758
23 2.38951469489845
};
\node at (axis cs:4.5,-1.5)[
  anchor=base,
  text=black,
  rotate=0.0
]{ A};
\node at (axis cs:10,-1.5)[
  anchor=base,
  text=black,
  rotate=0.0
]{ B};
\node at (axis cs:14.5,-1.5)[
  anchor=base,
  text=black,
  rotate=0.0
]{ C};
\node at (axis cs:20.5,-1.5)[
  anchor=base,
  text=black,
  rotate=0.0
]{ D};
\node at (axis cs:31.5,-1.5)[
  anchor=base,
  text=black,
  rotate=0.0
]{ E};
\end{axis}

\end{tikzpicture}%
		\subcaption{Longitudinal acceleration.}
		\label{fig:example ego states lon acc}
	\end{minipage}%
	\caption{Longitudinal position and acceleration of the ego vehicle. The blue dots show the measurements and the orange dotted line show the result of the models that are employed for modeling the position.}
	\label{fig:example ego states pos acc}
\end{figure}


\section{Dynamic environment}

\Cref{fig:example dynamic environment} shows the longitudinal and lateral position of the pickup truck and the station wagon with respect to the ego vehicle. The station wagon accelerates and performs a lane change at approximately $t\approx10$ s. At $t\approx12$ s, the station wagon overtakes the ego vehicle. The ego vehicle overtakes the pickup truck at $t\approx25$ s, as can be seen by the relative longitudinal distance of the pickup truck in \cref{fig:example dynamic environment lon}, which becomes negative at $t\approx25$ s. The events related to the pickup truck and the station wagon are shown in Figure~\ref{fig:example events}. Here, it can be seen that the pickup truck is in the field of view of the ego vehicle's sensors from $t=1$ s until $t=34$ s. The station wagon is in the field of view of the ego vehicle's sensors from $t=0$ s until $t=18$ s.

\setlength\figureheight{150pt}
\setlength\figurewidth{0.49\linewidth}
\begin{figure}
	\centering
	\begin{minipage}[b]{0.499\linewidth}
		\centering
		% This file was created by matplotlib2tikz v0.6.14.
\begin{tikzpicture}

\definecolor{color0}{rgb}{0.172549019607843,0.627450980392157,0.172549019607843}
\definecolor{color3}{rgb}{0.549019607843137,0.337254901960784,0.294117647058824}
\definecolor{color1}{rgb}{0.83921568627451,0.152941176470588,0.156862745098039}
\definecolor{color2}{rgb}{0.580392156862745,0.403921568627451,0.741176470588235}

\begin{axis}[
xlabel={$t$ [s]},
ylabel={Relative longitudinal distance [m]},
xmin=0, xmax=34.1600000858307,
ymin=-61.3809092368509, ymax=111.328903581739,
width=\figurewidth,
height=\figureheight,
tick align=outside,
tick pos=left,
xmajorgrids,
x grid style={lightgray!92.026143790849673!black},
ymajorgrids,
y grid style={lightgray!92.026143790849673!black},
clip marker paths
]
\addplot [semithick, color0, mark=*, mark size=1, mark options={solid}, only marks, forget plot]
table {%
0 -29.7074271603815
0.0400002002716064 -29.7193473728894
0.0800001621246338 -29.7383384251862
0.120000123977661 -29.7460490026369
0.160000085830688 -29.7664243002673
0.200000047683716 -29.779038823357
0.240000009536743 -29.7981872784185
0.28000020980835 -29.8295946814196
0.320000171661377 -29.8472175355928
0.360000133514404 -29.8765110027307
0.400000095367432 -29.8976779480472
0.440000057220459 -29.9161635570545
0.480000019073486 -29.9500390967969
0.519999980926514 -29.9626746751619
0.56000018119812 -29.9931910108844
0.600000143051147 -30.0158738229457
0.640000104904175 -30.0470681455827
0.680000066757202 -30.0739041547095
0.720000028610229 -30.1103165653785
0.759999990463257 -30.1499582959623
0.800000190734863 -30.1749450570478
0.840000152587891 -30.2110858703618
0.880000114440918 -30.2457953815992
0.920000076293945 -30.263344395089
0.960000038146973 -30.2944718244144
1 -30.3275244008983
1.04000020027161 -30.3461533459049
1.08000016212463 -30.404412105634
1.12000012397766 -30.4299916279488
1.16000008583069 -30.469466645427
1.20000004768372 -30.5036813379666
1.24000000953674 -30.5424287676051
1.28000020980835 -30.5992346562234
1.32000017166138 -30.6254737172967
1.3600001335144 -30.6676030564413
1.40000009536743 -30.7052177632922
1.44000005722046 -30.7366472687481
1.48000001907349 -30.7958631867077
1.51999998092651 -30.8290239459493
1.56000018119812 -30.8624309720108
1.60000014305115 -30.9005839226102
1.64000010490417 -30.936000716034
1.6800000667572 -30.9846167561773
1.72000002861023 -31.0065385651487
1.75999999046326 -31.0374855482733
1.80000019073486 -31.0730068300436
1.84000015258789 -31.1073708202912
1.88000011444092 -31.1585357137046
1.92000007629395 -31.1704042716519
1.96000003814697 -31.1984069697191
2 -31.2331737318946
2.04000020027161 -31.2569818936063
2.08000016212463 -31.3134767408137
2.12000012397766 -31.33457098291
2.16000008583069 -31.3772719417411
2.20000004768372 -31.4094436982869
2.24000000953674 -31.4405307212692
2.28000020980835 -31.4909659151326
2.32000017166138 -31.5102945930848
2.3600001335144 -31.5454310228088
2.40000009536743 -31.5738279901016
2.44000005722046 -31.6058762154462
2.48000001907349 -31.6507218298466
2.51999998092651 -31.6632214247602
2.56000018119812 -31.6973162838203
2.60000014305115 -31.7268398443266
2.64000010490417 -31.7571284448313
2.6800000667572 -31.808816608831
2.72000002861023 -31.8266725322683
2.75999999046326 -31.8658204385556
2.80000019073486 -31.8940627079883
2.84000015258789 -31.924219830471
2.88000011444092 -31.9840174570636
2.92000007629395 -32.0051623854943
2.96000003814697 -32.0588416641222
3 -32.0921224670674
3.04000020027161 -32.1367741825361
3.08000016212463 -32.196525812773
3.12000012397766 -32.2202876176016
3.16000008583069 -32.2727025826334
3.20000004768372 -32.3074194994151
3.24000000953674 -32.3521726259696
3.28000020980835 -32.4058262929539
3.32000017166138 -32.4296208923552
3.3600001335144 -32.4786247010088
3.40000009536743 -32.5130321161469
3.44000005722046 -32.5584644286209
3.48000001907349 -32.6141678829208
3.51999998092651 -32.6406701762025
3.56000018119812 -32.6823762718086
3.60000014305115 -32.7264484987882
3.64000010490417 -32.7713429856267
3.6800000667572 -32.8255580331006
3.72000002861023 -32.856951402975
3.75999999046326 -32.910074862717
3.80000019073486 -32.9485975998832
3.84000015258789 -32.9895311112195
3.88000011444092 -33.0398992823993
3.92000007629395 -33.0691890743856
3.96000003814697 -33.0977602159146
4 -33.159594566323
4.04000020027161 -33.1731866178852
4.08000016212463 -33.2148813720869
4.12000012397766 -33.2496221080928
4.16000008583069 -33.2803923345346
4.20000004768372 -33.3199595746482
4.24000000953674 -33.3390741809362
4.28000020980835 -33.3788283571612
4.32000017166138 -33.4023925829551
4.3600001335144 -33.4338549808817
4.40000009536743 -33.471955628771
4.44000005722046 -33.4928417260071
4.48000001907349 -33.5457550454139
4.51999998092651 -33.567108358624
4.56000018119812 -33.6140201087146
4.60000014305115 -33.6598147006916
4.64000010490417 -33.6878860990491
4.6800000667572 -33.7355324311775
4.72000002861023 -33.7683442772741
4.75999999046326 -33.8261440587339
4.80000019073486 -33.8687390755385
4.84000015258789 -33.9099162060938
4.88000011444092 -33.9738719872839
4.92000007629395 -33.9980793437735
4.96000003814697 -34.0767215251435
5 -34.1119637704469
5.04000020027161 -34.1593912526369
5.08000016212463 -34.214012101389
5.12000012397766 -34.2728830271117
5.16000008583069 -34.3172987296166
5.20000004768372 -34.383580664924
5.24000000953674 -34.4472430902279
5.28000020980835 -34.4886011693106
5.32000017166138 -34.5505265284828
5.3600001335144 -34.5957848910039
5.40000009536743 -34.6659730994397
5.44000005722046 -34.714066504248
5.48000001907349 -34.7630324400961
5.51999998092651 -34.8323117799755
5.56000018119812 -34.8743756620715
5.60000014305115 -34.9418214232555
5.64000010490417 -34.9858726325892
5.6800000667572 -35.0257822271687
5.72000002861023 -35.0993768488479
5.75999999046326 -35.1364252221247
5.80000019073486 -35.2288924718887
5.84000015258789 -35.2767301638487
5.88000011444092 -35.320289510335
5.92000007629395 -35.3877188390488
5.96000003814697 -35.4358495067754
6 -35.5199657559551
6.04000020027161 -35.60914191743
6.08000016212463 -35.6483753038101
6.12000012397766 -35.7333624830098
6.16000008583069 -35.7971334856502
6.20000004768372 -35.8847336545605
6.24000000953674 -35.9162921739353
6.28000020980835 -35.9859252078058
6.32000017166138 -36.0500405582061
6.3600001335144 -36.0750500777667
6.40000009536743 -36.1402647152117
6.44000005722046 -36.2106582522938
6.48000001907349 -36.2388524432699
6.51999998092651 -36.3051267675819
6.56000018119812 -36.3660438160696
6.60000014305115 -36.4123215717282
6.64000010490417 -36.4695835284092
6.6800000667572 -36.5204593584458
6.72000002861023 -36.5878880129167
6.75999999046326 -36.6522965958156
6.80000019073486 -36.7036051560499
6.84000015258789 -36.7577059540854
6.88000011444092 -36.8058453285303
6.92000007629395 -36.8648800318606
6.96000003814697 -36.9153492368514
7 -36.9676276877672
7.04000020027161 -37.0200439129221
7.08000016212463 -37.0696535416864
7.12000012397766 -37.1215291295975
7.16000008583069 -37.1879415296498
7.20000004768372 -37.2439704962035
7.24000000953674 -37.3016181066505
7.28000020980835 -37.3486455192688
7.32000017166138 -37.4160809300884
7.3600001335144 -37.4559836321223
7.40000009536743 -37.5210997451741
7.44000005722046 -37.5828120242459
7.48000001907349 -37.6020625552355
7.51999998092651 -37.6591875650829
7.56000018119812 -37.6927522353089
7.60000014305115 -37.7192409635682
7.64000010490417 -37.787252095015
7.6800000667572 -37.8113144948238
7.72000002861023 -37.848399161061
7.75999999046326 -37.8717833188257
7.80000019073486 -37.8933181017856
7.84000015258789 -37.9485304696154
7.88000011444092 -37.9624663773448
7.92000007629395 -38.0146831032653
7.96000003814697 -38.0409629800179
8 -38.0730814660183
8.04000020027161 -38.0862488593975
8.08000016212463 -38.099714781576
8.12000012397766 -38.1352397092269
8.16000008583069 -38.1369820298005
8.20000004768372 -38.147280783267
8.24000000953674 -38.1719449717675
8.28000020980835 -38.1563860035712
8.32000017166138 -38.1864213254939
8.3600001335144 -38.1782663117428
8.40000009536743 -38.1817153648408
8.44000005722046 -38.1742719525719
8.48000001907349 -38.161119165954
8.51999998092651 -38.1688152827337
8.56000018119812 -38.1701919391071
8.60000014305115 -38.1656855714955
8.64000010490417 -38.1476529812517
8.6800000667572 -38.1142517326825
8.72000002861023 -38.0963731100655
8.75999999046326 -38.0762722647596
8.80000019073486 -38.0255003242382
8.84000015258789 -38.0151752920829
8.88000011444092 -37.9693337524677
8.92000007629395 -37.9217376187516
8.96000003814697 -37.8897895912669
9 -37.812877766115
9.04000020027161 -37.7924157517209
9.08000016212463 -37.7151076064274
9.12000012397766 -37.640522485126
9.16000008583069 -37.5990262113537
9.20000004768372 -37.4940963246008
9.24000000953674 -37.4606192931769
9.28000020980835 -37.3371298010206
9.32000017166138 -37.23826584788
9.3600001335144 -37.1515125689275
9.40000009536743 -36.9918119513386
9.44000005722046 -36.9285336874655
9.48000001907349 -36.7962827252013
9.51999998092651 -36.6931405823179
9.56000018119812 -36.5694319435843
9.60000014305115 -36.4851381752378
9.64000010490417 -36.2599097684451
9.6800000667572 -36.1117586445798
9.72000002861023 -35.9809530615439
9.75999999046326 -35.8077820722556
9.80000019073486 -35.7035651437564
9.84000015258789 -35.4261585182394
9.88000011444092 -35.2410680008943
9.92000007629395 -35.0992599954934
9.96000003814697 -34.8726392893768
10 -34.7454597982378
10.0400002002716 -34.4411446770082
10.0800001621246 -34.233127927575
10.1200001239777 -34.0313373131121
10.1600000858307 -33.8354171908941
10.2000000476837 -33.5654751741386
10.2400000095367 -33.3117328715616
10.2800002098083 -33.0746374579539
10.3200001716614 -32.7310973859676
10.3600001335144 -32.5877009884753
10.4000000953674 -32.2912747910868
10.4400000572205 -32.0059770649514
10.4800000190735 -31.74191626227
10.5199999809265 -31.4421755285694
10.5600001811981 -31.2038473013163
10.6000001430511 -30.8992015238127
10.6400001049042 -30.5676398735959
10.6800000667572 -30.2787464374487
10.7200000286102 -29.9570306091809
10.7599999904633 -29.6933969814163
10.8000001907349 -29.3992908669352
10.8400001525879 -29.0010395113823
10.8800001144409 -28.6816084341845
10.9200000762939 -28.2293537422502
10.960000038147 -28.0533175649962
11 -27.7481268427546
11.0400002002716 -27.1939342429196
11.0800001621246 -26.9676291924116
11.1200001239777 -26.4778014131789
11.1600000858307 -26.2981569862113
11.2000000476837 -25.7823998891727
11.2400000095367 -25.4005123981951
11.2800002098083 -25.1704446725089
11.3200001716614 -24.6509466389034
11.3600001335144 -24.4665139782737
11.4000000953674 -23.9171570582475
11.4400000572205 -23.5083246894446
11.4800000190735 -23.2748806240961
11.5199999809265 -22.9254589927186
11.5600001811981 -22.5327539030077
11.6000001430511 -22.1384555607128
11.6400001049042 -21.7594986876302
11.6800000667572 -21.3516546520277
11.7200000286102 -20.9537587907071
11.7599999904633 -20.5428801713952
11.8000001907349 -20.0988092930475
11.8400001525879 -19.7071988682437
11.8800001144409 -19.2228427529953
11.9200000762939 -18.6044093580276
11.960000038147 -18.4035304260287
12 -17.9787021192078
12.0400002002716 -17.5429201174975
12.0800001621246 -17.079700822349
12.1200001239777 -16.6419597842887
12.1600000858307 -16.1375779392056
12.2000000476837 -15.7486996354583
12.2400000095367 -15.4589798666839
12.2800002098083 -14.8264947746738
12.3200001716614 -14.1703211419299
12.3600001335144 -13.8968332323839
12.4000000953674 -13.7114346185263
12.4400000572205 -12.9626582077162
12.4800000190735 -12.4685804669971
12.5199999809265 -12.2724841421514
12.5600001811981 -11.509945696087
12.6000001430511 -11.0268636638066
12.6400001049042 -10.7084681175693
12.6800000667572 -10.0204804755667
12.7200000286102 -9.3287742265511
12.7599999904633 -8.91321427856565
12.8000001907349 -8.21998669377172
12.8400001525879 -7.80847770170476
12.8800001144409 -7.268529564657
12.9200000762939 -6.72207962568973
12.960000038147 -6.23143474607969
13 -5.53360740659106
13.0400002002716 -4.8357759077234
13.0800001621246 -4.13794856823478
13.1200001239777 -3.44012122874616
13.1600000858307 -2.74229388925754
13.2000000476837 -2.04446654976891
13.2400000095367 -1.34663921028029
13.2800002098083 -0.648807711412632
13.3200001716614 0.0490196280759907
13.3600001335144 0.746846967564613
13.4000000953674 1.44467430705324
13.4400000572205 2.14250164654186
13.4800000190735 2.84032898603048
13.5199999809265 3.5381563255191
13.5600001811981 4.23598782438676
13.6000001430511 4.93381516387539
13.6400001049042 5.63164250336401
13.6800000667572 6.32946984285263
13.7200000286102 6.59663515279499
13.7599999904633 6.91507421624192
13.8000001907349 7.36358850169745
13.8400001525879 7.51360365835171
13.8800001144409 7.94167528904836
13.9200000762939 8.37009874303658
13.960000038147 8.68058540081438
14 8.91829565265107
14.0400002002716 9.67761481379603
14.0800001621246 10.3013219123295
14.1200001239777 10.6365294901971
14.1600000858307 11.1594149825978
14.2000000476837 11.7254737591356
14.2400000095367 12.1136385812315
14.2800002098083 12.5756518553208
14.3200001716614 13.0625135830778
14.3600001335144 13.4497263736266
14.4000000953674 13.9732587770159
14.4400000572205 14.3082878775567
14.4800000190735 14.9950598698833
14.5199999809265 15.652072886156
14.5600001811981 15.9875154633537
14.6000001430511 16.3510255465499
14.6400001049042 16.9892150072556
14.6800000667572 17.459709168912
14.7200000286102 18.0131879948458
14.7599999904633 18.4210395850077
14.8000001907349 18.8464314271132
14.8400001525879 19.3309100337337
14.8800001144409 20.0135110500887
14.9200000762939 20.2123697165607
14.960000038147 21.0479691289111
15 21.3648939073973
15.2000000476837 24.3686649939609
15.2400000095367 24.4334943564572
15.2800002098083 24.5685119330574
15.3200001716614 24.6704316239793
15.3600001335144 25.498243723534
15.4000000953674 25.5499004733701
15.4400000572205 26.2889026424637
15.4800000190735 26.4228098014428
15.5199999809265 26.5489864090632
15.5600001811981 26.6389803841303
15.6000001430511 26.9724396863803
15.6400001049042 27.1662413898084
15.6800000667572 27.2506874450501
15.7200000286102 28.6347522694105
15.7599999904633 29.8581646424063
15.8000001907349 30.2339136086321
15.8400001525879 30.3084825929363
15.8800001144409 31.4282980904154
15.9200000762939 32.3824495788267
15.960000038147 32.700553968125
16 33.6328035879069
16.0400002002716 33.8673750668913
16.0800001621246 34.381195706681
16.1200001239777 35.4193410273729
16.1600000858307 36.0757201501856
16.2000000476837 36.8633581596714
16.2400000095367 37.1473289805945
16.2800002098083 38.0630464370497
16.3200001716614 38.761957506611
16.3600001335144 39.1961462417075
16.4000000953674 39.937083341023
16.4400000572205 40.2649042567045
16.4800000190735 40.9478108444055
16.5199999809265 41.6263547744275
16.5600001811981 41.8948139790791
16.6000001430511 42.1654710781986
16.6400001049042 42.7387004257816
16.6800000667572 43.8411866273163
16.7200000286102 44.6673733483331
16.7599999904633 45.1244529573905
16.8000001907349 45.9967583665875
16.8400001525879 46.3654057657714
16.8800001144409 47.1395076419012
16.9200000762939 47.8641584779743
16.960000038147 48.2118127102403
17 48.9268941407809
17.0400002002716 49.2571265880215
17.0800001621246 49.8908053678169
17.1200001239777 50.1287872850517
17.1600000858307 51.0050614973097
17.2000000476837 51.6235918735074
17.2400000095367 52.3617076193077
17.2800002098083 52.7698168235875
17.3200001716614 53.5343237969464
17.3600001335144 53.9851337655946
17.4000000953674 54.1722581298818
17.4400000572205 55.2373630401016
17.4800000190735 55.8044304237519
17.5199999809265 56.0278847427908
17.5600001811981 57.0066121870695
17.6000001430511 57.6472276676432
17.6400001049042 58.1377461491975
17.6800000667572 58.8690445852444
17.7200000286102 59.737045744434
};
\addplot [semithick, color2, mark=*, mark size=1, mark options={solid}, only marks, forget plot]
table {%
0.800000190734863 101.481304680266
0.840000152587891 101.526985366871
0.880000114440918 101.565141825275
0.920000076293945 101.603370900992
0.960000038146973 101.646926877722
1 101.667385767189
1.04000020027161 101.69353465608
1.08000016212463 101.769457757408
1.12000012397766 101.813049728076
1.16000008583069 101.881585933119
1.20000004768372 101.943780187674
1.24000000953674 102.000734001884
1.28000020980835 102.077222044298
1.32000017166138 102.126680899333
1.3600001335144 102.110322336566
1.40000009536743 102.079564758151
1.44000005722046 102.055635798928
1.48000001907349 102.096734356823
1.51999998092651 102.126532141361
1.56000018119812 102.161250392581
1.60000014305115 102.18813955109
1.64000010490417 102.216541846394
1.6800000667572 102.252364719265
1.72000002861023 102.273045256319
1.75999999046326 102.336399366579
1.80000019073486 102.387245738459
1.84000015258789 102.435107366582
1.88000011444092 102.456775613236
1.92000007629395 102.55990562101
1.96000003814697 102.597046040619
2 102.62334959074
2.04000020027161 102.63779822287
2.08000016212463 102.726229786051
2.12000012397766 102.787255714891
2.16000008583069 102.805661899247
2.20000004768372 102.99257542323
2.24000000953674 103.01288461376
2.28000020980835 103.111777668471
2.32000017166138 103.194413163299
2.3600001335144 103.278176584017
2.40000009536743 103.38958907356
2.44000005722046 103.412386269027
2.48000001907349 103.47845754453
2.51999998092651 103.412045813178
2.56000018119812 103.347810031864
2.60000014305115 103.362673922806
2.64000010490417 103.289602193405
2.6800000667572 103.138496493411
2.72000002861023 103.065028071531
2.75999999046326 102.987044654885
2.80000019073486 102.909754278122
2.84000015258789 102.832580607483
2.88000011444092 102.815989859731
2.92000007629395 102.730714452638
2.96000003814697 102.581560204068
3 102.511840364392
3.04000020027161 102.436903193262
3.08000016212463 102.355229003415
3.12000012397766 102.332721774999
3.16000008583069 102.2535749293
3.20000004768372 102.157414712025
3.24000000953674 102.086958287555
3.28000020980835 101.999710693784
3.32000017166138 101.940805974398
3.3600001335144 101.832651298677
3.40000009536743 101.748329591424
3.44000005722046 101.656420665249
3.48000001907349 101.58479085585
3.51999998092651 101.549968598059
3.56000018119812 101.489744052611
3.60000014305115 101.432703987195
3.64000010490417 101.370124337356
3.6800000667572 101.296042011641
3.72000002861023 101.249227451573
3.75999999046326 101.167825882119
3.80000019073486 101.107207611258
3.84000015258789 101.035188781065
3.88000011444092 100.932492798865
3.92000007629395 100.884001395445
3.96000003814697 100.805232805578
4 100.701300281637
4.04000020027161 100.646773203041
4.08000016212463 100.537578849735
4.12000012397766 100.479427699776
4.16000008583069 100.378905875108
4.20000004768372 100.311561022285
4.24000000953674 100.23184603062
4.28000020980835 100.107443910925
4.32000017166138 100.086792307966
4.3600001335144 99.9634554600161
4.40000009536743 99.9118715765653
4.44000005722046 99.8326502830478
4.48000001907349 99.6991523945235
4.51999998092651 99.6752130427885
4.56000018119812 99.5582558588048
4.60000014305115 99.4875620480634
4.64000010490417 99.4227320332084
4.6800000667572 99.2846585254902
4.72000002861023 99.2659606480502
4.75999999046326 99.1239863613027
4.80000019073486 99.0389680428561
4.84000015258789 98.9748164943285
4.88000011444092 98.8400444662293
4.92000007629395 98.8021551418897
4.96000003814697 98.6503618025472
5 98.5737063134002
5.04000020027161 98.4724011675244
5.08000016212463 98.3350817488554
5.12000012397766 98.2544304569456
5.16000008583069 98.1617105926671
5.20000004768372 98.0490401065435
5.24000000953674 97.9434214909088
5.28000020980835 97.7951467589428
5.32000017166138 97.7160388562133
5.3600001335144 97.5758704656073
5.40000009536743 97.4677594501027
5.44000005722046 97.3651035750117
5.48000001907349 97.2280500001016
5.51999998092651 97.1267791856626
5.56000018119812 96.9988390636681
5.60000014305115 96.8766536203057
5.64000010490417 96.7502318796742
5.6800000667572 96.595016244979
5.72000002861023 96.4909765345074
5.75999999046326 96.3506675396111
5.80000019073486 96.2231530993176
5.84000015258789 96.0940445282467
5.88000011444092 95.9318985765749
5.92000007629395 95.8201664087719
5.96000003814697 95.6880559596902
6 95.5617145112883
6.04000020027161 95.4319133641111
6.08000016212463 95.290892333087
6.12000012397766 95.1832825687161
6.16000008583069 95.0511965021797
6.20000004768372 94.9054138146948
6.24000000953674 94.7717042434579
6.28000020980835 94.6209399415311
6.32000017166138 94.4793964985602
6.3600001335144 94.3239479667263
6.40000009536743 94.1716351935465
6.44000005722046 94.0197089062658
6.48000001907349 93.8496162846404
6.51999998092651 93.7080691100982
6.56000018119812 93.5053054968048
6.60000014305115 93.3972468432639
6.64000010490417 93.2392201299881
6.6800000667572 93.081127059726
6.72000002861023 92.9370776744418
6.75999999046326 92.7424178325964
6.80000019073486 92.6112349464693
6.84000015258789 92.4501744763838
6.88000011444092 92.2602529025735
6.92000007629395 92.1146002937967
6.96000003814697 91.901477039919
7 91.7567898048219
7.04000020027161 91.5933802124764
7.08000016212463 91.3856459812141
7.12000012397766 91.2318763106305
7.16000008583069 91.05425331474
7.20000004768372 90.8780010400769
7.24000000953674 90.6965430504169
7.28000020980835 90.4899699575835
7.32000017166138 90.3190719911581
7.3600001335144 90.1314329415472
7.40000009536743 89.8678534456667
7.44000005722046 89.7105932357827
7.48000001907349 89.5598101428404
7.51999998092651 89.385180428204
7.56000018119812 89.1950403277624
7.60000014305115 89.0904584861419
7.64000010490417 88.7740559649865
7.6800000667572 88.6523217211015
7.72000002861023 88.483337743306
7.75999999046326 88.2812435431842
7.80000019073486 88.179917384301
7.84000015258789 87.8431384234918
7.88000011444092 87.6960592417145
7.92000007629395 87.5249195184424
7.96000003814697 87.3316586626315
8 87.2234866174949
8.04000020027161 86.7690075841983
8.08000016212463 86.5583631163954
8.12000012397766 86.3758981975843
8.16000008583069 85.9997479503018
8.20000004768372 85.8797037677396
8.24000000953674 85.4146080198789
8.28000020980835 85.19328984811
8.32000017166138 84.9937010739686
8.3600001335144 84.7025286749194
8.40000009536743 84.5753696146949
8.44000005722046 84.1079965556528
8.48000001907349 83.8759391888416
8.51999998092651 83.6741330908808
8.56000018119812 83.3917022214828
8.60000014305115 83.2648741314861
8.64000010490417 82.9434520948016
8.6800000667572 82.6705405555949
8.72000002861023 82.3275562403887
8.75999999046326 82.1909267317787
8.80000019073486 81.8554675869018
8.84000015258789 81.7712238912536
8.88000011444092 81.4894668793258
8.92000007629395 81.252011492903
8.96000003814697 81.111811736153
9 80.7889219056851
9.04000020027161 80.6952749607754
9.08000016212463 80.4256673427735
9.12000012397766 80.1731179189155
9.16000008583069 80.023768318968
9.20000004768372 79.8262967206283
9.24000000953674 79.8892582716871
9.28000020980835 79.7639551922948
9.32000017166138 79.6147651942138
9.3600001335144 79.4572805592179
9.40000009536743 79.2871550613963
9.44000005722046 79.1624307891125
9.48000001907349 78.9950910817261
9.51999998092651 78.7732583720481
9.56000018119812 78.6711034929067
9.60000014305115 78.6076642039334
9.64000010490417 78.2965774110562
9.6800000667572 78.1482090560094
9.72000002861023 77.968358724529
9.75999999046326 77.8510652732966
9.80000019073486 77.7550080436522
9.84000015258789 77.483606921287
9.88000011444092 77.3648147486856
9.92000007629395 77.199954486081
9.96000003814697 77.0297278903963
10 76.9939979722676
10.0400002002716 76.6693078498174
10.0800001621246 76.5444326317156
10.1200001239777 76.3546367141134
10.1600000858307 76.1859152552315
10.2000000476837 75.9753319519295
10.2400000095367 75.8666503380427
10.2800002098083 75.7161943208976
10.3200001716614 75.4708107303068
10.3600001335144 75.3935057592043
10.4000000953674 75.163933548456
10.4400000572205 75.0476571020008
10.4800000190735 74.8952086077334
10.5199999809265 74.7158731643012
10.5600001811981 74.5807673227282
10.6000001430511 74.3446630600574
10.6400001049042 74.200188338189
10.6800000667572 74.0300311362025
10.7200000286102 73.8134511400185
10.7599999904633 73.6937341530866
10.8000001907349 73.5295471981026
10.8400001525879 73.3416382460964
10.8800001144409 73.1947327940961
10.9200000762939 72.9524152560025
10.960000038147 72.8766942431157
11 72.7442821386794
11.0400002002716 72.5081499953412
11.0800001621246 72.4115475297112
11.1200001239777 72.2131571529972
11.1600000858307 72.1400693315518
11.2000000476837 71.8960517075357
11.2400000095367 71.7618605381394
11.2800002098083 71.6672934896087
11.3200001716614 71.4594231316205
11.3600001335144 71.387142577576
11.4000000953674 71.1243966418333
11.4400000572205 71.0271575902225
11.4800000190735 70.8976694302601
11.5199999809265 70.8004433525639
11.5600001811981 70.6441063854381
11.6000001430511 70.4869642190315
11.6400001049042 70.3567480318907
11.6800000667572 70.1536921317584
11.7200000286102 70.0321304683785
11.7599999904633 69.9182711179819
11.8000001907349 69.7433549425223
11.8400001525879 69.596418219513
11.8800001144409 69.3925740738487
11.9200000762939 69.2369178161025
11.960000038147 69.1652093230332
12 68.9793267816203
12.0400002002716 68.8227919031087
12.0800001621246 68.6743533082499
12.1200001239777 68.5240155722659
12.1600000858307 68.3239246499761
12.2000000476837 68.1867124200489
12.2400000095367 68.1207254970323
12.2800002098083 67.8795839918312
12.3200001716614 67.7031552670123
12.3600001335144 67.5720209652227
12.4000000953674 67.5090173814588
12.4400000572205 67.2942849645933
12.4800000190735 67.0767863790734
12.5199999809265 67.0145269256827
12.5600001811981 66.7290996267184
12.6000001430511 66.6042031637789
12.6400001049042 66.5182594476209
12.6800000667572 66.2929317666913
12.7200000286102 66.0934819637423
12.7599999904633 65.9537550670393
12.8000001907349 65.7573855118571
12.8400001525879 65.6800711266769
12.8800001144409 65.5067331991486
12.9200000762939 65.3642591699263
12.960000038147 65.1874737130911
13 65.0314839285384
13.0400002002716 64.8260016394543
13.0800001621246 64.7241622982256
13.1200001239777 64.5517780405971
13.1600000858307 64.4496662423462
13.2000000476837 64.2750228351524
13.2400000095367 64.1197544686565
13.2800002098083 63.9956006658267
13.3200001716614 63.8239235176934
13.3600001335144 63.684230734556
13.4000000953674 63.646412660466
13.4400000572205 63.4113680738192
13.4800000190735 63.2575593062375
13.5199999809265 63.2219683607582
13.5600001811981 62.9787905826124
13.6000001430511 62.8097530966952
13.6400001049042 62.6886928021486
13.6800000667572 62.5056263512561
13.7200000286102 62.2962698809915
13.7599999904633 62.1990800204276
13.8000001907349 61.9741356745308
13.8400001525879 61.8384067664811
13.8800001144409 61.6900228698796
13.9200000762939 61.520578837406
13.960000038147 61.3712835800306
14 61.2882388707039
14.0400002002716 61.0777614707949
14.0800001621246 60.8721208536554
14.1200001239777 60.7426215103533
14.1600000858307 60.5652529110394
14.2000000476837 60.3577575218915
14.2400000095367 60.2275598712367
14.2800002098083 60.0633195360733
14.3200001716614 59.8775263221596
14.3600001335144 59.7498937537566
14.4000000953674 59.5479537440933
14.4400000572205 59.4550682828612
14.4800000190735 59.2441864623197
14.5199999809265 59.0239782900699
14.5600001811981 58.9227898398112
14.6000001430511 58.8075831533715
14.6400001049042 58.6272747536641
14.6800000667572 58.4883576822285
14.7200000286102 58.3151069097894
14.7599999904633 58.1817803648228
14.8000001907349 58.0446374164621
14.8400001525879 57.90368694676
14.8800001144409 57.7011295642951
14.9200000762939 57.6491269185244
14.960000038147 57.3918662189517
15 57.2738482121367
15.0400002002716 57.140767299783
15.0800001621246 56.9352796307194
15.1200001239777 56.8749033971799
15.1600000858307 56.6815295200704
15.2000000476837 56.4715277780888
15.2400000095367 56.3701543543903
15.2800002098083 56.1547464667055
15.3200001716614 55.9108216813875
15.3600001335144 55.8163490902298
15.4000000953674 55.701743822161
15.4400000572205 55.4955314710987
15.4800000190735 55.3445586924627
15.5199999809265 55.1796220770011
15.5600001811981 55.0720333679255
15.6000001430511 54.8234931573988
15.6400001049042 54.6227902485371
15.6800000667572 54.5294824312205
15.7200000286102 54.4552960028795
15.7599999904633 54.1995251637345
15.8000001907349 53.9896440220637
15.8400001525879 53.9449902595388
15.8800001144409 53.7395372224728
15.9200000762939 53.557330572432
15.960000038147 53.4161600219286
16 52.8933454087928
16.0400002002716 52.8287269836746
16.0800001621246 52.5958988120601
16.1200001239777 52.4423079533371
16.1600000858307 52.3535474912387
16.2000000476837 52.2108871194068
16.2400000095367 52.1207778991502
16.2800002098083 51.903704573846
16.3200001716614 51.760798637948
16.3600001335144 51.6290850177247
16.4000000953674 51.4615979602077
16.4400000572205 51.3711298913495
16.4800000190735 51.1557696345062
16.5199999809265 50.9565505869596
16.5600001811981 50.8678627867193
16.6000001430511 50.7931877469491
16.6400001049042 50.60271399456
16.6800000667572 50.400335798693
16.7200000286102 50.2500810495967
16.7599999904633 50.1143226950062
16.8000001907349 49.6605169290233
16.8400001525879 49.5427721819597
16.8800001144409 49.1769244162297
16.9200000762939 48.9661154405767
16.960000038147 48.7757367114227
17 48.5644383420695
17.0400002002716 48.3650144847143
17.0800001621246 48.2285070788039
17.1200001239777 48.1587131239721
17.1600000858307 47.8532749736496
17.2000000476837 47.5918577888442
17.2400000095367 47.440558965025
17.2800002098083 47.2987053716188
17.3200001716614 47.0434672121464
17.3600001335144 46.9339045326087
17.4000000953674 46.8686612129131
17.4400000572205 46.6488654722471
17.4800000190735 46.4288592162411
17.5199999809265 46.3696913238691
17.5600001811981 46.1644681866837
17.6000001430511 45.9707507779003
17.6400001049042 45.9121507764248
17.6800000667572 45.7037861452136
17.7200000286102 45.5738659826839
17.7599999904633 45.4974162556209
17.8000001907349 45.2500980906207
17.8400001525879 45.2068670938843
17.8800001144409 45.0974219083437
17.9200000762939 44.962168946462
17.960000038147 44.902468921613
18 44.8578663595617
18.0400002002716 44.7483312594868
18.0800001621246 44.5303715610207
18.1200001239777 44.4072990742334
18.1600000858307 44.2425721938253
18.2000000476837 44.0736742230511
18.2400000095367 43.9250716992774
18.2800002098083 43.4125830179619
18.3200001716614 43.2627285149829
18.3600001335144 43.1015450170635
18.4000000953674 42.9462275140231
18.4400000572205 42.7644965609688
18.4800000190735 42.5838978472748
18.5199999809265 42.5514266770879
18.5600001811981 42.268589014504
18.6000001430511 42.0533921655169
18.6400001049042 41.9679917284866
18.6800000667572 41.8392801879272
18.7200000286102 41.6548477970991
18.7599999904633 41.108161887707
18.8000001907349 40.7470343574023
18.8400001525879 40.4287551458619
18.8800001144409 40.1496807864478
18.9200000762939 39.8330183470425
18.960000038147 39.6826766453596
19 39.5896465481746
19.0400002002716 39.180948770947
19.0800001621246 38.9435419878337
19.1200001239777 38.6127631311138
19.1600000858307 38.5069438402097
19.2000000476837 38.2187979798509
19.2400000095367 37.9880523154225
19.2800002098083 37.8191972757377
19.3200001716614 37.4393978520202
19.3600001335144 37.3447040925294
19.4000000953674 37.2440134860899
19.4400000572205 36.8250454999161
19.4800000190735 36.6382115582674
19.5199999809265 36.3718098585359
19.5600001811981 36.2261817778926
19.6000001430511 35.9622111300469
19.6400001049042 35.8186304817591
19.6800000667572 35.502190601821
19.7200000286102 35.2647618670781
19.7599999904633 35.0647019009266
19.8000001907349 34.8113134998785
19.8400001525879 34.6043074554782
19.8800001144409 34.3549088083109
19.9200000762939 34.1001584362384
19.960000038147 33.9066141441508
20 33.6516875993439
20.0400002002716 33.4624807351775
20.0800001621246 33.2125688848901
20.1200001239777 32.9692411026317
20.1600000858307 32.7548883619074
20.2000000476837 32.5032443090495
20.2400000095367 32.3004009337401
20.2800002098083 32.0532929702767
20.3200001716614 31.8401740098016
20.3600001335144 31.6151487420575
20.4000000953674 31.3665451064917
20.4400000572205 31.1662384056654
20.4800000190735 30.9163665707802
20.5199999809265 30.7153229648211
20.5600001811981 30.4622102289341
20.6000001430511 30.2183949631981
20.6400001049042 30.0250388625391
20.6800000667572 29.7687860190053
20.7200000286102 29.5547277751211
20.7599999904633 29.298235489634
20.8000001907349 29.0485994120736
20.8400001525879 28.8450235084147
20.8800001144409 28.569980872051
20.9200000762939 28.4427271924342
20.960000038147 28.1593901164015
21 28.0129054721838
21.0400002002716 27.7393458578481
21.0800001621246 27.3870792166035
21.1200001239777 27.2283612260144
21.1600000858307 26.9626020073993
21.2000000476837 26.8179353694704
21.2400000095367 26.56275790158
21.2800002098083 26.1927447305952
21.3200001716614 26.012915401594
21.3600001335144 25.7670369710286
21.4000000953674 25.604121147735
21.4400000572205 25.3481422794357
21.4800000190735 24.985332189417
21.5199999809265 24.7180189203482
21.5600001811981 24.6019234601226
21.6000001430511 24.2369073490081
21.6400001049042 24.0194671260888
21.6800000667572 23.7517484831496
21.7200000286102 23.4756911986788
21.7599999904633 23.3234344469056
21.8000001907349 23.1603400349213
21.8400001525879 22.8276357945306
21.8800001144409 22.4754804450949
21.9200000762939 22.2051117929677
21.960000038147 21.9938474312949
22 21.721095581921
22.0400002002716 21.4512012303585
22.0800001621246 21.4938177093354
22.1200001239777 20.9611231369909
22.1600000858307 20.754933347127
22.2000000476837 20.4812174621875
22.2400000095367 20.214181801628
22.2800002098083 20.009337681513
22.3200001716614 19.7405844179157
22.3600001335144 19.5288031747223
22.4000000953674 19.2550208639459
22.4400000572205 18.9829085847778
22.4800000190735 18.8021288543441
22.5199999809265 18.4949580946304
22.5600001811981 18.2908743615117
22.6000001430511 18.0139553841254
22.6400001049042 17.7405577121663
22.6800000667572 17.5543233667195
22.7200000286102 17.2531851706717
22.7599999904633 16.996866964746
22.8000001907349 16.783870561294
22.8400001525879 16.5042519236868
22.8800001144409 16.3530786450374
22.9200000762939 16.0192735586043
22.960000038147 15.8489053919311
23 15.5407619015896
23.0400002002716 15.2617572443505
23.0800001621246 15.1525235334029
23.1200001239777 14.9577532499134
23.1600000858307 14.4972395834266
23.2000000476837 14.377653780386
23.2400000095367 14.1342848793338
23.2800002098083 13.764351130525
23.3200001716614 13.70190527052
23.3600001335144 13.3203086557933
23.4000000953674 13.0949248431934
23.4400000572205 12.9470117639921
23.4800000190735 12.5233736908904
23.5199999809265 12.4360788038102
23.5600001811981 12.1227882688472
23.6000001430511 11.8367790409102
23.6400001049042 11.5486776019152
23.6800000667572 11.4516330398765
23.7200000286102 11.2115306646938
23.7599999904633 10.9231616360375
23.8000001907349 10.5325445839007
23.8400001525879 10.3937985609482
23.8800001144409 10.0400192288544
23.9200000762939 9.85277089686679
23.960000038147 9.54638944603357
24 9.41215320500305
24.0400002002716 9.01416069670813
24.0800001621246 8.91639514231247
24.1200001239777 8.56138548014133
24.1600000858307 8.30973078305942
24.2000000476837 8.19010770008754
24.2400000095367 7.83452820983257
24.2800002098083 7.73887402655782
24.3200001716614 7.44058631563712
24.3600001335144 7.25586748671594
24.4000000953674 7.19955539676812
24.4400000572205 7.04521880163156
24.4800000190735 6.95760035088824
24.5199999809265 6.72991129099319
24.5600001811981 6.57797092521832
24.6000001430511 6.54692020296898
24.6400001049042 6.05203488708139
24.6800000667572 5.72135322686042
24.7200000286102 5.39067156663944
24.7599999904633 5.05998990641847
24.8000001907349 4.72930627517933
24.8400001525879 4.39862461495835
24.8800001144409 4.06794295473738
24.9200000762939 3.7372612945164
24.960000038147 3.40657963429543
25 3.07589797407446
25.0400002002716 2.74521434283531
25.0800001621246 2.41453268261434
25.1200001239777 2.08385102239337
25.1600000858307 1.75316936217239
25.2000000476837 1.42248770195142
25.2400000095367 1.09180604173044
25.2800002098083 0.761122410491302
25.3200001716614 0.430440750270328
25.3600001335144 0.0997590900493535
25.4000000953674 -0.230922570171621
25.4400000572205 -0.561604230392594
25.4800000190735 -0.892285890613568
25.5199999809265 -1.22296755083454
25.5600001811981 -1.55365118207368
25.6000001430511 -1.88433284229466
25.6400001049042 -2.21501450251563
25.6800000667572 -2.54569616273661
25.7200000286102 -2.87637782295758
25.7599999904633 -3.20705948317855
25.8000001907349 -3.5377431144177
25.8400001525879 -3.86842477463867
25.8800001144409 -4.19910643485965
25.9200000762939 -4.52978809508062
25.960000038147 -4.86046975530159
26 -5.19115141552257
26.0400002002716 -5.52183504676171
26.0800001621246 -5.85251670698268
26.1200001239777 -6.18319836720366
26.1600000858307 -6.34269787828816
26.2000000476837 -6.4551393305901
26.2400000095367 -6.75037590070315
26.2800002098083 -6.90046827996957
26.3200001716614 -7.19344752401412
26.3600001335144 -7.34442746254354
26.4000000953674 -7.41793076970134
26.4400000572205 -7.62233715359798
26.4800000190735 -7.84928645275249
26.5199999809265 -8.04018979420289
26.5600001811981 -8.18588997182269
26.6000001430511 -8.29410446614202
26.6400001049042 -8.65567044169256
26.6800000667572 -8.81020059913317
26.7200000286102 -9.44169753260576
26.7599999904633 -9.57321300304284
26.8000001907349 -9.66905491330908
26.8400001525879 -9.92303286529568
26.8800001144409 -10.1837382522117
26.9200000762939 -10.437642173385
26.960000038147 -10.5901330325814
27 -10.714304641655
27.0400002002716 -11.115113244512
27.0800001621246 -11.1807890889995
27.1200001239777 -11.5929128578919
27.1600000858307 -11.7531271840307
27.2000000476837 -11.8920918047334
27.2400000095367 -12.3448769610441
27.2800002098083 -12.4047798967331
27.3200001716614 -12.8432828591212
27.3600001335144 -13.0121988515566
27.4000000953674 -13.1575324731275
27.4400000572205 -13.5198671809085
27.4800000190735 -13.6593364139644
27.5199999809265 -14.1228665389553
27.5600001811981 -14.3009511535074
27.6000001430511 -14.4450129154975
27.6400001049042 -14.8259638117888
27.6800000667572 -14.9741990941238
27.7200000286102 -15.4201923588407
27.7599999904633 -15.6063091848428
27.8000001907349 -15.744839267667
27.8400001525879 -16.1226214781709
27.8800001144409 -16.26811134932
27.9200000762939 -16.7081152942865
27.960000038147 -16.9013549478186
28 -17.031722869433
28.0400002002716 -17.4857008491854
28.0800001621246 -17.5482300134117
28.1200001239777 -17.989597254742
28.1600000858307 -18.1812828324637
28.2000000476837 -18.3161401467132
28.2400000095367 -18.7027211175937
28.2800002098083 -18.8243539524337
28.3200001716614 -19.2614023603674
28.3600001335144 -19.4496491320879
28.4000000953674 -19.5795754933843
28.4400000572205 -19.9885586535365
28.4800000190735 -20.0871460088183
28.5199999809265 -20.5337730135634
28.5600001811981 -20.7258404268814
28.6000001430511 -20.8391770099352
28.6400001049042 -21.2792300030942
28.6800000667572 -21.3372227602304
28.7200000286102 -21.7875850943765
28.7599999904633 -21.966646818948
28.8000001907349 -22.0789943004893
28.8400001525879 -22.5296290901115
28.8800001144409 -22.5779580321941
28.9200000762939 -23.0307706867843
28.960000038147 -23.2094415304346
29 -23.329274779735
29.0400002002716 -23.7829151760325
29.0800001621246 -23.8514600270719
29.1200001239777 -24.29606071323
29.1600000858307 -24.4764859031711
29.2000000476837 -24.6006154368442
29.2400000095367 -25.04397066196
29.2800002098083 -25.1032559630439
29.3200001716614 -25.5465036126388
29.3600001335144 -25.7271658477785
29.4000000953674 -25.8584042771254
29.4400000572205 -26.3298960640313
29.4800000190735 -26.3803117962343
29.5199999809265 -26.8196697242493
29.5600001811981 -27.000737984521
29.6000001430511 -27.120171080629
29.6400001049042 -27.5609671461189
29.6800000667572 -27.613976538194
29.7200000286102 -28.0602445297445
29.7599999904633 -28.2402723911637
29.8000001907349 -28.350380686883
29.8400001525879 -28.7974325652121
29.8800001144409 -28.8449194284676
29.9200000762939 -29.2919462002501
29.960000038147 -29.4713486230939
30 -29.5891572474447
30.0400002002716 -30.0384627457734
30.0800001621246 -30.085822151912
30.1200001239777 -30.5276554665434
30.1600000858307 -30.7051624945307
30.2000000476837 -30.8220865840685
30.2400000095367 -31.2604995585916
30.2800002098083 -31.4663039977349
30.3200001716614 -31.50723830469
30.3600001335144 -31.8893539981473
30.4000000953674 -32.2303377878834
30.4400000572205 -32.4776080970605
30.4800000190735 -32.5241237622886
30.5199999809265 -32.9659232401609
30.5600001811981 -33.1411500857121
30.6000001430511 -33.2644369610498
30.6400001049042 -33.696888092114
30.6800000667572 -33.9012092473549
30.7200000286102 -34.179165681313
30.7599999904633 -34.3537292364963
30.8000001907349 -34.6670679562812
30.8400001525879 -34.7327897953492
30.8800001144409 -35.1076711571714
30.9200000762939 -35.4016810404628
30.960000038147 -35.5708997145193
31 -35.68662248896
31.0400002002716 -36.1101107898867
31.0800001621246 -36.3245547379538
31.1200001239777 -36.3940119977815
31.1600000858307 -36.7792782806591
31.2000000476837 -36.8941622725961
31.2400000095367 -37.2641448524082
31.2800002098083 -37.5294697985701
31.3200001716614 -37.6197988614749
31.3600001335144 -37.9410825236919
31.4000000953674 -38.3351562655953
31.4400000572205 -38.460242380108
31.4800000190735 -38.7578077735398
31.5199999809265 -39.0553079623751
31.5600001811981 -39.2124785794003
31.6000001430511 -39.339861883207
31.6400001049042 -39.6392979916636
31.6800000667572 -39.9714506718101
31.7200000286102 -40.2120387440227
31.7599999904633 -40.3800448817892
31.8000001907349 -40.7808287118187
31.8400001525879 -40.8529843019987
31.8800001144409 -41.1867113863664
31.9200000762939 -41.3927126456638
31.960000038147 -41.6663102513758
32 -41.9992190241483
32.0400002002716 -42.2862690839338
32.0800001621246 -42.2896692464965
32.1200001239777 -42.5738616783146
32.1600000858307 -42.9095866755324
32.2000000476837 -43.1674840022533
32.2400000095367 -43.2791464588245
32.2800002098083 -43.6763818520012
32.3200001716614 -44.0221587463711
32.3600001335144 -44.0924282124433
32.4000000953674 -44.364746152427
32.4400000572205 -44.5912162997502
32.4800000190735 -44.8292777579081
32.5199999809265 -45.0631080417843
32.5600001811981 -45.3207879184338
32.6000001430511 -45.5526869969035
32.6400001049042 -45.8052508141391
32.6800000667572 -46.083107164588
32.7200000286102 -46.1749218554123
32.7599999904633 -46.5542458971668
32.8000001907349 -46.6748024754779
32.8400001525879 -46.9905359527656
32.8800001144409 -47.3234813646541
32.9200000762939 -47.5254904218091
32.960000038147 -47.8058956548521
33 -48.0321085357282
33.0400002002716 -48.2845930197091
33.0800001621246 -48.5443714751727
33.1200001239777 -48.7856099574819
33.1600000858307 -49.0503223747673
33.2000000476837 -49.2860321341523
33.2400000095367 -49.5399599587145
33.2800002098083 -49.8401721027712
33.3200001716614 -49.9311127161327
33.3600001335144 -50.3093517807902
33.4000000953674 -50.4274854639716
33.4400000572205 -50.7881927215003
33.4800000190735 -51.0582099315816
33.5199999809265 -51.3752998567343
33.5600001811981 -51.5594607958938
33.6000001430511 -51.9085602364794
33.6400001049042 -52.0470561753627
33.6800000667572 -52.2931560794677
33.7200000286102 -52.5886182672257
33.7599999904633 -52.7895713949219
33.8000001907349 -53.1437683703589
33.8400001525879 -53.2811793235214
33.8800001144409 -53.5304631996423
33.9200000762939 -53.5304631996423
33.960000038147 -53.5304631996423
34 -53.5304631996423
34.0400002002716 -53.5304631996423
34.0800001621246 -53.5304631996423
34.1200001239777 -53.5304631996423
34.1600000858307 -53.5304631996423
};
\path [draw=black, fill=black] (axis cs:6,90)
--(axis cs:6.5,85)
--(axis cs:6.125,85)
--(axis cs:6.125,55)
--(axis cs:5.875,55)
--(axis cs:5.875,85)
--(axis cs:5.5,85)
--cycle;

\path [draw=black, fill=black] (axis cs:6,-30)
--(axis cs:5.5,-25)
--(axis cs:5.875,-25)
--(axis cs:5.875,0)
--(axis cs:6.125,0)
--(axis cs:6.125,-25)
--(axis cs:6.5,-25)
--cycle;

\addplot [line width=2.4000000000000004pt, color1, dashed, forget plot]
table {%
0 -29.7188931389674
0.0400002002716064 -29.7299563409853
0.0800001621246338 -29.7411997067783
0.120000123977661 -29.7533279461824
0.160000085830688 -29.7667942653549
0.200000047683716 -29.78089545525
0.240000009536743 -29.7954320841488
0.28000020980835 -29.8107226668905
0.320000171661377 -29.8266145371051
0.360000133514404 -29.8435831895172
0.400000095367432 -29.8612173987661
0.440000057220459 -29.8792910559265
0.480000019073486 -29.8984000859081
0.519999980926514 -29.9181969884245
0.56000018119812 -29.9378888004887
0.600000143051147 -29.9586271579341
0.640000104904175 -29.9809722186025
0.680000066757202 -30.0036080929521
0.720000028610229 -30.0261793658401
0.759999990463257 -30.0491229499174
0.800000190734863 -30.0727676251873
0.840000152587891 -30.0967012373611
0.880000114440918 -30.1204339951422
0.920000076293945 -30.1449705067457
0.960000038146973 -30.1705194735039
1 -30.1960332089379
1.04000020027161 -30.2226529458099
1.08000016212463 -30.2506743427523
1.12000012397766 -30.2790858547478
1.16000008583069 -30.3081599824256
1.20000004768372 -30.3374458924741
1.24000000953674 -30.3667706765184
1.28000020980835 -30.3964421972041
1.32000017166138 -30.4276625305402
1.3600001335144 -30.4589606067176
1.40000009536743 -30.4908940866163
1.44000005722046 -30.5229066418293
1.48000001907349 -30.5548583033759
1.51999998092651 -30.5874071023747
1.56000018119812 -30.6214014755213
1.60000014305115 -30.6558450356474
1.64000010490417 -30.6900410166422
1.6800000667572 -30.7241311668848
1.72000002861023 -30.7587892330752
1.75999999046326 -30.7945972090148
1.80000019073486 -30.8305036649767
1.84000015258789 -30.8662867778163
1.88000011444092 -30.9023954351012
1.92000007629395 -30.939712919949
1.96000003814697 -30.9773674603892
2 -31.0151675563924
2.04000020027161 -31.0531163806854
2.08000016212463 -31.0913630290855
2.12000012397766 -31.1303082083296
2.16000008583069 -31.1690068189263
2.20000004768372 -31.20774495773
2.24000000953674 -31.2474772824826
2.28000020980835 -31.2881010864492
2.32000017166138 -31.328181017626
2.3600001335144 -31.369283837434
2.40000009536743 -31.4106404982149
2.44000005722046 -31.4517258743984
2.48000001907349 -31.4934874444162
2.51999998092651 -31.5357277066921
2.56000018119812 -31.5773942107844
2.60000014305115 -31.6196359297934
2.64000010490417 -31.6627700405243
2.6800000667572 -31.7053379542075
2.72000002861023 -31.748335028955
2.75999999046326 -31.7917505977884
2.80000019073486 -31.8348820314408
2.84000015258789 -31.8789320151247
2.88000011444092 -31.923359919059
2.92000007629395 -31.9675297909875
2.96000003814697 -32.0123794028405
3 -32.0574352176028
3.04000020027161 -32.1024283934086
3.08000016212463 -32.1476527069657
3.12000012397766 -32.1924483062808
3.16000008583069 -32.2381474228078
3.20000004768372 -32.2837373659379
3.24000000953674 -32.3293289259164
3.28000020980835 -32.3753128285893
3.32000017166138 -32.4212036462442
3.3600001335144 -32.4592673242483
3.40000009536743 -32.5050920371814
3.44000005722046 -32.551517298032
3.48000001907349 -32.5979449910911
3.51999998092651 -32.6436784814627
3.56000018119812 -32.6893456115955
3.60000014305115 -32.7342927432693
3.64000010490417 -32.7794284517455
3.6800000667572 -32.8247678869039
3.72000002861023 -32.8708062970172
3.75999999046326 -32.9173511741383
3.80000019073486 -32.9628429776992
3.84000015258789 -33.0217852501104
3.88000011444092 -33.0730370543073
3.92000007629395 -33.1215302165186
3.96000003814697 -33.1702130788799
4 -33.2187705760904
4.04000020027161 -33.2672521660474
4.08000016212463 -33.3158798314681
4.12000012397766 -33.3640687446004
4.16000008583069 -33.413080970087
4.20000004768372 -33.4623609274466
4.24000000953674 -33.5111815383242
4.28000020980835 -33.5601704987456
4.32000017166138 -33.6087291434336
4.3600001335144 -33.6582035889056
4.40000009536743 -33.7086284591084
4.44000005722046 -33.7587357636379
4.48000001907349 -33.8092909008177
4.51999998092651 -33.8602163767029
4.56000018119812 -33.9104551436356
4.60000014305115 -33.9602149857272
4.64000010490417 -34.0102712813008
4.6800000667572 -34.0595618402416
4.72000002861023 -34.1094765000282
4.75999999046326 -34.1589279279602
4.80000019073486 -34.2087467840773
4.84000015258789 -34.2588738511894
4.88000011444092 -34.3096253149394
4.92000007629395 -34.3595248525727
4.96000003814697 -34.4091738029831
5 -34.4596382304195
5.04000020027161 -34.510051647409
5.08000016212463 -34.5596717998596
5.12000012397766 -34.6100142418731
5.16000008583069 -34.6598356758821
5.20000004768372 -34.7095505231636
5.24000000953674 -34.7594245983285
5.28000020980835 -34.8091479736693
5.32000017166138 -34.8597812923927
5.3600001335144 -34.9108997405343
5.40000009536743 -34.9619114460077
5.44000005722046 -35.012712703222
5.48000001907349 -35.0635139202775
5.51999998092651 -35.1140386895103
5.56000018119812 -35.1638787886297
5.60000014305115 -35.2136727264368
5.64000010490417 -35.2638796462667
5.6800000667572 -35.3131849811394
5.72000002861023 -35.3624117420911
5.75999999046326 -35.41246656023
5.80000019073486 -35.4624775136113
5.84000015258789 -35.5131054113043
5.88000011444092 -35.564313966377
5.92000007629395 -35.6155791128782
5.96000003814697 -35.6663992122161
6 -35.7170631068987
6.04000020027161 -35.7678505871746
6.08000016212463 -35.818263324465
6.12000012397766 -35.8701123070442
6.16000008583069 -35.9196865313809
6.20000004768372 -35.969084559418
6.24000000953674 -36.018369180037
6.28000020980835 -36.0677739171766
6.32000017166138 -36.1173218876912
6.3600001335144 -36.1673519508237
6.40000009536743 -36.2172556112757
6.44000005722046 -36.2666630826898
6.48000001907349 -36.3157289909648
6.51999998092651 -36.3647926215127
6.56000018119812 -36.4139701526667
6.60000014305115 -36.4632179974224
6.64000010490417 -36.5121893692333
6.6800000667572 -36.5605265477663
6.72000002861023 -36.6082225151977
6.75999999046326 -36.6552661830385
6.80000019073486 -36.702604210177
6.84000015258789 -36.7496355455635
6.88000011444092 -36.7967507853433
6.92000007629395 -36.8442942416474
6.96000003814697 -36.8917431624572
7 -36.9381884650527
7.04000020027161 -36.9840257807805
7.08000016212463 -37.0295941149252
7.12000012397766 -37.0742821491065
7.16000008583069 -37.1177326955777
7.20000004768372 -37.160925224589
7.24000000953674 -37.2033953462056
7.28000020980835 -37.244819201771
7.32000017166138 -37.2857257201783
7.3600001335144 -37.325516581038
7.40000009536743 -37.3654158583043
7.44000005722046 -37.4052747150782
7.48000001907349 -37.4438445141569
7.51999998092651 -37.4813501532385
7.56000018119812 -37.5179983706579
7.60000014305115 -37.5532756891371
7.64000010490417 -37.5868489676343
7.6800000667572 -37.6191734689775
7.72000002861023 -37.6508982455054
7.75999999046326 -37.6810749318967
7.80000019073486 -37.7096938476282
7.84000015258789 -37.7373698531326
7.88000011444092 -37.7642174209035
7.92000007629395 -37.7900940968963
7.96000003814697 -37.8147540802152
8 -37.8383397184766
8.04000020027161 -37.8603338082544
8.08000016212463 -37.8802792019204
8.12000012397766 -37.8985402810716
8.16000008583069 -37.9158260004466
8.20000004768372 -37.9315331652459
8.24000000953674 -37.9442846229136
8.28000020980835 -37.9553907017183
8.32000017166138 -37.9648760816569
8.3600001335144 -37.9705257162266
8.40000009536743 -37.9716108073267
8.44000005722046 -37.9690704723479
8.48000001907349 -37.9617984915756
8.51999998092651 -37.9491842414429
8.56000018119812 -37.9334883182979
8.60000014305115 -37.9146903349072
8.64000010490417 -37.8917562236944
8.6800000667572 -37.8650619844981
8.72000002861023 -37.8351236944606
8.75999999046326 -37.8002980688337
8.80000019073486 -37.761130469867
8.84000015258789 -37.7167650033463
8.88000011444092 -37.6671386003961
8.92000007629395 -37.6123557220074
8.96000003814697 -37.5505932142281
9 -37.4838478676474
9.04000020027161 -37.4117540899989
9.08000016212463 -37.3345321074775
9.12000012397766 -37.2527519864792
9.16000008583069 -37.1664856947391
9.20000004768372 -37.0749866664422
9.24000000953674 -36.9792316795759
9.28000020980835 -36.8782455112196
9.32000017166138 -36.7719483427981
9.3600001335144 -36.660439164239
9.40000009536743 -36.5451999201506
9.44000005722046 -36.4198624717382
9.48000001907349 -36.2916153775127
9.51999998092651 -36.1625755319952
9.56000018119812 -36.0287222928359
9.60000014305115 -35.890168287735
9.64000010490417 -35.7468003044487
9.6800000667572 -35.5988831835993
9.72000002861023 -35.447273049114
9.75999999046326 -35.2905741975939
9.80000019073486 -35.1297702756238
9.84000015258789 -34.9648442664693
9.88000011444092 -34.7946991999288
9.92000007629395 -34.6193052207576
9.96000003814697 -34.439131312072
10 -34.2541185019251
10.0400002002716 -34.0640385965227
10.0800001621246 -33.8698366940526
10.1200001239777 -33.6553297692756
10.1600000858307 -33.4474959585272
10.2000000476837 -33.237471066401
10.2400000095367 -33.0225184184428
10.2800002098083 -32.8025488747119
10.3200001716614 -32.5779124666641
10.3600001335144 -32.3481764405966
10.4000000953674 -32.1142898460657
10.4400000572205 -31.8743998250167
10.4800000190735 -31.6292110225473
10.5199999809265 -31.3785994291177
10.5600001811981 -31.1231872712592
10.6000001430511 -30.8792270835557
10.6400001049042 -30.6331333950744
10.6800000667572 -30.3649801025207
10.7200000286102 -30.0917225137079
10.7599999904633 -29.8129063037376
10.8000001907349 -29.5297489690965
10.8400001525879 -29.2411685225155
10.8800001144409 -28.9470928298366
10.9200000762939 -28.6478473399682
10.960000038147 -28.3429305891538
11 -28.0329702322223
11.0400002002716 -27.7171549238476
11.0800001621246 -27.395326604608
11.1200001239777 -27.0690414917607
11.1600000858307 -26.7383646979205
11.2000000476837 -26.4021243944844
11.2400000095367 -26.0605019317856
11.2800002098083 -25.7133163891195
11.3200001716614 -25.360445619182
11.3600001335144 -25.002200562557
11.4000000953674 -24.638283264491
11.4400000572205 -24.2688907406558
11.4800000190735 -23.8946273118927
11.5199999809265 -23.5147698131029
11.5600001811981 -23.1288586413156
11.6000001430511 -22.7364755500646
11.6400001049042 -22.3380911250798
11.6800000667572 -21.9344135477036
11.7200000286102 -21.5254500353819
11.7599999904633 -21.1106676904657
11.8000001907349 -20.689129993824
11.8400001525879 -20.29706282257
11.8800001144409 -19.8640017509812
11.9200000762939 -19.4243193013899
11.960000038147 -18.9791681990435
12 -18.5287150879937
12.0400002002716 -18.0725029798952
12.0800001621246 -17.6104655563942
12.1200001239777 -17.1429110009244
12.1600000858307 -16.6699749198931
12.2000000476837 -16.1912963054992
12.2400000095367 -15.7054675027534
12.2800002098083 -15.2131750129306
12.3200001716614 -14.7138768417118
12.3600001335144 -14.2095792255251
12.4000000953674 -13.6978115715972
12.4400000572205 -13.179520582162
12.4800000190735 -12.6543907475131
12.5199999809265 -12.1225540712096
12.5600001811981 -11.584035230646
12.6000001430511 -11.0397221572369
12.6400001049042 -10.4895868694578
12.6800000667572 -9.933095934613
12.7200000286102 -9.37020839303646
12.7599999904633 -8.80109403413127
12.8000001907349 -8.22483974919149
12.8400001525879 -7.64114440126028
12.8800001144409 -7.05104279400427
12.9200000762939 -6.45465243026399
12.960000038147 -5.85060706217353
13 -5.23925471207057
13.0400002002716 -4.62149392747352
13.0800001621246 -3.99563893770535
13.1200001239777 -3.36218795751825
13.1600000858307 -2.72167172545472
13.2000000476837 -2.02333188773975
13.2400000095367 -1.47429107590696
13.2800002098083 -0.925342423273158
13.3200001716614 -0.376433323753442
13.3600001335144 0.171123961537185
13.4000000953674 0.718733085198463
13.4400000572205 1.26526814659516
13.4800000190735 1.81135029260447
13.5199999809265 2.35725936841064
13.5600001811981 2.90216660249036
13.6000001430511 3.44770807447346
13.6400001049042 3.99348754269006
13.6800000667572 4.53927865494391
13.7200000286102 5.08476626977801
13.7599999904633 5.63016415363298
13.8000001907349 6.1757693812579
13.8400001525879 6.72150995431999
13.8800001144409 7.26667416943683
13.9200000762939 7.81166418847609
13.960000038147 8.35744923760418
14 8.9025868554062
14.0400002002716 9.44769869378443
14.0800001621246 9.99285118839083
14.1200001239777 10.5377816516648
14.1600000858307 11.0836967334246
14.2000000476837 11.6285997311623
14.2400000095367 12.1738728108323
14.2800002098083 12.7198306486966
14.3200001716614 13.2654431479518
14.3600001335144 13.8105008968487
14.4000000953674 14.3552822839232
14.4400000572205 14.9005741769356
14.4800000190735 15.4461069554744
14.5199999809265 15.9921543134437
14.5600001811981 16.5379013998782
14.6000001430511 17.0833472963177
14.6400001049042 17.6290455360846
14.6800000667572 18.1757509891891
14.7200000286102 18.7226558340471
14.7599999904633 19.2692106163673
14.8000001907349 19.8160475314199
14.8400001525879 20.3629392388697
14.8800001144409 20.9102117508727
14.9200000762939 21.4578481938898
14.960000038147 22.0053472193294
15 22.5528425907148
15.2000000476837 25.2899095826299
15.2400000095367 25.8377236079452
15.2800002098083 26.3855831254082
15.3200001716614 26.9330683431346
15.3600001335144 27.4800058034118
15.4000000953674 28.0274731093809
15.4400000572205 28.5748163464334
15.4800000190735 29.1219529713711
15.5199999809265 29.6684568284541
15.5600001811981 30.2158317122956
15.6000001430511 30.7633607296686
15.6400001049042 31.3102488882495
15.6800000667572 31.8569314564365
15.7200000286102 32.4038060012691
15.7599999904633 32.951076337087
15.8000001907349 33.4984552516116
15.8400001525879 34.0458504458165
15.8800001144409 34.5934387558902
15.9200000762939 35.1403920420385
15.960000038147 35.6871319056818
16 36.2348834702941
16.0400002002716 36.7833425578203
16.0800001621246 37.3312715881442
16.1200001239777 37.8806158064033
16.1600000858307 38.4297826294229
16.2000000476837 38.9794293522664
16.2400000095367 39.5301103399228
16.2800002098083 40.082624553994
16.3200001716614 40.6350632341094
16.3600001335144 41.1870866521758
16.4000000953674 41.7397438571715
16.4400000572205 42.2921193874646
16.4800000190735 42.8451485132065
16.5199999809265 43.39727776784
16.5600001811981 43.9503692262001
16.6000001430511 44.502402537165
16.6400001049042 45.056495029347
16.6800000667572 45.6096962532241
16.7200000286102 46.1620915321461
16.7599999904633 46.7140196689211
16.8000001907349 47.2668847594364
16.8400001525879 47.8197911967044
16.8800001144409 48.3717706872225
16.9200000762939 48.9230596652657
16.960000038147 49.4758804491005
17 50.0271016965841
17.0400002002716 50.5774130986268
17.0800001621246 51.1242836598685
17.1200001239777 51.672545309317
17.1600000858307 52.2201959140254
17.2000000476837 52.7670063488809
17.2400000095367 53.3131671285719
17.2800002098083 53.8577651885671
17.3200001716614 54.4016310368785
17.3600001335144 54.9443054556316
17.4000000953674 55.4838825124497
17.4400000572205 56.022849208036
17.4800000190735 56.5609616948444
17.5199999809265 57.0969222566835
17.5600001811981 57.6309200385135
17.6000001430511 58.1645269295756
17.6400001049042 58.6974319173026
17.6800000667572 59.2278827586706
17.7200000286102 59.7572335684629
};
\addplot [line width=2.4000000000000004pt, color3, dashed, forget plot]
table {%
0.800000190734863 101.465466608788
0.840000152587891 101.508171380934
0.880000114440918 101.550091975127
0.920000076293945 101.590320374476
0.960000038146973 101.628320310422
1 101.665256728601
1.04000020027161 101.700516717164
1.08000016212463 101.733437676589
1.12000012397766 101.764888781866
1.16000008583069 101.79428611963
1.20000004768372 101.822255750803
1.24000000953674 101.849203872278
1.28000020980835 101.874684713105
1.32000017166138 101.897256684029
1.3600001335144 101.918547127497
1.40000009536743 101.937959382701
1.44000005722046 101.955878449043
1.48000001907349 101.972811025616
1.51999998092651 101.987941380819
1.56000018119812 102.000239689067
1.60000014305115 102.010907037484
1.64000010490417 102.020632319241
1.6800000667572 102.02944605594
1.72000002861023 102.036526461145
1.75999999046326 102.041009300919
1.80000019073486 102.0443810469
1.84000015258789 102.04666503889
1.88000011444092 102.047492395032
1.92000007629395 102.046190770439
1.96000003814697 102.043525903541
2 102.039479665938
2.04000020027161 102.034052242092
2.08000016212463 102.027257814753
2.12000012397766 102.018745044625
2.16000008583069 102.009337300196
2.20000004768372 101.998655608808
2.24000000953674 101.985919152105
2.28000020980835 101.971142302881
2.32000017166138 101.955478324773
2.3600001335144 101.936062535231
2.40000009536743 101.909116730549
2.44000005722046 101.891533257563
2.48000001907349 101.871982693068
2.51999998092651 101.850689447015
2.56000018119812 101.827798227201
2.60000014305115 101.801651186686
2.64000010490417 101.773416958882
2.6800000667572 101.744575313137
2.72000002861023 101.714197669936
2.75999999046326 101.682207923448
2.80000019073486 101.649584403147
2.84000015258789 101.614986458969
2.88000011444092 101.583862175153
2.92000007629395 101.547840617267
2.96000003814697 101.507149993773
3 101.464923174395
3.04000020027161 101.421887481755
3.08000016212463 101.378541212529
3.12000012397766 101.334831993266
3.16000008583069 101.288921905912
3.20000004768372 101.241839616412
3.24000000953674 101.193494494269
3.28000020980835 101.143474145806
3.32000017166138 101.097779504246
3.3600001335144 101.057485521361
3.40000009536743 101.002884391101
3.44000005722046 100.946411491978
3.48000001907349 100.888623304287
3.51999998092651 100.830221972261
3.56000018119812 100.770595516136
3.60000014305115 100.710319752151
3.64000010490417 100.648519359272
3.6800000667572 100.585147828171
3.72000002861023 100.51971097126
3.75999999046326 100.452440118411
3.80000019073486 100.384786781866
3.84000015258789 100.320395143145
3.88000011444092 100.247478852722
3.92000007629395 100.17519098234
3.96000003814697 100.101481527236
4 100.026666396513
4.04000020027161 99.950694439651
4.08000016212463 99.8733430170946
4.12000012397766 99.7951874182703
4.16000008583069 99.7149741684241
4.20000004768372 99.6332550944455
4.24000000953674 99.5507631001514
4.28000020980835 99.4668732589905
4.32000017166138 99.3821839403125
4.3600001335144 99.2953405750387
4.40000009536743 99.2063115574046
4.44000005722046 99.1163581801867
4.48000001907349 99.0247170600869
4.51999998092651 98.9314682965128
4.56000018119812 98.8376659727511
4.60000014305115 98.7431052516281
4.64000010490417 98.647009366826
4.6800000667572 98.5504308456814
4.72000002861023 98.4519766159655
4.75999999046326 98.3527347877912
4.80000019073486 98.2518569398588
4.84000015258789 98.1494159738837
4.88000011444092 98.0450890777364
4.92000007629395 97.9403433066317
4.96000003814697 97.8345927441907
5 97.717402050284
5.04000020027161 97.6082556560407
5.08000016212463 97.4986318688928
5.12000012397766 97.3870094696572
5.16000008583069 97.2746308166943
5.20000004768372 97.1610885992189
5.24000000953674 97.0461313492432
5.28000020980835 96.9301000968196
5.32000017166138 96.8118969700399
5.3600001335144 96.6919414709391
5.40000009536743 96.5708736394519
5.44000005722046 96.4487900451215
5.48000001907349 96.3254685197967
5.51999998092651 96.2012153358897
5.56000018119812 96.0762507496213
5.60000014305115 95.9500476973486
5.64000010490417 95.8221755093968
5.6800000667572 95.6939564506665
5.72000002861023 95.5645201887134
5.75999999046326 95.4329666181412
5.80000019073486 95.3001604243693
5.84000015258789 95.1654702918253
5.88000011444092 95.0289222761752
5.92000007629395 94.8910501482853
5.96000003814697 94.7523857089783
6 94.6127048628777
6.04000020027161 94.4716812654369
6.08000016212463 94.3299192366003
6.12000012397766 94.1880417136584
6.16000008583069 94.0445106193647
6.20000004768372 93.8999096838418
6.24000000953674 93.7541062121897
6.28000020980835 93.6069221272582
6.32000017166138 93.4582837621784
6.3600001335144 93.3078527423384
6.40000009536743 93.1561946346525
6.44000005722046 93.0037151282287
6.48000001907349 92.8502430071796
6.51999998092651 92.6954435381169
6.56000018119812 92.5392398938169
6.60000014305115 92.3816342915312
6.64000010490417 92.2229757259688
6.6800000667572 92.0635606458248
6.72000002861023 91.9034085441897
6.75999999046326 91.7425144331308
6.80000019073486 91.5799079518556
6.84000015258789 91.4161285351568
6.88000011444092 91.2508014745571
6.92000007629395 91.0834981580226
6.96000003814697 90.9146793316134
7 90.7452076992629
7.04000020027161 90.5746560652683
7.08000016212463 90.4026669216031
7.12000012397766 90.2298164431686
7.16000008583069 90.0564471836372
7.20000004768372 89.8815189298149
7.24000000953674 89.705462727994
7.28000020980835 89.5285091828133
7.32000017166138 89.3501681226308
7.3600001335144 89.1709799281471
7.40000009536743 88.989699075209
7.44000005722046 88.8064419076188
7.48000001907349 88.622470498655
7.51999998092651 88.4374692915834
7.56000018119812 88.251221015129
7.60000014305115 88.0641743190748
7.64000010490417 87.8765914336536
7.6800000667572 87.688003159597
7.72000002861023 87.4977572326571
7.75999999046326 87.3067446896694
7.80000019073486 87.1149455768267
7.84000015258789 86.9216942904709
7.88000011444092 86.7268355173789
7.92000007629395 86.5304957870994
7.96000003814697 86.3327945979764
8 86.1336612265625
8.04000020027161 85.9335711445442
8.08000016212463 85.7329440909613
8.12000012397766 85.5314228909331
8.16000008583069 85.3282581009228
8.20000004768372 85.1239921698082
8.24000000953674 84.9198888653146
8.28000020980835 84.7146857756439
8.32000017166138 84.5082790681873
8.3600001335144 84.3028765433739
8.40000009536743 84.0991606789504
8.44000005722046 83.8961530200729
8.48000001907349 83.6950185302418
8.51999998092651 83.4962802985132
8.56000018119812 83.2975518952508
8.60000014305115 83.098905324023
8.64000010490417 82.9012753914575
8.6800000667572 82.7043213768006
8.72000002861023 82.5075169093375
8.75999999046326 82.3125709171429
8.80000019073486 82.1188022875684
8.84000015258789 81.9271504063619
8.88000011444092 81.7374803191095
8.92000007629395 81.5464787071878
8.96000003814697 81.3409456003816
9 81.1612145760248
9.04000020027161 80.9834388807422
9.08000016212463 80.8071908803188
9.12000012397766 80.6318435882067
9.16000008583069 80.4573231624023
9.20000004768372 80.2844263820603
9.24000000953674 80.111818829007
9.28000020980835 79.9405930539112
9.32000017166138 79.7707260404285
9.3600001335144 79.6022106134816
9.40000009536743 79.4335290944813
9.44000005722046 79.2762151472035
9.48000001907349 79.1112427750271
9.51999998092651 78.943618714578
9.56000018119812 78.7768507143628
9.60000014305115 78.6107847821658
9.64000010490417 78.4455010029867
9.6800000667572 78.2807005433926
9.72000002861023 78.115506524454
9.75999999046326 77.9512765798501
9.80000019073486 77.7869901798834
9.84000015258789 77.6226243553022
9.88000011444092 77.4592511865325
9.92000007629395 77.2968628358098
9.96000003814697 77.1352009782151
10 76.9744147855574
10.0400002002716 76.8144860810798
10.0800001621246 76.6543216333883
10.1200001239777 76.5077079895385
10.1600000858307 76.3488649113857
10.2000000476837 76.18947483516
10.2400000095367 76.0303798628399
10.2800002098083 75.8716093818074
10.3200001716614 75.7127949072004
10.3600001335144 75.5542893313577
10.4000000953674 75.3951832455615
10.4400000572205 75.2372158086712
10.4800000190735 75.0795757644566
10.5199999809265 74.9223155463665
10.5600001811981 74.7648132029937
10.6000001430511 74.6077911152214
10.6400001049042 74.4477513542248
10.6800000667572 74.2915626057954
10.7200000286102 74.1355438423443
10.7599999904633 73.9801253219129
10.8000001907349 73.823994981082
10.8400001525879 73.6682645151115
10.8800001144409 73.5129208039836
10.9200000762939 73.3576417647607
10.960000038147 73.202929888692
11 73.0481572151907
11.0400002002716 72.8940880218246
11.0800001621246 72.7408464375047
11.1200001239777 72.5869587887586
11.1600000858307 72.4321769877472
11.2000000476837 72.2776578458731
11.2400000095367 72.1232083661357
11.2800002098083 71.9690830841409
11.3200001716614 71.8153201639372
11.3600001335144 71.6615515864351
11.4000000953674 71.5081057483543
11.4400000572205 71.354702894441
11.4800000190735 71.2007368542867
11.5199999809265 71.0468006114144
11.5600001811981 70.8932914813136
11.6000001430511 70.7404882260844
11.6400001049042 70.5879813362569
11.6800000667572 70.4351017591217
11.7200000286102 70.2818067638911
11.7599999904633 70.1286980746389
11.8000001907349 69.9761166552294
11.8400001525879 69.8231873655727
11.8800001144409 69.6702712183523
11.9200000762939 69.5179826019143
11.960000038147 69.365064028374
12 69.2112867654504
12.0400002002716 69.0572446822591
12.0800001621246 68.9029657257043
12.1200001239777 68.7481064561243
12.1600000858307 68.5924261540286
12.2000000476837 68.4362901577738
12.2400000095367 68.2810399433327
12.2800002098083 68.1258743941194
12.3200001716614 67.9714420653347
12.3600001335144 67.8155036959318
12.4000000953674 67.6605814965305
12.4400000572205 67.5058002994636
12.4800000190735 67.3515307252465
12.5199999809265 67.1975048776493
12.5600001811981 67.0435096645865
12.6000001430511 66.8888138334459
12.6400001049042 66.7333994231666
12.6800000667572 66.5777597124291
12.7200000286102 66.4218781809627
12.7599999904633 66.2654992446979
12.8000001907349 66.1094914094813
12.8400001525879 65.9539124669883
12.8800001144409 65.7979407905959
12.9200000762939 65.6412822948327
12.960000038147 65.4851791260443
13 65.3295228182616
13.0400002002716 65.1733966381897
13.0800001621246 65.0181669106104
13.1200001239777 64.8636319622695
13.1600000858307 64.7090828087404
13.2000000476837 64.5541453454782
13.2400000095367 64.3989154496912
13.2800002098083 64.2436707280104
13.3200001716614 64.0883749838987
13.3600001335144 63.9322583409428
13.4000000953674 63.7761475246937
13.4400000572205 63.6192935353966
13.4800000190735 63.4620249582441
13.5199999809265 63.3046554887386
13.5600001811981 63.1466607588448
13.6000001430511 62.9892958745526
13.6400001049042 62.8322530036825
13.6800000667572 62.6752072157405
13.7200000286102 62.5177905331673
13.7599999904633 62.3601198963861
13.8000001907349 62.2027852075524
13.8400001525879 62.0455947859602
13.8800001144409 61.8877562595208
13.9200000762939 61.729884853707
13.960000038147 61.5727759976336
14 61.4150052804567
14.0400002002716 61.2571397285544
14.0800001621246 61.0992793260796
14.1200001239777 60.9411479851152
14.1600000858307 60.7840246211574
14.2000000476837 60.6262087558196
14.2400000095367 60.4685730460752
14.2800002098083 60.3114901395628
14.3200001716614 60.1540972748559
14.3600001335144 59.9961041533824
14.4000000953674 59.8378721859086
14.4400000572205 59.6800288809867
14.4800000190735 59.5225227092641
14.5199999809265 59.3654275396787
14.5600001811981 59.2080042050245
14.6000001430511 59.0502348920216
14.6400001049042 58.8926696121296
14.6800000667572 58.7356518129402
14.7200000286102 58.5785679158959
14.7599999904633 58.4211323402833
14.8000001907349 58.2636333404352
14.8400001525879 58.1061449368038
14.8800001144409 57.9488288674554
14.9200000762939 57.7918089239485
14.960000038147 57.6347275010522
15 57.4773470840657
15.0400002002716 57.3195471916877
15.0800001621246 57.1621018832065
15.1200001239777 57.004722407619
15.1600000858307 56.8470292178756
15.2000000476837 56.6897481587284
15.2400000095367 56.5329115037562
15.2800002098083 56.3763729267182
15.3200001716614 56.2194413248137
15.3600001335144 56.0621310635828
15.4000000953674 55.9053281749818
15.4400000572205 55.7484374195664
15.4800000190735 55.5912148346689
15.5199999809265 55.4335795052116
15.5600001811981 55.2767308966675
15.6000001430511 55.1202215562291
15.6400001049042 54.9631376879994
15.6800000667572 54.8058658470374
15.7200000286102 54.6486175817357
15.7599999904633 54.4916688866989
15.8000001907349 54.3348273846328
15.8400001525879 54.1781925174416
15.8800001144409 54.0215532697723
15.9200000762939 53.8645283368053
15.960000038147 53.7073452924015
16 53.5506587706211
16.0400002002716 53.3945779251125
16.0800001621246 53.2377651999576
16.1200001239777 53.0821161765203
16.1600000858307 52.9261648536891
16.2000000476837 52.7704641841648
16.2400000095367 52.6154037578739
16.2800002098083 52.4605870670757
16.3200001716614 52.3062258060836
16.3600001335144 52.1520533777766
16.4000000953674 51.9977982253957
16.4400000572205 51.8436445541677
16.4800000190735 51.6896252070198
16.5199999809265 51.5355181585452
16.5600001811981 51.3808141778506
16.6000001430511 51.2256426458942
16.6400001049042 51.0718430881181
16.6800000667572 50.9171879543837
16.7200000286102 50.7620637554124
16.7599999904633 50.6074441353987
16.8000001907349 50.4542410423314
16.8400001525879 50.3008424082106
16.8800001144409 50.1463480567084
16.9200000762939 49.991889410051
16.960000038147 49.8364584874565
17 49.680755345217
17.0400002002716 49.5241363723502
17.0800001621246 49.3658149151723
17.1200001239777 49.2075458681001
17.1600000858307 49.0502429314958
17.2000000476837 48.8908844969691
17.2400000095367 48.7297844372532
17.2800002098083 48.5680727321286
17.3200001716614 48.4055478419073
17.3600001335144 48.2417579380372
17.4000000953674 48.0760834446151
17.4400000572205 47.9099712774263
17.4800000190735 47.7445262442941
17.5199999809265 47.5758185565421
17.5600001811981 47.405764809612
17.6000001430511 47.2338839207132
17.6400001049042 47.0615859915252
17.6800000667572 46.8858760394434
17.7200000286102 46.7114654614059
17.7599999904633 46.5357732718061
17.8000001907349 46.3577969681155
17.8400001525879 46.1770224997191
17.8800001144409 45.9939779319207
17.9200000762939 45.8088067516682
17.960000038147 45.6247661465719
18 45.4389638195626
18.0400002002716 45.2532075783406
18.0800001621246 45.0652248616047
18.1200001239777 44.8788711213539
18.1600000858307 44.6866347105773
18.2000000476837 44.4922387757162
18.2400000095367 44.2971713221496
18.2800002098083 44.0997178207817
18.3200001716614 43.9005299901639
18.3600001335144 43.7001070461897
18.4000000953674 43.5000557146376
18.4400000572205 43.3001815317167
18.4800000190735 43.0989442319801
18.5199999809265 42.8980331793653
18.5600001811981 42.6947782746714
18.6000001430511 42.4921151442457
18.6400001049042 42.2890667176372
18.6800000667572 42.0882581862528
18.7200000286102 41.885270231771
18.7599999904633 41.6800936391941
18.8000001907349 41.4753134734698
18.8400001525879 41.2683281656195
18.8800001144409 41.0596030721356
18.9200000762939 40.8491342421977
18.960000038147 40.63333577016
19 40.4169038567798
19.0400002002716 40.2018712976426
19.0800001621246 39.9839735906491
19.1200001239777 39.7655748841976
19.1600000858307 39.5450075121882
19.2000000476837 39.3265391825446
19.2400000095367 39.1038292206461
19.2800002098083 38.8819368553974
19.3200001716614 38.6587068100507
19.3600001335144 38.434416495177
19.4000000953674 38.2102950502687
19.4400000572205 37.9860085354236
19.4800000190735 37.7621601464763
19.5199999809265 37.5372004794173
19.5600001811981 37.3115851677587
19.6000001430511 37.0854205358464
19.6400001049042 36.8565161641218
19.6800000667572 36.628258094558
19.7200000286102 36.3991636975544
19.7599999904633 36.1728613590257
19.8000001907349 35.9450418932745
19.8400001525879 35.7174673439664
19.8800001144409 35.4891479261291
19.9200000762939 35.2607498453736
19.960000038147 35.0302667526709
20 34.797891505878
20.0400002002716 34.5647762542711
20.0800001621246 34.3326455584301
20.1200001239777 34.1008253248913
20.1600000858307 33.8696069970241
20.2000000476837 33.6387757384778
20.2400000095367 33.4063957793103
20.2800002098083 33.1723540265775
20.3200001716614 32.9385698180619
20.3600001335144 32.7046653930613
20.4000000953674 32.4691006619885
20.4400000572205 32.2351494695508
20.4800000190735 32.0012811365015
20.5199999809265 31.7682605112782
20.5600001811981 31.5340312545496
20.6000001430511 31.2998609444236
20.6400001049042 31.0653043568497
20.6800000667572 30.8293653942445
20.7200000286102 30.5936716234625
20.7599999904633 30.3586078081221
20.8000001907349 30.1227285543418
20.8400001525879 29.8861403761844
20.8800001144409 29.6496261817034
20.9200000762939 29.413439977471
20.960000038147 29.1762909206114
21 28.9393289275502
21.0400002002716 28.7022825813321
21.0800001621246 28.4643979863655
21.1200001239777 28.2261197332291
21.1600000858307 27.987112663702
21.2000000476837 27.7469072457258
21.2400000095367 27.5063743349271
21.2800002098083 27.2657489787819
21.3200001716614 27.0242966352622
21.3600001335144 26.7832853552066
21.4000000953674 26.5416657448216
21.4400000572205 26.300141845967
21.4800000190735 26.0578793760505
21.5199999809265 25.817045363988
21.5600001811981 25.5762341878491
21.6000001430511 25.3345923240945
21.6400001049042 25.0926570579492
21.6800000667572 24.8501573432208
21.7200000286102 24.6069673409354
21.7599999904633 24.3636732593147
21.8000001907349 24.1193555215649
21.8400001525879 23.8746458946486
21.8800001144409 23.6298754814919
21.9200000762939 23.3849374971739
21.960000038147 23.1399097407925
22 22.8949193445956
22.0400002002716 22.6491203505718
22.0800001621246 22.4038365642698
22.1200001239777 22.1587283130219
22.1600000858307 21.9133153150657
22.2000000476837 21.6683673328099
22.2400000095367 21.4222541488439
22.2800002098083 21.1757613280753
22.3200001716614 20.9286143751897
22.3600001335144 20.6809377793
22.4000000953674 20.4331468412434
22.4400000572205 20.1854313947352
22.4800000190735 19.9381302143106
22.5199999809265 19.6897192409059
22.5600001811981 19.441552433052
22.6000001430511 19.1935446112685
22.6400001049042 18.9452564227879
22.6800000667572 18.6961902983461
22.7200000286102 18.4474586984681
22.7599999904633 18.1986775220064
22.8000001907349 17.949887873525
22.8400001525879 17.70098297255
22.8800001144409 17.4517825332636
22.9200000762939 17.202289534821
22.960000038147 16.9535230777206
23 16.7039720263089
23.0400002002716 16.4545381575863
23.0800001621246 16.2050122028874
23.1200001239777 15.9548622987454
23.1600000858307 15.7048234752558
23.2000000476837 15.4549346808338
23.2400000095367 15.2051046483577
23.2800002098083 14.955120579584
23.3200001716614 14.7050584746453
23.3600001335144 14.4552137385411
23.4000000953674 14.2048617999041
23.4400000572205 13.9540893575258
23.4800000190735 13.7039293383133
23.5199999809265 13.4516092889239
23.5600001811981 13.1993863861105
23.6000001430511 12.9465123234725
23.6400001049042 12.69350372124
23.6800000667572 12.4405551971526
23.7200000286102 12.1872920285993
23.7599999904633 11.9334940906324
23.8000001907349 11.6799448338252
23.8400001525879 11.426291113364
23.8800001144409 11.1731234817774
23.9200000762939 10.9198860101382
23.960000038147 10.6668188715893
24 10.4146965541029
24.0400002002716 10.1636406218949
24.0800001621246 9.91344282755913
24.1200001239777 9.6631160583321
24.1600000858307 9.41340511002221
24.2000000476837 9.16296972336683
24.2400000095367 8.91224697782263
24.2800002098083 8.6613473053971
24.3200001716614 8.41084024104772
24.3600001335144 8.16044470829353
24.4000000953674 7.90966275926215
24.4400000572205 7.65870030456552
24.4800000190735 7.40827563899264
24.5199999809265 7.15797719064307
24.5600001811981 6.90801385414846
24.6000001430511 6.65765349164203
24.6400001049042 6.40764621441656
24.6800000667572 6.15797722523948
24.7200000286102 5.90872236761659
24.7599999904633 5.65970831361437
24.8000001907349 5.41036750494398
24.8400001525879 5.16131861787894
24.8800001144409 4.91192559128296
24.9200000762939 4.66250596482899
24.960000038147 4.41329970590703
25 4.1636741208305
25.0400002002716 3.91396756286075
25.0800001621246 3.6644408763731
25.1200001239777 3.41046217096077
25.1600000858307 3.16057783074613
25.2000000476837 2.91072241549958
25.2400000095367 2.66037380271954
25.2800002098083 2.4103527264704
25.3200001716614 2.15995148349976
25.3600001335144 1.91001296444877
25.4000000953674 1.65990240585027
25.4400000572205 1.41032819933659
25.4800000190735 1.16061433611834
25.5199999809265 0.910811568931973
25.5600001811981 0.66148801482618
25.6000001430511 0.411402739427066
25.6400001049042 0.160781639393686
25.6800000667572 -0.0898112608159694
25.7200000286102 -0.340664168378237
25.7599999904633 -0.59153821833448
25.8000001907349 -0.842725440563792
25.8400001525879 -1.09356091729119
25.8800001144409 -1.34471368372387
25.9200000762939 -1.59587085991672
25.960000038147 -1.84697025984159
26 -2.09830544098833
26.0400002002716 -2.34945707323683
26.0800001621246 -2.60022234182376
26.1200001239777 -2.85054307501551
26.1600000858307 -3.10067257407531
26.2000000476837 -3.35152228025431
26.2400000095367 -3.60290706315831
26.2800002098083 -3.85485035732425
26.3200001716614 -4.10689346841187
26.3600001335144 -4.35916383306359
26.4000000953674 -4.61140687276691
26.4400000572205 -4.86368395901019
26.4800000190735 -5.11621882441978
26.5199999809265 -5.36867396369211
26.5600001811981 -5.62075975996893
26.6000001430511 -5.87264712197534
26.6400001049042 -6.12431815446832
26.6800000667572 -6.37618390975538
26.7200000286102 -6.62801394705092
26.7599999904633 -6.87935049609781
26.8000001907349 -7.13063141663258
26.8400001525879 -7.38185382191955
26.8800001144409 -7.63310604086204
26.9200000762939 -7.88434526008133
26.960000038147 -8.13629612969886
27 -8.38799030963377
27.0400002002716 -8.63961120406166
27.0800001621246 -8.89193691961668
27.1200001239777 -9.14404363468748
27.1600000858307 -9.39566893000909
27.2000000476837 -9.64729622211332
27.2400000095367 -9.89908071316457
27.2800002098083 -10.1504419896779
27.3200001716614 -10.4020868747503
27.3600001335144 -10.6526775786703
27.4000000953674 -10.9031656901268
27.4400000572205 -11.1530968028841
27.4800000190735 -11.4031116811071
27.5199999809265 -11.6531185646257
27.5600001811981 -11.903024547133
27.6000001430511 -12.1531702473933
27.6400001049042 -12.4035431429459
27.6800000667572 -12.6536518891218
27.7200000286102 -12.9037818024708
27.7599999904633 -13.1537432814631
27.8000001907349 -13.4041263577938
27.8400001525879 -13.6546924414542
27.8800001144409 -13.9051902490042
27.9200000762939 -14.1570200899525
27.960000038147 -14.4079210852829
28 -14.6581620410596
28.0400002002716 -14.9077834504288
28.0800001621246 -15.1572916172504
28.1200001239777 -15.4068213184548
28.1600000858307 -15.6562591984014
28.2000000476837 -15.9062961191131
28.2400000095367 -16.1564870935656
28.2800002098083 -16.4064019286268
28.3200001716614 -16.6567447449977
28.3600001335144 -16.9068841790189
28.4000000953674 -17.1573678577616
28.4400000572205 -17.4081428097771
28.4800000190735 -17.6590326828916
28.5199999809265 -17.9099635948383
28.5600001811981 -18.1607080649002
28.6000001430511 -18.4112955939156
28.6400001049042 -18.6619079006554
28.6800000667572 -18.9127360815621
28.7200000286102 -19.1634117427168
28.7599999904633 -19.4134454544847
28.8000001907349 -19.6638689645441
28.8400001525879 -19.9145966654544
28.8800001144409 -20.1643008876558
28.9200000762939 -20.4148265806106
28.960000038147 -20.6651612582664
29 -20.9161327302933
29.0400002002716 -21.1678452767611
29.0800001621246 -21.4196657600986
29.1200001239777 -21.6710868265815
29.1600000858307 -21.9227981843048
29.2000000476837 -22.174428535416
29.2400000095367 -22.4255641885559
29.2800002098083 -22.6764721829951
29.3200001716614 -22.9279361742118
29.3600001335144 -23.1792142937129
29.4000000953674 -23.4304529264594
29.4400000572205 -23.6816277972663
29.4800000190735 -23.9323202454243
29.5199999809265 -24.1829121901347
29.5600001811981 -24.4334753146694
29.6000001430511 -24.6840727549135
29.6400001049042 -24.9349498332542
29.6800000667572 -25.1852717997397
29.7200000286102 -25.4353035832072
29.7599999904633 -25.6851541800636
29.8000001907349 -25.9356122078494
29.8400001525879 -26.1861004903034
29.8800001144409 -26.435725346515
29.9200000762939 -26.6862659981062
29.960000038147 -26.9355077142518
30 -27.1851403836143
30.0400002002716 -27.4343687275215
30.0800001621246 -27.6837120231568
30.1200001239777 -27.9331002692602
30.1600000858307 -28.1830084620029
30.2000000476837 -28.4326941087759
30.2400000095367 -28.6825677363671
30.2800002098083 -28.9311672565118
30.3200001716614 -29.1794634032031
30.3600001335144 -29.4269262364131
30.4000000953674 -29.6757638811973
30.4400000572205 -29.9236873212456
30.4800000190735 -30.1731422541724
30.5199999809265 -30.4215908067139
30.5600001811981 -30.6689437779959
30.6000001430511 -30.9169328361235
30.6400001049042 -31.1646489972985
30.6800000667572 -31.4120994050819
30.7200000286102 -31.6591438336593
30.7599999904633 -31.9068416640432
30.8000001907349 -32.155454239994
30.8400001525879 -32.4011850624117
30.8800001144409 -32.6480014672964
30.9200000762939 -32.895720033516
30.960000038147 -33.142910812059
31 -33.3904483715059
31.0400002002716 -33.6397445152825
31.0800001621246 -33.8885898949994
31.1200001239777 -34.135976004499
31.1600000858307 -34.3823958939774
31.2000000476837 -34.6288573277577
31.2400000095367 -34.873695470996
31.2800002098083 -35.1200161690857
31.3200001716614 -35.3694530481699
31.3600001335144 -35.6193213597913
31.4000000953674 -35.8693258612159
31.4400000572205 -36.1176556133975
31.4800000190735 -36.368372687031
31.5199999809265 -36.616107428335
31.5600001811981 -36.864436065903
31.6000001430511 -37.1223447582784
31.6400001049042 -37.3829906944785
31.6800000667572 -37.6323460354048
31.7200000286102 -37.8829340169007
31.7599999904633 -38.1334249245977
31.8000001907349 -38.3811259864001
31.8400001525879 -38.6255686624233
31.8800001144409 -38.8721771605604
31.9200000762939 -39.1184131717691
31.960000038147 -39.3657804213787
32 -39.6133783333926
32.0400002002716 -39.8653647030356
32.0800001621246 -40.1152547104709
32.1200001239777 -40.3468182635226
32.1600000858307 -40.5982964491093
32.2000000476837 -40.8508694061365
32.2400000095367 -41.1018683634634
32.2800002098083 -41.3481569348728
32.3200001716614 -41.5988502368218
32.3600001335144 -41.8497175864068
32.4000000953674 -42.1016204946384
32.4400000572205 -42.3563765135033
32.4800000190735 -42.6067553166379
32.5199999809265 -42.8605561609665
32.5600001811981 -43.1123823217141
32.6000001430511 -43.3644422338657
32.6400001049042 -43.6148984073859
32.6800000667572 -43.8705664539516
32.7200000286102 -44.124004625709
32.7599999904633 -44.3830507666645
32.8000001907349 -44.6429121899369
32.8400001525879 -44.891883672455
32.8800001144409 -45.1132104586108
32.9200000762939 -45.3658799832957
32.960000038147 -45.6183144355328
33 -45.8728416185282
33.0400002002716 -46.1244282063641
33.0800001621246 -46.3780250628728
33.1200001239777 -46.629640021988
33.1600000858307 -46.8811776996769
33.2000000476837 -47.1317261861055
33.2400000095367 -47.3838468187955
33.2800002098083 -47.6370728531147
33.3200001716614 -47.8895552130781
33.3600001335144 -48.1767597663419
33.4000000953674 -48.445251124502
33.4400000572205 -48.700369210031
33.4800000190735 -48.9560792606711
33.5199999809265 -49.2091274871409
33.5600001811981 -49.4630959422897
33.6000001430511 -49.7160740555733
33.6400001049042 -49.9723448197502
33.6800000667572 -50.2273517707698
33.7200000286102 -50.4833576371101
33.7599999904633 -50.7373973388056
33.8000001907349 -50.9909787714524
33.8400001525879 -51.2448373823363
33.8800001144409 -51.498799091827
33.9200000762939 -51.7519371681672
33.960000038147 -52.0052327324858
34 -52.2602762255826
34.0400002002716 -52.5148682300343
34.0800001621246 -52.7700201273581
34.1200001239777 -53.0245602068897
34.1600000858307 -53.2779622756575
};
\node at (axis cs:6,55)[
  anchor=north,
  text=black,
  rotate=0.0
]{ Pickup truck};
\node at (axis cs:6,0)[
  anchor=south,
  text=black,
  rotate=0.0
]{ Station wagon};
\end{axis}

\end{tikzpicture}%
		\subcaption{Longitudinal distance of other vehicles with respect to ego vehicle.}
		\label{fig:example dynamic environment lon}
	\end{minipage}%
	\hfill
	\begin{minipage}[b]{0.49\linewidth}
		\centering
		% This file was created by matplotlib2tikz v0.6.14.
\begin{tikzpicture}

\definecolor{color3}{rgb}{0.549019607843137,0.337254901960784,0.294117647058824}
\definecolor{color0}{rgb}{0.172549019607843,0.627450980392157,0.172549019607843}
\definecolor{color2}{rgb}{0.580392156862745,0.403921568627451,0.741176470588235}
\definecolor{color1}{rgb}{0.83921568627451,0.152941176470588,0.156862745098039}

\begin{axis}[
xlabel={$t$ [s]},
ylabel={Relative lateral distance [m]},
xmin=0, xmax=34.1600000858307,
ymin=-3.6944648464502, ymax=3.79876173773052,
width=\figurewidth,
height=\figureheight,
tick align=outside,
tick pos=left,
xmajorgrids,
x grid style={lightgray!92.026143790849673!black},
ymajorgrids,
y grid style={lightgray!92.026143790849673!black},
clip marker paths
]
\addplot [semithick, color0, mark=*, mark size=1, mark options={solid}, only marks, forget plot]
table {%
0 -0.0417787396782297
0.0400002002716064 -0.0434833324633892
0.0800001621246338 -0.00793896066897679
0.120000123977661 -0.0290788596134219
0.160000085830688 -0.00511650671826534
0.200000047683716 0.026785123248344
0.240000009536743 0.0205590435513228
0.28000020980835 -0.0226561519444399
0.320000171661377 -0.00188009407999191
0.360000133514404 -0.00377719282237902
0.400000095367432 -0.0155240147251251
0.440000057220459 -0.0306548872097677
0.480000019073486 -0.0606875543387061
0.519999980926514 -0.0349109763945282
0.56000018119812 -0.0274273183054372
0.600000143051147 -0.0255048803860961
0.640000104904175 -0.0142116114614845
0.680000066757202 -0.0144869500468184
0.720000028610229 -0.00634319719655396
0.759999990463257 0.00264264047317925
0.800000190734863 0.0105393193251953
0.840000152587891 0.0138097292343392
0.880000114440918 0.0291535571020424
0.920000076293945 0.0366716656258401
0.960000038146973 0.0297297342137315
1 0.0217351281834865
1.04000020027161 0.0103010571788247
1.08000016212463 -0.012837636547751
1.12000012397766 -0.000666284051467028
1.16000008583069 -0.0196001067388685
1.20000004768372 -0.00508295675099915
1.24000000953674 -0.0302094671738396
1.28000020980835 -0.0652997387717573
1.32000017166138 -0.0518677885327045
1.3600001335144 -0.062191293287375
1.40000009536743 -0.0583230235330692
1.44000005722046 -0.0637485272072151
1.48000001907349 -0.0848910782785489
1.51999998092651 -0.0887539614874596
1.56000018119812 -0.0915527800488931
1.60000014305115 -0.094033839429671
1.64000010490417 -0.0986166093511816
1.6800000667572 -0.101533324659034
1.72000002861023 -0.097418266876103
1.75999999046326 -0.0951203112193729
1.80000019073486 -0.0768142142537029
1.84000015258789 -0.0758488769996015
1.88000011444092 -0.0657215850608496
1.92000007629395 -0.0596173475464898
1.96000003814697 -0.0774694105485675
2 -0.0819576342397531
2.04000020027161 -0.0737450374304574
2.08000016212463 -0.0632697980881775
2.12000012397766 -0.065146546111649
2.16000008583069 -0.0379296996072994
2.20000004768372 -0.0298124907636027
2.24000000953674 -0.0104287031121793
2.28000020980835 -0.000783934746713261
2.32000017166138 -0.0146727622936489
2.3600001335144 -0.0116230187023605
2.40000009536743 -0.00800980925905722
2.44000005722046 -0.0569500180955829
2.48000001907349 -0.0194943770969402
2.51999998092651 0.0151885988311344
2.56000018119812 0.0656892286827005
2.60000014305115 0.072081029845608
2.64000010490417 0.0845924477338064
2.6800000667572 0.105742618639308
2.72000002861023 0.119406130685862
2.75999999046326 0.12824628236362
2.80000019073486 0.139938514157885
2.84000015258789 0.148019094859861
2.88000011444092 0.25583690341537
2.92000007629395 0.306554447711603
2.96000003814697 0.278564315163056
3 0.244764365743769
3.04000020027161 0.234652405155853
3.08000016212463 0.237902782307596
3.12000012397766 0.235483204363214
3.16000008583069 0.239052413177167
3.20000004768372 0.232617816658026
3.24000000953674 0.260339993436241
3.28000020980835 0.258706864517433
3.32000017166138 0.3574681748221
3.3600001335144 0.901479815125867
3.40000009536743 0.918641901317467
3.44000005722046 0.952841176995706
3.48000001907349 0.986183120187692
3.51999998092651 0.988627684322505
3.56000018119812 1.00186064011514
3.60000014305115 1.01180667626192
3.64000010490417 1.01591165870319
3.6800000667572 1.04432666445691
3.72000002861023 1.06150147302882
3.75999999046326 1.09942249772943
3.80000019073486 1.11890528969187
3.84000015258789 0.636547295423912
3.88000011444092 0.405911449278643
3.92000007629395 0.389694720289353
3.96000003814697 0.41780990901604
4 0.414559403397376
4.04000020027161 0.420269918803582
4.08000016212463 0.449963792167126
4.12000012397766 0.418271949134033
4.16000008583069 0.429857636237069
4.20000004768372 0.407949681378507
4.24000000953674 0.415058137103386
4.28000020980835 0.431460320618351
4.32000017166138 0.413362124001608
4.3600001335144 0.415399411612331
4.40000009536743 0.395302178564078
4.44000005722046 0.415159251091263
4.48000001907349 0.40905699661753
4.51999998092651 0.392595507550862
4.56000018119812 0.253553725662647
4.60000014305115 0.250807804710817
4.64000010490417 0.263448388963936
4.6800000667572 0.283164182167046
4.72000002861023 0.260610309470746
4.75999999046326 0.264162076704724
4.80000019073486 0.249604925137301
4.84000015258789 0.225354232271675
4.88000011444092 0.22634417982975
4.92000007629395 0.198150775817807
4.96000003814697 0.197521672397681
5 0.182313265544689
5.04000020027161 0.194446681398852
5.08000016212463 0.221427069799978
5.12000012397766 0.189127232670235
5.16000008583069 0.209346408108924
5.20000004768372 0.163559505577604
5.24000000953674 0.139405583163256
5.28000020980835 0.16605103007943
5.32000017166138 0.119682918201199
5.3600001335144 0.131886410955644
5.40000009536743 0.112349399868415
5.44000005722046 0.11047751927303
5.48000001907349 0.144443713438279
5.51999998092651 0.117395918805469
5.56000018119812 0.1614418734515
5.60000014305115 0.12118195280291
5.64000010490417 0.168949677401619
5.6800000667572 0.213932440600319
5.72000002861023 0.18063949003915
5.75999999046326 0.220967184431904
5.80000019073486 0.191980929452678
5.84000015258789 0.205817819563089
5.88000011444092 0.234895748189027
5.92000007629395 0.1917900253487
5.96000003814697 0.201873463051521
6 0.191549293366018
6.04000020027161 0.165585265412575
6.08000016212463 0.214685869194678
6.12000012397766 0.329367520623184
6.16000008583069 0.337115141512145
6.20000004768372 0.325935692683958
6.24000000953674 0.372339743198021
6.28000020980835 0.388521744711257
6.32000017166138 0.393964485258088
6.3600001335144 0.456375310335519
6.40000009536743 0.469509591666124
6.44000005722046 0.45917285117698
6.48000001907349 0.527237683952156
6.51999998092651 0.502872396074831
6.56000018119812 0.524079771777374
6.60000014305115 0.523329068632831
6.64000010490417 0.527918935909131
6.6800000667572 0.545962045217118
6.72000002861023 0.502727470836427
6.75999999046326 0.482184477662969
6.80000019073486 0.477095554564596
6.84000015258789 0.444204777600845
6.88000011444092 0.44850776699175
6.92000007629395 0.416270276152863
6.96000003814697 0.420725740179746
7 0.414874046023105
7.04000020027161 0.38603477883713
7.08000016212463 0.399460643621401
7.12000012397766 0.392432457259056
7.16000008583069 0.383393069805033
7.20000004768372 0.37759507821646
7.24000000953674 0.360152991313975
7.28000020980835 0.369336407770393
7.32000017166138 0.318961385583633
7.3600001335144 0.338294565931529
7.40000009536743 0.351662758020797
7.44000005722046 0.334317194704018
7.48000001907349 0.351272703445039
7.51999998092651 0.356151038601924
7.56000018119812 0.370647339093179
7.60000014305115 0.361366313709705
7.64000010490417 0.358361932543915
7.6800000667572 0.377377270780715
7.72000002861023 0.367017494780768
7.75999999046326 0.385762241170437
7.80000019073486 0.377993316734781
7.84000015258789 0.385367301479262
7.88000011444092 0.42649928924518
7.92000007629395 0.376952841485367
7.96000003814697 0.32511495293492
8 0.313859984558447
8.04000020027161 0.303630483685739
8.08000016212463 0.354531774151929
8.12000012397766 0.298820160892349
8.16000008583069 0.284601621841473
8.20000004768372 0.27365181123016
8.24000000953674 0.255467651706338
8.28000020980835 0.298369025216015
8.32000017166138 0.225502205215161
8.3600001335144 0.23583690821237
8.40000009536743 0.224470184450044
8.44000005722046 0.220576961146684
8.48000001907349 0.284973198074928
8.51999998092651 0.210160351948411
8.56000018119812 0.232056477010676
8.60000014305115 0.22463754242515
8.64000010490417 0.246830359009954
8.6800000667572 0.320988924900878
8.72000002861023 0.235447281501163
8.75999999046326 0.242505953399173
8.80000019073486 0.308828695113704
8.84000015258789 0.226868064053963
8.88000011444092 0.298516583395571
8.92000007629395 0.112239341909362
8.96000003814697 0.233319506270085
9 0.222930825871416
9.04000020027161 0.173799730495064
9.08000016212463 0.133034145241127
9.12000012397766 0.271600448852642
9.16000008583069 0.240262057773767
9.20000004768372 0.272490445382458
9.24000000953674 0.290509172202269
9.28000020980835 0.203546656494682
9.32000017166138 0.240886270162535
9.3600001335144 0.292646349281046
9.40000009536743 0.24357226985176
9.44000005722046 0.0457886916948231
9.48000001907349 0.113493234988786
9.51999998092651 -0.000772053611664832
9.56000018119812 0.112654417957984
9.60000014305115 0.19868405246894
9.64000010490417 0.289832259579926
9.6800000667572 0.336780099632129
9.72000002861023 0.546307860143487
9.75999999046326 0.605127321445179
9.80000019073486 0.709152583080343
9.84000015258789 0.86787252442825
9.88000011444092 0.855295164046075
9.92000007629395 1.01077547057974
9.96000003814697 1.03499509359713
10 1.09273377718478
10.0400002002716 1.1245075426224
10.0800001621246 1.04513123852704
10.1200001239777 0.590629446494773
10.1600000858307 0.490265626919553
10.2000000476837 0.525338806170876
10.2400000095367 0.582509165031396
10.2800002098083 0.550542282891613
10.3200001716614 0.569727304410648
10.3600001335144 0.58622729699181
10.4000000953674 0.600210523025847
10.4400000572205 0.745514249594854
10.4800000190735 0.744749887615746
10.5199999809265 0.891990517277949
10.5600001811981 0.899691473428397
10.6000001430511 1.20455610658641
10.6400001049042 1.70099379394201
10.6800000667572 1.76319378604241
10.7200000286102 1.8848609444214
10.7599999904633 1.942461616096
10.8000001907349 2.00791692021982
10.8400001525879 2.08900244115099
10.8800001144409 2.15631064645281
10.9200000762939 2.24277276393423
10.960000038147 2.26034880712028
11 2.2847860362985
11.0400002002716 2.38502518267754
11.0800001621246 2.41149425484073
11.1200001239777 2.46287144341076
11.1600000858307 2.45554859483347
11.2000000476837 2.50325144504693
11.2400000095367 2.4748176006217
11.2800002098083 2.53893538470829
11.3200001716614 2.6149336120236
11.3600001335144 2.58903833083674
11.4000000953674 2.6709464333538
11.4400000572205 2.71302412942567
11.4800000190735 2.71737212973987
11.5199999809265 2.73644748425465
11.5600001811981 2.78274488624952
11.6000001430511 2.81908712925223
11.6400001049042 2.84631088708817
11.6800000667572 2.87858715226997
11.7200000286102 2.8526665423972
11.7599999904633 2.82163607458537
11.8000001907349 2.81946235112765
11.8400001525879 2.77045666703232
11.8800001144409 2.81113021465049
11.9200000762939 2.83063097523508
11.960000038147 2.8547743756347
12 2.85652138845991
12.0400002002716 2.83735527463517
12.0800001621246 2.84826293840178
12.1200001239777 2.94650899192099
12.1600000858307 2.96834716525122
12.2000000476837 2.9459678874127
12.2400000095367 2.95552809754005
12.2800002098083 2.98052639741588
12.3200001716614 3.03308143254575
12.3600001335144 3.03516685628485
12.4000000953674 3.03830755409455
12.4400000572205 3.05198260959917
12.4800000190735 3.0678785048125
12.5199999809265 3.07061290230589
12.5600001811981 3.08839716534139
12.6000001430511 3.0832046156668
12.6400001049042 3.09135883909449
12.6800000667572 3.10858451990355
12.7200000286102 3.03941106173655
12.7599999904633 3.0286950647477
12.8000001907349 3.0340884180457
12.8400001525879 2.99966137050916
12.8800001144409 3.00007675451334
12.9200000762939 2.92460981991743
12.960000038147 2.91652296508919
13 2.84387097629045
13.0400002002716 2.82922891089638
13.0800001621246 2.784352756705
13.1200001239777 2.7617945740745
13.1600000858307 2.75657701446277
13.2000000476837 2.63835897994763
13.2400000095367 2.61485943274389
13.2800002098083 2.59233925987106
13.3200001716614 2.54632233644069
13.3600001335144 2.53502060367273
13.4000000953674 2.53048062053379
13.4400000572205 2.4933923496161
13.4800000190735 2.49556359346183
13.5199999809265 2.49158551094374
13.5600001811981 2.52205181480324
13.6000001430511 2.55991862959851
13.6400001049042 2.55413321029144
13.6800000667572 2.54124768509059
13.7200000286102 2.60696714951397
13.7599999904633 2.60402189107454
13.8000001907349 2.66285799233406
13.8400001525879 2.71306910117497
13.8800001144409 2.71683005828622
13.9200000762939 2.72031786075228
13.960000038147 2.73824721379009
14 2.74174426664553
14.0400002002716 2.78187433127749
14.0800001621246 2.86799681689796
14.1200001239777 2.87925312794838
14.1600000858307 2.89610280642867
14.2000000476837 2.91435251281274
14.2400000095367 2.92266929287403
14.2800002098083 2.93546697106999
14.3200001716614 3.00334255019609
14.3600001335144 3.01768822302289
14.4000000953674 3.03556724485131
14.4400000572205 3.01297996432891
14.4800000190735 3.01002518731099
14.5199999809265 3.0251960717066
14.5600001811981 3.00817321327189
14.6000001430511 3.01049629972815
14.6400001049042 2.99215120632159
14.6800000667572 2.99225186649551
14.7200000286102 2.95555593367695
14.7599999904633 2.9528534374746
14.8000001907349 2.95507923768087
14.8400001525879 2.99711772569004
14.8800001144409 2.99898927695971
14.9200000762939 3.00566460827698
14.960000038147 3.01015581961001
15 3.01644627112053
15.2000000476837 3.3620217705862
15.2400000095367 3.36236510006473
15.2800002098083 3.35264191549625
15.3200001716614 3.45816052935867
15.3600001335144 3.28908788760309
15.4000000953674 3.38583134435336
15.4400000572205 3.25156133938931
15.4800000190735 3.23405095981175
15.5199999809265 3.32269457553571
15.5600001811981 3.31887655673157
15.6000001430511 3.40343196799003
15.6400001049042 3.40716664924312
15.6800000667572 3.40956343709971
15.7200000286102 3.31672020465732
15.7599999904633 3.24866368285574
15.8000001907349 3.32397850709351
15.8400001525879 3.32650928617088
15.8800001144409 3.20824210039003
15.9200000762939 3.18228634701947
15.960000038147 3.15800436307625
16 3.2586327824088
16.0400002002716 3.25299428130253
16.0800001621246 3.22557641523254
16.1200001239777 3.18870393876074
16.1600000858307 3.21108888850197
16.2000000476837 3.16667358794859
16.2400000095367 3.16149233014769
16.2800002098083 3.12760682312541
16.3200001716614 3.02746670675659
16.3600001335144 2.99932713581388
16.4000000953674 3.0794154043405
16.4400000572205 3.07811237480725
16.4800000190735 3.17230242186235
16.5199999809265 3.13040485931683
16.5600001811981 3.16553648913305
16.6000001430511 3.32523983633196
16.6400001049042 3.33552316814808
16.6800000667572 3.20596796864897
16.7200000286102 3.15341218549578
16.7599999904633 3.14067369762768
16.8000001907349 3.12112796439612
16.8400001525879 3.11060822932094
16.8800001144409 3.15898096982907
16.9200000762939 3.09511129442034
16.960000038147 3.23231214923136
17 3.19656782795671
17.0400002002716 3.23766812528487
17.0800001621246 3.20983969689052
17.1200001239777 3.25479426494558
17.1600000858307 3.18412109014656
17.2000000476837 3.18202625105208
17.2400000095367 3.16884716106373
17.2800002098083 3.20181181156275
17.3200001716614 3.19369118608081
17.3600001335144 3.20997527150776
17.4000000953674 3.24944315752324
17.4400000572205 3.12758829404063
17.4800000190735 3.10249432758436
17.5199999809265 3.13525716009609
17.5600001811981 3.02941396537664
17.6000001430511 3.04261737822625
17.6400001049042 3.05885295114235
17.6800000667572 3.08026320160039
17.7200000286102 2.98320854652992
};
\addplot [semithick, color2, mark=*, mark size=1, mark options={solid}, only marks, forget plot]
table {%
0.800000190734863 0.566310207694819
0.840000152587891 0.610843449352141
0.880000114440918 0.646629711214613
0.920000076293945 0.675662457956674
0.960000038146973 0.677383853894884
1 0.652233842837334
1.04000020027161 0.617705696642899
1.08000016212463 0.59707644187842
1.12000012397766 0.582773435789414
1.16000008583069 0.57319667237221
1.20000004768372 0.544899207724968
1.24000000953674 0.520278560671475
1.28000020980835 0.486819697340015
1.32000017166138 0.471235098019191
1.3600001335144 0.482251080844038
1.40000009536743 0.498766569171103
1.44000005722046 0.525970282627713
1.48000001907349 0.546432838878246
1.51999998092651 0.570399635333319
1.56000018119812 0.603940263122975
1.60000014305115 0.637042374110034
1.64000010490417 0.671570818616942
1.6800000667572 0.768573049015447
1.72000002861023 0.854839414125309
1.75999999046326 0.921061127208674
1.80000019073486 0.972855269388712
1.84000015258789 1.02464311800119
1.88000011444092 1.06445564214345
1.92000007629395 0.997485319848947
1.96000003814697 0.905785809928621
2 0.826862849189619
2.04000020027161 0.833231959205346
2.08000016212463 0.651094119210421
2.12000012397766 0.258213124322324
2.16000008583069 0.236399503264688
2.20000004768372 0.0448681660358898
2.24000000953674 0.0231399156933656
2.28000020980835 -0.0463956041838169
2.32000017166138 0.0291771100656393
2.3600001335144 0.043865510420214
2.40000009536743 0.0832582982314672
2.44000005722046 -0.0523143689140556
2.48000001907349 0.0410986709152589
2.51999998092651 0.0119072781358935
2.56000018119812 -0.0342830780656945
2.60000014305115 -0.0409483506911916
2.64000010490417 0.00288127999770402
2.6800000667572 0.0280583743270773
2.72000002861023 -0.00967951797540012
2.75999999046326 -0.0235366700101965
2.80000019073486 0.0118969778108895
2.84000015258789 0.0404327908569302
2.88000011444092 -0.0160557699592126
2.92000007629395 0.00551272847477687
2.96000003814697 0.0458368679025721
3 0.00325816038120176
3.04000020027161 0.0808821824303544
3.08000016212463 0.104864896536889
3.12000012397766 0.0942851146853216
3.16000008583069 0.12244417347231
3.20000004768372 0.142543943031819
3.24000000953674 0.158265555786577
3.28000020980835 0.178455439635281
3.32000017166138 0.0464456730218337
3.3600001335144 1.71684116383779
3.40000009536743 1.80984270687744
3.44000005722046 1.8959152879875
3.48000001907349 1.96452916292364
3.51999998092651 2.01682161590802
3.56000018119812 2.07905052524954
3.60000014305115 2.14780912078468
3.64000010490417 2.2216339958221
3.6800000667572 2.27966431102876
3.72000002861023 2.35061349917311
3.75999999046326 2.42485896977757
3.80000019073486 2.48595951961801
3.84000015258789 1.35608615843019
3.88000011444092 0.800633217619586
3.92000007629395 0.806348912787322
3.96000003814697 0.793682823714538
4 0.781876202946178
4.04000020027161 0.779983413198837
4.08000016212463 0.743868245406269
4.12000012397766 0.714927957546375
4.16000008583069 0.673641919155063
4.20000004768372 0.640113584474457
4.24000000953674 0.609827272549596
4.28000020980835 0.619925098724548
4.32000017166138 0.593513720024116
4.3600001335144 0.585118610765496
4.40000009536743 0.542491445026684
4.44000005722046 0.514798994022574
4.48000001907349 0.495125997341028
4.51999998092651 0.455079345551271
4.56000018119812 0.411902997749079
4.60000014305115 0.414135237819253
4.64000010490417 0.392015050219518
4.6800000667572 0.381675412418476
4.72000002861023 0.373460900595211
4.75999999046326 0.356430183844504
4.80000019073486 0.353605795312068
4.84000015258789 0.338985483920712
4.88000011444092 0.297862774116112
4.92000007629395 0.270229610660648
4.96000003814697 0.234806377942342
5 0.197459938713907
5.04000020027161 0.160225007734228
5.08000016212463 0.129040816772171
5.12000012397766 0.10975342962631
5.16000008583069 0.102104780014038
5.20000004768372 0.070453461075896
5.24000000953674 0.0478691878701982
5.28000020980835 0.0317514027692696
5.32000017166138 0.0180137507933136
5.3600001335144 0.0146181817805523
5.40000009536743 -0.0146888413641624
5.44000005722046 -0.0213805046728804
5.48000001907349 -0.00648853465014628
5.51999998092651 0.00530702145781812
5.56000018119812 0.00225931900450016
5.60000014305115 -0.00310089537004378
5.64000010490417 -0.00748847930732433
5.6800000667572 0.030737689679129
5.72000002861023 0.0600183831049996
5.75999999046326 0.0508558302315239
5.80000019073486 0.0936188338125774
5.84000015258789 0.102305691474107
5.88000011444092 0.105565447949982
5.92000007629395 0.0765239492002127
5.96000003814697 0.112394011334227
6 0.120328411551254
6.04000020027161 0.110537871326804
6.08000016212463 0.138945929325808
6.12000012397766 0.166163563309374
6.16000008583069 0.183933603125409
6.20000004768372 0.1953600528994
6.24000000953674 0.197118693448992
6.28000020980835 0.195238790931309
6.32000017166138 0.199424579114837
6.3600001335144 0.197564247717759
6.40000009536743 0.190046991038282
6.44000005722046 0.183229822318331
6.48000001907349 0.152979201536525
6.51999998092651 0.125625637089855
6.56000018119812 0.0841914873662651
6.60000014305115 0.0510819772389325
6.64000010490417 0.0504747655169298
6.6800000667572 0.00932592477742102
6.72000002861023 -0.0129447702698416
6.75999999046326 -0.0164690498439564
6.80000019073486 -0.0262103437160532
6.84000015258789 -0.0329670587294698
6.88000011444092 -0.0384263744283208
6.92000007629395 -0.0262605696872765
6.96000003814697 -0.0351780229515891
7 -0.0432309925837881
7.04000020027161 -0.0151981113250474
7.08000016212463 0.0338100476016889
7.12000012397766 0.0357679562669317
7.16000008583069 0.0592956410198714
7.20000004768372 0.0931031224978316
7.24000000953674 0.109539192761189
7.28000020980835 0.146419504304959
7.32000017166138 0.177827589876486
7.3600001335144 0.208639028888138
7.40000009536743 0.241338746802459
7.44000005722046 0.242672236542932
7.48000001907349 0.289665332091754
7.51999998092651 0.321895144934148
7.56000018119812 0.328060234949125
7.60000014305115 0.34673592934652
7.64000010490417 0.389457517074822
7.6800000667572 0.427037136242893
7.72000002861023 0.457571806512278
7.75999999046326 0.469860417641607
7.80000019073486 0.477150907636018
7.84000015258789 0.50638483287477
7.88000011444092 0.517117167656501
7.92000007629395 0.481972901194229
7.96000003814697 0.520853194232856
8 0.555521718227795
8.04000020027161 0.524745312629409
8.08000016212463 0.534539829863825
8.12000012397766 0.570375525037023
8.16000008583069 0.516393241180434
8.20000004768372 0.537096404126594
8.24000000953674 0.482310099671144
8.28000020980835 0.501713631853749
8.32000017166138 0.523627180729363
8.3600001335144 0.487648713697725
8.40000009536743 0.510649260771908
8.44000005722046 0.470145698223579
8.48000001907349 0.481780408790086
8.51999998092651 0.499816324526014
8.56000018119812 0.489557610179938
8.60000014305115 0.513996777222232
8.64000010490417 0.513186447871466
8.6800000667572 0.535450528900637
8.72000002861023 0.53446231575872
8.75999999046326 0.546404819940854
8.80000019073486 0.564148460926211
8.84000015258789 0.551930000078432
8.88000011444092 0.58428951203691
8.92000007629395 0.764741188347422
8.96000003814697 1.55375805788906
9 1.49972280383137
9.04000020027161 1.42908013881005
9.08000016212463 1.34565933230997
9.12000012397766 1.28410794596156
9.16000008583069 1.20190537542625
9.20000004768372 1.05106507492809
9.24000000953674 0.97219118376518
9.28000020980835 0.85186529657332
9.32000017166138 0.793467436620649
9.3600001335144 0.711929168654996
9.40000009536743 0.636589809799067
9.44000005722046 0.0111924741370204
9.48000001907349 0.184432560387042
9.51999998092651 0.0928070148279042
9.56000018119812 -0.0035604359525262
9.60000014305115 0.00764646137104876
9.64000010490417 0.109509650398619
9.6800000667572 0.183795393463004
9.72000002861023 0.28361730148163
9.75999999046326 0.369584729844607
9.80000019073486 0.436554177496068
9.84000015258789 0.53993090817934
9.88000011444092 0.60936556681791
9.92000007629395 0.701576671484731
9.96000003814697 0.766588155874757
10 0.795320818150348
10.0400002002716 0.822163739214917
10.0800001621246 0.819465444767499
10.1200001239777 0.160630832338604
10.1600000858307 0.368721974828637
10.2000000476837 0.429991782721702
10.2400000095367 0.491752522703567
10.2800002098083 0.55212876541671
10.3200001716614 0.57109841778446
10.3600001335144 0.623836037549689
10.4000000953674 0.67635406648016
10.4400000572205 0.702848258560234
10.4800000190735 0.738574248868264
10.5199999809265 0.787722261072312
10.5600001811981 0.829154987026429
10.6000001430511 0.390461506995752
10.6400001049042 0.088210622555838
10.6800000667572 0.118442428606877
10.7200000286102 0.0862487188791029
10.7599999904633 0.109056080858271
10.8000001907349 0.0999966986860275
10.8400001525879 0.112223766509256
10.8800001144409 0.125940889207166
10.9200000762939 0.0939719427990252
10.960000038147 0.0954343909285349
11 0.0495611679202714
11.0400002002716 0.0499953368987809
11.0800001621246 0.0522129403243237
11.1200001239777 0.000892564743278264
11.1600000858307 0.00416056760696292
11.2000000476837 -0.0114628756722205
11.2400000095367 -0.0287236270662935
11.2800002098083 -0.0303494237736359
11.3200001716614 0.00908528571427872
11.3600001335144 0.00101057931345877
11.4000000953674 0.031681532710757
11.4400000572205 0.0330757205978438
11.4800000190735 0.0441293978755079
11.5199999809265 0.0546038547184303
11.5600001811981 0.0586091412894697
11.6000001430511 0.0628140335914317
11.6400001049042 0.0870766563838677
11.6800000667572 0.0898512604110598
11.7200000286102 0.0952771450046767
11.7599999904633 0.0934841837491885
11.8000001907349 0.0864083913513956
11.8400001525879 0.0849029112771754
11.8800001144409 0.0901387050616518
11.9200000762939 0.109423038511248
11.960000038147 0.117090565579414
12 0.128960814528231
12.0400002002716 0.146061296707827
12.0800001621246 0.140733140537947
12.1200001239777 0.154080838830084
12.1600000858307 0.133019293641508
12.2000000476837 0.126440674501128
12.2400000095367 0.11672144334535
12.2800002098083 0.0787771936971545
12.3200001716614 0.0256632256961259
12.3600001335144 0.0261733034384334
12.4000000953674 0.0303289529064699
12.4400000572205 -0.000944586122840076
12.4800000190735 -0.00968484800864877
12.5199999809265 -0.000273297049830694
12.5600001811981 -0.0052414452892844
12.6000001430511 0.0147797617976586
12.6400001049042 0.0146422796429012
12.6800000667572 0.0333753775917582
12.7200000286102 0.0594680698933437
12.7599999904633 0.0433029877477595
12.8000001907349 0.0758864451535606
12.8400001525879 0.132975106045697
12.8800001144409 0.152004155587209
12.9200000762939 0.193625398094346
12.960000038147 0.193292114022355
13 0.234876265889778
13.0400002002716 0.244397281672815
13.0800001621246 0.292911429532265
13.1200001239777 0.345180930434846
13.1600000858307 0.341011107825739
13.2000000476837 0.35596996360535
13.2400000095367 0.353265448015626
13.2800002098083 0.384979407064744
13.3200001716614 0.435744258130932
13.3600001335144 0.411237814555455
13.4000000953674 0.436464773778761
13.4400000572205 0.443470334011553
13.4800000190735 0.454673119389233
13.5199999809265 0.463302464134554
13.5600001811981 0.425825911697083
13.6000001430511 0.44239121962162
13.6400001049042 0.397786357296248
13.6800000667572 0.386133315776477
13.7200000286102 0.346914837932545
13.7599999904633 0.33891940100642
13.8000001907349 0.341350930766195
13.8400001525879 0.309790035946524
13.8800001144409 0.29750845610216
13.9200000762939 0.291639044370678
13.960000038147 0.288427938829084
14 0.274808572329144
14.0400002002716 0.25634174722573
14.0800001621246 0.257866649496806
14.1200001239777 0.251856026769563
14.1600000858307 0.244557503850695
14.2000000476837 0.23708568710016
14.2400000095367 0.233878568140892
14.2800002098083 0.23902251210239
14.3200001716614 0.197539551329336
14.3600001335144 0.208320735808366
14.4000000953674 0.196750895956934
14.4400000572205 0.204102732860377
14.4800000190735 0.212694612034852
14.5199999809265 0.218067670237961
14.5600001811981 0.211442809596015
14.6000001430511 0.175364470214746
14.6400001049042 0.173946624046427
14.6800000667572 0.168364948373697
14.7200000286102 0.168050266272658
14.7599999904633 0.17601600013847
14.8000001907349 0.181933060884221
14.8400001525879 0.177373849603686
14.8800001144409 0.197899952538733
14.9200000762939 0.213707804393029
14.960000038147 0.271238492358582
15 0.251444165503166
15.0400002002716 0.301896236658166
15.0800001621246 0.371651485747634
15.1200001239777 0.346971511720756
15.1600000858307 0.39226866970485
15.2000000476837 0.380482063803515
15.2400000095367 0.399623485315372
15.2800002098083 0.437694281824329
15.3200001716614 0.429529490736138
15.3600001335144 0.444422087994565
15.4000000953674 0.424034117137449
15.4400000572205 0.466133375252341
15.4800000190735 0.518911727566767
15.5199999809265 0.481066425741159
15.5600001811981 0.509581586610643
15.6000001430511 0.466325235525048
15.6400001049042 0.488396030737598
15.6800000667572 0.494905862370299
15.7200000286102 0.473714082063661
15.7599999904633 0.486474816098836
15.8000001907349 0.533609134022427
15.8400001525879 0.541887694552721
15.8800001144409 0.570135719502637
15.9200000762939 0.515307093370158
15.960000038147 0.516284628827609
16 0.330152038730749
16.0400002002716 0.367435154613507
16.0800001621246 0.327668306125615
16.1200001239777 0.34750467826179
16.1600000858307 0.352935054294176
16.2000000476837 0.391503327390787
16.2400000095367 0.432523304960908
16.2800002098083 0.406801371463973
16.3200001716614 0.45341647350826
16.3600001335144 0.505841648368469
16.4000000953674 0.467727895771039
16.4400000572205 0.506102889100949
16.4800000190735 0.473141046614069
16.5199999809265 0.479638922135009
16.5600001811981 0.525803814491428
16.6000001430511 0.515817534573067
16.6400001049042 0.514998743522826
16.6800000667572 0.480735315134913
16.7200000286102 0.478211176518587
16.7599999904633 0.477970710889447
16.8000001907349 0.274700255046721
16.8400001525879 0.260458180035826
16.8800001144409 0.0784179343925006
16.9200000762939 0.0813456511950165
16.960000038147 0.0119598054522513
17 0.00689240438001953
17.0400002002716 0.0203330215755822
17.0800001621246 -0.0136044825773348
17.1200001239777 0.0546710789614252
17.1600000858307 0.0640013333894688
17.2000000476837 0.193876301265352
17.2400000095367 0.12437566054101
17.2800002098083 0.0964831749665264
17.3200001716614 0.0776572298173204
17.3600001335144 0.137382696135798
17.4000000953674 0.0842097948356671
17.4400000572205 0.0489896024131397
17.4800000190735 0.116478444927321
17.5199999809265 0.0354244869242144
17.5600001811981 0.0598586379706848
17.6000001430511 0.0075663268172757
17.6400001049042 -0.0277578415043747
17.6800000667572 -0.11066865723301
17.7200000286102 -0.0933330142122363
17.7599999904633 -0.132659020156006
17.8000001907349 -0.1190174236814
17.8400001525879 -0.135165294037518
17.8800001144409 -0.176077792953201
17.9200000762939 -0.233190318349368
17.960000038147 -0.224460237723129
18 -0.276026780436551
18.0400002002716 -0.27100310888099
18.0800001621246 -0.281742996626321
18.1200001239777 -0.261303698134667
18.1600000858307 -0.305963607402613
18.2000000476837 -0.366999641402713
18.2400000095367 -0.415860436203306
18.2800002098083 -0.381925922617564
18.3200001716614 -0.449515713244959
18.3600001335144 -0.479008649191778
18.4000000953674 -0.526065153924264
18.4400000572205 -0.555062775086706
18.4800000190735 -0.563840935419121
18.5199999809265 -0.650808802236333
18.5600001811981 -0.645287875088315
18.6000001430511 -0.674056840978301
18.6400001049042 -0.696454080488008
18.6800000667572 -0.705973488418374
18.7200000286102 -0.687588581678176
18.7599999904633 -0.651796222040873
18.8000001907349 -0.669318613447536
18.8400001525879 -0.668670633539845
18.8800001144409 -0.699659836064121
18.9200000762939 -0.684550010908187
18.960000038147 -0.774139808445955
19 -0.847483846314773
19.0400002002716 -0.789853684681743
19.0800001621246 -0.882705806092499
19.1200001239777 -0.887991217056227
19.1600000858307 -1.00670539696924
19.2000000476837 -1.02696917033419
19.2400000095367 -1.10683648923957
19.2800002098083 -1.17721374071009
19.3200001716614 -1.14842622928867
19.3600001335144 -1.26307464292016
19.4000000953674 -1.33711854363836
19.4400000572205 -1.37062088606972
19.4800000190735 -1.45597983487223
19.5199999809265 -1.51829022184806
19.5600001811981 -1.59490033597692
19.6000001430511 -1.71039401653401
19.6400001049042 -1.79332094032618
19.6800000667572 -1.90688166443455
19.7200000286102 -1.93277090809491
19.7599999904633 -1.98157763298714
19.8000001907349 -2.01415135951048
19.8400001525879 -2.03603167152783
19.8800001144409 -2.08145048839823
19.9200000762939 -2.13427038389052
19.960000038147 -2.17986351736825
20 -2.2215365356194
20.0400002002716 -2.25179126968644
20.0800001621246 -2.27575958698248
20.1200001239777 -2.29509360091733
20.1600000858307 -2.35356574807719
20.2000000476837 -2.37988602596902
20.2400000095367 -2.39586358905932
20.2800002098083 -2.39791396949491
20.3200001716614 -2.41858756009156
20.3600001335144 -2.43708588184386
20.4000000953674 -2.44596525147163
20.4400000572205 -2.47621739818731
20.4800000190735 -2.49832537918992
20.5199999809265 -2.5289021242192
20.5600001811981 -2.54101956334867
20.6000001430511 -2.51667965594035
20.6400001049042 -2.52324005282374
20.6800000667572 -2.53152789311519
20.7200000286102 -2.56837064707082
20.7599999904633 -2.58717134270261
20.8000001907349 -2.60497636494775
20.8400001525879 -2.65142507249725
20.8800001144409 -2.66978526305454
20.9200000762939 -2.72750643483382
20.960000038147 -2.75167688992973
21 -2.77789235085584
21.0400002002716 -2.81564277616127
21.0800001621246 -2.82378367018273
21.1200001239777 -2.88394419058188
21.1600000858307 -2.92664984135118
21.2000000476837 -2.96763101280468
21.2400000095367 -3.00754755581002
21.2800002098083 -3.01767951747502
21.3200001716614 -3.04959447355103
21.3600001335144 -3.06483539860865
21.4000000953674 -3.08108535471511
21.4400000572205 -3.09988038177406
21.4800000190735 -3.07938192962744
21.5199999809265 -3.07450179160106
21.5600001811981 -3.11126867081998
21.6000001430511 -3.0772164112275
21.6400001049042 -3.09095531866039
21.6800000667572 -3.10269302185509
21.7200000286102 -3.10620439626484
21.7599999904633 -3.14846644360885
21.8000001907349 -3.15143292877913
21.8400001525879 -3.17539500684569
21.8800001144409 -3.16336295867154
21.9200000762939 -3.16684817754278
21.960000038147 -3.19276571448159
22 -3.16221074430849
22.0400002002716 -3.14962689474149
22.0800001621246 -3.19963011182661
22.1200001239777 -3.15126987289506
22.1600000858307 -3.17046951271812
22.2000000476837 -3.17009791109698
22.2400000095367 -3.18408447135549
22.2800002098083 -3.20265986033878
22.3200001716614 -3.18135070467814
22.3600001335144 -3.23107140953676
22.4000000953674 -3.20049992058705
22.4400000572205 -3.205351533558
22.4800000190735 -3.24689020341992
22.5199999809265 -3.24618794564699
22.5600001811981 -3.31296933309177
22.6000001430511 -3.32238731323191
22.6400001049042 -3.32085445451615
22.6800000667572 -3.35386363807835
22.7200000286102 -3.33440467245594
22.7599999904633 -3.27004764585176
22.8000001907349 -3.31705378775831
22.8400001525879 -3.32534795773888
22.8800001144409 -3.35238068703109
22.9200000762939 -3.30895890740336
22.960000038147 -3.30847977971329
23 -3.31484166481725
23.0400002002716 -3.31039253919606
23.0800001621246 -3.30921887967748
23.1200001239777 -3.31496718737862
23.1600000858307 -3.30344205879509
23.2000000476837 -3.3262146420195
23.2400000095367 -3.27484787419676
23.2800002098083 -3.26676867533237
23.3200001716614 -3.27966046112407
23.3600001335144 -3.25184571148178
23.4000000953674 -3.25812513888116
23.4400000572205 -3.23622625108198
23.4800000190735 -3.20261363734703
23.5199999809265 -3.2068762297282
23.5600001811981 -3.17399343737053
23.6000001430511 -3.17085537683501
23.6400001049042 -3.09102151167569
23.6800000667572 -3.07990074763016
23.7200000286102 -3.09147149598831
23.7599999904633 -3.04601365509262
23.8000001907349 -3.00431183095612
23.8400001525879 -2.96059977383607
23.8800001144409 -2.92070586405343
23.9200000762939 -2.88615294547282
23.960000038147 -2.84515825858977
24 -2.81476280561682
24.0400002002716 -2.76147381586448
24.0800001621246 -2.74818468608631
24.1200001239777 -2.71746692050952
24.1600000858307 -2.6698666124587
24.2000000476837 -2.63814799391607
24.2400000095367 -2.5860266896065
24.2800002098083 -2.57211902024349
24.3200001716614 -2.51950719128724
24.3600001335144 -2.58940700535638
24.4000000953674 -2.59120146933747
24.4400000572205 -2.56571001634112
24.4800000190735 -2.55655367996439
24.5199999809265 -2.53298952777565
24.5600001811981 -2.50554421944333
24.6000001430511 -2.4825804451454
24.6400001049042 -2.43864632084909
24.6800000667572 -2.43142804747703
24.7200000286102 -2.41183840668535
24.7599999904633 -2.39964292712566
24.8000001907349 -2.39825341311332
24.8400001525879 -2.46718946023636
24.8800001144409 -2.46491657907907
24.9200000762939 -2.42284496173631
24.960000038147 -2.40576869199577
25 -2.39544185452559
25.0400002002716 -2.35186107948322
25.0800001621246 -2.34811620717847
25.1200001239777 -2.32657782026529
25.1600000858307 -2.31207744852484
25.2000000476837 -2.31167722606531
25.2400000095367 -2.29977998642911
25.2800002098083 -2.29588080764236
25.3200001716614 -2.27847447474836
25.3600001335144 -2.2741632278136
25.4000000953674 -2.27419167934333
25.4400000572205 -2.26377251310878
25.4800000190735 -2.2628438239207
25.5199999809265 -2.2697114643276
25.5600001811981 -2.24988932495678
25.6000001430511 -2.23938165513604
25.6400001049042 -2.23462903161494
25.6800000667572 -2.22760710015989
25.7200000286102 -2.22802254104448
25.7599999904633 -2.2445938752599
25.8000001907349 -2.24200788333704
25.8400001525879 -2.23557840532604
25.8800001144409 -2.22249489206524
25.9200000762939 -2.21677273716648
25.960000038147 -2.22533793896382
26 -2.23404312104643
26.0400002002716 -2.2593390499648
26.0800001621246 -2.27615922686597
26.1200001239777 -2.28648302737267
26.1600000858307 -2.30341907296632
26.2000000476837 -2.34453838697522
26.2400000095367 -2.4281880745939
26.2800002098083 -2.43865109641466
26.3200001716614 -2.48650210017652
26.3600001335144 -2.44970634277806
26.4000000953674 -2.45475809371296
26.4400000572205 -2.46755026655577
26.4800000190735 -2.47681788892842
26.5199999809265 -2.48396338592078
26.5600001811981 -2.49725018920557
26.6000001430511 -2.50490773581862
26.6400001049042 -2.49390751057243
26.6800000667572 -2.50288433766719
26.7200000286102 -2.4513887598724
26.7599999904633 -2.45362101184234
26.8000001907349 -2.47598063118207
26.8400001525879 -2.50348544767923
26.8800001144409 -2.52222594099045
26.9200000762939 -2.53797841460131
26.960000038147 -2.54510518718761
27 -2.57157083690572
27.0400002002716 -2.62029681031889
27.0800001621246 -2.62464812958259
27.1200001239777 -2.64189720240698
27.1600000858307 -2.6526236577082
27.2000000476837 -2.67282090842203
27.2400000095367 -2.65491264456947
27.2800002098083 -2.6574144228119
27.3200001716614 -2.70728465858581
27.3600001335144 -2.71422010483035
27.4000000953674 -2.73367314242464
27.4400000572205 -2.75768024577716
27.4800000190735 -2.77047956984109
27.5199999809265 -2.74682595589388
27.5600001811981 -2.75181138442918
27.6000001430511 -2.76049805157132
27.6400001049042 -2.73587135978063
27.6800000667572 -2.71002586255824
27.7200000286102 -2.73210615906794
27.7599999904633 -2.73611180861448
27.8000001907349 -2.71526808339848
27.8400001525879 -2.71963398659206
27.8800001144409 -2.74113967170981
27.9200000762939 -2.80385364035568
27.960000038147 -2.81192461418436
28 -2.82151117847386
28.0400002002716 -2.8662003204457
28.0800001621246 -2.86789730066448
28.1200001239777 -2.89002520328185
28.1600000858307 -2.8973569880417
28.2000000476837 -2.88877312670165
28.2400000095367 -2.92543677157336
28.2800002098083 -2.87852247512787
28.3200001716614 -2.87398743011126
28.3600001335144 -2.88075960259602
28.4000000953674 -2.83739239807695
28.4400000572205 -2.81173675389736
28.4800000190735 -2.81765386803174
28.5199999809265 -2.78147043421476
28.5600001811981 -2.78150852027185
28.6000001430511 -2.75318300056105
28.6400001049042 -2.76884247961643
28.6800000667572 -2.75838117614253
28.7200000286102 -2.76060105664844
28.7599999904633 -2.75837731917145
28.8000001907349 -2.76579487820908
28.8400001525879 -2.87378495935205
28.8800001144409 -2.8760473519881
28.9200000762939 -2.90582607072331
28.960000038147 -2.91425939096074
29 -2.93657455293848
29.0400002002716 -2.94575688980919
29.0800001621246 -2.97916976518797
29.1200001239777 -3.01724646396958
29.1600000858307 -3.02416281957772
29.2000000476837 -2.99568116818603
29.2400000095367 -3.02261579426044
29.2800002098083 -3.00669704736426
29.3200001716614 -3.05927785604151
29.3600001335144 -3.06683659375732
29.4000000953674 -3.0684727393434
29.4400000572205 -3.0380287586231
29.4800000190735 -3.0375686355351
29.5199999809265 -3.04340467587728
29.5600001811981 -3.04687883838274
29.6000001430511 -2.98890304163213
29.6400001049042 -2.9990769212831
29.6800000667572 -2.95547896174324
29.7200000286102 -2.98736340999783
29.7599999904633 -2.98522332973768
29.8000001907349 -2.99763381204247
29.8400001525879 -3.01471429135446
29.8800001144409 -3.0108347675208
29.9200000762939 -3.03038114623164
29.960000038147 -3.02784021447056
30 -3.00634249675779
30.0400002002716 -3.04254666053674
30.0800001621246 -3.03064264155103
30.1200001239777 -3.07279285756794
30.1600000858307 -3.07773873836244
30.2000000476837 -3.06462669168446
30.2400000095367 -3.12608503955543
30.2800002098083 -3.12523114745277
30.3200001716614 -3.09766465266305
30.3600001335144 -3.14790265819908
30.4000000953674 -3.26171812992882
30.4400000572205 -3.23586516421113
30.4800000190735 -3.22879795246832
30.5199999809265 -3.24008360942349
30.5600001811981 -3.23119952952039
30.6000001430511 -3.22169658525557
30.6400001049042 -3.23104687855408
30.6800000667572 -3.21673438359682
30.7200000286102 -3.20422933458226
30.7599999904633 -3.19849504325427
30.8000001907349 -3.15316621136524
30.8400001525879 -3.11502467137735
30.8800001144409 -3.09366118605176
30.9200000762939 -3.12221302091524
30.960000038147 -3.10746115952173
31 -3.11950373313499
31.0400002002716 -3.11420236875809
31.0800001621246 -3.1462762997198
31.1200001239777 -3.12190182512929
31.1600000858307 -3.1413995377537
31.2000000476837 -3.07391465867788
31.2400000095367 -3.07354459741406
31.2800002098083 -3.01698151800646
31.3200001716614 -2.99326740638999
31.3600001335144 -2.97976676209942
31.4000000953674 -2.9505046102769
31.4400000572205 -2.92490161649248
31.4800000190735 -2.90618219454163
31.5199999809265 -2.86002318940931
31.5600001811981 -2.90033390738575
31.6000001430511 -2.96704108478009
31.6400001049042 -3.07472301259415
31.6800000667572 -3.06883054768239
31.7200000286102 -3.00946945013467
31.7599999904633 -2.97002807907616
31.8000001907349 -2.95274972630053
31.8400001525879 -2.85992727628906
31.8800001144409 -2.78052903604534
31.9200000762939 -2.70911716382429
31.960000038147 -2.64946456537842
32 -2.51962526279854
32.0400002002716 -2.46365620557724
32.0800001621246 -2.47920235855583
32.1200001239777 -2.26170462784135
32.1600000858307 -2.20704535742009
32.2000000476837 -2.24299002479363
32.2400000095367 -2.20396445839646
32.2800002098083 -2.09618641617809
32.3200001716614 -1.9770392294933
32.3600001335144 -1.97561668216405
32.4000000953674 -1.9571304304276
32.4400000572205 -1.92008592337765
32.4800000190735 -1.84891699494613
32.5199999809265 -1.82512187252416
32.5600001811981 -1.71903130640654
32.6000001430511 -1.62936754554125
32.6400001049042 -1.59266188297517
32.6800000667572 -1.62128089127122
32.7200000286102 -1.60972519106358
32.7599999904633 -1.58332909982517
32.8000001907349 -1.64517557763296
32.8400001525879 -1.55431943700714
32.8800001144409 -1.09901995845896
32.9200000762939 -1.01675410008137
32.960000038147 -0.933309716740302
33 -0.898074741824657
33.0400002002716 -0.829436605363087
33.0800001621246 -0.783736095627514
33.1200001239777 -0.742654390037021
33.1600000858307 -0.70801887724927
33.2000000476837 -0.601050299094892
33.2400000095367 -0.522436921926209
33.2800002098083 -0.381034988096312
33.3200001716614 -0.350820764074171
33.3600001335144 -0.683945456779577
33.4000000953674 -0.879792398929308
33.4400000572205 -0.792223194796082
33.4800000190735 -0.760605237564636
33.5199999809265 -0.691684602925099
33.5600001811981 -0.668958795737696
33.6000001430511 -0.575949471987297
33.6400001049042 -0.661034602890749
33.6800000667572 -0.70473204934717
33.7200000286102 -0.685390736011053
33.7599999904633 -0.675678718792236
33.8000001907349 -0.615669240678282
33.8400001525879 -0.602849580566696
33.8800001144409 -0.565048967036272
33.9200000762939 -0.577432030661821
33.960000038147 -0.555647910091302
34 -0.634970574233768
34.0400002002716 -0.680970577002872
34.0800001621246 -0.736431421807733
34.1200001239777 -0.775654138014481
34.1600000858307 -0.800359213922397
};
\path [draw=black, fill=black] (axis cs:15,0.12)
--(axis cs:15.5,-0.1)
--(axis cs:15.125,-0.1)
--(axis cs:15.125,-2)
--(axis cs:14.875,-2)
--(axis cs:14.875,-0.1)
--(axis cs:14.5,-0.1)
--cycle;

\path [draw=black, fill=black] (axis cs:17,2.97)
--(axis cs:17.5,2.75)
--(axis cs:17.125,2.75)
--(axis cs:17.125,2.2)
--(axis cs:16.875,2.2)
--(axis cs:16.875,2.75)
--(axis cs:16.5,2.75)
--cycle;

\addplot [line width=2.4000000000000004pt, color1, dashed, forget plot]
table {%
0 -0.0417787396782297
0.0400002002716064 -0.0234894145660781
0.0800001621246338 -0.00316769069747583
0.120000123977661 0.0171410027993888
0.160000085830688 0.0353008948896494
0.200000047683716 0.0519132876052167
0.240000009536743 0.047620933155939
0.28000020980835 0.0354367334666362
0.320000171661377 0.0243375702749589
0.360000133514404 0.0147874162005871
0.400000095367432 0.00626683768990029
0.440000057220459 -0.00117534729387254
0.480000019073486 -0.00759592921282741
0.519999980926514 -0.0134284232414302
0.56000018119812 -0.0182270835009956
0.600000143051147 -0.022438744343384
0.640000104904175 -0.026070568590947
0.680000066757202 -0.0291511821203337
0.720000028610229 -0.0319400999575896
0.759999990463257 -0.0338169274746991
0.800000190734863 -0.0352436446804821
0.840000152587891 -0.0363389452063536
0.880000114440918 -0.0371832000288373
0.920000076293945 -0.0378755526499222
0.960000038146973 -0.0379489570409571
1 -0.0377584947677796
1.04000020027161 -0.0376390233325493
1.08000016212463 -0.0374952551139339
1.12000012397766 -0.0372917578207389
1.16000008583069 -0.0369938784342557
1.20000004768372 -0.0366344301219262
1.24000000953674 -0.0360929457393068
1.28000020980835 -0.0354371027776929
1.32000017166138 -0.0345109548601747
1.3600001335144 -0.0334581229483179
1.40000009536743 -0.0324192791695158
1.44000005722046 -0.030980956112695
1.48000001907349 -0.0293790907626932
1.51999998092651 -0.0272088269687246
1.56000018119812 -0.024882385056818
1.60000014305115 -0.0226886086533483
1.64000010490417 -0.0190998139744281
1.6800000667572 -0.0151483238709784
1.72000002861023 -0.0110182941941803
1.75999999046326 -0.00714158725049169
1.80000019073486 -0.00356854737946884
1.84000015258789 0.00048544445960752
1.88000011444092 0.00451394310810775
1.92000007629395 0.00841177825825727
1.96000003814697 0.0119801034553173
2 0.01527436559098
2.04000020027161 0.0187924828262535
2.08000016212463 0.0222216665186688
2.12000012397766 0.0254353989836296
2.16000008583069 0.0283560527972835
2.20000004768372 0.0310911488127918
2.24000000953674 0.0341602267221054
2.28000020980835 0.037264868811495
2.32000017166138 0.0403389177592532
2.3600001335144 0.0432238920390264
2.40000009536743 0.0463886045314703
2.44000005722046 0.0496490639713337
2.48000001907349 0.0526877993519101
2.51999998092651 0.0564129153244892
2.56000018119812 0.0602846008974836
2.60000014305115 0.0638569429020855
2.64000010490417 0.0677044802291149
2.6800000667572 0.0715338351926627
2.72000002861023 0.0756254784338014
2.75999999046326 0.079603058963355
2.80000019073486 0.0832873003763861
2.84000015258789 0.0874524012518538
2.88000011444092 0.0916925884473081
2.92000007629395 0.0960618042740521
2.96000003814697 0.100200412245092
3 0.104271661953581
3.04000020027161 0.108167817374537
3.08000016212463 0.111783227709938
3.12000012397766 0.115684704289017
3.16000008583069 0.119522929461469
3.20000004768372 0.12315851979343
3.24000000953674 0.125141227311931
3.28000020980835 0.124029761602394
3.32000017166138 0.122598006925633
3.3600001335144 0.121279407271869
3.40000009536743 0.119340532654102
3.44000005722046 0.117189462631772
3.48000001907349 0.127221066263764
3.51999998092651 0.113301125155132
3.56000018119812 0.11144265952714
3.60000014305115 0.109868516394461
3.64000010490417 0.107719769339863
3.6800000667572 0.106500928672304
3.72000002861023 0.105042468454708
3.75999999046326 0.103877521716361
3.80000019073486 0.102937771611001
3.84000015258789 0.102171915711101
3.88000011444092 0.101579580501347
3.92000007629395 0.101226654302596
3.96000003814697 0.101113989199034
4 0.101616293426719
4.04000020027161 0.102418837022119
4.08000016212463 0.103231067324208
4.12000012397766 0.104120570963666
4.16000008583069 0.105118444100313
4.20000004768372 0.1061130722062
4.24000000953674 0.107099502125028
4.28000020980835 0.10808039216406
4.32000017166138 0.10905367990172
4.3600001335144 0.110000619620419
4.40000009536743 0.110710271456434
4.44000005722046 0.121753809426045
4.48000001907349 0.111943376392363
4.51999998092651 0.112505432724912
4.56000018119812 0.113037328530031
4.60000014305115 0.113606122874255
4.64000010490417 0.114212016721308
4.6800000667572 0.114937984492269
4.72000002861023 0.115942378212514
4.75999999046326 0.117000394895368
4.80000019073486 0.118307752468225
4.84000015258789 0.119605674340697
4.88000011444092 0.120902193469763
4.92000007629395 0.122538104963462
4.96000003814697 0.124289500251855
5 0.126108929838909
5.04000020027161 0.127843713065531
5.08000016212463 0.129489474225519
5.12000012397766 0.131066693075334
5.16000008583069 0.132579986193214
5.20000004768372 0.13395889045711
5.24000000953674 0.135024695265243
5.28000020980835 0.136097275109658
5.32000017166138 0.137301891448054
5.3600001335144 0.138522282896608
5.40000009536743 0.140322133246502
5.44000005722046 0.142360937875605
5.48000001907349 0.144572170213079
5.51999998092651 0.146722021325224
5.56000018119812 0.148727839038772
5.60000014305115 0.150742032117834
5.64000010490417 0.152816424101772
5.6800000667572 0.154949766866354
5.72000002861023 0.157208720762795
5.75999999046326 0.159367421435761
5.80000019073486 0.161433522444049
5.84000015258789 0.163369160557194
5.88000011444092 0.164821837764834
5.92000007629395 0.166124342488959
5.96000003814697 0.167369968847014
6 0.168116567065514
6.04000020027161 0.168625163666823
6.08000016212463 0.168751701543232
6.12000012397766 0.168478845217233
6.16000008583069 0.168307856437836
6.20000004768372 0.168231728600646
6.24000000953674 0.168188884458664
6.28000020980835 0.167959200938365
6.32000017166138 0.167825517347502
6.3600001335144 0.167778793190581
6.40000009536743 0.167686063493615
6.44000005722046 0.167577660315265
6.48000001907349 0.167414043164538
6.51999998092651 0.167352124229287
6.56000018119812 0.16730211208751
6.60000014305115 0.16726784797416
6.64000010490417 0.166864112795472
6.6800000667572 0.165765939475376
6.72000002861023 0.164868805072815
6.75999999046326 0.164104203344044
6.80000019073486 0.163175634034676
6.84000015258789 0.162404237764292
6.88000011444092 0.161376542242743
6.92000007629395 0.160366243110981
6.96000003814697 0.159332653747468
7 0.158396458797799
7.04000020027161 0.157743728819583
7.08000016212463 0.157670882961741
7.12000012397766 0.157998181690912
7.16000008583069 0.158570906013077
7.20000004768372 0.159516228425704
7.24000000953674 0.160643503265796
7.28000020980835 0.162220036107218
7.32000017166138 0.163808533736528
7.3600001335144 0.165476506757328
7.40000009536743 0.167105373620904
7.44000005722046 0.168665782559703
7.48000001907349 0.169805000798356
7.51999998092651 0.170751800959562
7.56000018119812 0.171827140200371
7.60000014305115 0.172935017713235
7.64000010490417 0.174222904089218
7.6800000667572 0.175631408882665
7.72000002861023 0.176995424582933
7.75999999046326 0.178355224513547
7.80000019073486 0.179779153121134
7.84000015258789 0.181321531430517
7.88000011444092 0.183416695002654
7.92000007629395 0.185557119325337
7.96000003814697 0.188003976170535
8 0.190525844739645
8.04000020027161 0.193066122560267
8.08000016212463 0.195395145382109
8.12000012397766 0.197570964148006
8.16000008583069 0.199726411527884
8.20000004768372 0.201854983180803
8.24000000953674 0.203986200810348
8.28000020980835 0.206255365030284
8.32000017166138 0.208366027970262
8.3600001335144 0.210387484551534
8.40000009536743 0.212325539830309
8.44000005722046 0.214381526209379
8.48000001907349 0.216517628578634
8.51999998092651 0.218616770703851
8.56000018119812 0.221095809107547
8.60000014305115 0.22355554916839
8.64000010490417 0.226048185659467
8.6800000667572 0.228557681507201
8.72000002861023 0.230702647590885
8.75999999046326 0.23243481748732
8.80000019073486 0.234081004254391
8.84000015258789 0.23558011349421
8.88000011444092 0.236788258690557
8.92000007629395 0.237924488952279
8.96000003814697 0.238731923484593
9 0.239499077388328
9.04000020027161 0.239734043719992
9.08000016212463 0.239799866402045
9.12000012397766 0.239988545359118
9.16000008583069 0.240262057773767
9.20000004768372 0.27249044538246
9.24000000953674 0.289991062240143
9.28000020980835 0.306487324447838
9.32000017166138 0.322004096110039
9.3600001335144 0.336267088026604
9.40000009536743 0.350096351134862
9.44000005722046 0.363051527791928
9.48000001907349 0.375882951802808
9.51999998092651 0.388226280791758
9.56000018119812 0.400754764070436
9.60000014305115 0.414008879288613
9.64000010490417 0.427690903121654
9.6800000667572 0.441974399828963
9.72000002861023 0.457064922618433
9.75999999046326 0.473005982408575
9.80000019073486 0.490280442656957
9.84000015258789 0.509064261184364
9.88000011444092 0.529485993180569
9.92000007629395 0.552115731425192
9.96000003814697 0.576684401166912
10 0.60325270940883
10.0400002002716 0.649831130856617
10.0800001621246 0.663265913770234
10.1200001239777 0.697282600223065
10.1600000858307 0.734081533815605
10.2000000476837 0.773597182126251
10.2400000095367 0.815988606103267
10.2800002098083 0.861317943809147
10.3200001716614 0.909598423165704
10.3600001335144 0.961256891333492
10.4000000953674 1.01645285377695
10.4400000572205 1.07552925798337
10.4800000190735 1.13846920296723
10.5199999809265 1.20554806746357
10.5600001811981 1.27517714672044
10.6000001430511 1.35045689811077
10.6400001049042 1.42734502594035
10.6800000667572 1.5065466494216
10.7200000286102 1.58753967702282
10.7599999904633 1.66973949662613
10.8000001907349 1.75282800954276
10.8400001525879 1.8364926698402
10.8800001144409 1.92004994454084
10.9200000762939 2.00314700538797
10.960000038147 2.08522734424601
11 2.16606889393365
11.0400002002716 2.2451298367563
11.0800001621246 2.32178394087486
11.1200001239777 2.39562664107728
11.1600000858307 2.46621406025363
11.2000000476837 2.55660417158619
11.2400000095367 2.59551971083539
11.2800002098083 2.65369595965046
11.3200001716614 2.70647836614852
11.3600001335144 2.75369088275477
11.4000000953674 2.79472028564417
11.4400000572205 2.82932207898036
11.4800000190735 2.85686386213842
11.5199999809265 2.87708908193203
11.5600001811981 2.88963753824534
11.6000001430511 2.89418151611971
11.6400001049042 2.89020334658292
11.6800000667572 2.87707720121421
11.7200000286102 2.85438721110032
11.7599999904633 2.82163607458536
11.8000001907349 2.81946235112765
11.8400001525879 2.82116591494254
11.8800001144409 2.82292899929882
11.9200000762939 2.82487931454895
11.960000038147 2.82687799052953
12 2.82901084490146
12.0400002002716 2.83118772363196
12.0800001621246 2.83336431404865
12.1200001239777 2.83554041684541
12.1600000858307 2.83770421958513
12.2000000476837 2.83941622049006
12.2400000095367 2.84091721275971
12.2800002098083 2.84239674481811
12.3200001716614 2.84384764503737
12.3600001335144 2.84521260282666
12.4000000953674 2.84626790164339
12.4400000572205 2.84747487840261
12.4800000190735 2.84850920790293
12.5199999809265 2.8494915962145
12.5600001811981 2.85025217788525
12.6000001430511 2.85099510747219
12.6400001049042 2.85197948850158
12.6800000667572 2.85318728562048
12.7200000286102 2.8545628921675
12.7599999904633 2.85601886621953
12.8000001907349 2.8575184955966
12.8400001525879 2.85890121794457
12.8800001144409 2.85989187934172
12.9200000762939 2.86090450652759
12.960000038147 2.86173581345884
13 2.86263661943759
13.0400002002716 2.86365100062622
13.0800001621246 2.86510819226026
13.1200001239777 2.86678465209858
13.1600000858307 2.86855735242891
13.2000000476837 2.87037406523006
13.2400000095367 2.87229955612768
13.2800002098083 2.87434111080212
13.3200001716614 2.87639503129837
13.3600001335144 2.8784233679331
13.4000000953674 2.88032740506776
13.4400000572205 2.88234655692363
13.4800000190735 2.88488259176671
13.5199999809265 2.88750131518647
13.5600001811981 2.89008614787922
13.6000001430511 2.89265031655973
13.6400001049042 2.89572143181652
13.6800000667572 2.89886507095544
13.7200000286102 2.90189527304969
13.7599999904633 2.90485947248228
13.8000001907349 2.90775999261128
13.8400001525879 2.91064154555121
13.8800001144409 2.9135789669666
13.9200000762939 2.91669274001231
13.960000038147 2.92104549711275
14 2.92549185329583
14.0400002002716 2.9310390710285
14.0800001621246 2.93652863347197
14.1200001239777 2.94194056627251
14.1600000858307 2.9481758322133
14.2000000476837 2.95432151438774
14.2400000095367 2.96028160901867
14.2800002098083 2.96602889185049
14.3200001716614 2.97117268010714
14.3600001335144 2.97577392363073
14.4000000953674 2.98027180157868
14.4400000572205 2.98479270070921
14.4800000190735 2.9891602958126
14.5199999809265 2.99374931884239
14.5600001811981 2.99818193068313
14.6000001430511 3.00241745373083
14.6400001049042 3.0061846823938
14.6800000667572 3.00978859099482
14.7200000286102 3.01305770668934
14.7599999904633 3.01603644197872
14.8000001907349 3.01867789969351
14.8400001525879 3.02114224926212
14.8800001144409 3.02332188552257
14.9200000762939 3.02522324269404
14.960000038147 3.02714463113932
15 3.02902239115438
15.2000000476837 3.03581285876219
15.2400000095367 3.03682919669528
15.2800002098083 3.03647918557522
15.3200001716614 3.03588694612308
15.3600001335144 3.03560266712244
15.4000000953674 3.03519843673356
15.4400000572205 3.03466744154129
15.4800000190735 3.03419394789522
15.5199999809265 3.03361529292566
15.5600001811981 3.03314342764902
15.6000001430511 3.03218547747298
15.6400001049042 3.03154820456549
15.6800000667572 3.03083075783232
15.7200000286102 3.03019726795694
15.7599999904633 3.02973953767429
15.8000001907349 3.02993220761341
15.8400001525879 3.03059178639963
15.8800001144409 3.032015592186
15.9200000762939 3.03396743956902
15.960000038147 3.0362630955797
16 3.03912061617393
16.0400002002716 3.0419153660061
16.0800001621246 3.04472339185263
16.1200001239777 3.04807617225551
16.1600000858307 3.05143227117412
16.2000000476837 3.05530167094487
16.2400000095367 3.05923011438077
16.2800002098083 3.06321889489764
16.3200001716614 3.06716322145259
16.3600001335144 3.07085821029738
16.4000000953674 3.07493434208533
16.4400000572205 3.07912470784107
16.4800000190735 3.08352294577925
16.5199999809265 3.08816542401901
16.5600001811981 3.0928191337742
16.6000001430511 3.09779896709111
16.6400001049042 3.10257366390395
16.6800000667572 3.10763294493057
16.7200000286102 3.11255837134667
16.7599999904633 3.11724023618146
16.8000001907349 3.12234074492009
16.8400001525879 3.12758011258413
16.8800001144409 3.13321495302348
16.9200000762939 3.13822630652863
16.960000038147 3.13753400168899
17 3.13608046020745
17.0400002002716 3.13498617028145
17.0800001621246 3.13327969580245
17.1200001239777 3.12992708185703
17.1600000858307 3.12667921847227
17.2000000476837 3.12170348773769
17.2400000095367 3.11752597903024
17.2800002098083 3.11183076637403
17.3200001716614 3.10599601329874
17.3600001335144 3.09492581492654
17.4000000953674 3.08307202517713
17.4400000572205 3.07334654246205
17.4800000190735 3.0628729101452
17.5199999809265 3.05207030549858
17.5600001811981 3.04111577457144
17.6000001430511 3.02891162055312
17.6400001049042 3.01603343285441
17.6800000667572 2.99881223059587
17.7200000286102 2.98320854652992
};
\addplot [line width=2.4000000000000004pt, color3, dashed, forget plot]
table {%
0.800000190734863 0.566310207694819
0.840000152587891 0.563927639889975
0.880000114440918 0.561796117788519
0.920000076293945 0.559816497888462
0.960000038146973 0.558455826218455
1 0.55735902121266
1.04000020027161 0.5561912176962
1.08000016212463 0.555047718635843
1.12000012397766 0.553963948650066
1.16000008583069 0.552974560757577
1.20000004768372 0.552046741790935
1.24000000953674 0.551300958894582
1.28000020980835 0.550669526904506
1.32000017166138 0.550308407543052
1.3600001335144 0.550073972175936
1.40000009536743 0.549825548675766
1.44000005722046 0.549976604453615
1.48000001907349 0.550291202524645
1.51999998092651 0.551174199039641
1.56000018119812 0.552213365999857
1.60000014305115 0.553119875124355
1.64000010490417 0.555421402524303
1.6800000667572 0.558085625348781
1.72000002861023 0.560928387746607
1.75999999046326 0.563517827411323
1.80000019073486 0.565803592330656
1.84000015258789 0.56857031689076
1.88000011444092 0.571311548260288
1.92000007629395 0.573922116131465
1.96000003814697 0.576203174049553
2 0.578210168906244
2.04000020027161 0.580441011189827
2.08000016212463 0.58258292760327
2.12000012397766 0.584509392789259
2.16000008583069 0.58614277932394
2.20000004768372 0.587590608060477
2.24000000953674 0.589372418690818
2.28000020980835 0.591189785828517
2.32000017166138 0.592976567497303
2.3600001335144 0.594574274498105
2.40000009536743 0.596451719711576
2.44000005722046 0.598424911872467
2.48000001907349 0.600176379974072
2.51999998092651 0.602614228667679
2.56000018119812 0.605198639288983
2.60000014305115 0.607483714014613
2.64000010490417 0.61004398406267
2.6800000667572 0.612586071747246
2.72000002861023 0.615390447709412
2.75999999046326 0.618080760959994
2.80000019073486 0.620477727421334
2.84000015258789 0.62335556101783
2.88000011444092 0.626308480934312
2.92000007629395 0.629390429482084
2.96000003814697 0.632241770174151
3 0.635025752603668
3.04000020027161 0.637634633072934
3.08000016212463 0.639962776129363
3.12000012397766 0.64257698542947
3.16000008583069 0.64512794332295
3.20000004768372 0.647476266375938
3.24000000953674 0.648171706615468
3.28000020980835 0.64577296595424
3.32000017166138 0.643053943998507
3.3600001335144 0.640448077065771
3.40000009536743 0.637221935169032
3.44000005722046 0.633783597867729
3.48000001907349 0.642527934220749
3.51999998092651 0.627320725833145
3.56000018119812 0.624174985253463
3.60000014305115 0.621313574841812
3.64000010490417 0.617877560508242
3.6800000667572 0.615371452561711
3.72000002861023 0.612625725065143
3.75999999046326 0.610173511047823
3.80000019073486 0.607946485990773
3.84000015258789 0.605893362811901
3.88000011444092 0.604013760323175
3.92000007629395 0.602373566845452
3.96000003814697 0.600973634462917
4 0.600188671411631
4.04000020027161 0.59970394005534
4.08000016212463 0.599228903078457
4.12000012397766 0.598831139438943
4.16000008583069 0.598541745296617
4.20000004768372 0.598249106123532
4.24000000953674 0.597948268763388
4.28000020980835 0.59764188385073
4.32000017166138 0.597327904309418
4.3600001335144 0.596987576749144
4.40000009536743 0.596409961306188
4.44000005722046 0.606166231996826
4.48000001907349 0.595068531684172
4.51999998092651 0.594343320737749
4.56000018119812 0.593587941591178
4.60000014305115 0.592869468656429
4.64000010490417 0.59218809522451
4.6800000667572 0.591626795716499
4.72000002861023 0.591343922157773
4.75999999046326 0.591114671561654
4.80000019073486 0.591134754182821
4.84000015258789 0.591145408776321
4.88000011444092 0.591154660626415
4.92000007629395 0.591503304841141
4.96000003814697 0.591967432850562
5 0.592499595158644
5.04000020027161 0.592947103433576
5.08000016212463 0.593305597314591
5.12000012397766 0.593595548885434
5.16000008583069 0.593821574724342
5.20000004768372 0.593913211709266
5.24000000953674 0.593691749238427
5.28000020980835 0.593477054131152
5.32000017166138 0.593394403190575
5.3600001335144 0.593327527360158
5.40000009536743 0.593840110431079
5.44000005722046 0.59459164778121
5.48000001907349 0.595515612839712
5.51999998092651 0.596378196672885
5.56000018119812 0.597096739434742
5.60000014305115 0.597823665234832
5.64000010490417 0.598610789939798
5.6800000667572 0.599456865425409
5.72000002861023 0.600428552042877
5.75999999046326 0.601299985436871
5.80000019073486 0.602078811493469
5.84000015258789 0.602727182327641
5.88000011444092 0.602892592256309
5.92000007629395 0.602907829701462
5.96000003814697 0.602866188780546
6 0.602325519720073
6.04000020027161 0.601546841369692
6.08000016212463 0.600386111967129
6.12000012397766 0.598825988362157
6.16000008583069 0.597367732303788
6.20000004768372 0.596004337187626
6.24000000953674 0.594674225766672
6.28000020980835 0.593157267294683
6.32000017166138 0.591736316424847
6.3600001335144 0.590402324988954
6.40000009536743 0.589022328013016
6.44000005722046 0.587626657555694
6.48000001907349 0.586175773125995
6.51999998092651 0.584826586911772
6.56000018119812 0.583489299818304
6.60000014305115 0.582167768425982
6.64000010490417 0.580476765968322
6.6800000667572 0.578091325369254
6.72000002861023 0.575906923687721
6.75999999046326 0.573855054679978
6.80000019073486 0.571639210418919
6.84000015258789 0.569580546869564
6.88000011444092 0.567265584069042
6.92000007629395 0.564968017658308
6.96000003814697 0.562647161015823
7 0.560423698787181
7.04000020027161 0.558483693857275
7.08000016212463 0.557123580720461
7.12000012397766 0.55616361217066
7.16000008583069 0.555449069213853
7.20000004768372 0.555107124347508
7.24000000953674 0.554947131908628
7.28000020980835 0.555236389798359
7.32000017166138 0.555537620148697
7.3600001335144 0.555918325890525
7.40000009536743 0.556259925475129
7.44000005722046 0.556533067134955
7.48000001907349 0.556385018094637
7.51999998092651 0.556044550976871
7.56000018119812 0.555832615265989
7.60000014305115 0.555653225499881
7.64000010490417 0.555653844596892
7.6800000667572 0.555775082111367
7.72000002861023 0.555851830532663
7.75999999046326 0.555924363184304
7.80000019073486 0.556061016840202
7.84000015258789 0.556316127870613
7.88000011444092 0.557124024163776
7.92000007629395 0.557977181207488
7.96000003814697 0.559136770773713
8 0.560371372063851
8.04000020027161 0.561624374932784
8.08000016212463 0.562666130475653
8.12000012397766 0.563554681962578
8.16000008583069 0.564422862063484
8.20000004768372 0.565264166437431
8.24000000953674 0.566108116788004
8.28000020980835 0.56709000605625
8.32000017166138 0.567913401717255
8.3600001335144 0.568647591019555
8.40000009536743 0.569298379019357
8.44000005722046 0.570067098119456
8.48000001907349 0.570915933209738
8.51999998092651 0.571727808055983
8.56000018119812 0.572919571507989
8.60000014305115 0.57409204428986
8.64000010490417 0.575297413501965
8.6800000667572 0.576519642070727
8.72000002861023 0.577377340875438
8.75999999046326 0.577822243492902
8.80000019073486 0.578181155308282
8.84000015258789 0.578392997269129
8.88000011444092 0.578313875186503
8.92000007629395 0.578162838169254
8.96000003814697 0.577683005422596
9 0.577162892047358
9.04000020027161 0.576110583427332
9.08000016212463 0.574889138830413
9.12000012397766 0.573790550508514
9.16000008583069 0.572776795644191
9.20000004768372 0.57182208536616
9.24000000953674 0.570855994111366
9.28000020980835 0.569824442798667
9.32000017166138 0.568593387227091
9.3600001335144 0.566729314620851
9.40000009536743 0.564893159455378
9.44000005722046 0.56248544762589
9.48000001907349 0.560097396475496
9.51999998092651 0.557205547166552
9.56000018119812 0.554323939860843
9.60000014305115 0.551834118657573
9.64000010490417 0.549279152598448
9.6800000667572 0.546673489480972
9.72000002861023 0.544063566051141
9.75999999046326 0.54133377676557
9.80000019073486 0.53880774691803
9.84000015258789 0.536502553116121
9.88000011444092 0.534387519592495
9.92000007629395 0.532873622664875
9.96000003814697 0.53153267112004
10 0.530266255499196
10.0400002002716 0.546925549187508
10.0800001621246 0.528198039854712
10.1200001239777 0.527649981150819
10.1600000858307 0.527322600214427
10.2000000476837 0.526991248162036
10.2400000095367 0.526655869480012
10.2800002098083 0.526219203610849
10.3200001716614 0.525535906794667
10.3600001335144 0.524873446159609
10.4000000953674 0.524232210708219
10.4400000572205 0.52379603146589
10.4800000190735 0.523388890985208
10.5199999809265 0.523228539236435
10.5600001811981 0.522074161923519
10.6000001430511 0.523476872174834
10.6400001049042 0.523843772231425
10.6800000667572 0.524329791026581
10.7200000286102 0.524862646749479
10.7599999904633 0.525307537003112
10.8000001907349 0.525795678861039
10.8400001525879 0.526465323014089
10.8800001144409 0.5270822559427
10.9200000762939 0.527743459111048
10.960000038147 0.528342234104429
11 0.529106323462411
11.0400002002716 0.529943259380987
11.0800001621246 0.530677554474344
11.1200001239777 0.531353981030273
11.1600000858307 0.531978471659719
11.2000000476837 0.556058809265848
11.2400000095367 0.532767539329965
11.2800002098083 0.533289424039148
11.3200001716614 0.533420410755758
11.3600001335144 0.53343390396562
11.4000000953674 0.533166489564569
11.4400000572205 0.53282348143712
11.4800000190735 0.532222288679252
11.5199999809265 0.531556167825498
11.5600001811981 0.530914673708877
11.6000001430511 0.530420060561538
11.6400001049042 0.53000436711532
11.6800000667572 0.529491574670342
11.7200000286102 0.528915624034234
11.7599999904633 0.528229023271985
11.8000001907349 0.527339178792165
11.8400001525879 0.526501261177906
11.8800001144409 0.525722864105031
11.9200000762939 0.525131697926006
11.960000038147 0.524588892477429
12 0.524180265420195
12.0400002002716 0.523815647573118
12.0800001621246 0.523450756560656
12.1200001239777 0.52308537792826
12.1600000858307 0.522707699238823
12.2000000476837 0.521878218714596
12.2400000095367 0.52083772955509
12.2800002098083 0.519775765035912
12.3200001716614 0.518685183826015
12.3600001335144 0.517508660186145
12.4000000953674 0.516022477573725
12.4400000572205 0.514687972903787
12.4800000190735 0.513180820974952
12.5199999809265 0.511621727857366
12.5600001811981 0.50984081295053
12.6000001430511 0.508042261108313
12.6400001049042 0.506485160708548
12.6800000667572 0.505151476398292
12.7200000286102 0.50398560151616
12.7599999904633 0.502900094139036
12.8000001907349 0.501858226938522
12.8400001525879 0.500699467857341
12.8800001144409 0.499148647825328
12.9200000762939 0.497619793582047
12.960000038147 0.495909619084141
13 0.494268943633728
13.0400002002716 0.492741828244786
13.0800001621246 0.49165753844967
13.1200001239777 0.490792516858827
13.1600000858307 0.490023735760005
13.2000000476837 0.489298967132004
13.2400000095367 0.48868297660046
13.2800002098083 0.488183034697325
13.3200001716614 0.487695473764414
13.3600001335144 0.487182328969994
13.4000000953674 0.48654488467549
13.4400000572205 0.486022555102209
13.4800000190735 0.486017108516137
13.5199999809265 0.486094350506732
13.5600001811981 0.486137686621905
13.6000001430511 0.486160373873264
13.6400001049042 0.486690007700896
13.6800000667572 0.487292165410659
13.7200000286102 0.487780886075753
13.7599999904633 0.488203604079186
13.8000001907349 0.488562627630611
13.8400001525879 0.488902699141377
13.8800001144409 0.489298639127612
13.9200000762939 0.489870930744165
13.960000038147 0.491682206415451
14 0.493587081169375
14.0400002002716 0.496592802324463
14.0800001621246 0.49954088333878
14.1200001239777 0.502411334710166
14.1600000858307 0.506105119221799
14.2000000476837 0.509709319967081
14.2400000095367 0.513127933168859
14.2800002098083 0.516333719423099
14.3200001716614 0.518936026250591
14.3600001335144 0.520995788345029
14.4000000953674 0.522952184863815
14.4400000572205 0.524931602565191
14.4800000190735 0.526757716239426
14.5199999809265 0.528805257840063
14.5600001811981 0.530696373103217
14.6000001430511 0.532390414721763
14.6400001049042 0.533616161955578
14.6800000667572 0.534678589127438
14.7200000286102 0.535406223392801
14.7599999904633 0.53584347725303
14.8000001907349 0.535943438390238
14.8400001525879 0.535866306529692
14.8800001144409 0.535504461360984
14.9200000762939 0.534864337103298
14.960000038147 0.534244244119427
15 0.53358052270533
15.0400002002716 0.532662227247694
15.0800001621246 0.531738019324988
15.1200001239777 0.530434322229276
15.1600000858307 0.529142164816211
15.2000000476837 0.527663568018939
15.2400000095367 0.526138424522874
15.2800002098083 0.523246916825231
15.3200001716614 0.520113195943936
15.3600001335144 0.517287435514134
15.4000000953674 0.514341723696098
15.4400000572205 0.51126924707468
15.4800000190735 0.508254271999453
15.5199999809265 0.505134135600738
15.5600001811981 0.502120773746518
15.6000001430511 0.49862134214132
15.6400001049042 0.495442587804676
15.6800000667572 0.492183659642341
15.7200000286102 0.489008688337813
15.7599999904633 0.486009476626005
15.8000001907349 0.483660649987539
15.8400001525879 0.481778747344605
15.8800001144409 0.48066107170182
15.9200000762939 0.480071437655686
15.960000038147 0.479825612237212
16 0.480141651402282
16.0400002002716 0.480394904656874
16.0800001621246 0.480661449074243
16.1200001239777 0.481472748047967
16.1600000858307 0.482287365537429
16.2000000476837 0.483615283879024
16.2400000095367 0.485002245885759
16.2800002098083 0.486449529825051
16.3200001716614 0.487852374950846
16.3600001335144 0.489005882366482
16.4000000953674 0.490540532725273
16.4400000572205 0.492189417051856
16.4800000190735 0.494046173560887
16.5199999809265 0.496147170371491
16.5600001811981 0.498259383549096
16.6000001430511 0.500697735436852
16.6400001049042 0.502930950820535
16.6800000667572 0.505448750417994
16.7200000286102 0.507832695404942
16.7599999904633 0.509973078810582
16.8000001907349 0.512532090971624
16.8400001525879 0.515229977206515
16.8800001144409 0.518323336216702
16.9200000762939 0.520793208292698
16.960000038147 0.517559422023902
17 0.513564399113204
17.0400002002716 0.509928612609625
17.0800001621246 0.505680656701469
17.1200001239777 0.499786561326895
17.1600000858307 0.493997216512978
17.2000000476837 0.486480004349246
17.2400000095367 0.479761014212639
17.2800002098083 0.471524304978844
17.3200001716614 0.463148070474397
17.3600001335144 0.449536390673046
17.4000000953674 0.435141119494483
17.4400000572205 0.422874155350239
17.4800000190735 0.409859041604239
17.5199999809265 0.396514955528463
17.5600001811981 0.383018928023737
17.6000001430511 0.368273292576262
17.6400001049042 0.352853623448394
17.6800000667572 0.333090939760707
17.7200000286102 0.314945774265597
17.7599999904633 0.292931715366997
17.8000001907349 0.272382267142471
17.8400001525879 0.249909796719576
17.8800001144409 0.227007533258429
17.9200000762939 0.205556134603106
17.960000038147 0.183123015003599
18 0.161866119933623
18.0400002002716 0.136477825554885
18.0800001621246 0.10859364002154
18.1200001239777 0.0820231452406671
18.1600000858307 0.0551840037930995
18.2000000476837 0.0281780294691293
18.2400000095367 -0.00166156217087476
18.2800002098083 -0.0294411269390502
18.3200001716614 -0.0647262780977721
18.3600001335144 -0.100047232038578
18.4000000953674 -0.133356194565386
18.4400000572205 -0.168643587611325
18.4800000190735 -0.20022827116965
18.5199999809265 -0.231791290968456
18.5600001811981 -0.263616250127143
18.6000001430511 -0.295076251382716
18.6400001049042 -0.32779320112527
18.6800000667572 -0.360207325313561
18.7200000286102 -0.40085453669958
18.7599999904633 -0.44070215252467
18.8000001907349 -0.481750767354482
18.8400001525879 -0.520361587622894
18.8800001144409 -0.555477954414314
18.9200000762939 -0.594487755733066
18.960000038147 -0.63305777529849
19 -0.676109740801417
19.0400002002716 -0.717576484920324
19.0800001621246 -0.757689262558619
19.1200001239777 -0.796505524920395
19.1600000858307 -0.834101763723671
19.2000000476837 -0.872172195942563
19.2400000095367 -0.913016155435836
19.2800002098083 -0.968049491227328
19.3200001716614 -1.01991097991935
19.3600001335144 -1.06765599933218
19.4000000953674 -1.1180306563882
19.4400000572205 -1.16629438798686
19.4800000190735 -1.21530488378435
19.5199999809265 -1.26282637580694
19.5600001811981 -1.30827650525853
19.6000001430511 -1.3483711088362
19.6400001049042 -1.3877826525578
19.6800000667572 -1.42670953515843
19.7200000286102 -1.46196574961523
19.7599999904633 -1.49793501433718
19.8000001907349 -1.54051604284473
19.8400001525879 -1.5799879611256
19.8800001144409 -1.62026129007031
19.9200000762939 -1.65852213254711
19.960000038147 -1.69807372700624
20 -1.73579085751822
20.0400002002716 -1.78144333065173
20.0800001621246 -1.82696101900612
20.1200001239777 -1.86554507952709
20.1600000858307 -1.90319968209045
20.2000000476837 -1.93972804944744
20.2400000095367 -1.97513656869397
20.2800002098083 -2.01177486442198
20.3200001716614 -2.04732632861467
20.3600001335144 -2.08936807891758
20.4000000953674 -2.12920894768779
20.4400000572205 -2.17014513101829
20.4800000190735 -2.20857703921015
20.5199999809265 -2.24874836858047
20.5600001811981 -2.28833880398407
20.6000001430511 -2.32398723712702
20.6400001049042 -2.35779262868662
20.6800000667572 -2.39064231061415
20.7200000286102 -2.42490908481529
20.7599999904633 -2.46052811872392
20.8000001907349 -2.49493405280961
20.8400001525879 -2.52879016170607
20.8800001144409 -2.56051054590224
20.9200000762939 -2.59280653394196
20.960000038147 -2.6222437538717
21 -2.6493792230323
21.0400002002716 -2.68040269635519
21.0800001621246 -2.71216477842703
21.1200001239777 -2.7442752944313
21.1600000858307 -2.77340037722636
21.2000000476837 -2.80056889123675
21.2400000095367 -2.82618478350346
21.2800002098083 -2.84921539358482
21.3200001716614 -2.87201483468714
21.3600001335144 -2.89643900154388
21.4000000953674 -2.91935013798375
21.4400000572205 -2.94074032305327
21.4800000190735 -2.96085979618567
21.5199999809265 -2.98073681292527
21.5600001811981 -2.99899311492534
21.6000001430511 -3.01748447248947
21.6400001049042 -3.03606762567086
21.6800000667572 -3.05296805890715
21.7200000286102 -3.06780476079598
21.7599999904633 -3.0814757779189
21.8000001907349 -3.09411753194364
21.8400001525879 -3.10574577956912
21.8800001144409 -3.11647873874857
21.9200000762939 -3.12586667343871
21.960000038147 -3.1343722840652
22 -3.14219651637647
22.0400002002716 -3.14867125157525
22.0800001621246 -3.15464672327943
22.1200001239777 -3.16005534907116
22.1600000858307 -3.16506436122899
22.2000000476837 -3.16960233300277
22.2400000095367 -3.17346310079766
22.2800002098083 -3.17698262871556
22.3200001716614 -3.18022028826922
22.3600001335144 -3.18318574715157
22.4000000953674 -3.18569926685406
22.4400000572205 -3.18734445343163
22.4800000190735 -3.18889711749647
22.5199999809265 -3.19029333838841
22.5600001811981 -3.19115479067629
22.6000001430511 -3.19191818456762
22.6400001049042 -3.19096453474222
22.6800000667572 -3.18954611464741
22.7200000286102 -3.18824063325985
22.7599999904633 -3.18625309557559
22.8000001907349 -3.18421658952843
22.8400001525879 -3.18178541444866
22.8800001144409 -3.17917107832383
22.9200000762939 -3.17612119904749
22.960000038147 -3.17219201616058
23 -3.16857081015958
23.0400002002716 -3.16515725619533
23.0800001621246 -3.16115810188355
23.1200001239777 -3.15739879649366
23.1600000858307 -3.1527030006046
23.2000000476837 -3.14811909138904
23.2400000095367 -3.14373820703735
23.2800002098083 -3.13677307923692
23.3200001716614 -3.12980436707365
23.3600001335144 -3.1230111220286
23.4000000953674 -3.11638853139072
23.4400000572205 -3.11044626452857
23.4800000190735 -3.10465819940002
23.5199999809265 -3.09872369727516
23.5600001811981 -3.09313368842693
23.6000001430511 -3.08761979197185
23.6400001049042 -3.08266926256685
23.6800000667572 -3.07768587405984
23.7200000286102 -3.07279755586332
23.7599999904633 -3.06819058357103
23.8000001907349 -3.06385226643817
23.8400001525879 -3.06017879579008
23.8800001144409 -3.05650434515481
23.9200000762939 -3.05283166234428
23.960000038147 -3.04953571408226
24 -3.04625645159661
24.0400002002716 -3.04241159189465
24.0800001621246 -3.03884764818621
24.1200001239777 -3.03541160452394
24.1600000858307 -3.03227214109481
24.2000000476837 -3.02930256644047
24.2400000095367 -3.02656384632114
24.2800002098083 -3.02398035216542
24.3200001716614 -3.02168990780746
24.3600001335144 -3.01999081570264
24.4000000953674 -3.01861795719642
24.4400000572205 -3.01744409467171
24.4800000190735 -3.01725621068551
24.5199999809265 -3.01859008113163
24.5600001811981 -3.01993128343023
24.6000001430511 -3.02102215142922
24.6400001049042 -3.02229521646069
24.6800000667572 -3.0235127758207
24.7200000286102 -3.0250542217152
24.7599999904633 -3.02663353472275
24.8000001907349 -3.02813112776014
24.8400001525879 -3.03131457513377
24.8800001144409 -3.03476407440966
24.9200000762939 -3.03863570358227
24.960000038147 -3.04284474806492
25 -3.04722629528103
25.0400002002716 -3.05212834414081
25.0800001621246 -3.05671280009892
25.1200001239777 -3.06101180433907
25.1600000858307 -3.06494335731499
25.2000000476837 -3.06861409720777
25.2400000095367 -3.07206865959385
25.2800002098083 -3.07538054795567
25.3200001716614 -3.07922833770367
25.3600001335144 -3.0830691049182
25.4000000953674 -3.08645949521743
25.4400000572205 -3.08964065330216
25.4800000190735 -3.09268463359324
25.5199999809265 -3.09573090250841
25.5600001811981 -3.09879765826611
25.6000001430511 -3.10155757056799
25.6400001049042 -3.10380737601632
25.6800000667572 -3.10588661248871
25.7200000286102 -3.1076445877
25.7599999904633 -3.10926472539818
25.8000001907349 -3.1107860064564
25.8400001525879 -3.11151635221304
25.8800001144409 -3.11193606990889
25.9200000762939 -3.11123573900057
25.960000038147 -3.11012607009645
26 -3.10921695826847
26.0400002002716 -3.10813323274568
26.0800001621246 -3.10695589521532
26.1200001239777 -3.10585963500949
26.1600000858307 -3.1047741822136
26.2000000476837 -3.10367425916209
26.2400000095367 -3.10199112543039
26.2800002098083 -3.10040470990576
26.3200001716614 -3.09842523185473
26.3600001335144 -3.09660275143389
26.4000000953674 -3.09491983696476
26.4400000572205 -3.09210404374657
26.4800000190735 -3.08904975209152
26.5199999809265 -3.08598222497896
26.5600001811981 -3.08295351500312
26.6000001430511 -3.07989391487829
26.6400001049042 -3.0768073492051
26.6800000667572 -3.07398176848043
26.7200000286102 -3.07142041446374
26.7599999904633 -3.06884085519035
26.8000001907349 -3.0664312114736
26.8400001525879 -3.06443046921649
26.8800001144409 -3.06259046814647
26.9200000762939 -3.06088428597908
26.960000038147 -3.05927089701911
27 -3.057863201863
27.0400002002716 -3.05614577467182
27.0800001621246 -3.05434368041142
27.1200001239777 -3.05255039938368
27.1600000858307 -3.05076923499717
27.2000000476837 -3.04912681307888
27.2400000095367 -3.04752941532641
27.2800002098083 -3.0461495349667
27.3200001716614 -3.04502254463527
27.3600001335144 -3.04519891389882
27.4000000953674 -3.04554464894153
27.4400000572205 -3.0467402093181
27.4800000190735 -3.04819348534165
27.5199999809265 -3.05020248015308
27.5600001811981 -3.05228497304395
27.6000001430511 -3.05413729221519
27.6400001049042 -3.05603369881925
27.6800000667572 -3.05795171264271
27.7200000286102 -3.05980518403021
27.7599999904633 -3.0616272723732
27.8000001907349 -3.06332656566938
27.8400001525879 -3.06506498571582
27.8800001144409 -3.06670845008976
27.9200000762939 -3.06819324484233
27.960000038147 -3.07142120880214
28 -3.0747425711339
28.0400002002716 -3.07892630048381
28.0800001621246 -3.08315042161306
28.1200001239777 -3.08676719839462
28.1600000858307 -3.09083348408661
28.2000000476837 -3.09518268124913
28.2400000095367 -3.09929184238516
28.2800002098083 -3.10331192038176
28.3200001716614 -3.10690328534861
28.3600001335144 -3.11011354009615
28.4000000953674 -3.1131590884315
28.4400000572205 -3.11524478376099
28.4800000190735 -3.11721636773006
28.5199999809265 -3.11916297325803
28.5600001811981 -3.12098541192642
28.6000001430511 -3.12266619658216
28.6400001049042 -3.12413676233627
28.6800000667572 -3.1254747602067
28.7200000286102 -3.12691185707381
28.7599999904633 -3.12846578872467
28.8000001907349 -3.12989501329388
28.8400001525879 -3.13126186887142
28.8800001144409 -3.13267266867584
28.9200000762939 -3.13395580594061
28.960000038147 -3.13471372949924
29 -3.13538757328855
29.0400002002716 -3.13490932103054
29.0800001621246 -3.1341252739791
29.1200001239777 -3.13341849006276
29.1600000858307 -3.13232679834441
29.2000000476837 -3.13122018636376
29.2400000095367 -3.12997079926449
29.2800002098083 -3.12863161921327
29.3200001716614 -3.12766667516988
29.3600001335144 -3.12713027711098
29.4000000953674 -3.1266910232181
29.4400000572205 -3.12629975284805
29.4800000190735 -3.12583140628429
29.5199999809265 -3.12544686715797
29.5600001811981 -3.12390693799057
29.6000001430511 -3.12227947671975
29.6400001049042 -3.11978546323652
29.6800000667572 -3.11714347716993
29.7200000286102 -3.11485986613387
29.7599999904633 -3.11271153060996
29.8000001907349 -3.11064119611469
29.8400001525879 -3.10852454494824
29.8800001144409 -3.10626817302645
29.9200000762939 -3.10426862792158
29.960000038147 -3.102470316615
30 -3.10080797207185
30.0400002002716 -3.09947724489471
30.0800001621246 -3.09852574292678
30.1200001239777 -3.09769334645661
30.1600000858307 -3.09665787860721
30.2000000476837 -3.09564611681417
30.2400000095367 -3.0944041163228
30.2800002098083 -3.09322931651988
30.3200001716614 -3.09221760720683
30.3600001335144 -3.09091134842154
30.4000000953674 -3.08927227027325
30.4400000572205 -3.08761952013195
30.4800000190735 -3.08567962711249
30.5199999809265 -3.08391170246022
30.5600001811981 -3.08215708845827
30.6000001430511 -3.08046221756638
30.6400001049042 -3.07844508410792
30.6800000667572 -3.07642796688965
30.7200000286102 -3.07455139759923
30.7599999904633 -3.07154034670631
30.8000001907349 -3.06842096207526
30.8400001525879 -3.06495137941348
30.8800001144409 -3.06127082902975
30.9200000762939 -3.05795891665044
30.960000038147 -3.05451936904564
31 -3.05032685485218
31.0400002002716 -3.0458846782397
31.0800001621246 -3.03942147273481
31.1200001239777 -3.03337294953614
31.1600000858307 -3.02621415403253
31.2000000476837 -3.01921699007315
31.2400000095367 -3.01257162831714
31.2800002098083 -3.00489113274947
31.3200001716614 -2.99745877256903
31.3600001335144 -2.98760299835945
31.4000000953674 -2.97698141416926
31.4400000572205 -2.96584412151992
31.4800000190735 -2.95404659434919
31.5199999809265 -2.94226596984242
31.5600001811981 -2.92936315032436
31.6000001430511 -2.91695392775296
31.6400001049042 -2.89905599374622
31.6800000667572 -2.88168217302807
31.7200000286102 -2.86509627952803
31.7599999904633 -2.84596053977222
31.8000001907349 -2.82661375479393
31.8400001525879 -2.80698480522511
31.8800001144409 -2.78839463494095
31.9200000762939 -2.76614228427842
31.960000038147 -2.74142163191376
32 -2.71894206978041
32.0400002002716 -2.69695768531981
32.0800001621246 -2.67326320789268
32.1200001239777 -2.65080879597698
32.1600000858307 -2.62572252576284
32.2000000476837 -2.6020264114522
32.2400000095367 -2.56920242720012
32.2800002098083 -2.53735767543013
32.3200001716614 -2.50207906979168
32.3600001335144 -2.46802778932719
32.4000000953674 -2.43880097392883
32.4400000572205 -2.40988958326224
32.4800000190735 -2.38184194627937
32.5199999809265 -2.35352341977618
32.5600001811981 -2.32471803490088
32.6000001430511 -2.29536513174775
32.6400001049042 -2.260115290012
32.6800000667572 -2.23197929838065
32.7200000286102 -2.19604967731882
32.7599999904633 -2.15758113755654
32.8000001907349 -2.12269741627009
32.8400001525879 -2.08866960478418
32.8800001144409 -2.05553259784881
32.9200000762939 -2.01525945631235
32.960000038147 -1.97671398049917
33 -1.93924104769494
33.0400002002716 -1.90277044243617
33.0800001621246 -1.86712145972105
33.1200001239777 -1.83278821936665
33.1600000858307 -1.79745270023536
33.2000000476837 -1.76215683394709
33.2400000095367 -1.7274161989365
33.2800002098083 -1.68559335896546
33.3200001716614 -1.64719181793382
33.3600001335144 -1.59536739468671
33.4000000953674 -1.54347700936637
33.4400000572205 -1.49439716594584
33.4800000190735 -1.44719190199054
33.5199999809265 -1.40686581037722
33.5600001811981 -1.3645624196192
33.6000001430511 -1.32289170731313
33.6400001049042 -1.28580426174661
33.6800000667572 -1.24778905473918
33.7200000286102 -1.20982473069319
33.7599999904633 -1.17062406807163
33.8000001907349 -1.1339732081369
33.8400001525879 -1.09528476817552
33.8800001144409 -1.05709751020293
33.9200000762939 -1.02131671136402
33.960000038147 -0.951625583804719
34 -0.943213589384272
34.0400002002716 -0.905293540085858
34.0800001621246 -0.869052156940205
34.1200001239777 -0.834685359899977
34.1600000858307 -0.800359213922397
};
\node at (axis cs:13.5,-2)[
  anchor=north,
  text=black,
  rotate=0.0
]{ Pickup truck};
\node at (axis cs:17,2.2)[
  anchor=north,
  text=black,
  rotate=0.0
]{ Station wagon};
\end{axis}

\end{tikzpicture}%
		\subcaption{Lateral distance of other vehicles with respect to ego vehicle.}
		\label{fig:example dynamic environment lat}
	\end{minipage}%
	\caption{States of the dynamic environment.}
	\label{fig:example dynamic environment}
\end{figure}

The station wagon starts with acceleration. The parameters for this activity are shown in \cref{tab:parameters acceleration station wagon}. After accelerating, the station wagon starts cruising, for which the parameters are presented in \cref{tab:parameters station wagon cruising}. The parameters for the lane change of the station wagon are presented in \cref{tab:parameters lane change station wagon} and the parameters for the straight-driving activities are presented in \cref{tab:parameters station wagon straight 1} and \cref{tab:parameters station wagon straight 2}. The pickup truck is driving straight and at a constant speed. The normalization parameters are presented in \cref{tab:parameters pickup}.

\begin{table}
	\centering
	\caption{Parameters of the acceleration activity of the station wagon.}
	\label{tab:parameters acceleration station wagon}
	\begin{minipage}[b]{.33\linewidth}
		\centering
		\subcaption{Normalization parameters.}
		\begin{tabular}{cr}
    \toprule
    Parameter & Value \\ \otoprule
    $x_{\textup{start}}$    & $\unit[-2.33 \cdot 10^{1}]{m}$ \\
    $x_{\textup{end}}$      & $\unit[2.12 \cdot 10^{2}]{m}$ \\
    $t_{\textup{start}}$    & $\unit[0.00]{s}$ \\
    $t_{\textup{duration}}$ & $\unit[1.32 \cdot 10^{1}]{s}$ \\
    \bottomrule
\end{tabular}
%
	\end{minipage}%
	\begin{minipage}[b]{.5\linewidth}
		\centering
		\subcaption{Spline coefficients.}
		\begin{tabular}{crr}
    \toprule
    Coefficients & Spline 1 & Spline 2 \\ \otoprule
    $a$ & $2.63 \cdot 10^{-2}$ & $8.86 \cdot 10^{-1}$ \\
    $b$ & $2.58 \cdot 10^{-1}$ & $-1.03$ \\
    $c$ & $6.08 \cdot 10^{-1}$ & $1.25$ \\
    $d$ & $3.59 \cdot 10^{-16}$ & $-1.07 \cdot 10^{-1}$ \\
    \bottomrule
\end{tabular}
%
	\end{minipage}%
\end{table}

\begin{table}
	\centering
	\caption{Parameters of the lane change activity of the station wagon.}
	\label{tab:parameters lane change station wagon}
	\begin{minipage}[b]{.33\linewidth}
		\centering
		\subcaption{Normalization parameters.}
		\begin{tabular}{cr}
    \toprule
    Parameter & Value \\ \otoprule
    $x_{\textup{start}}$    & $\unit[2.94 \cdot 10^{-1}]{m}$ \\
    $x_{\textup{end}}$      & $\unit[2.86]{m}$ \\
    $t_{\textup{start}}$    & $\unit[9.20]{s}$ \\
    $t_{\textup{duration}}$ & $\unit[2.56]{s}$ \\
    \bottomrule
\end{tabular}
%
	\end{minipage}%
	\begin{minipage}[b]{.5\linewidth}
		\centering
		\subcaption{Spline coefficients.}
		\begin{tabular}{crr}
    \toprule
    Coefficients & Spline 1 & Spline 2 \\ \otoprule
    $a$ & $2.71$ & $-7.65$ \\
    $b$ & $-8.75 \cdot 10^{-1}$ & $1.47 \cdot 10^{1}$ \\
    $c$ & $4.64 \cdot 10^{-1}$ & $-7.30$ \\
    $d$ & $7.29 \cdot 10^{-16}$ & $1.29$ \\
    \bottomrule
\end{tabular}
%
	\end{minipage}%
\end{table}

\begin{table}
	\centering
	\caption{Normalization parameters of the cruising and straight-driving activities of the station wagon.}
	\label{tab:parameters station wagon cruising and straight}
	\begin{minipage}[b]{.33\linewidth}
		\centering
		\subcaption{Cruising.}
		\label{tab:parameters station wagon cruising}
		\begin{tabular}{cr}
    \toprule
    Parameter & Value \\ \otoprule
    $x_{\textup{start}}$    & $\unit[2.14 \cdot 10^{2}]{m}$ \\
    $x_{\textup{end}}$      & $\unit[3.53 \cdot 10^{2}]{m}$ \\
    $t_{\textup{start}}$    & $\unit[1.32 \cdot 10^{1}]{s}$ \\
    $t_{\textup{duration}}$ & $\unit[4.52]{s}$ \\
    \bottomrule
\end{tabular}
%
	\end{minipage}%
	\begin{minipage}[b]{.33\linewidth}
		\centering
		\subcaption{Straight before lane change.}
		\label{tab:parameters station wagon straight 1}
		\begin{tabular}{cr}
    \toprule
    Parameter & Value \\ \otoprule
    $x_{\textup{start}}$    & $\unit[5.20 \cdot 10^{-2}]{m}$ \\
    $x_{\textup{end}}$      & $\unit[2.61 \cdot 10^{-1}]{m}$ \\
    $t_{\textup{start}}$    & $\unit[0.00]{s}$ \\
    $t_{\textup{duration}}$ & $\unit[9.16]{s}$ \\
    \bottomrule
\end{tabular}
%
	\end{minipage}%
	\begin{minipage}[b]{.33\linewidth}
		\centering
		\subcaption{Straight after lane change.}
		\label{tab:parameters station wagon straight 2}
		\begin{tabular}{cr}
    \toprule
    Parameter & Value \\ \otoprule
    $x_{\textup{start}}$    & $\unit[2.86]{m}$ \\
    $x_{\textup{end}}$      & $\unit[3.18]{m}$ \\
    $t_{\textup{start}}$    & $\unit[1.18 \cdot 10^{1}]{s}$ \\
    $t_{\textup{duration}}$ & $\unit[5.92]{s}$ \\
    \bottomrule
\end{tabular}
%
	\end{minipage}%
\end{table}

\begin{table}
	\centering
	\caption{Normalization parameters of the cruising and straight-driving activities of the pickup truck.}
	\label{tab:parameters pickup}
	\begin{minipage}[b]{.33\linewidth}
		\centering
		\subcaption{Cruising.}
		\label{tab:parameters pickup cruising}
		\begin{tabular}{cr}
    \toprule
    Parameter & Value \\ \otoprule
    $x_{\textup{start}}$    & $\unit[1.17 \cdot 10^{2}]{m}$ \\
    $x_{\textup{end}}$      & $\unit[5.56 \cdot 10^{2}]{m}$ \\
    $t_{\textup{start}}$    & $\unit[8.00 \cdot 10^{-1}]{s}$ \\
    $t_{\textup{duration}}$ & $\unit[3.34 \cdot 10^{1}]{s}$ \\
    \bottomrule
\end{tabular}
%
	\end{minipage}%
	\begin{minipage}[b]{.33\linewidth}
		\centering
		\subcaption{Straight driving.}
		\label{tab:parameters pickup straight}
		\begin{tabular}{cr}
    \toprule
    Parameter & Value \\ \otoprule
    $x_{\textup{start}}$    & $\unit[6.72 \cdot 10^{-1}]{m}$ \\
    $x_{\textup{end}}$      & $\unit[3.60 \cdot 10^{-1}]{m}$ \\
    $t_{\textup{start}}$    & $\unit[8.00 \cdot 10^{-1}]{s}$ \\
    $t_{\textup{duration}}$ & $\unit[3.34 \cdot 10^{1}]{s}$ \\
    \bottomrule
\end{tabular}
%
	\end{minipage}%
\end{table}


\section{Scenarios}

Scenarios are constructed by combining the activities of all road users (ego + dynamic environment) with the static environment. Since currently we don't have enough reliable information on the static environment, we focus on the ego and its dynamic environment for constructing the first scenarios. As an example, we take the 'ego overtaking 1' scenario. The first relevant event is the start of the lane change of the ego vehicle, indicated by the red vertical line in \cref{fig:example events}. The scenario ends at the end of the second lane change. The relevant activities are therefore G,H,I for the lateral ego activities and accordingly C,D,E for the longitudinal ego activities. Similarly, the scenario contains the activities of the pickup (cruising + lane following). 

\printbibliography

\end{document} 